% Miss Assistant — LaTeX Template for Print-Ready Textbook Sample
% Based on LAYOUT.md specs: 17×24cm, Insper brand colors

\documentclass[11pt,twoside,openright]{book}

% === Page Geometry ===
\usepackage[
  paperwidth=17cm,
  paperheight=24cm,
  top=2.5cm,
  bottom=2.5cm,
  inner=3cm,
  outer=2cm
]{geometry}

% === Language ===
\usepackage[brazil]{babel}

% === Fonts ===
\usepackage{fontspec}
\setmainfont{DejaVu Serif}
\setsansfont{Inter}
\setmonofont{DejaVu Sans Mono}[Scale=0.85]

% === Colors (from LAYOUT.md) ===
\usepackage{xcolor}
\definecolor{insper}{HTML}{C8102E}
\definecolor{brasilgreen}{HTML}{2E7D32}
\definecolor{brasilbg}{HTML}{E8F5E9}
\definecolor{mundoblue}{HTML}{1B3A5C}
\definecolor{mundobg}{HTML}{E3EBF2}
\definecolor{exampurple}{HTML}{6A1B9A}
\definecolor{exambg}{HTML}{F3E5F5}
\definecolor{warnorange}{HTML}{E65100}
\definecolor{warnbg}{HTML}{FFF3E0}
\definecolor{defblue}{HTML}{1565C0}
\definecolor{defbg}{HTML}{E3F2FD}
\definecolor{tipgreen}{HTML}{43A047}
\definecolor{tipbg}{HTML}{E8F5E9}
\definecolor{gray}{HTML}{757575}
\definecolor{graybg}{HTML}{F5F5F5}
\definecolor{bodytext}{HTML}{333333}

% === Math ===
\usepackage{amsmath,amssymb,amsfonts,mathtools}
\usepackage{unicode-math}

% === Graphics ===
\usepackage{graphicx}
\usepackage{float}

% === Tables ===
\usepackage{booktabs}
\usepackage{longtable}
\usepackage{array}
\usepackage{multirow}

% === Boxes (tcolorbox) ===
\usepackage[most]{tcolorbox}

% Base box style
\tcbset{
  basebox/.style={
    enhanced,
    breakable,
    left=8pt,
    right=8pt,
    top=6pt,
    bottom=6pt,
    boxrule=0pt,
    frame hidden,
    borderline west={3pt}{0pt}{#1},
    fonttitle=\sffamily\bfseries\color{#1},
    coltitle=#1,
    colback=#1!5!white,
    coltext=bodytext,
    attach title to upper={\par\smallskip},
    before upper={\setlength{\parindent}{0pt}\setlength{\parskip}{4pt}},
  }
}

% === Box environments ===

% Definition
\newtcolorbox{definitionbox}[1][]{
  basebox=defblue,
  title={#1},
}

% Theorem
\newtcolorbox{theorembox}[1][]{
  basebox=defblue,
  title={#1},
}

% Proof
\newtcolorbox{proofbox}[1][]{
  basebox=gray,
  title={#1},
  colback=graybg,
}

% Intuição Econômica
\newtcolorbox{intuitionbox}[1][]{
  basebox=insper,
  title={#1},
  colback=insper!5!white,
}

% Tip
\newtcolorbox{tipbox}[1][]{
  basebox=tipgreen,
  title={#1},
  colback=tipbg,
}

% Note
\newtcolorbox{notebox}[1][]{
  basebox=defblue,
  title={#1},
  colback=defbg,
}

% Warning / Erro Comum
\newtcolorbox{warningbox}[1][]{
  basebox=warnorange,
  title={#1},
  colback=warnbg,
}

% Example
\newtcolorbox{examplebox}[1][]{
  basebox=exampurple,
  title={#1},
  colback=exambg,
}

% Success (answer)
\newtcolorbox{successbox}[1][]{
  basebox=tipgreen,
  title={#1},
  colback=tipbg,
}

% Info
\newtcolorbox{infobox}[1][]{
  basebox=defblue,
  title={#1},
  colback=defbg,
}

% Box Brasil
\newtcolorbox{boxbrasilbox}[1][]{
  basebox=brasilgreen,
  title={\raisebox{-1pt}{🇧🇷}\ #1},
  colback=brasilbg,
}

% Box Mundo
\newtcolorbox{boxmundobox}[1][]{
  basebox=mundoblue,
  title={\raisebox{-1pt}{🌍}\ #1},
  colback=mundobg,
}

% Exercício Resolvido
\newtcolorbox{exresolvidobox}[1][]{
  basebox=exampurple,
  title={#1},
  colback=exambg,
}

% Classroom Activity
\newtcolorbox{classroombox}[1][]{
  basebox=insper,
  title={#1},
  colback=insper!5!white,
}

% Pesquisa em Ação
\newtcolorbox{pesquisabox}[1][]{
  basebox=mundoblue,
  title={#1},
  colback=mundobg,
}

% === Figure placeholder ===
\newcommand{\figurePlaceholder}[1]{%
  \begin{center}
    \begin{tcolorbox}[
      width=0.85\textwidth,
      colback=graybg,
      colframe=gray,
      boxrule=0.5pt,
      arc=2pt,
      halign=center,
      valign=center,
      height=5cm
    ]
      \centering
      \vfill
      {\sffamily\color{gray}\large [Figura Interativa]}\\[6pt]
      {\sffamily\small\color{gray} Versão interativa disponível online}\\[4pt]
      {\ttfamily\tiny\color{gray} #1}
      \vfill
    \end{tcolorbox}
  \end{center}
}

% === Headers and Footers ===
\usepackage{fancyhdr}
\pagestyle{fancy}
\fancyhf{}
\fancyhead[LE]{\small\sffamily\textcolor{insper}{\leftmark}}
\fancyhead[RO]{\small\sffamily\textcolor{insper}{\rightmark}}
\fancyfoot[LE,RO]{\small\sffamily\thepage}
\renewcommand{\headrulewidth}{0.4pt}
\renewcommand{\headrule}{\hbox to\headwidth{\color{insper}\leaders\hrule height \headrulewidth\hfill}}

% Chapter pages: plain style
\fancypagestyle{plain}{
  \fancyhf{}
  \fancyfoot[C]{\small\sffamily\thepage}
  \renewcommand{\headrulewidth}{0pt}
}

% === Chapter Styling ===
\usepackage{titlesec}

\titleformat{\chapter}[display]
  {\sffamily\bfseries}
  {\textcolor{insper}{\fontsize{48}{52}\selectfont\thechapter}}
  {10pt}
  {\fontsize{24}{28}\selectfont\textcolor{insper}}

\titleformat{\section}
  {\Large\sffamily\bfseries\color{insper!80!black}}
  {\thesection}{0.8em}{}

\titleformat{\subsection}
  {\large\sffamily\bfseries\color{insper!60!black}}
  {\thesubsection}{0.6em}{}

\titleformat{\subsubsection}
  {\normalsize\sffamily\bfseries\color{bodytext}}
  {\thesubsubsection}{0.5em}{}

% === Hyperlinks ===
\usepackage{hyperref}
\hypersetup{
  colorlinks=true,
  linkcolor=insper,
  urlcolor=mundoblue,
  citecolor=brasilgreen,
  pdfauthor={Darcio Genicolo Martins},
  pdftitle={Microeconomia: Do Zero ao Equilíbrio (e Além) — Amostra},
  pdflang={pt-BR},
}

% === Lists ===
\usepackage{enumitem}
\setlist{nosep, leftmargin=*}

% === Body text color ===
\color{bodytext}

% === Line spacing ===
\usepackage{setspace}
\setstretch{1.2}

% === Pandoc-specific ===
\providecommand{\tightlist}{\setlength{\itemsep}{0pt}\setlength{\parskip}{0pt}}

% Allow pandoc code blocks

% === Document ===
\begin{document}

% ============ COVER PAGE ============
\begin{titlepage}
\begin{center}
\vspace*{3cm}

{\fontsize{32}{38}\selectfont\sffamily\bfseries\textcolor{insper}{Microeconomia}}

\vspace{0.4cm}

{\fontsize{16}{20}\selectfont\sffamily Do Zero ao Equilíbrio (e Além)}

\vspace{2cm}

{\Large\sffamily Capítulos Selecionados para Avaliação Editorial}

\vspace{1.5cm}

{\large\sffamily Prof. Darcio Genicolo Martins}

{\normalsize\sffamily Insper — Instituto de Ensino e Pesquisa}

\vspace{0.5cm}

{\normalsize\sffamily São Paulo, Março de 2026}

\vspace{3cm}

\rule{0.4\textwidth}{0.5pt}

\vspace{0.3cm}

{\small\itshape Documento confidencial — para avaliação editorial}

\vspace{1cm}

{\small\sffamily Capítulos incluídos:}\\[6pt]
{\small Cap.\ 3 — Mais É Melhor — Será?}\\
{\small Cap.\ 15 — Sozinho, Feliz e Cobrando Caro}\\
{\small Cap.\ 20 — Fumaça, Caronas e Tragédias}

\end{center}
\end{titlepage}

% ============ TOC ============
\tableofcontents
\newpage

% ============ CONTENT ============
\chapter{Capítulo 3 --- Mais É Melhor ---
Será?}\label{capuxedtulo-3-mais-uxe9-melhor-seruxe1}

Netflix, Disney+ ou HBO Max? Você só pode assinar um. A escolha parece
simples --- até que você percebe que prefere Netflix a Disney+, Disney+
a HBO Max, mas HBO Max a Netflix (por causa daquela \emph{única} série).
Parabéns: suas preferências são intransitivas, e a microeconomia acabou
de quebrar. Ou será que não?

Você entra no supermercado com R\$ 200 e sai com um carrinho cheio de
escolhas. Arroz ou macarrão? Cerveja ou refrigerante? Marca A ou marca
B? Cada decisão revela algo sobre suas preferências --- algo que você
talvez nem consiga articular, mas que a microeconomia pretende capturar
com precisão matemática. A teoria das preferências e da utilidade é o
alicerce de tudo que vem depois: curvas de demanda, análise de
bem-estar, efeitos de políticas públicas. Se este capítulo for bem
assimilado, o resto do livro é consequência. Se não for\ldots{} bem, vai
ser uma longa jornada. (E aquele dilema do streaming? Voltamos a ele na
Seção 3.1, quando discutirmos por que a transitividade é tão importante
--- e o que acontece quando ela falha.)

Neste capítulo, formalizamos as preferências individuais por meio de
axiomas --- regras mínimas de coerência que, quando satisfeitas,
permitem representar o comportamento do consumidor por uma função de
utilidade. A partir dessa função, construímos progressivamente as
ferramentas analíticas centrais: curvas de indiferença, que oferecem uma
representação geométrica das preferências; a taxa marginal de
substituição (TMS), que quantifica a disposição do consumidor a trocar
um bem por outro; e as famílias específicas de funções de utilidade ---
Cobb-Douglas, CES, substitutos e complementos perfeitos, quase-linear
--- que dão conteúdo empírico ao arcabouço teórico.

O capítulo segue a abordagem axiomática presente em Mas-Colell, Whinston
e Green (1995), complementada pela exposição mais intuitiva de Varian
(2015) e Nicholson e Snyder (2017). Para uma abordagem particularmente
acessível, com ênfase em intuição econômica e exemplos aplicados, ver
Besanko e Braeutigam (2014, Caps. 3--4). Ao final do capítulo, o leitor
terá em mãos todo o aparato conceitual necessário para resolver o
problema de otimização do consumidor --- tema central do Capítulo 4,
onde as preferências aqui formalizadas se encontram com a restrição
orçamentária para determinar as escolhas ótimas.

\begin{theorembox}{Roteiro do capítulo}

{\def\LTcaptype{none} % do not increment counter
\begin{longtable}[]{@{}
  >{\centering\arraybackslash}p{(\linewidth - 6\tabcolsep) * \real{0.3125}}
  >{\raggedright\arraybackslash}p{(\linewidth - 6\tabcolsep) * \real{0.1875}}
  >{\raggedright\arraybackslash}p{(\linewidth - 6\tabcolsep) * \real{0.1875}}
  >{\centering\arraybackslash}p{(\linewidth - 6\tabcolsep) * \real{0.3125}}@{}}
\toprule\noalign{}
\begin{minipage}[b]{\linewidth}\centering
Seção
\end{minipage} & \begin{minipage}[b]{\linewidth}\raggedright
Pergunta-guia
\end{minipage} & \begin{minipage}[b]{\linewidth}\raggedright
O que você vai aprender
\end{minipage} & \begin{minipage}[b]{\linewidth}\centering
Página
\end{minipage} \\
\midrule\noalign{}
\endhead
\bottomrule\noalign{}
\endlastfoot
\textbf{3.1--3.3} & O que torna um consumidor ``coerente''? & Axiomas,
Teorema de Debreu, curvas de indiferença & Axiomas \\
\textbf{3.4--3.5} & A que ``preço pessoal'' você troca um bem por outro?
& TMS, utilidade marginal, TMS = UMg₁/UMg₂ & TMS \\
\textbf{3.6--3.7} & Que ``personagem'' descreve melhor suas
preferências? & Cobb-Douglas, CES, Leontief, quase-linear, homotéticas,
transformações monotônicas & Preferências \\
\textbf{Exercícios} & Teste e aplique & Revisão, exercícios, ANPEC,
atividades de sala & Exercícios \\
\textbf{Pesquisa} & O que a fronteira diz? & Artigos seminais com dados
reais & Pesquisa \\
\end{longtable}
}

\end{theorembox}

\chapter{3.1--3.3 Axiomas, Utilidade e Curvas de
Indiferença}\label{axiomas-utilidade-e-curvas-de-indiferenuxe7a}

\section{3.1 As Regras Mínimas da Sanidade: Axiomas da
Escolha}\label{as-regras-muxednimas-da-sanidade-axiomas-da-escolha}

Antes de construir qualquer modelo de escolha, precisamos responder a
uma pergunta incômoda: o que \emph{exatamente} estamos supondo sobre o
consumidor? Não vamos supor que ele é gênio, nem que usa planilha Excel
no supermercado. Vamos supor apenas que ele é \emph{coerente} --- que
consegue comparar opções e não se contradiz.\footnote{No sketch
  \emph{Dead Parrot} de Monty Python, o cliente e o vendedor discordam
  se o papagaio está morto ou ``descansando''. O cliente insiste: ``This
  parrot is no more! It has ceased to be!'' O vendedor: ``No, no, he's
  resting!'' O axioma da completude exige que ambos \emph{consigam}
  comparar as cestas ``papagaio vivo'' e ``papagaio morto'' --- o que
  claramente fazem. O problema é que discordam. A boa notícia: a teoria
  do consumidor não exige que todos concordem sobre o ranking; exige
  apenas que \emph{cada} consumidor seja coerente consigo mesmo. Essas
  exigências mínimas de bom senso são formalizadas como \textbf{axiomas}
  das preferências. A beleza da coisa: desse ponto de partida modesto,
  toda a teoria do consumidor pode ser construída. É como erguer um
  arranha-céu a partir de quatro tijolos bem colocados.}

Seja \(X \subseteq \mathbb{R}^n_+\) o \textbf{conjunto de consumo}, isto
é, o conjunto de todas as cestas de bens fisicamente disponíveis ao
consumidor. Uma cesta \(\mathbf{x} = (x_1, x_2, \ldots, x_n)\)
especifica a quantidade de cada um dos \(n\) bens --- por exemplo, em um
modelo com dois bens, \(\mathbf{x} = (x_1, x_2)\) poderia representar
quilos de arroz e litros de leite. O conjunto de consumo inclui todas as
combinações não negativas dessas quantidades, formando o quadrante
positivo do espaço \(\mathbb{R}^n\).

Uma \textbf{relação de preferência} \(\succsim\) sobre \(X\) indica que,
dadas duas cestas \(\mathbf{x}\) e \(\mathbf{y}\), o consumidor
considera \(\mathbf{x}\) pelo menos tão boa quanto \(\mathbf{y}\). A
notação \(\mathbf{x} \succsim \mathbf{y}\) se lê ``\(\mathbf{x}\) é
fracamente preferida a \(\mathbf{y}\)''. Note que a relação \(\succsim\)
é \textbf{primitiva} --- ela é o ponto de partida da teoria, não algo
derivado de uma função de utilidade. Tomamos como dado que o consumidor
\emph{possui} preferências; o que os axiomas fazem é impor estrutura
sobre essas preferências, de modo a torná-las matematicamente tratáveis.
A função de utilidade, como veremos na próxima seção, é uma
\emph{consequência} dos axiomas impostos sobre \(\succsim\), não uma
hipótese adicional.

A partir de \(\succsim\), definimos duas relações derivadas que
utilizaremos extensivamente:

\begin{itemize}
\tightlist
\item
  \textbf{Preferência estrita}: \(\mathbf{x} \succ \mathbf{y}\) se e
  somente se \(\mathbf{x} \succsim \mathbf{y}\) e não
  \(\mathbf{y} \succsim \mathbf{x}\). O consumidor considera
  \(\mathbf{x}\) estritamente melhor que \(\mathbf{y}\).
\item
  \textbf{Indiferença}: \(\mathbf{x} \sim \mathbf{y}\) se e somente se
  \(\mathbf{x} \succsim \mathbf{y}\) e
  \(\mathbf{y} \succsim \mathbf{x}\). O consumidor é igualmente
  satisfeito com ambas as cestas.
\end{itemize}

Para que as preferências do consumidor sejam ``bem comportadas'' e
passíveis de análise formal, exigimos um conjunto de axiomas. Cada um
deles captura uma exigência de coerência que, embora possa parecer óbvia
isoladamente, desempenha um papel específico e insubstituível na
construção do edifício teórico. Abandonar qualquer um deles --- como
ilustra o caso das preferências lexicográficas --- pode comprometer
resultados fundamentais, como a própria existência de uma função de
utilidade.

\begin{definitionbox}{Axioma 1 — Completude}

Para quaisquer \(\mathbf{x}, \mathbf{y} \in X\), vale
\(\mathbf{x} \succsim \mathbf{y}\) ou \(\mathbf{y} \succsim \mathbf{x}\)
(ou ambos). Em outras palavras, o consumidor é sempre capaz de comparar
duas cestas quaisquer.

\end{definitionbox}

\begin{definitionbox}{Axioma 2 — Transitividade}

Para quaisquer \(\mathbf{x}, \mathbf{y}, \mathbf{z} \in X\), se
\(\mathbf{x} \succsim \mathbf{y}\) e \(\mathbf{y} \succsim \mathbf{z}\),
então \(\mathbf{x} \succsim \mathbf{z}\). A transitividade garante
consistência interna nas escolhas.

\end{definitionbox}

\begin{definitionbox}{Axioma 3 — Continuidade}

Para todo \(\mathbf{y} \in X\), os conjuntos
\(\{\mathbf{x} \in X : \mathbf{x} \succsim \mathbf{y}\}\) (conjunto de
pelo menos tão bom quanto) e
\(\{\mathbf{x} \in X : \mathbf{y} \succsim \mathbf{x}\}\) (conjunto de
no máximo tão bom quanto) são fechados. A continuidade exclui ``saltos''
nas preferências --- pequenas mudanças na cesta não causam inversões
bruscas no ordenamento.

\end{definitionbox}

\begin{definitionbox}{Axioma 4 — Monotonicidade (não saciedade local)}

Se \(\mathbf{x} \geq \mathbf{y}\) (componente a componente) e
\(\mathbf{x} \neq \mathbf{y}\), então \(\mathbf{x} \succ \mathbf{y}\).
``Mais é melhor'' --- o consumidor sempre prefere ter mais de pelo menos
um bem, tudo o mais constante.

\end{definitionbox}

Uma nota de precaução que desenvolveremos na próxima seção: a função de
utilidade que construiremos será \emph{ordinal} --- os números que ela
atribui importam apenas pela ordem, não pela magnitude. Dizer que
\(u(A) = 10\) e \(u(B) = 5\) significa que A é preferido a B, não que A
é ``duas vezes melhor''.

\begin{tipbox}{Observação sobre preferências lexicográficas}

As preferências lexicográficas satisfazem completude, transitividade e
monotonicidade, mas \textbf{não} satisfazem continuidade. Por isso, não
podem ser representadas por uma função de utilidade contínua. Este é um
caso clássico que ilustra a importância do axioma de continuidade.

\end{tipbox}

\begin{infobox}{Exemplo — Racionalidade nas escolhas de consumo no Brasil}

Os axiomas de completude e transitividade podem parecer abstratos, mas
refletem exigências mínimas de consistência nas escolhas cotidianas. A
Pesquisa de Orçamentos Familiares (POF 2017-2018) do
\href{https://www.ibge.gov.br}{IBGE} mostra que famílias brasileiras,
mesmo as de baixa renda, alocam seus orçamentos de forma estável e
internamente consistente: quando a renda aumenta, a participação
relativa da alimentação cai de maneira suave e previsível (de 22\% para
famílias com renda até R\$ 1.908 a 7,6\% para famílias acima de R\$
23.850). Essa regularidade empírica é compatível com preferências que
satisfazem os axiomas acima --- especialmente transitividade e
monotonicidade --- e justifica o uso do arcabouço axiomático como ponto
de partida para modelar o consumidor.

\end{infobox}

\begin{intuitionbox}{Intuição Econômica}

\textbf{Em uma frase:} Os axiomas de preferência são apenas regras
mínimas de coerência --- exigem que o consumidor saiba comparar opções e
não se contradiga.

\textbf{Pense assim:} Quando você vai à feira e decide que prefere
banana a maçã, e maçã a pera, seria estranho voltar e dizer que prefere
pera a banana. Completude é conseguir comparar qualquer par de frutas;
transitividade é não entrar em contradição. A monotonicidade diz que,
entre dois sacos de frutas, você prefere o maior --- desde que goste de
todas.

\textbf{Por que isso importa:} Sem essas regras mínimas, não seria
possível prever o comportamento do consumidor nem formular políticas
públicas baseadas em escolhas racionais.

\end{intuitionbox}

\begin{notebox}{Conexão com a teoria da preferência revelada}

Os axiomas sobre \(\succsim\) são postulados sobre preferências
\emph{internas} do consumidor --- algo que, em princípio, não podemos
observar diretamente. Existe, porém, uma contrapartida observacional: o
\textbf{Axioma Generalizado da Preferência Revelada (GARP)}, que traduz
completude e transitividade em condições testáveis sobre \emph{escolhas
de consumo observadas}. Se as compras de um consumidor satisfazem o
GARP, seus dados são consistentes com a maximização de alguma função de
utilidade bem comportada. Na seção ``Pesquisa em Ação'' ao final deste
capítulo, o artigo de Choi et al.~(2014) testa exatamente essa hipótese
em uma amostra representativa --- com resultados surpreendentes sobre a
relação entre consistência das escolhas e acumulação de riqueza.

\end{notebox}

Antes de prosseguir para a função de utilidade, vale uma pausa para quem
quer saber: esses axiomas funcionam na prática? O Box a seguir examina a
evidência experimental --- e a resposta é surpreendente.

\begin{boxmundobox}{Box Mundo 3.1 — Violações de transitividade: os experimentos de Tversky e replicações internacionais}

\textbf{Contexto:} O axioma da transitividade --- se \(A \succsim B\) e
\(B \succsim C\), então \(A \succsim C\) --- é um dos pilares da teoria
da escolha racional. Sem ele, como discutido na Seção 3.1, não é
possível construir uma função de utilidade que represente as
preferências do consumidor. Porém, desde os anos 1960, psicólogos e
economistas comportamentais têm documentado violações sistemáticas desse
axioma em contextos experimentais. O trabalho mais influente é o de Amos
Tversky (1969), que demonstrou que indivíduos frequentemente exibem
preferências intransitivas quando as alternativas diferem em múltiplas
dimensões e as diferenças em alguma dimensão são ``imperceptíveis''.

\textbf{Dados:} Tversky (1969) apresentou a participantes do experimento
escolhas entre pares de apostas que variavam em probabilidade de ganho e
valor do prêmio. Quando a diferença de probabilidades entre duas apostas
era pequena (por exemplo, 7/24 vs.~8/24), os participantes tendiam a
escolher a aposta com prêmio maior --- mas quando confrontados com
comparações entre extremos da cadeia, a probabilidade dominava, gerando
ciclos de intransitividade. A taxa de intransitividade observada chegou
a 17\% das tríades de escolha. Estudos subsequentes replicaram o
fenômeno em diversos contextos e culturas: Regenwetter et al.~(2011)
conduziram meta-análises com dados de laboratórios nos EUA, Europa e
Ásia e encontraram que, embora a frequência de intransitividade varie
entre populações (de 5\% a 25\% das tríades), o padrão é robusto e não
desaparece com incentivos monetários reais. Na Alemanha, Birnbaum e
Schmidt (2015) replicaram o resultado com mais de 400 participantes e
diferentes tipos de loterias.

\textbf{Análise:} Essas violações não invalidam a teoria da escolha
racional, mas delimitam seu domínio de aplicação --- uma questão central
para a verificação de modelos (Capítulo 1). A resposta da teoria
econômica tem sido dupla. Primeiro, modelos de utilidade aleatória
(\emph{random utility}), como o de McFadden (1974, Nobel 2000),
incorporam um componente estocástico que pode gerar intransitividade
aparente mesmo quando as preferências subjacentes são transitivas --- as
violações refletem ``erros'' de percepção, não preferências genuinamente
cíclicas. Segundo, a economia comportamental desenvolveu modelos
alternativos, como a \emph{regret theory} de Loomes e Sugden (1982), que
relaxam a transitividade mantendo outras propriedades. Do ponto de vista
prático, a evidência sugere que a transitividade é uma aproximação
excelente para decisões familiares e repetidas, mas pode falhar em
escolhas complexas e não habituais --- exatamente a distinção entre
modelos descritivos e normativos que permeia todo este livro.

\textbf{Fonte:} Tversky, A. (1969). Intransitivity of preferences.
\emph{Psychological Review}, 76(1), 31--48. Regenwetter, M.; Dana, J.;
Davis-Stober, C. P. (2011). Transitivity of preferences.
\emph{Psychological Review}, 118(1), 42--56. Birnbaum, M. H.; Schmidt,
U. (2015). Testing transitivity in choice under risk. \emph{Theory and
Decision}, 79(4), 631--660.

\end{boxmundobox}

Com essa ressalva empírica em mente --- os axiomas são uma excelente
aproximação para decisões cotidianas, mas podem falhar em escolhas
complexas ---, prossigamos para a recompensa teórica: a função de
utilidade.

\begin{center}\rule{0.5\linewidth}{0.5pt}\end{center}

\section{3.2 Do ``Prefiro'' ao Número: Função de
Utilidade}\label{do-prefiro-ao-nuxfamero-funuxe7uxe3o-de-utilidade}

Ótimo: o consumidor é coerente, prefere mais a menos e não dá saltos.
Mas como \emph{usar} essas preferências para resolver um problema de
otimização? Comparar cestas duas a duas, sem uma escala numérica, é como
tentar navegar sem mapa --- dá para fazer, mas é dolorosamente lento. O
que precisamos é de uma \textbf{função} que atribua um número a cada
cesta, de modo que cestas melhores recebam números maiores. Se tivermos
isso, todo o arsenal do Capítulo 2 --- derivadas, Lagrange, condições de
ótimo --- entra em campo.

A pergunta, portanto, é: podemos traduzir a relação de preferência
\(\succsim\) em uma \textbf{função de utilidade}
\(u: X \to \mathbb{R}\)? A resposta é sim --- e é exatamente para isso
que precisávamos dos axiomas. A passagem de ``eu prefiro A a B'' para
``\(u(A) > u(B)\)'' é o passe de mágica que transforma a microeconomia
de uma disciplina verbal em uma ciência analítica.

A definição formal de representação é a seguinte. Uma \textbf{função de
utilidade} \(u: X \to \mathbb{R}\) \textbf{representa} a relação de
preferência \(\succsim\) se, para todo \(\mathbf{x}, \mathbf{y} \in X\):

{[} \mathbf{x} \succsim \mathbf{y} \iff u(\mathbf{x})
\geq u(\mathbf{y}). \label{eq:3.2.1} \tag{3.2.1} {]}

A condição expressa na equação (\ref{eq:3.2.1}) exige apenas que a
função preserve o ordenamento das preferências --- ela não atribui
significado aos valores numéricos em si, apenas à sua ordenação
relativa. A existência de tal função, contudo, não é trivial: não é para
qualquer relação de preferência que conseguimos encontrar uma
representação numérica. O teorema a seguir fornece condições suficientes
e mostra que os axiomas da Seção 3.1 são exatamente o que precisamos
para garantir essa representação.

\begin{theorembox}{Teorema 3.1 — Existência da Função de Utilidade (Debreu, 1954)}

Se a relação de preferência \(\succsim\) sobre
\(X \subseteq \mathbb{R}^n_+\) é \textbf{completa}, \textbf{transitiva},
\textbf{contínua} e \textbf{monótona}, então existe uma função de
utilidade contínua \(u: X \to \mathbb{R}\) que representa \(\succsim\).

\end{theorembox}

\begin{infobox}{Prêmio Nobel — Gerard Debreu (1983)}

\textbf{Gerard Debreu} (1921--2004) foi um economista e matemático
franco-americano, formado pela École Normale Supérieure de Paris. Atuou
na Cowles Commission e na Universidade da Califórnia, Berkeley.

\textbf{Por que ganhou o Nobel:} Premiado por sua contribuição à teoria
do equilíbrio geral e por ter incorporado novos métodos analíticos à
teoria econômica. Em \emph{Theory of Value} (1959), Debreu forneceu uma
axiomatização rigorosa da teoria do consumidor e demonstrou a existência
de equilíbrio geral competitivo sob condições precisas. O Teorema 3.1
deste capítulo é uma de suas contribuições mais fundamentais.

\textbf{Conexão com este capítulo:} O teorema de existência da função de
utilidade garante que os axiomas da Seção 3.1 --- completude,
transitividade, continuidade e monotonicidade --- são suficientes para
construir toda a maquinaria analítica que se segue: curvas de
indiferença, TMS e otimização do consumidor.

\end{infobox}

A demonstração completa pode ser encontrada em Mas-Colell, Whinston e
Green (1995, Proposição 3.C.1). A ideia central é elegante: para cada
cesta \(\mathbf{x}\), encontramos o escalar \(t\) tal que
\(\mathbf{x} \sim t \cdot \mathbf{1}\), onde
\(\mathbf{1} = (1, 1, \ldots, 1)\) é a cesta com uma unidade de cada
bem. Esse escalar \(t\) é o ``nível de utilidade'' atribuído a
\(\mathbf{x}\). A monotonicidade garante que tal \(t\) é único ---
cestas melhores correspondem a valores maiores de \(t\) --- e a
continuidade garante que a função resultante é contínua, sem saltos que
gerariam cestas impossíveis de classificar. Cada axioma desempenha,
portanto, um papel preciso na construção.

\begin{notebox}{Utilidade é ordinal, não cardinal}

A função de utilidade atribui números a cestas apenas para preservar o
\textbf{ordenamento}. Os valores absolutos não têm significado econômico
intrínseco. Se \(u(\mathbf{x}) = 10\) e \(u(\mathbf{y}) = 5\), sabemos
apenas que \(\mathbf{x} \succ \mathbf{y}\), \textbf{não} que
\(\mathbf{x}\) é ``duas vezes melhor'' que \(\mathbf{y}\).

\end{notebox}

\begin{intuitionbox}{Intuição Econômica}

\textbf{Em uma frase:} Os números da utilidade funcionam como a ordem de
chegada numa corrida --- só importa quem chegou antes, não por quantos
segundos.

\textbf{Pense assim:} Imagine que você dá nota 10 para um prato de
feijoada e nota 5 para uma salada. Isso não significa que a feijoada é
``duas vezes mais gostosa'' --- apenas que você a prefere. Se um amigo
desse nota 100 e 50, as preferências seriam as mesmas. O ranking é o que
importa, não o placar.

\textbf{Por que isso importa:} Essa propriedade ordinal nos liberta de
medir ``felicidade'' em unidades absolutas --- basta saber o que o
consumidor prefere para analisar suas escolhas. No Capítulo 7, veremos
que a introdução de incerteza exige uma utilidade \textbf{cardinal}
(VNM): a curvatura da função passa a ter significado econômico,
determinando a aversão ao risco do agente.

\end{intuitionbox}

\begin{center}\rule{0.5\linewidth}{0.5pt}\end{center}

\section{3.3 O Mapa do ``Tanto Faz'': Curvas de
Indiferença}\label{o-mapa-do-tanto-faz-curvas-de-indiferenuxe7a}

Números são úteis, mas humanos pensam em imagens. A \textbf{curva de
indiferença} é onde a microeconomia ganha olhos. A ideia: pegue todas as
cestas que dão ao consumidor a \emph{mesma} satisfação e ligue-as por
uma curva. Pronto --- você tem um ``mapa topográfico'' das preferências,
onde cada curva é uma altitude de felicidade. Cestas em curvas mais
altas são preferidas; cestas na mesma curva são equivalentes.

A forma das curvas conta uma história: se são quase retas, os bens são
facilmente substituíveis (café e chá); se são em ``L'', devem ser
consumidos juntos (pé esquerdo e pé direito de sapato). No Capítulo 4,
quando sobrepormos essas curvas à reta orçamentária, o ponto ótimo do
consumidor emergirá como o ponto de tangência --- uma das imagens mais
elegantes de toda a ciência econômica.

\begin{definitionbox}{Curva de indiferença}

A curva de indiferença associada ao nível de utilidade \(\bar{u}\) é o
conjunto \(\{\mathbf{x} \in X : u(\mathbf{x}) = \bar{u}\}\). Trata-se de
uma \textbf{curva de nível} da função de utilidade.

\end{definitionbox}

\begin{tipbox}{Propriedades das curvas de indiferença}

Sob os axiomas da Seção 3.1:

\begin{enumerate}
\def\labelenumi{\arabic{enumi}.}
\tightlist
\item
  \textbf{Cobrem todo o espaço de consumo}: pela completude, toda cesta
  pertence a alguma curva de indiferença.
\item
  \textbf{Não se cruzam}: se duas curvas se cruzassem em um ponto, a
  transitividade seria violada.
\item
  \textbf{Possuem inclinação negativa}: pela monotonicidade, manter a
  utilidade constante exige compensar um aumento em \(x_1\) com uma
  redução em \(x_2\).
\item
  \textbf{Cestas em curvas mais altas (a nordeste) são preferidas}:
  consequência direta da monotonicidade.
\end{enumerate}

O \textbf{mapa de indiferença} é a família de todas as curvas de
indiferença. Ele oferece uma representação visual completa das
preferências do consumidor no espaço bidimensional.

\end{tipbox}

\figurePlaceholder{/micro-book/graficos/cap03/curvas-indiferenca.html}

\textbf{Figura 3.1 --- Mapa de curvas de indiferença.} Arraste o ponto
sobre a curva para ver a TMS. Use os sliders para alterar os parâmetros
e o menu para trocar o tipo de preferência (Cobb-Douglas, substitutos
perfeitos, complementos perfeitos, CES, quase-linear).

\begin{tipbox}{Convexidade estrita}

Se as preferências forem \textbf{estritamente convexas} --- isto é, se
\(\mathbf{x} \sim \mathbf{y}\) e \(\mathbf{x} \neq \mathbf{y}\)
implicarem \(t\mathbf{x} + (1-t)\mathbf{y} \succ \mathbf{x}\) para todo
\(t \in (0,1)\) --- então as curvas de indiferença são estritamente
convexas em relação à origem. Isso reflete a ideia de que consumidores
preferem variedade: uma cesta ``mista'' é preferida a cestas extremas. A
seção sobre a TMS (a seguir) mostrará que essa propriedade equivale a
uma TMS decrescente.

\end{tipbox}

\begin{warningbox}{Erro Comum}

\textbf{Confundir preferências convexas com função de utilidade
côncava.} São conceitos relacionados, mas distintos:

\begin{itemize}
\tightlist
\item
  \textbf{Preferências convexas} significam que os \emph{conjuntos
  superiores} \(\{\mathbf{x} : u(\mathbf{x}) \geq \bar{u}\}\) são
  convexos --- geometricamente, as curvas de indiferença são
  ``abauladas'' em direção à origem.
\item
  \textbf{Função côncava} é uma condição sobre a própria função \(u\),
  mais forte que a convexidade das preferências.
\end{itemize}

Toda função de utilidade côncava gera preferências convexas, mas a
recíproca é falsa: como a utilidade é ordinal, podemos sempre encontrar
uma transformação monotônica que torne a função convexa (não côncava)
sem alterar as preferências. O Exercício 3.5, com \(u = x_1^2 + x_2^2\),
ilustra o caso oposto: uma função convexa que gera preferências
\emph{côncavas} (curvas de indiferença circulares, côncavas em relação à
origem).

\end{warningbox}

\chapter{3.4--3.5 TMS e Utilidade
Marginal}\label{tms-e-utilidade-marginal}

\section{3.4 Quanto Cerveja Vale um Café? A Taxa Marginal de
Substituição}\label{quanto-cerveja-vale-um-cafuxe9-a-taxa-marginal-de-substituiuxe7uxe3o}

As curvas de indiferença mostram \emph{o quê} o consumidor aceita; a TMS
diz \emph{a que preço}. Imagine que você tem 10 fatias de pizza e 2
cervejas. Alguém te oferece trocar 1 pizza por 1 cerveja. Você aceita?
Provavelmente sim --- pizza está sobrando, cerveja está escassa. Mas e
se já tiver 10 cervejas e 2 pizzas? Agora a troca é desvantajosa. A
\textbf{taxa marginal de substituição (TMS)} captura exatamente essa
lógica: quantas unidades do bem 2 você topa ceder por uma unidade a mais
do bem 1, \emph{sem perder satisfação}.

A TMS é, talvez, o conceito mais operacional de toda a teoria do
consumidor. No próximo capítulo, veremos que o consumidor está no ótimo
quando sua TMS (o quanto ele \emph{quer} trocar) se iguala à razão de
preços (o quanto o mercado \emph{cobra} para trocar). Se o valor
subjetivo e o preço de mercado divergem, há um negócio a ser feito --- e
o consumidor faz.

\begin{classroombox}{Atividade — Sua TMS pessoal}

\textbf{Passo 1 (individual, 2 min):} Você tem 10 fatias de pizza e 2
latas de cerveja. Quantas fatias de pizza você toparia trocar por 1
cerveja extra? Anote um número.

\textbf{Passo 2 (individual, 1 min):} Agora inverta: você tem 2 fatias
de pizza e 10 cervejas. Quantas fatias toparia trocar por 1 cerveja?
Anote.

\textbf{Passo 3 (em dupla, 3 min):} Compare com o colega. (a) Os números
de vocês são iguais? Provavelmente não --- preferências são subjetivas.
(b) O número do Passo 2 é menor que o do Passo 1? Quase certamente sim.
Esse é o fenômeno da \textbf{TMS decrescente}.

\textbf{Debrief (turma):} O professor formaliza: a TMS caiu porque pizza
ficou escassa e cerveja abundante. A disposição a trocar não é fixa ---
depende da composição da cesta. Isso é a convexidade estrita em ação.

\end{classroombox}

\begin{definitionbox}{Taxa marginal de substituição}

A \textbf{taxa marginal de substituição} do bem 1 pelo bem 2 é definida
como o valor absoluto da inclinação da curva de indiferença no ponto
\((x_1, x_2)\):

{[} \text{TMS}\emph{\{12\} = -\frac{dx_2}{dx_1}\bigg\textbar{}}\{u =
\bar\{u\}\} \label{eq:3.4.2} \tag{3.4.2} {]}

A TMS mede a quantidade do bem 2 que o consumidor está disposto a abrir
mão para obter uma unidade adicional do bem 1, mantendo o nível de
utilidade constante.

\end{definitionbox}

\begin{infobox}{Prêmio Nobel — John R. Hicks (1972)}

\textbf{John Richard Hicks} (1904--1989) foi um economista britânico,
formado na Universidade de Oxford. Foi professor em Manchester, Oxford e
na London School of Economics. Dividiu o Nobel com Kenneth Arrow.

\textbf{Por que ganhou o Nobel:} Premiado por suas contribuições
pioneiras à teoria do equilíbrio geral e à teoria do bem-estar. Em
\emph{Value and Capital} (1939), Hicks axiomatizou a teoria do
consumidor, introduzindo a taxa marginal de substituição como ferramenta
central, e desenvolveu o conceito de demanda compensada (hicksiana), que
separa efeitos renda e substituição.

\textbf{Conexão com este capítulo:} A taxa marginal de substituição
(TMS) formalizada por Hicks é o conceito central deste capítulo. A
condição de tangência entre a curva de indiferença e a restrição
orçamentária --- que exige igualdade entre TMS e razão de preços --- é a
base da teoria da escolha do consumidor desenvolvida aqui.

\end{infobox}

\textbf{Interpretação econômica.} A TMS é a \textbf{taxa de troca
subjetiva} do consumidor --- o ``preço pessoal'' que ele atribui ao bem
1 em termos do bem 2. Se \(\text{TMS}_{12} = 3\), o consumidor está
disposto a abrir mão de até 3 unidades do bem 2 para obter 1 unidade
adicional do bem 1, permanecendo no mesmo nível de satisfação. Note que
essa taxa é puramente subjetiva e pode diferir radicalmente da razão de
preços do mercado --- é precisamente essa diferença que cria incentivos
para o consumidor realocar seu consumo.

\textbf{TMS decrescente.} Sob preferências estritamente convexas, a TMS
é decrescente ao longo da curva de indiferença: à medida que o
consumidor adquire mais do bem 1, cada unidade adicional torna-se
relativamente menos valiosa em relação ao bem 2, que se torna cada vez
mais escasso na cesta. O consumidor, portanto, está disposto a
sacrificar cada vez menos do bem 2 para obter mais uma unidade do bem 1.
Esta propriedade --- matematicamente equivalente à convexidade estrita
das curvas de indiferença --- reflete a ideia intuitiva de que os
consumidores valorizam a diversidade: cestas equilibradas são preferidas
a cestas extremas.

\begin{intuitionbox}{Intuição Econômica}

\textbf{Em uma frase:} A TMS mede ``quanto do bem 2 você toparia trocar
por mais uma unidade do bem 1'' --- é o seu preço pessoal de troca.

\textbf{Pense assim:} Imagine que você tem muito arroz e pouco feijão em
casa. Você toparia trocar bastante arroz por um pouco de feijão. Mas à
medida que ganha feijão e perde arroz, cada porção adicional de feijão
vale menos para você, e cada porção de arroz que abre mão dói mais. Essa
``taxa de troca pessoal'' que vai caindo é a TMS decrescente.

\textbf{Por que isso importa:} No capítulo seguinte veremos que o
consumidor otimiza quando sua TMS iguala a razão de preços do mercado
--- é o ponto onde o ``preço pessoal'' coincide com o ``preço de
mercado''.

\end{intuitionbox}

\begin{infobox}{Exemplo — TMS e escolhas alimentares no Brasil}

Considere um consumidor brasileiro escolhendo entre alimentação dentro
de casa (\(x_1\)) e alimentação fora de casa (\(x_2\)). Segundo a POF
2017-2018 do IBGE, a despesa média per capita com alimentação no
domicílio era de aproximadamente R\$ 136 mensais, contra R\$ 73 com
alimentação fora. Uma família com muito gasto em alimentação domiciliar
(cesta ``extrema'') teria uma TMS alta: estaria disposta a abrir mão de
várias refeições caseiras por uma refeição fora. À medida que aumenta a
alimentação fora de casa, a TMS diminui --- o consumidor valoriza cada
vez menos uma refeição adicional fora. Essa TMS decrescente é a
manifestação empírica da convexidade estrita das preferências.

\end{infobox}

\begin{center}\rule{0.5\linewidth}{0.5pt}\end{center}

\section{3.5 A Última Fatia de Pizza: Utilidade Marginal e
TMS}\label{a-uxfaltima-fatia-de-pizza-utilidade-marginal-e-tms}

Na seção anterior, a TMS era uma inclinação que ``víamos'' no gráfico.
Agora queremos \emph{calculá-la} --- sem desenhar nada. O truque: a TMS
é simplesmente a razão entre as \textbf{utilidades marginais} dos dois
bens. A utilidade marginal mede o ``quanto a mais de satisfação'' que
uma unidade extra de um bem proporciona (nosso velho amigo \emph{ceteris
paribus}). Se a próxima fatia de pizza te dá prazer 3 e a próxima
cerveja te dá prazer 6, você toparia trocar até 2 pizzas por 1 cerveja
--- e a TMS é exatamente 2. Vejamos isso com rigor.

Formalmente, a \textbf{utilidade marginal} do bem \(i\) é a derivada
parcial da função de utilidade em relação à quantidade desse bem:

{[} \text{UMg}\_i = \frac{\partial u}{\partial x_i}. \label{eq:3.5.3}
\tag{3.5.3} {]}

Se a utilidade marginal é positiva --- o que é garantido pela
monotonicidade ---, o consumidor sempre se beneficia de uma unidade
adicional do bem \(i\), mantendo tudo o mais constante. Contudo, é
importante lembrar que, como a utilidade é ordinal (Seção 3.2), o valor
numérico da utilidade marginal \emph{em si} não possui significado
econômico absoluto --- ele depende da escala escolhida para representar
as preferências. O que \emph{tem} significado é a razão entre utilidades
marginais, como veremos imediatamente.

\begin{warningbox}{Cuidado: utilidade marginal não tem significado isolado}

A utilidade marginal \(\text{UMg}_i\) muda de valor quando aplicamos uma
transformação monotônica à função de utilidade (Seção 3.7). Se
\(\hat{u} = \ln(u)\), a utilidade marginal passa de
\(\partial u / \partial x_i\) para
\(\frac{1}{u} \cdot \partial u / \partial x_i\) --- um valor
completamente diferente. Por isso, afirmações como ``a utilidade
marginal do bem 1 é 5'' não possuem conteúdo econômico: o número 5
depende arbitrariamente da escala escolhida.

O que \textbf{tem} significado econômico é a \textbf{razão}
\(\text{UMg}_1 / \text{UMg}_2\), que é a TMS --- invariante sob
transformações monotônicas. Não confunda, ainda, ``utilidade marginal
decrescente'' (propriedade cardinal, sem significado ordinal) com ``TMS
decrescente'' (propriedade ordinal, com significado econômico preciso:
preferências convexas).

\end{warningbox}

A relação fundamental entre a utilidade marginal, definida pela equação
(\ref{eq:3.5.3}), e a TMS, definida pela equação (\ref{eq:3.4.2}),
revela que a inclinação da curva de indiferença pode ser inteiramente
expressa em termos de derivadas da função de utilidade. Essa ponte entre
geometria e cálculo --- entre a inclinação visual da curva e as
derivadas parciais da função --- é dada pela proposição a seguir.

\begin{theorembox}{Proposição 3.2 — TMS como razão de utilidades marginais}

Se \(u(x_1, x_2)\) é diferenciável e \(\text{UMg}_2 > 0\), então:

{[} \text{TMS}\_\{12\} = \frac{\text{UMg}_1}{\text{UMg}_2} =
\frac{\partial u / \partial x_1}{\partial u / \partial x_2}
\label{eq:3.5.4} \tag{3.5.4} {]}

\end{theorembox}

\begin{proofbox}{Demonstração}

Considere o consumidor sobre uma curva de indiferença, de modo que
\(u(x_1, x_2) = \bar{u}\). Tomando o \textbf{diferencial total} da
função de utilidade ao longo dessa curva:

{[} du = \frac{\partial u}{\partial x_1} dx\_1 +
\frac{\partial u}{\partial x_2} dx\_2 = 0, {]}

pois o nível de utilidade é constante (\(du = 0\)) ao longo da curva de
indiferença. Reorganizando:

{[} \frac{\partial u}{\partial x_2} dx\_2 =
-\frac{\partial u}{\partial x_1} dx\_1, {]}

{[} \frac{dx_2}{dx_1} =
-\frac{\partial u / \partial x_1}{\partial u / \partial x_2}. {]}

Como a TMS é definida como o valor absoluto (com sinal positivo) da
inclinação da curva de indiferença:

{[} \text{TMS}\emph{\{12\} = -\frac{dx_2}{dx_1}\bigg\textbar{}}\{u =
\bar\{u\}\} =
\frac{\partial u / \partial x_1}{\partial u / \partial x_2} =
\frac{\text{UMg}_1}{\text{UMg}_2}. \qquad \blacksquare {]}

\end{proofbox}

\figurePlaceholder{/micro-book/graficos/cap03/tms-ponto.html}

\textbf{Figura 3.2 --- Taxa Marginal de Substituição (TMS).} Arraste o
ponto P ao longo da curva de indiferença para ver a reta tangente e o
cálculo da \(\text{TMS} = \text{UMg}_1/\text{UMg}_2\) em tempo real.
Selecione entre Cobb-Douglas, linear, Leontief, CES e quase-linear.

\begin{exresolvidobox}{Exercício Resolvido 3.1}

\textbf{Enunciado:} Um consumidor tem preferências representadas por
\(u(x_1, x_2) = x_1^{2/5} \, x_2^{3/5}\). Calcule a TMS no ponto
\((x_1, x_2) = (10, 15)\) e interprete o resultado.

\textbf{Dados:} \(a = 2/5\), \(b = 3/5\), \(x_1 = 10\), \(x_2 = 15\).

\textbf{Resolução:}

\textbf{Passo 1 --- Cálculo das utilidades marginais}

{[} \text{UMg}\_1 = \frac{\partial u}{\partial x_1} = \frac{2}{5} ,
x\_1\^{}\{-3/5\} , x\_2\^{}\{3/5\} {]}

{[} \text{UMg}\_2 = \frac{\partial u}{\partial x_2} = \frac{3}{5} ,
x\_1\^{}\{2/5\} , x\_2\^{}\{-2/5\} {]}

\textbf{Passo 2 --- Cálculo da TMS}

{[} \text{TMS}\_\{12\} = \frac{\text{UMg}_1}{\text{UMg}_2} =
\frac{\frac{2}{5} \, x_1^{-3/5} \, x_2^{3/5}}{\frac{3}{5} \, x_1^{2/5} \, x_2^{-2/5}}
= \frac{2}{3} \cdot \frac{x_2}{x_1} {]}

Note que, para qualquer Cobb-Douglas \(u = x_1^a x_2^b\), a TMS assume a
forma geral \(\text{TMS}_{12} = \frac{a}{b} \cdot \frac{x_2}{x_1}\).

\textbf{Passo 3 --- Avaliação no ponto dado}

{[} \text{TMS}\_\{12\}(10, 15) = \frac{2}{3} \cdot \frac{15}{10} =
\frac{2}{3} \cdot 1{,}5 = 1 {]}

\textbf{Resultado:} \(\text{TMS}_{12} = 1\) no ponto \((10, 15)\).

\textbf{Interpretação econômica:} No ponto \((10, 15)\), o consumidor
está disposto a trocar exatamente 1 unidade do bem 2 por 1 unidade
adicional do bem 1, mantendo-se indiferente. Se pensarmos no bem 1 como
transporte e no bem 2 como alimentação fora de casa no orçamento de uma
família brasileira, a TMS igual a 1 significa que o consumidor valoriza
igualmente uma unidade adicional de cada bem naquela composição de
cesta. Se o preço relativo diferir de 1, haverá incentivo para realocar
o consumo --- tema do Capítulo 4.

\end{exresolvidobox}

\chapter{3.6--3.7 Formas Funcionais e Transformações
Monotônicas}\label{formas-funcionais-e-transformauxe7uxf5es-monotuxf4nicas}

\section{3.6 O Cardápio das Preferências: Cobb-Douglas, CES e
Companhia}\label{o-carduxe1pio-das-preferuxeancias-cobb-douglas-ces-e-companhia}

Até aqui, a função de utilidade era abstrata --- uma \(u\) genérica sem
rosto. Agora é hora de conhecer as protagonistas pelo nome. Cada família
de funções de utilidade captura um tipo diferente de consumidor: o que
substitui facilmente (Cobb-Douglas), o que não substitui de jeito nenhum
(Leontief), o que trata os bens como idênticos (substitutos perfeitos),
e o camaleão que se adapta a qualquer grau de substituição (CES).
Conhecer essas ``personagens'' é essencial porque elas reaparecerão em
\emph{todo} capítulo daqui para frente --- da derivação de demandas à
análise de comércio internacional.

\subsection{3.6.1 Cobb-Douglas}\label{cobb-douglas}

A função Cobb-Douglas é, sem dúvida, a forma funcional mais utilizada em
microeconomia --- tanto pela sua notável tratabilidade analítica quanto
por suas propriedades econômicas intuitivas e empiricamente relevantes.
Proposta originalmente por Charles Cobb e Paul Douglas (1928) no
contexto da teoria da produção, ela rapidamente migrou para a teoria do
consumidor, onde se tornou o ``cavalo de batalha'' de inúmeros modelos
teóricos e aplicados. A função é definida como:

{[} u(x\_1, x\_2) = x\_1\^{}a , x\_2\^{}b, \quad a, b \textgreater{} 0.
\label{eq:3.6.5} \tag{3.6.5} {]}

Geometricamente, as curvas de indiferença da Cobb-Douglas são hipérboles
convexas, estritamente decrescentes e assintóticas aos eixos
coordenados. A assíntota implica que o consumidor nunca deseja
quantidades nulas de qualquer bem: por maior que seja a quantidade de
\(x_1\), uma pequena quantidade positiva de \(x_2\) sempre é necessária
para manter a utilidade positiva. A TMS é:

{[} \text{TMS}\_\{12\} = \frac{a \, x_2}{b \, x_1}. \label{eq:3.6.6}
\tag{3.6.6} {]}

Observe que a TMS depende da razão \(x_2/x_1\): quando o consumidor tem
muito do bem 2 e pouco do bem 1, a TMS é alta --- ele valoriza muito uma
unidade adicional do bem 1. À medida que adquire mais do bem 1, a TMS
cai, refletindo a TMS decrescente discutida na Seção 3.4.

A função Cobb-Douglas é extremamente conveniente por gerar funções de
demanda com forma fechada, o que simplifica consideravelmente os
cálculos. A participação de cada bem na despesa total é constante:
\(\frac{a}{a+b}\) para o bem 1 e \(\frac{b}{a+b}\) para o bem 2 --- uma
propriedade notável que independe tanto dos preços quanto da renda. A
elasticidade de substituição é \(\sigma = 1\), indicando que, quando o
preço relativo de um bem sobe 1\%, a razão entre as quantidades
consumidas se ajusta exatamente 1\% na direção oposta.

\subsection{3.6.2 Substitutos Perfeitos}\label{substitutos-perfeitos}

No extremo oposto da substituibilidade, considere bens que o consumidor
troca livremente entre si a uma taxa fixa --- como, para muitos
consumidores, manteiga e margarina, ou gasolina comum e aditivada. Para
esse consumidor, o que importa é apenas a quantidade total de ``gordura
para cozinhar'' ou de ``combustível'', não a composição específica da
cesta. Nesse caso, a função de utilidade assume a forma linear:

{[} u(x\_1, x\_2) = a x\_1 + b x\_2, \quad a, b \textgreater{} 0.
\label{eq:3.6.7} \tag{3.6.7} {]}

As curvas de indiferença são \textbf{linhas retas} com inclinação
\(-a/b\), paralelas entre si no espaço \((x_1, x_2)\). Curvas de
indiferença mais altas (maiores níveis de utilidade) estão mais
afastadas da origem. A TMS é constante em todos os pontos:

{[} \text{TMS}\_\{12\} = \frac{a}{b}. \label{eq:3.6.8} \tag{3.6.8} {]}

O consumidor troca os bens a uma taxa fixa, independente da composição
da cesta --- não há TMS decrescente, pois a disposição a substituir
nunca se altera. A elasticidade de substituição é \(\sigma = \infty\),
refletindo a perfeita substituibilidade.

\begin{warningbox}{Substitutos perfeitos geram soluções de canto}

Diferentemente das demais formas funcionais, a solução ótima com
substitutos perfeitos tipicamente \textbf{não é interior}: o consumidor
gasta toda a renda no bem que oferece maior ``valor por unidade
monetária'' --- isto é, maior razão \(a_i/p_i\). Formalmente:

\begin{itemize}
\tightlist
\item
  Se \(a/p_1 > b/p_2\): \(x_1^* = m/p_1\), \(x_2^* = 0\) ---
  especialização total no bem 1.
\item
  Se \(a/p_1 < b/p_2\): \(x_1^* = 0\), \(x_2^* = m/p_2\) ---
  especialização total no bem 2.
\item
  Se \(a/p_1 = b/p_2\): qualquer cesta na reta orçamentária é ótima
  (caso-faca).
\end{itemize}

Na prática, isso significa que a condição de tangência
\(\text{TMS} = p_1/p_2\) não se aplica: como a TMS é constante e a reta
orçamentária também tem inclinação constante, as duas nunca se
``encontram'' em um ponto interior (a menos do caso-faca). Essa é uma
exceção importante à regra geral de otimização do Capítulo 4.

\end{warningbox}

\subsection{3.6.3 Complementos Perfeitos}\label{complementos-perfeitos}

Se os substitutos perfeitos representam a máxima disposição a trocar
entre bens, os complementos perfeitos ocupam o extremo oposto do
espectro de substituibilidade: os bens só têm valor quando consumidos
juntos, em proporções fixas.\footnote{Em \emph{Monty Python and the Holy
  Grail}, os Cavaleiros Que Dizem Ni exigem um arbusto
  (\emph{shrubbery}) --- e nenhum substituto é aceito. ``We want\ldots{}
  a shrubbery!'' A demanda é perfeitamente inelástica ao preço: não
  importa quanto custe, nada além de um arbusto satisfaz a restrição. É
  o complemento perfeito levado ao absurdo: sem o arbusto, a utilidade é
  zero, independente de quanto mais se tenha de qualquer outra coisa.
  Pense em um pé esquerdo e um pé direito de sapato --- um sem o outro é
  inútil. Ter dez pés esquerdos e apenas um pé direito não é melhor do
  que ter um par, pois os pés extras não podem ser aproveitados. A
  função de utilidade que captura essa rigidez é:}

{[} u(x\_1, x\_2) = \min\{a x\_1, , b x\_2\}, \quad a, b \textgreater{}
0. \label{eq:3.6.9} \tag{3.6.9} {]}

As curvas de indiferença têm formato de \textbf{L} (ângulo reto), com
vértice na reta \(a x_1 = b x_2\), que define a proporção ótima de
consumo. A TMS é indefinida no vértice (onde a curva faz uma ``quina''),
zero nos segmentos horizontais (o consumidor não valoriza mais do bem 1
sem mais do bem 2) e infinita nos segmentos verticais (o consumidor não
valoriza mais do bem 2 sem mais do bem 1). A elasticidade de
substituição é \(\sigma = 0\), indicando que nenhum ajuste de preços
relativos pode induzir o consumidor a alterar a proporção entre os bens.
Exemplos clássicos incluem sapato esquerdo e sapato direito, CPU e
monitor, e, em contextos mais modernos, cartucho de impressora e a
impressora correspondente.

\subsection{3.6.4 CES (Elasticidade de Substituição
Constante)}\label{ces-elasticidade-de-substituiuxe7uxe3o-constante}

Os três casos anteriores --- Cobb-Douglas (\ref{eq:3.6.5}), substitutos
perfeitos (\ref{eq:3.6.7}) e complementos perfeitos (\ref{eq:3.6.9}) ---
podem parecer categorias isoladas e desconexas, mas na verdade são
membros de uma mesma família: a função de utilidade com
\textbf{Elasticidade de Substituição Constante} (CES, do inglês
\emph{Constant Elasticity of Substitution}). Introduzida por Arrow,
Chenery, Minhas e Solow (1961) no contexto da teoria da produção, a CES
permite capturar \emph{qualquer} grau de substituibilidade entre os bens
por meio de um único parâmetro \(\rho\), que governa a curvatura das
curvas de indiferença. Ao variar \(\rho\) continuamente, transitamos
suavemente dos complementos perfeitos aos substitutos perfeitos,
passando pela Cobb-Douglas como caso intermediário --- o que torna a CES
uma ferramenta extremamente versátil tanto na teoria quanto na estimação
empírica.

{[} u(x\_1, x\_2) = \left(x\_1\^{}\{\rho\} +
x\_2\textsuperscript{\{\rho\}\right)}\{1/\rho\}, \quad \rho \leq 1, ;
\rho \neq 0. \label{eq:3.6.10} \tag{3.6.10} {]}

A elasticidade de substituição --- que mede a sensibilidade da razão
\(x_1/x_2\) a variações no preço relativo --- é determinada diretamente
pelo parâmetro \(\rho\):

{[} \sigma = \frac{1}{1 - \rho}. \label{eq:3.6.11} \tag{3.6.11} {]}

Essa relação é a chave para entender a versatilidade da CES: ao variar
um único parâmetro, percorremos todo o espectro de substituibilidade. A
função CES engloba como casos especiais as três formas funcionais
anteriores:

\begin{itemize}
\tightlist
\item
  \(\rho \to 0\): Cobb-Douglas (\(\sigma = 1\)). A passagem ao limite
  requer a regra de L'Hôpital aplicada ao logaritmo da função.
\item
  \(\rho = 1\): Substitutos perfeitos (\(\sigma = \infty\)). A utilidade
  se reduz a uma combinação linear dos bens.
\item
  \(\rho \to -\infty\): Complementos perfeitos (\(\sigma = 0\)). A
  utilidade converge para o mínimo das quantidades ponderadas.
\end{itemize}

Para valores intermediários de \(\rho\), a CES produz curvas de
indiferença com curvatura intermediária --- nem tão ``redondas'' quanto
as da Cobb-Douglas, nem tão ``angulosas'' quanto as dos complementos
perfeitos. Isso a torna ideal para estimação empírica, pois os dados
podem ``escolher'' o grau de substituibilidade que melhor se ajusta ao
comportamento observado.

A TMS para a CES assume uma forma compacta e elegante:

{[} \text{TMS}\_\{12\} = \left(\frac{x_1}{x_2}\right)\^{}\{\rho - 1\}.
\label{eq:3.6.12} \tag{3.6.12} {]}

Note que, quando \(\rho < 1\) (o caso economicamente relevante), o
expoente \(\rho - 1\) é negativo, de modo que a TMS é decrescente em
\(x_1/x_2\) --- confirmando a convexidade das curvas de indiferença.

\figurePlaceholder{/micro-book/graficos/cap03/ces-continua.html}

\textbf{Figura 3.3 --- CES Contínua: de Leontief a Substitutos
Perfeitos.} Arraste o slider de \(\rho\) para observar a transformação
contínua das curvas de indiferença: de ângulos retos
(\(\rho \to -\infty\), complementos perfeitos) a hipérboles suaves
(\(\rho = 0\), Cobb-Douglas) a retas (\(\rho \to 1\), substitutos
perfeitos). A elasticidade de substituição \(\sigma = 1/(1-\rho)\) é
exibida em tempo real.

\begin{examplebox}{Exemplo numérico: CES com diferentes valores de ρ}

Considere o ponto \((x_1, x_2) = (4, 8)\). A TMS da CES é
\(\text{TMS}_{12} = (x_1/x_2)^{\rho-1} = (4/8)^{\rho-1} = (0{,}5)^{\rho-1}\).
Veja como ela varia com \(\rho\):

{\def\LTcaptype{none} % do not increment counter
\begin{longtable}[]{@{}
  >{\raggedright\arraybackslash}p{(\linewidth - 6\tabcolsep) * \real{0.2500}}
  >{\raggedright\arraybackslash}p{(\linewidth - 6\tabcolsep) * \real{0.2500}}
  >{\raggedright\arraybackslash}p{(\linewidth - 6\tabcolsep) * \real{0.2500}}
  >{\raggedright\arraybackslash}p{(\linewidth - 6\tabcolsep) * \real{0.2500}}@{}}
\toprule\noalign{}
\begin{minipage}[b]{\linewidth}\raggedright
\(\rho\)
\end{minipage} & \begin{minipage}[b]{\linewidth}\raggedright
\(\sigma = \frac{1}{1-\rho}\)
\end{minipage} & \begin{minipage}[b]{\linewidth}\raggedright
\(\text{TMS}_{12}(4,8)\)
\end{minipage} & \begin{minipage}[b]{\linewidth}\raggedright
Interpretação
\end{minipage} \\
\midrule\noalign{}
\endhead
\bottomrule\noalign{}
\endlastfoot
\(0{,}5\) & \(2\) & \((0{,}5)^{-0,5} = 1{,}41\) & Substituição fácil:
troca 1,41 unidades de \(x_2\) por 1 de \(x_1\) \\
\(0\) (Cobb-Douglas) & \(1\) & \((0{,}5)^{-1} = 2\) & Substituição
moderada \\
\(-1\) & \(0{,}5\) & \((0{,}5)^{-2} = 4\) & Substituição difícil: exige
4 unidades de \(x_2\) por 1 de \(x_1\) \\
\(-5\) & \(0{,}17\) & \((0{,}5)^{-6} = 64\) & Quase complementares:
compensação altíssima \\
\end{longtable}
}

À medida que \(\rho\) diminui (e \(\sigma\) cai), a TMS cresce
explosivamente para cestas desequilibradas --- o consumidor resiste cada
vez mais a trocar o bem escasso. Para \(\rho \to -\infty\), a TMS
diverge para infinito fora do vértice, recuperando os complementos
perfeitos.

\end{examplebox}

\begin{boxbrasilbox}{Box Brasil 3.1 — Etanol versus gasolina: preferências CES reveladas nas bombas brasileiras}

\textbf{Contexto:} O Brasil oferece um laboratório natural único para
observar a elasticidade de substituição entre bens de consumo. Desde a
introdução dos veículos flex-fuel em 2003, motoristas brasileiros
escolhem, a cada abastecimento, entre etanol hidratado e gasolina comum
--- bens que são substitutos próximos, mas não perfeitos (o etanol rende
aproximadamente 70\% da quilometragem por litro em relação à gasolina).
A regra prática difundida entre consumidores --- ``abasteça etanol se o
preço for inferior a 70\% do preço da gasolina'' --- é, na essência, uma
aplicação intuitiva da condição de otimização com preferências CES.

\textbf{Dados:} Segundo a ANP (Agência Nacional do Petróleo), a
participação do etanol no consumo total de combustíveis leves oscila
fortemente com o preço relativo. Em São Paulo, onde o etanol é mais
barato (média de 65--70\% do preço da gasolina em 2023), a participação
do etanol atingiu cerca de 45\% do volume total. No Rio de Janeiro e nos
estados do Sul, onde a razão de preços frequentemente supera 75\%, a
participação do etanol caiu para 15--20\%. Estimativas econométricas de
Santos (2013, \emph{Energy Economics}) e Salvo e Huse (2013,
\emph{Journal of Environmental Economics and Management}) encontraram
elasticidades de substituição entre etanol e gasolina na faixa de
\(\sigma \approx 3\) a \(5\) para proprietários de veículos flex ---
valores que situam esses combustíveis na região intermediária do
espectro CES, entre a Cobb-Douglas (\(\sigma = 1\)) e os substitutos
perfeitos (\(\sigma = \infty\)).

\textbf{Análise:} O caso brasileiro do etanol ilustra com precisão como
o parâmetro \(\rho\) da função CES se manifesta em decisões reais. Se
etanol e gasolina fossem substitutos perfeitos (\(\sigma \to \infty\)),
observaríamos soluções de canto puras: 100\% etanol ou 100\% gasolina em
cada estado, sem transição gradual. O fato de a participação variar
\emph{suavemente} com o preço relativo --- em vez de saltar
descontinuamente --- revela que \(\sigma\) é alto, mas finito:
diferenças na autonomia, na disponibilidade de postos, na percepção de
desempenho do motor e em hábitos de consumo impedem a substituição
perfeita. A curva de indiferença entre etanol e gasolina é convexa, mas
relativamente ``achatada'', refletindo alta substituibilidade com
fricções.

\textbf{Para refletir:} Se o governo aumentar a tributação sobre a
gasolina em 10\%, que mudança percentual na razão etanol/gasolina
consumida você esperaria observar com \(\sigma = 4\)? Compare com o que
ocorreria se \(\sigma = 1\) (Cobb-Douglas) ou \(\sigma = \infty\)
(substitutos perfeitos).

\end{boxbrasilbox}

\subsection{3.6.5 Quase-linear}\label{quase-linear}

As formas funcionais anteriores compartilham uma propriedade importante:
são todas homotéticas (como veremos na Seção 3.6.6), o que significa que
a proporção entre os bens consumidos não se altera quando a renda varia
--- o consumidor simplesmente ``escala'' sua cesta ótima. Na prática,
porém, existem situações em que a demanda por um bem específico é
essencialmente insensível à renda: por exemplo, o consumo de sal de
cozinha dificilmente se altera quando uma família recebe um aumento
salarial. Nesses casos, todo o acréscimo de renda é direcionado aos
demais bens. A função de utilidade quase-linear captura exatamente essa
assimetria:

{[} u(x\_1, x\_2) = v(x\_1) + x\_2, \quad v' \textgreater{} 0, ; v'\,'
\textless{} 0. \label{eq:3.6.13} \tag{3.6.13} {]}

A TMS depende apenas de \(x_1\) --- e não de \(x_2\), o que é uma
peculiaridade notável:

{[} \text{TMS}\_\{12\} = v'(x\_1). \label{eq:3.6.14} \tag{3.6.14} {]}

Como \(v'' < 0\), a TMS é decrescente em \(x_1\): à medida que o
consumidor obtém mais do bem 1, sua disposição a trocar o bem 2 por
unidades adicionais do bem 1 diminui --- preservando a convexidade das
curvas de indiferença.

\begin{notebox}{TMS independente de x₂: a assinatura da quase-linearidade}

A propriedade mais distintiva da utilidade quase-linear é que a TMS
depende \emph{exclusivamente} de \(x_1\) --- a quantidade do bem 2 é
irrelevante para a taxa de troca subjetiva do consumidor. Em todas as
demais formas funcionais deste capítulo (Cobb-Douglas, CES, complementos
perfeitos), a TMS depende da composição \emph{completa} da cesta. Na
quase-linear, o consumidor que possui 10 ou 1.000 unidades de \(x_2\)
atribui exatamente o mesmo valor marginal relativo ao bem 1, desde que
\(x_1\) seja o mesmo. É essa independência que elimina o efeito renda
sobre \(x_1\) e gera curvas de indiferença que são translações
verticais.

\end{notebox}

As curvas de indiferença são \textbf{translações verticais} umas das
outras: possuem a mesma forma, apenas deslocadas paralelamente ao eixo
\(x_2\). Essa propriedade geométrica tem uma implicação econômica direta
e poderosa: não há efeito renda sobre o bem 1 (para soluções
interiores), pois a demanda por \(x_1\) depende apenas dos preços, não
da renda. Todo acréscimo de renda é absorvido pelo bem 2, que funciona
como um ``numerário'' --- um bem residual que absorve as variações de
poder aquisitivo. A utilidade quase-linear é particularmente útil em
análises de equilíbrio parcial e em modelos de organização industrial,
onde se deseja isolar o mercado de um bem específico sem que efeitos
renda contaminem a análise.

\begin{theorembox}{Proposição — Propriedades da utilidade quase-linear}

Seja \(u(x_1, x_2) = v(x_1) + x_2\) com \(v' > 0\) e \(v'' < 0\). Então,
para soluções interiores:

\begin{enumerate}
\def\labelenumi{\arabic{enumi}.}
\tightlist
\item
  \textbf{Efeito renda nulo sobre \(x_1\)}: a demanda marshalliana por
  \(x_1\) depende apenas de \(p_1/p_2\), não da renda \(I\).
\item
  \textbf{Translação vertical}: fixado \(x_1\), a curva de indiferença
  de nível \(\bar{u}\) é \(x_2 = \bar{u} - v(x_1)\). Ao passar para o
  nível \(\bar{u}' > \bar{u}\), a curva se desloca verticalmente em
  \(\bar{u}' - \bar{u}\).
\item
  \textbf{Medidas de bem-estar coincidem}: variação compensatória,
  variação equivalente e variação do excedente do consumidor são iguais
  --- \(VC = VE = \Delta EC\) --- porque a demanda hicksiana por \(x_1\)
  independe do nível de utilidade (ver Cap. 5, §5.8.4).
\item
  \textbf{Função dispêndio}:
  \(E(\mathbf{p}, \bar{u}) = c(p_1, p_2) + p_2 \bar{u}\), onde
  \(c(\cdot)\) é uma função apenas dos preços. A função dispêndio é
  \textbf{linear} em \(\bar{u}\).
\end{enumerate}

\end{theorembox}

\begin{intuitionbox}{Intuição Econômica}

\textbf{Em uma frase:} Na utilidade quase-linear, o bem 2 funciona como
``dinheiro'' --- todo aumento de renda vai para \(x_2\), sem afetar a
escolha de \(x_1\).

\textbf{Pense assim:} Imagine que você decide quanto café tomar por dia
baseado apenas no preço do café, e todo o resto do seu orçamento vai
para ``outros gastos'' (\(x_2\)). Se você ganha um aumento de salário,
continua tomando a mesma quantidade de café --- a renda extra vai toda
para os outros gastos. Essa é a essência da quase-linearidade.

\textbf{Por que isso importa:} Essa propriedade torna a utilidade
quase-linear a preferida em modelos de organização industrial e análises
de equilíbrio parcial: como \(VC = VE = \Delta EC\), a medida de
bem-estar é única e não ambígua.

\end{intuitionbox}

\textbf{Quase-linear vs.~homotética.} É importante distinguir claramente
a utilidade quase-linear das preferências homotéticas, pois ambas são
amplamente utilizadas mas possuem implicações opostas sobre o efeito
renda. A utilidade quase-linear \textbf{não} é homotética (exceto no
caso trivial em que \(v\) é linear). Enquanto as preferências
homotéticas geram curvas de indiferença que se expandem radialmente a
partir da origem --- mantendo constante a razão \(x_1/x_2\) ao longo do
caminho de expansão da renda ---, as curvas da utilidade quase-linear se
deslocam por translação vertical: a razão \(x_1/x_2\) muda com a renda,
pois todo o aumento vai para \(x_2\), violando a condição das funções
homotéticas (ver §3.6.6).

\subsection{3.6.6 Funções
homotéticas}\label{funuxe7uxf5es-homotuxe9ticas}

Após examinar a quase-linearidade --- onde todo o efeito renda recai
sobre um único bem ---, é natural perguntar: existe uma classe de
preferências em que a composição \emph{relativa} da cesta permanece
inalterada quando a renda varia? Ou seja, preferências tais que ricos e
pobres, ao enfrentarem os mesmos preços, escolham as mesmas
\emph{proporções} entre os bens, diferindo apenas na \emph{escala} de
consumo? A resposta são as preferências \textbf{homotéticas}, que ocupam
um papel central tanto na teoria do consumidor quanto nos modelos de
equilíbrio geral e macroeconômico. A homoteticidade é a propriedade que
permite agregar consumidores heterogêneos em um ``consumidor
representativo'' --- uma simplificação poderosa, embora nem sempre
realista, como discutiremos adiante.

\begin{definitionbox}{Função homotética}

Uma função de utilidade \(u(x_1, x_2)\) é \textbf{homotética} se pode
ser escrita como uma transformação monotônica crescente de uma função
homogênea de grau 1:

{[} u(x\_1, x\_2) = g!\big(h(x\_1, x\_2)\big), \quad g' \textgreater{}
0, \quad h(\lambda x\_1, \lambda x\_2) = \lambda , h(x\_1, x\_2)
;;\forall, \lambda \textgreater{} 0. {]}

\end{definitionbox}

\textbf{Propriedade fundamental.} Para preferências homotéticas, a TMS
depende apenas da \textbf{razão} \(x_1/x_2\), e não das quantidades
absolutas:

{[} \text{TMS}\_\{12\}(x\_1, x\_2) = \phi!\left(\frac{x_1}{x_2}\right).
\label{eq:3.6.15} \tag{3.6.15} {]}

A implicação geométrica é elegante: ao longo de qualquer raio que parte
da origem (\(x_2 = k \cdot x_1\)), a TMS é constante --- todas as curvas
de indiferença cruzam esse raio com a mesma inclinação. Isso significa
que as curvas de indiferença são \textbf{expansões radiais} umas das
outras: escalar uma curva de indiferença a partir da origem produz outra
curva de indiferença. Visualmente, o mapa de indiferença possui uma
simetria radial que contrasta com o deslocamento puramente vertical da
utilidade quase-linear. Essa simetria é precisamente o que garante que o
caminho de expansão da renda --- o conjunto de cestas ótimas para
diferentes níveis de renda, a preços fixos --- seja uma reta que passa
pela origem.

\begin{theorembox}{Proposição — Propriedades das preferências homotéticas}

Se \(u\) é homotética, então:

\begin{enumerate}
\def\labelenumi{\arabic{enumi}.}
\tightlist
\item
  \textbf{Caminho de expansão da renda linear}: a reta que passa pela
  origem e pela cesta ótima contém todas as cestas ótimas para
  diferentes níveis de renda (preços fixos). A razão \(x_1^*/x_2^*\) é
  constante em \(I\).
\item
  \textbf{Curvas de Engel lineares}:
  \(x_i^*(I) = \alpha_i(\mathbf{p}) \cdot I\), onde \(\alpha_i\) depende
  apenas dos preços.
\item
  \textbf{Elasticidade-renda unitária}: \(\varepsilon_{x_i, I} = 1\)
  para todo bem \(i\). Todos os bens são normais e nem de luxo nem de
  necessidade.
\item
  \textbf{Participação constante na despesa}: \(p_i x_i^* / I\) é
  constante para variações da renda.
\item
  \textbf{Função dispêndio separável}:
  \(E(\mathbf{p}, \bar{u}) = b(\mathbf{p}) \cdot e(\bar{u})\), onde
  \(b\) depende apenas dos preços e \(e\) apenas da utilidade.
\end{enumerate}

\end{theorembox}

\textbf{Exemplos.} A maioria das funções de utilidade estudadas neste
capítulo é homotética --- o que não é coincidência, pois a
homoteticidade simplifica enormemente a análise. Vale, portanto,
classificar cada uma delas:

\begin{itemize}
\tightlist
\item
  \textbf{Cobb-Douglas} \(u = x_1^a x_2^b\): é homotética, pois é
  homogênea de grau \(a+b\). A razão ótima é
  \(x_1^*/x_2^* = (a p_2)/(b p_1)\), constante em \(I\) --- a
  participação de cada bem na despesa total não se altera com a renda.
\item
  \textbf{CES} \(u = (x_1^\rho + x_2^\rho)^{1/\rho}\): homotética
  (homogênea de grau 1). A razão ótima depende apenas de \(p_1/p_2\) e
  \(\rho\), nunca da renda.
\item
  \textbf{Substitutos perfeitos} e \textbf{complementos perfeitos}:
  ambos são homotéticos (homogêneos de grau 1). No caso dos substitutos,
  a solução de canto faz com que toda a renda vá para um único bem, mas
  a \emph{proporção} não depende de \(I\).
\item
  \textbf{Quase-linear} \(u = v(x_1) + x_2\): \textbf{não} é homotética.
  A TMS depende de \(x_1\) isoladamente, não da razão \(x_1/x_2\). A
  razão ótima \(x_1^*/x_2^*\) varia com a renda, pois todo o acréscimo
  de renda é absorvido por \(x_2\).
\end{itemize}

\begin{boxbrasilbox}{Box Brasil 3.2 — Lei de Engel e a POF: evidência contra a homoteticidade}

A \textbf{Lei de Engel} (1857) afirma que a participação da alimentação
no orçamento familiar \emph{cai} à medida que a renda \emph{sobe}. Essa
regularidade empírica --- uma das mais robustas de toda a economia ---
implica que a elasticidade-renda da alimentação é menor que 1 (bem
necessário), violando a previsão de elasticidade unitária das
preferências homotéticas.

Os dados da \textbf{POF 2017-2018} (IBGE) confirmam a Lei de Engel para
o Brasil com notável clareza:

{\def\LTcaptype{none} % do not increment counter
\begin{longtable}[]{@{}
  >{\raggedright\arraybackslash}p{(\linewidth - 6\tabcolsep) * \real{0.2500}}
  >{\raggedright\arraybackslash}p{(\linewidth - 6\tabcolsep) * \real{0.2500}}
  >{\raggedright\arraybackslash}p{(\linewidth - 6\tabcolsep) * \real{0.2500}}
  >{\raggedright\arraybackslash}p{(\linewidth - 6\tabcolsep) * \real{0.2500}}@{}}
\toprule\noalign{}
\begin{minipage}[b]{\linewidth}\raggedright
Faixa de renda familiar (R\$/mês)
\end{minipage} & \begin{minipage}[b]{\linewidth}\raggedright
Participação da alimentação (\%)
\end{minipage} & \begin{minipage}[b]{\linewidth}\raggedright
Participação de transporte (\%)
\end{minipage} & \begin{minipage}[b]{\linewidth}\raggedright
Participação de educação (\%)
\end{minipage} \\
\midrule\noalign{}
\endhead
\bottomrule\noalign{}
\endlastfoot
Até 1.908 & 22,0 & 11,5 & 1,6 \\
1.908 a 2.862 & 18,5 & 14,2 & 2,0 \\
5.724 a 9.540 & 12,1 & 18,8 & 3,7 \\
Acima de 23.850 & 7,6 & 15,4 & 6,2 \\
\end{longtable}
}

A alimentação segue o padrão engeliano clássico (de 22\% para 7,6\%),
enquanto educação exibe comportamento oposto --- bem de luxo, com
participação crescente na renda. Transporte apresenta uma relação não
monotônica: cresce nas faixas intermediárias e recua nas mais altas
(possivelmente refletindo a troca de transporte público por automóvel
próprio já amortizado).

\textbf{Implicação teórica:} preferências homotéticas --- incluindo a
Cobb-Douglas --- são uma aproximação razoável \emph{dentro} de uma faixa
de renda estreita, mas não capturam a variação \emph{entre} faixas.
Modelos aplicados ao Brasil, como os que estimam impactos distributivos
de reformas tributárias, devem usar especificações não homotéticas (como
o AIDS de Deaton e Muellbauer) para capturar esses padrões.

\textbf{Fonte:} IBGE, Pesquisa de Orçamentos Familiares 2017-2018.

\end{boxbrasilbox}

\begin{intuitionbox}{Intuição Econômica}

\textbf{Em uma frase:} Com preferências homotéticas, ricos e pobres
gastam a mesma \emph{proporção} da renda em cada bem --- só a escala
muda.

\textbf{Pense assim:} Uma família homotética que gasta 30\% da renda com
alimentação e 70\% com outros bens manterá essa proporção se sua renda
dobrar, triplicar ou cair pela metade. O caminho de expansão da renda é
uma reta que sai da origem --- escalar a cesta ótima é como ``dar zoom''
na mesma cesta.

\textbf{Por que isso importa:} Essa propriedade permite agregar
consumidores com rendas diferentes em um ``consumidor representativo''
--- base de grande parte dos modelos macroeconômicos. Contudo, os dados
da POF (ver Box Brasil, §3.6.6) mostram que preferências reais raramente
são homotéticas: a participação da alimentação cai com a renda (Lei de
Engel), evidenciando preferências não homotéticas na prática.

\end{intuitionbox}

A \hyperref[tabela-3-1]{Tabela 3.1} a seguir sintetiza as principais
formas funcionais de utilidade discutidas neste capítulo, reunindo em
uma única referência a expressão da função, a TMS, o formato das curvas
de indiferença e a elasticidade de substituição de cada caso.
Recomenda-se consultá-la sempre que surgir dúvida sobre as propriedades
de uma forma funcional específica.

{\def\LTcaptype{none} % do not increment counter
\begin{longtable}[]{@{}
  >{\raggedright\arraybackslash}p{(\linewidth - 8\tabcolsep) * \real{0.2000}}
  >{\raggedright\arraybackslash}p{(\linewidth - 8\tabcolsep) * \real{0.2000}}
  >{\raggedright\arraybackslash}p{(\linewidth - 8\tabcolsep) * \real{0.2000}}
  >{\raggedright\arraybackslash}p{(\linewidth - 8\tabcolsep) * \real{0.2000}}
  >{\raggedright\arraybackslash}p{(\linewidth - 8\tabcolsep) * \real{0.2000}}@{}}
\toprule\noalign{}
\begin{minipage}[b]{\linewidth}\raggedright
Tipo
\end{minipage} & \begin{minipage}[b]{\linewidth}\raggedright
Função \(u(x_1, x_2)\)
\end{minipage} & \begin{minipage}[b]{\linewidth}\raggedright
TMS\(_{12}\)
\end{minipage} & \begin{minipage}[b]{\linewidth}\raggedright
Curvas de indiferença
\end{minipage} & \begin{minipage}[b]{\linewidth}\raggedright
Elasticidade de substituição (\(\sigma\))
\end{minipage} \\
\midrule\noalign{}
\endhead
\bottomrule\noalign{}
\endlastfoot
Cobb-Douglas & \(x_1^a x_2^b\) & \(\dfrac{a\, x_2}{b\, x_1}\) &
Hipérboles convexas & \(1\) \\
Substitutos perfeitos & \(ax_1 + bx_2\) & \(\dfrac{a}{b}\) (constante) &
Retas paralelas & \(\infty\) \\
Complementos perfeitos & \(\min\{ax_1, bx_2\}\) & Indefinida no vértice
& Ângulo reto (L) & \(0\) \\
CES & \((x_1^{\rho}+x_2^{\rho})^{1/\rho}\) &
\(\left(\frac{x_1}{x_2}\right)^{\rho-1}\) & Convexa (curvatura varia com
\(\rho\)) & \(\dfrac{1}{1-\rho}\) \\
Quase-linear & \(v(x_1) + x_2\) & \(v'(x_1)\) & Translações verticais &
Variável \\
Homotética (geral) & \(g(h(x_1,x_2))\), \(h\) homogênea grau 1 &
\(\phi(x_1/x_2)\) & Expansões radiais & Depende de \(h\) \\
\end{longtable}
}

\textbf{Tabela 3.1 --- Comparativa das funções de utilidade.}

\begin{tipbox}{Como escolher a forma funcional certa?}

Diante de tantas opções, como decidir qual função de utilidade usar em
um modelo ou exercício? A escolha depende do objetivo da análise:

\begin{itemize}
\tightlist
\item
  \textbf{Cobb-Douglas}: primeira escolha para modelos teóricos que
  exigem soluções analíticas em forma fechada. Ideal para exercícios e
  demonstrações pedagógicas. Limitação: elasticidade de substituição
  fixa em 1 e participação constante na despesa.
\item
  \textbf{CES}: preferida em trabalhos empíricos e modelos de comércio
  internacional (Armington), pois permite estimar a elasticidade de
  substituição \(\sigma\) a partir dos dados. Também útil quando se
  deseja generalidade teórica sem perder tratabilidade.
\item
  \textbf{Substitutos perfeitos}: para mercados em que os bens são
  essencialmente intercambiáveis (commodities homogêneas, marcas
  genéricas). Gera soluções de canto --- útil para ilustrar
  especialização.
\item
  \textbf{Complementos perfeitos (Leontief)}: para insumos que devem ser
  usados em proporções fixas. Aparece frequentemente na teoria da
  produção (Capítulo 10) e em modelos de equilíbrio geral computável.
\item
  \textbf{Quase-linear}: a escolha natural para análises de equilíbrio
  parcial e modelos de organização industrial, pois elimina efeitos
  renda e garante \(VC = VE = \Delta EC\). Ideal quando o foco é um
  mercado específico.
\end{itemize}

\textbf{Regra prática:} comece com a forma mais simples que capture o
fenômeno de interesse. Se a Cobb-Douglas for suficiente, não use a CES.

\end{tipbox}

\figurePlaceholder{/micro-book/graficos/cap03/tipos-preferencias.html}

\textbf{Figura 3.4 --- Comparação dos quatro tipos de preferências:
Cobb-Douglas (hipérboles convexas), substitutos perfeitos (retas),
complementos perfeitos (ângulo reto) e quase-linear (translações
verticais).}

\figurePlaceholder{/micro-book/graficos/cap03/funcoes-utilidade.html}

\textbf{Figura 3.5 --- Superfície 3D da função de utilidade.} Rotacione
e aplique zoom com o mouse. Use o menu para trocar entre Cobb-Douglas,
substitutos perfeitos (plano), complementos perfeitos (superfície em
cunha), CES e quase-linear. Ajuste os parâmetros nos sliders.

\figurePlaceholder{/micro-book/graficos/cap03/homotetica-vs-quaselinear.html}

\textbf{Figura 3.6 --- Homotética vs quase-linear: compare a expansão da
renda.} À esquerda (Cobb-Douglas), o caminho de expansão é um raio da
origem --- a razão \(x_1/x_2\) é constante. À direita (quase-linear), o
caminho é vertical --- \(x_1\) não muda com a renda. Ajuste \(I\) e
\(p_1/p_2\) nos sliders.

\begin{exresolvidobox}{Exercício Resolvido 3.2}

\textbf{Enunciado:} Considere a função CES
\(u(x_1, x_2) = (x_1^{\rho} + x_2^{\rho})^{1/\rho}\) com \(\rho = -1\).
(a) Calcule a elasticidade de substituição. (b) Derive a TMS. (c)
Compare as curvas de indiferença com os casos Cobb-Douglas e
complementos perfeitos.

\textbf{Dados:} \(\rho = -1\), logo
\(u(x_1, x_2) = (x_1^{-1} + x_2^{-1})^{-1}\).

\textbf{Resolução:}

\textbf{Passo 1 --- Elasticidade de substituição}

{[} \sigma = \frac{1}{1-\rho} = \frac{1}{1-(-1)} = \frac{1}{2} {]}

A elasticidade de substituição é \(1/2\), um valor entre 0 (complementos
perfeitos) e 1 (Cobb-Douglas). Isso indica substituibilidade baixa, mas
não nula.

\textbf{Passo 2 --- Cálculo da TMS}

Usando a fórmula geral da CES:

{[} \text{TMS}\_\{12\} = \left(\frac{x_1}{x_2}\right)\^{}\{\rho-1\} =
\left(\frac{x_1}{x_2}\right)\^{}\{-2\} =
\left(\frac{x_2}{x_1}\right)\^{}\{2\} {]}

\textbf{Passo 3 --- Comparação com outros casos}

Com \(\sigma = 1/2\), as curvas de indiferença são convexas e mais
``encurvadas'' do que as da Cobb-Douglas (\(\sigma = 1\)), mas sem o
ângulo reto dos complementos perfeitos (\(\sigma = 0\)). Intuitivamente,
o consumidor aceita trocar um bem pelo outro, mas exige compensações
crescentes de maneira mais acentuada do que no caso Cobb-Douglas.

\textbf{Resultado:} \(\sigma = 1/2\); \(\text{TMS}_{12} = (x_2/x_1)^2\).

\textbf{Interpretação econômica:} No contexto brasileiro, uma
elasticidade de substituição baixa como \(\sigma = 1/2\) pode descrever,
por exemplo, a relação entre energia elétrica e gás de cozinha para
cocção: o consumidor pode substituir parcialmente um pelo outro (fogão
elétrico vs.~fogão a gás), mas com dificuldade crescente --- refletindo
custos de troca de equipamentos e hábitos de consumo enraizados.

\end{exresolvidobox}

\begin{boxmundobox}{Box Mundo 3.2 — Elasticidades de substituição no comércio internacional: as elasticidades de Armington}

\textbf{Contexto:} A função CES apresentada na Seção 3.6.4 não é apenas
uma ferramenta teórica para modelar preferências individuais --- ela é a
base de praticamente todos os modelos de comércio internacional
utilizados por organizações como a OMC, o FMI e o Banco Mundial. Em
1969, Paul Armington propôs tratar bens importados de diferentes países
como substitutos imperfeitos, modelando as preferências dos compradores
(consumidores, firmas ou governos) por bens diferenciados por origem com
uma função CES. Nesse enquadramento, a elasticidade de substituição
\(\sigma\) mede o grau em que compradores estão dispostos a trocar, por
exemplo, aço brasileiro por aço chinês em resposta a variações de preços
relativos.

\textbf{Dados:} A estimação das chamadas ``elasticidades de Armington''
é uma indústria acadêmica por si só. O trabalho de referência de Broda e
Weinstein (2006) estimou elasticidades de substituição para mais de
10.000 categorias de produtos importados pelos Estados Unidos no período
1972--2001 e encontrou uma mediana de \(\sigma \approx 3{,}1\), com
variação enorme entre setores: commodities homogêneas como petróleo
bruto apresentam \(\sigma > 20\) (quase substitutos perfeitos entre
origens), enquanto bens diferenciados como automóveis e vinhos
apresentam \(\sigma\) entre 1 e 3 (substituição limitada). O modelo GTAP
(\emph{Global Trade Analysis Project}), utilizado pela OMC para avaliar
rodadas de negociação, adota elasticidades de Armington calibradas
setorialmente: \(\sigma = 5{,}2\) para grãos, \(\sigma = 2{,}8\) para
manufaturas e \(\sigma = 1{,}9\) para serviços. Hertel et al.~(2007)
mostram que os resultados de simulação são altamente sensíveis a esses
parâmetros: dobrar a elasticidade de Armington pode triplicar os ganhos
estimados de uma liberalização comercial.

\textbf{Análise:} A aplicação das elasticidades de Armington ao comércio
internacional ilustra três pontos centrais da Seção 3.6. Primeiro, a
flexibilidade da função CES: ao variar \(\sigma\), o mesmo arcabouço
analítico cobre desde bens quase homogêneos (\(\sigma \to \infty\),
substitutos perfeitos) até bens altamente diferenciados
(\(\sigma \to 0\), complementos perfeitos). Segundo, a importância
empírica do parâmetro \(\sigma\): políticas comerciais que parecem
benéficas sob elasticidades altas podem ser inócuas ou até prejudiciais
sob elasticidades baixas. Terceiro, a conexão entre microeconomia e
macroeconomia: preferências CES de consumidores individuais, quando
agregadas, determinam fluxos de comércio entre nações e os efeitos de
políticas tarifárias sobre o bem-estar global. Para o Brasil, as
elasticidades de Armington são particularmente relevantes na avaliação
de acordos comerciais como o Mercosul-UE, onde o grau de
substituibilidade entre produtos brasileiros e europeus determina os
ganhos e perdas setoriais esperados.

\textbf{Fonte:} Armington, P. S. (1969). A theory of demand for products
distinguished by place of production. \emph{IMF Staff Papers}, 16(1),
159--178. Broda, C.; Weinstein, D. E. (2006). Globalization and the
gains from variety. \emph{Quarterly Journal of Economics}, 121(2),
541--585. Hertel, T.; Hummels, D.; Ivanic, M.; Keeney, R. (2007). How
confident can we be of CGE-based assessments of free trade agreements?
\emph{Economic Modelling}, 24(4), 611--635.

\end{boxmundobox}

Em resumo: a Cobb-Douglas é a melhor amiga do microeconomista --- fácil
de usar, gera soluções fechadas, e nunca reclama. A CES é mais versátil
mas exige mais trabalho. A Leontief é inflexível --- literalmente: só
aceita proporções fixas. E a quase-linear é a favorita da análise de
bem-estar, porque isola efeito renda em um único bem. Escolher a forma
funcional certa para cada problema é metade da arte.

\begin{center}\rule{0.5\linewidth}{0.5pt}\end{center}

\section{3.7 Troque o Rótulo, Não o Ranking: Transformações
Monotônicas}\label{troque-o-ruxf3tulo-nuxe3o-o-ranking-transformauxe7uxf5es-monotuxf4nicas}

\emph{And now for something completely different} --- ou melhor,
\emph{exactly the same thing, but with different numbers}. Aqui vai uma
informação libertadora: a função de utilidade que você escolhe \emph{não
importa} --- desde que preserve o ranking. Se \(u\) diz que pizza é
melhor que salada, e você aplica logaritmo, \(\ln u\) \emph{continua}
dizendo que pizza é melhor que salada. Pode elevar ao cubo, multiplicar
por mil, somar uma constante --- nada muda na economia. É como medir
temperatura em Celsius ou Fahrenheit: os números são diferentes, mas a
água ferve no mesmo ponto.

Essa liberdade é um presente prático enorme: quando uma Cobb-Douglas
\(u = x_1^{0{,}3} x_2^{0{,}7}\) está dando dor de cabeça, basta tomar o
log e trabalhar com \(\ln u = 0{,}3 \ln x_1 + 0{,}7 \ln x_2\) --- que é
linear e muito mais tratável. As preferências são as mesmas; só a
embalagem mudou.

\begin{theorembox}{Proposição 3.3 — Invariância sob transformação monotônica}

Se \(u(\mathbf{x})\) representa \(\succsim\) e
\(f: \mathbb{R} \to \mathbb{R}\) é estritamente crescente, então
\(\hat{u}(\mathbf{x}) = f(u(\mathbf{x}))\) também representa
\(\succsim\).

\end{theorembox}

A demonstração é direta e vale a pena percorrer, pois ilustra a
simplicidade do argumento ordinal: como \(f\) é estritamente crescente,
a desigualdade \(u(\mathbf{x}) \geq u(\mathbf{y})\) é preservada pela
aplicação de \(f\), de modo que
\(f(u(\mathbf{x})) \geq f(u(\mathbf{y}))\). Reciprocamente, se
\(f(u(\mathbf{x})) \geq f(u(\mathbf{y}))\), a monotonicidade estrita de
\(f\) garante que \(u(\mathbf{x}) \geq u(\mathbf{y})\). Portanto,
\(\hat{u}\) preserva exatamente o mesmo ordenamento que \(u\).

\begin{intuitionbox}{Intuição Econômica}

\textbf{Em uma frase:} Você pode trocar a ``régua'' que mede a utilidade
sem alterar as preferências --- qualquer escala que preserve a ordem
serve.

\textbf{Pense assim:} É como converter temperatura de Celsius para
Fahrenheit: os números mudam, mas a água continua fervendo antes do
metal derreter. Da mesma forma, se você aplica logaritmo na utilidade
Cobb-Douglas para facilitar as contas, o consumidor continua preferindo
as mesmas cestas de antes.

\textbf{Por que isso importa:} Isso dá ao economista liberdade para
escolher a forma funcional que facilite os cálculos --- por exemplo,
trabalhar com \(\ln u\) em vez de \(u\) --- sem perder informação
econômica.

\end{intuitionbox}

\textbf{Exemplos de transformações monotônicas úteis.} Duas aplicações
da invariância ordinal são especialmente frequentes na prática:

\begin{itemize}
\tightlist
\item
  A Cobb-Douglas \(u = x_1^a x_2^b\) pode ser transformada em
  \(\hat{u} = a \ln x_1 + b \ln x_2\) via \(f(u) = \ln(u)\). O logaritmo
  transforma o produto em soma, eliminando os expoentes e simplificando
  consideravelmente as derivadas parciais e a resolução das condições de
  primeira ordem.
\item
  Qualquer Cobb-Douglas pode ser normalizada de modo que
  \(\hat{u} = \frac{a}{a+b} \ln x_1 + \frac{b}{a+b} \ln x_2\), com
  coeficientes somando 1. Essa normalização é conveniente porque os
  coeficientes passam a coincidir com as participações na despesa ---
  tornando a interpretação econômica imediata.
\end{itemize}

\figurePlaceholder{/micro-book/graficos/cap03/transformacao-monotonica.html}

\textbf{Figura 3.7 --- Transformação Monotônica e Invariância Ordinal.}
Esquerda: curvas de indiferença de \(U = x_1 x_2\). Direita: curvas da
transformada \(\hat{U} = f(U)\). As curvas são \textbf{idênticas} ---
apenas os rótulos numéricos mudam. Escolha entre \(\ln\), \(\sqrt{}\),
\(U^2\), \(U^3\) e \(-1/U\). Mova o ponto de referência para verificar
que a TMS é a mesma em ambos os painéis.

\begin{examplebox}{Transformação monotônica em ação: simplificando a CES}

A função CES \(u = (x_1^\rho + x_2^\rho)^{1/\rho}\) envolve uma raiz que
complica as derivadas parciais. Aplique a transformação monotônica
\(f(u) = u^\rho\) (estritamente crescente para \(u > 0\) quando
\(\rho > 0\), e estritamente decrescente --- mas aí tomamos
\(f(u) = -u^\rho\) --- para \(\rho < 0\)). O resultado é:

{[} \hat{u} = x\_1\^{}\rho + x\_2\^{}\rho {]}

Agora as CPOs ficam muito mais simples:

{[} \frac{\partial \hat{u}}{\partial x_1} = \rho , x\_1\^{}\{\rho-1\},
\qquad \frac{\partial \hat{u}}{\partial x_2} = \rho , x\_2\^{}\{\rho-1\}
{]}

E a TMS se obtém diretamente: \(\text{TMS}_{12} = (x_1/x_2)^{\rho-1}\)
--- o mesmo resultado da equação (\ref{eq:3.6.12}), mas com cálculo
muito mais rápido. Essa é a essência da invariância ordinal: simplificar
a álgebra sem alterar a economia.

\end{examplebox}

\begin{warningbox}{Erro Comum}

\textbf{Tratar a utilidade como grandeza cardinal --- dizer que ``Maria
tem utilidade 2 vezes maior que João''.}

A função de utilidade ordinal mede apenas \emph{ordem}, não
\emph{intensidade}. Se \(u(\text{pizza}) = 10\) e
\(u(\text{salada}) = 5\), isso significa apenas que pizza é preferida a
salada --- \textbf{não} que pizza é ``duas vezes melhor''. Qualquer
transformação monotônica (como \(\hat{u} = \ln u\)) preserva as
preferências mas muda a razão entre os valores numéricos. Comparações
interpessoais (``Maria gosta mais de pizza que João'') também são
impossíveis com utilidade ordinal --- um ponto que se torna crucial na
economia do bem-estar (Capítulo 17).

\end{warningbox}

\begin{tipbox}{Implicações práticas da ordinalidade}

A TMS é invariante sob transformações monotônicas. A utilidade marginal,
porém, \textbf{não é invariante}: ela muda com a transformação. Por
isso, a utilidade marginal não tem significado cardinal --- apenas a
razão entre utilidades marginais (a TMS) possui significado econômico
bem definido.

\end{tipbox}

\begin{exresolvidobox}{Exercício Resolvido 3.3}

\textbf{Enunciado:} Mostre que \(u(x_1, x_2) = \ln x_1 + 2\ln x_2\) e
\(v(x_1, x_2) = x_1 \cdot x_2^2\) representam as mesmas preferências, e
verifique que ambas produzem a mesma TMS.

\textbf{Resolução:}

\textbf{Passo 1 --- Identificar a transformação monotônica}

Note que \(u = \ln x_1 + 2\ln x_2 = \ln(x_1 \cdot x_2^2) = \ln(v)\).
Portanto, \(u = f(v)\) com \(f(v) = \ln(v)\), que é estritamente
crescente para \(v > 0\). Pela Proposição 3.3, \(u\) e \(v\) representam
as mesmas preferências.

\textbf{Passo 2 --- TMS pela função \(u\)}

{[} \text{TMS}\_\{12\}\^{}\{(u)\} =
\frac{\partial u/\partial x_1}{\partial u/\partial x_2} =
\frac{1/x_1}{2/x_2} = \frac{x_2}{2x_1} {]}

\textbf{Passo 3 --- TMS pela função \(v\)}

{[} \text{TMS}\_\{12\}\^{}\{(v)\} =
\frac{\partial v/\partial x_1}{\partial v/\partial x_2} =
\frac{x_2^2}{2x_1 x_2} = \frac{x_2}{2x_1} {]}

\textbf{Resultado:}
\(\text{TMS}_{12}^{(u)} = \text{TMS}_{12}^{(v)} = \dfrac{x_2}{2x_1}\).

\textbf{Interpretação econômica:} As utilidades marginais diferem ---
para \(u\), \(\text{UMg}_1 = 1/x_1\); para \(v\),
\(\text{UMg}_1 = x_2^2\) --- mas a TMS é idêntica. Isso confirma que a
TMS é a grandeza economicamente relevante: ela mede a taxa de troca
subjetiva do consumidor, independentemente da ``escala'' escolhida para
medir a utilidade. Na prática, o economista pode escolher a
representação que facilite os cálculos sem perda de conteúdo econômico.

\end{exresolvidobox}

\begin{center}\rule{0.5\linewidth}{0.5pt}\end{center}

As seções anteriores construíram progressivamente o arcabouço teórico
das preferências: partimos de axiomas mínimos de coerência, chegamos à
função de utilidade via o Teorema de Debreu, derivamos curvas de
indiferença e TMS como ferramentas geométricas e analíticas, exploramos
as principais formas funcionais --- Cobb-Douglas, substitutos e
complementos perfeitos, CES e quase-linear --- e, por fim, demonstramos
que a invariância ordinal confere liberdade ao analista na escolha da
representação. A progressão foi deliberada: cada conceito se construiu
sobre os anteriores, formando uma cadeia lógica que vai da abstração dos
axiomas à concretude das formas funcionais.

Antes de consolidar esses conceitos no resumo do capítulo, vale observar
como esse arcabouço se manifesta --- e, por vezes, é posto à prova ---
em um contexto histórico concreto do Brasil. O Box a seguir mostra que
as mudanças nos padrões de consumo após a estabilização monetária de
1994 podem ser interpretadas à luz da teoria das preferências
desenvolvida nestas páginas.

\begin{boxbrasilbox}{Box Brasil 3.3 — O Plano Real e a revolução no consumo}

A hiperinflação brasileira que antecedeu o Plano Real (julho de 1994)
distorcia profundamente as preferências reveladas dos consumidores. Com
taxas de inflação que chegaram a superar 2.000\% ao ano em 1993, o
comportamento de compra era dominado pela \textbf{antecipação de
consumo}: famílias corriam ao supermercado no dia do pagamento para
converter salários em bens antes que os preços subissem.

Dados da Pesquisa de Orçamentos Familiares (POF) do IBGE revelam
mudanças marcantes nos padrões de consumo ao longo das décadas:

\begin{itemize}
\tightlist
\item
  \textbf{Antes da estabilização (POF 1987-88)}: famílias de baixa renda
  concentravam até 53\% dos gastos em alimentação, com forte preferência
  por produtos estocáveis e não perecíveis --- um reflexo racional da
  tentativa de manter o valor real da renda.
\item
  \textbf{Após a estabilização (POF 1995-96 e 2002-03)}: a participação
  da alimentação caiu para cerca de 33\% nas famílias de menor renda,
  com diversificação para bens duráveis, serviços de lazer e educação.
\item
  \textbf{POF 2017-18}: a tendência de diversificação se manteve, com
  crescimento expressivo dos gastos com comunicação (telefonia móvel) e
  transporte.
\end{itemize}

Do ponto de vista da teoria das preferências, a estabilização monetária
não apenas alterou a \textbf{restrição orçamentária} (via ganho de renda
real, sobretudo para os mais pobres), mas também permitiu que as
\textbf{preferências subjacentes} se manifestassem sem a distorção
imposta pelo imposto inflacionário.

O conceito de TMS decrescente ganha concretude nesse contexto: à medida
que cestas de alimentação básica se tornaram acessíveis com menor fração
da renda, os consumidores puderam mover-se ao longo de suas curvas de
indiferença em direção a cestas mais diversificadas.

Vasconcellos e Garcia (2014) contextualizam os ciclos econômicos
brasileiros que moldaram essas transformações nas escolhas de consumo.

\textbf{Fonte}: IBGE, Pesquisa de Orçamentos Familiares (várias
edições); Barros, R. P. de; Foguel, M. N.; Ulyssea, G. (orgs.).
\emph{Desigualdade de renda no Brasil: uma análise da queda recente}.
Brasília: IPEA, 2007.

\end{boxbrasilbox}

\begin{boxmundobox}{Box Mundo 3.3 — Preferência revelada em big data: dados de scanner de supermercados e testes do GARP}

\textbf{Contexto:} A teoria do consumidor desenvolvida neste capítulo
repousa sobre axiomas --- completude, transitividade, continuidade,
monotonicidade --- que são hipóteses sobre as preferências dos agentes.
Mas como verificar empiricamente se consumidores reais se comportam de
acordo com esses axiomas? A resposta clássica vem da teoria da
preferência revelada (que será formalizada na Seção 5.9): se observarmos
as escolhas de um consumidor em diferentes situações de preço e renda,
podemos testar se essas escolhas são compatíveis com a maximização de
alguma função de utilidade. O critério operacional é o \textbf{Axioma
Generalizado da Preferência Revelada} (GARP, \emph{Generalized Axiom of
Revealed Preference}), formulado por Varian (1982): se as escolhas
observadas satisfazem o GARP, então existe uma função de utilidade que
as racionaliza. A revolução dos dados de scanner de supermercados ---
registros eletrônicos de cada item comprado por cada consumidor ---
abriu, a partir dos anos 2000, a possibilidade de testar essa hipótese
com milhões de observações individuais.

\textbf{Dados:} Echenique, Lee e Shum (2011, \emph{American Economic
Review}) utilizaram dados de scanner do painel Homescan da Nielsen, que
rastreia todas as compras de supermercado de uma amostra representativa
de famílias americanas ao longo de vários anos. Com observações de
preços e quantidades de centenas de produtos para cada família, os
autores testaram o GARP individualmente para mais de 500 famílias. O
resultado principal: apenas cerca de 5\% das famílias apresentaram
violações do GARP, e mesmo essas violações foram de magnitude pequena
--- medidas pelo índice de Afriat (1967), que quantifica o grau de
inconsistência, as violações raramente ultrapassaram 1\% do orçamento.
Dean e Martin (2016, \emph{American Economic Review}) confirmaram esses
resultados usando dados experimentais com escolhas sobre cestas reais de
bens, encontrando taxas de violação do GARP inferiores a 10\%. Na
Europa, Cherchye, De Rock e Vermeulen (2011, \emph{Review of Economics
and Statistics}) aplicaram testes de preferência revelada a microdados
belgas e holandeses, concluindo que o modelo coletivo de consumo
domiciliar (onde marido e mulher têm preferências distintas) racionaliza
os dados significativamente melhor do que o modelo unitário ---
evidência de que a unidade de decisão importa para a validade dos
axiomas.

\textbf{Análise:} Esses resultados fornecem suporte empírico notável
para os axiomas apresentados na Seção 3.1: na grande maioria dos casos,
as escolhas de consumo observadas em dados de alta frequência são
compatíveis com a existência de uma função de utilidade racional. A taxa
de violação de 5\% é consistente com a evidência de intransitividade
discutida no Box Mundo 3.1 --- violações existem, mas são a exceção, não
a regra. Do ponto de vista da Seção 3.7 (invariância ordinal), os testes
de preferência revelada são particularmente elegantes: eles verificam se
existe \emph{alguma} função de utilidade compatível com os dados, sem
impor uma forma funcional específica (Cobb-Douglas, CES, etc.) ---
exatamente porque a utilidade é ordinal e qualquer transformação
monotônica preserva o ordenamento. A explosão de dados de scanner e de
compras online oferece oportunidades sem precedentes para testar a
teoria do consumidor em escala, transformando os axiomas deste capítulo
de postulados filosóficos em hipóteses empiricamente verificáveis.

\textbf{Fonte:} Varian, H. R. (1982). The nonparametric approach to
demand analysis. \emph{Econometrica}, 50(4), 945--973. Echenique, F.;
Lee, S.; Shum, M. (2011). The money pump as a measure of revealed
preference violations. \emph{American Economic Review}, 101(4),
1645--1651. Cherchye, L.; De Rock, B.; Vermeulen, F. (2011). The
revealed preference approach to collective consumption behaviour.
\emph{Review of Economics and Statistics}, 93(1), 223--238.

\end{boxmundobox}

Três convicções para levar deste capítulo:

\begin{enumerate}
\def\labelenumi{\arabic{enumi}.}
\item
  \textbf{Os axiomas são modestos --- as consequências, não.} Exigir que
  o consumidor saiba comparar cestas e não se contradiga parece pouco.
  Mas desse ponto de partida modesto emerge todo o edifício: função de
  utilidade, curvas de indiferença, TMS, demandas, bem-estar. A
  desproporção entre premissa e consequência é o que torna a teoria
  elegante.
\item
  \textbf{A forma funcional importa --- e muito.} Cobb-Douglas, CES,
  Leontief e quase-linear não são caprichos algébricos: cada uma implica
  um comportamento de demanda, uma elasticidade de substituição e uma
  resposta à renda \emph{qualitativamente diferente}. Escolher a forma
  funcional errada é como usar o mapa errado --- tecnicamente bonito,
  mas te leva ao destino errado.
\item
  \textbf{A utilidade é ordinal --- e isso é liberdade, não limitação.}
  A invariância sob transformações monotônicas significa que você pode
  trocar a ``régua'' sem perder informação econômica. Logaritmo, raiz,
  cubo --- nada muda na economia. A TMS é a grandeza real; os números da
  utilidade são mera conveniência.
\end{enumerate}

\emph{Sabemos o que o consumidor quer. No próximo capítulo, descobrimos
o que ele pode --- e o que ele escolhe.}

\chapter{Exercícios e ANPEC}\label{exercuxedcios-e-anpec}

\begin{tipbox}{Dados para exercicios empiricos}

\begin{itemize}
\tightlist
\item
  \textbf{POF/IBGE (cestas de consumo):}
  \href{https://sidra.ibge.gov.br}{sidra.ibge.gov.br} --- Tabelas
  7169--7175 (POF 2017-2018). Explore como familias de diferentes faixas
  de renda alocam gastos entre alimentacao, transporte e lazer ---
  revelando preferencias diferentes.
\item
  \textbf{PNAD Continua (consumo e renda):}
  \href{https://www.ibge.gov.br/estatisticas/sociais/trabalho/9171-pesquisa-nacional-por-amostra-de-domicilios-continua-mensal.html}{ibge.gov.br/pnadc}
  --- Series trimestrais para analisar como a composicao do consumo muda
  com a renda (curvas de Engel).
\item
  \textbf{Notebook:} N1 --- Elasticidades de alimentos (POF 2017-2018,
  metodo de Deaton).
\end{itemize}

\end{tipbox}

\figurePlaceholder{/micro-book/graficos/cap03/webr-utilidade.html}

\textbf{WebR 3.1 --- Curvas de indiferença comparadas.} Troque entre
Cobb-Douglas, CES, quaselinear e substitutos perfeitos. Altere o
parâmetro ρ na CES para ver a transição de substitutos (isoquantas
retas) a complementos (formato de L).

\figurePlaceholder{/micro-book/graficos/cap03/webr-tms.html}

\textbf{WebR 3.2 --- TMS ao longo da curva de indiferença.} Calcule a
TMS em vários pontos de uma Cobb-Douglas e veja a tangente mudar.
Observe: quanto mais x você tem, menos y está disposto a trocar ---
convexidade em ação.

\figurePlaceholder{/micro-book/graficos/cap03/webr-preferencias.html}

\textbf{WebR 3.3 --- Zoológico de preferências.} Quatro famílias em 4
painéis: Cobb-Douglas (suave), substitutos perfeitos (retas), Leontief
(L) e quase-linear (paralelas). Cada forma de CI conta uma história
diferente sobre como o consumidor troca bens.

\figurePlaceholder{/micro-book/graficos/cap03/webr-transf-monotonica.html}

\textbf{WebR 3.4 --- Transformação monotônica preserva tudo.} Três
funções utilidade (U, ln U, U²) geram a \emph{mesma} curva de
indiferença e a \emph{mesma} TMS. A utilidade é ordinal: só o ranking
importa, não o número.

\begin{center}\rule{0.5\linewidth}{0.5pt}\end{center}

\begin{classroombox}{Atividade 3.1 — Desenhando suas próprias curvas de indiferença (20 min)}

\textbf{Objetivo:} Descobrir, por introspecção guiada, que preferências
pessoais geram curvas de indiferença --- e que a TMS varia ao longo
delas.

\textbf{Material:} Folha quadriculada, caneta colorida.

\textbf{Protocolo:}

\begin{enumerate}
\def\labelenumi{\arabic{enumi}.}
\tightlist
\item
  \textbf{Setup} (3 min): ``Eixo horizontal = horas de Netflix por
  semana. Eixo vertical = horas de academia por semana. Marque o ponto
  (10, 2) --- sua dotação atual.''
\item
  \textbf{Revelação de preferências} (7 min): ``Agora responda
  honestamente: se eu tirasse 2 horas de Netflix, quantas horas a mais
  de academia você precisaria para ficar \emph{igualmente feliz}?''
  Marque o novo ponto. Repita para tirar mais 2 horas. E mais 2. Conecte
  os pontos --- isso é sua curva de indiferença.
\item
  \textbf{Comparação em duplas} (4 min): Duplas comparam curvas.
  Pergunte: ``As curvas são iguais? Deveriam ser?'' → Não! Preferências
  são subjetivas. ``Qual é a TMS de cada um no ponto (10, 2)?''
\item
  \textbf{Debrief} (6 min):

  \begin{itemize}
  \tightlist
  \item
    ``A curva ficou convexa (curvada para a origem)? Por quê?'' → TMS
    decrescente: quando você já tem pouca Netflix, cada hora a menos dói
    mais.
  \item
    ``Alguém desenhou uma reta?'' → Substitutos perfeitos. ``Alguém fez
    um L?'' → Complementos perfeitos.
  \item
    ``Se eu multiplicar todos os números de utilidade por 2, a curva
    muda?'' → Não. Transformação monotônica (Seção 3.7).
  \item
    Conecte: o exercício que vocês fizeram é exatamente como Edgeworth
    imaginou as curvas em 1881 --- só que ele não tinha Netflix.
  \end{itemize}
\end{enumerate}

\textbf{Conexão com o conteúdo:} Seções 3.3 (curvas de indiferença), 3.4
(TMS), 3.6 (formas funcionais), 3.7 (transformações monotônicas).

\end{classroombox}

\begin{classroombox}{Atividade 3.2 — O teste da transitividade: flagrando incoerência (10 min)}

\textbf{Objetivo:} Testar se as preferências declaradas dos alunos
satisfazem os axiomas --- e mostrar que violações existem (e importam).

\textbf{Material:} 3 cartões por aluno com opções escritas (ou projeção
no slide).

\textbf{Protocolo:}

\begin{enumerate}
\def\labelenumi{\arabic{enumi}.}
\tightlist
\item
  \textbf{As opções} (2 min): Apresente três ``cestas'' de lazer para o
  fim de semana: (A) Cinema + pizza, (B) Futebol + churrasco, (C) Praia
  + açaí. Cada aluno vota nas comparações dois a dois: A vs.~B, B vs.~C,
  A vs.~C (levantar a mão ou usar Mentimeter).
\item
  \textbf{Contagem} (3 min): O professor registra as proporções.
  Tipicamente, com preferências heterogêneas na turma, o resultado
  \emph{agregado} pode violar transitividade: A \textgreater{} B, B
  \textgreater{} C, mas C \textgreater{} A (Paradoxo de Condorcet!).
\item
  \textbf{Debrief} (5 min):

  \begin{itemize}
  \tightlist
  \item
    ``Se a \emph{turma} fosse uma pessoa, ela seria irracional ---
    prefere A a B, B a C, e C a A. Isso é possível para um indivíduo?''
    → Não, se satisfaz o axioma de transitividade.
  \item
    ``Mas o \emph{grupo} violou. O que isso significa?'' → Agregação de
    preferências individuais racionais pode gerar intransitividade
    coletiva --- antecipe o Teorema de Arrow (Capítulo 14, Seção 14.12).
  \item
    Feche: ``Os axiomas do Capítulo 3 são sobre \emph{indivíduos}. Para
    \emph{sociedades}, a coisa complica --- e Arrow ganhou o Nobel por
    explicar por quê.''
  \end{itemize}
\end{enumerate}

\textbf{Conexão com o conteúdo:} Seções 3.1 (axiomas de completude e
transitividade), antecipação do Capítulo 14, Seção 14.12 (Escolha Social
e Teorema de Arrow).

\end{classroombox}

\begin{center}\rule{0.5\linewidth}{0.5pt}\end{center}

\section{Revisão Rápida}\label{revisuxe3o-ruxe1pida}

Teste seu entendimento dos conceitos centrais deste capítulo.

\begin{infobox}{1. O axioma de transitividade das preferências implica que:}

\begin{itemize}
\tightlist
\item
  \begin{enumerate}
  \def\labelenumi{(\alph{enumi})}
  \tightlist
  \item
    O consumidor sempre prefere mais a menos
  \end{enumerate}
\item
  \begin{enumerate}
  \def\labelenumi{(\alph{enumi})}
  \setcounter{enumi}{1}
  \tightlist
  \item
    Se \(A \succsim B\) e \(B \succsim C\), então \(A \succsim C\)
  \end{enumerate}
\item
  \begin{enumerate}
  \def\labelenumi{(\alph{enumi})}
  \setcounter{enumi}{2}
  \tightlist
  \item
    O consumidor consegue comparar quaisquer duas cestas
  \end{enumerate}
\item
  \begin{enumerate}
  \def\labelenumi{(\alph{enumi})}
  \setcounter{enumi}{3}
  \tightlist
  \item
    A utilidade marginal é sempre decrescente
  \end{enumerate}
\end{itemize}

??? success ``Resposta'' \textbf{(b)} Transitividade exige consistência
nas ordenações: se A é pelo menos tão bom quanto B, e B pelo menos tão
bom quanto C, então A deve ser pelo menos tão bom quanto C. A
alternativa (a) descreve monotonicidade; (c) descreve completude; (d) é
uma propriedade da função de utilidade, não um axioma sobre
preferências.

\end{infobox}

\begin{infobox}{2. Se $u(x_1, x_2) = x_1^{0{,}5} x_2^{0{,}5}$ e $v = 2u + 10$, é correto afirmar que:}

\begin{itemize}
\tightlist
\item
  \begin{enumerate}
  \def\labelenumi{(\alph{enumi})}
  \tightlist
  \item
    \(v\) representa preferências diferentes de \(u\), pois os valores
    numéricos mudam
  \end{enumerate}
\item
  \begin{enumerate}
  \def\labelenumi{(\alph{enumi})}
  \setcounter{enumi}{1}
  \tightlist
  \item
    \(v\) representa as mesmas preferências, pois é uma transformação
    monotônica crescente de \(u\)
  \end{enumerate}
\item
  \begin{enumerate}
  \def\labelenumi{(\alph{enumi})}
  \setcounter{enumi}{2}
  \tightlist
  \item
    \(v\) só representa as mesmas preferências se a renda do consumidor
    não mudar
  \end{enumerate}
\item
  \begin{enumerate}
  \def\labelenumi{(\alph{enumi})}
  \setcounter{enumi}{3}
  \tightlist
  \item
    \(v\) inverte a ordenação das cestas em relação a \(u\)
  \end{enumerate}
\end{itemize}

??? success ``Resposta'' \textbf{(b)} A função de utilidade é ordinal:
qualquer transformação monotônica crescente \(v = g(u)\) com \(g' > 0\)
preserva a ordenação das cestas e, portanto, representa as mesmas
preferências. A alternativa (a) confunde valores cardinais com ordinais;
(c) adiciona condição inexistente; (d) seria verdade apenas para
transformações decrescentes.

\end{infobox}

\begin{infobox}{3. A Taxa Marginal de Substituição (TMS) entre dois bens mede:}

\begin{itemize}
\tightlist
\item
  \begin{enumerate}
  \def\labelenumi{(\alph{enumi})}
  \tightlist
  \item
    A razão entre os preços dos dois bens no mercado
  \end{enumerate}
\item
  \begin{enumerate}
  \def\labelenumi{(\alph{enumi})}
  \setcounter{enumi}{1}
  \tightlist
  \item
    A quantidade de um bem que o consumidor está disposto a abrir mão
    para obter uma unidade adicional do outro, mantendo a utilidade
    constante
  \end{enumerate}
\item
  \begin{enumerate}
  \def\labelenumi{(\alph{enumi})}
  \setcounter{enumi}{2}
  \tightlist
  \item
    A variação percentual na demanda de um bem quando o preço do outro
    varia
  \end{enumerate}
\item
  \begin{enumerate}
  \def\labelenumi{(\alph{enumi})}
  \setcounter{enumi}{3}
  \tightlist
  \item
    O custo de oportunidade de produzir um bem em vez do outro
  \end{enumerate}
\end{itemize}

??? success ``Resposta'' \textbf{(b)} A TMS é a inclinação da curva de
indiferença e mede a taxa subjetiva de troca entre bens que mantém o
consumidor indiferente. Matematicamente,
\(\text{TMS}_{12} = -dx_2/dx_1|_{u=\bar{u}} = \text{UMg}_1/\text{UMg}_2\).
A alternativa (a) descreve a inclinação da restrição orçamentária; (c)
descreve elasticidade cruzada; (d) refere-se à produção.

\end{infobox}

\begin{infobox}{4. As curvas de indiferença de preferências estritamente convexas são:}

\begin{itemize}
\tightlist
\item
  \begin{enumerate}
  \def\labelenumi{(\alph{enumi})}
  \tightlist
  \item
    Retas com inclinação constante
  \end{enumerate}
\item
  \begin{enumerate}
  \def\labelenumi{(\alph{enumi})}
  \setcounter{enumi}{1}
  \tightlist
  \item
    Convexas em relação à origem (curvadas para fora)
  \end{enumerate}
\item
  \begin{enumerate}
  \def\labelenumi{(\alph{enumi})}
  \setcounter{enumi}{2}
  \tightlist
  \item
    Côncavas em relação à origem (curvadas para dentro, `arqueadas')
  \end{enumerate}
\item
  \begin{enumerate}
  \def\labelenumi{(\alph{enumi})}
  \setcounter{enumi}{3}
  \tightlist
  \item
    Círculos concêntricos
  \end{enumerate}
\end{itemize}

??? success ``Resposta'' \textbf{(c)} Preferências convexas implicam que
médias são preferidas a extremos, gerando curvas de indiferença côncavas
em relação à origem (convexas do ponto de vista do conjunto preferido).
A TMS é decrescente em valor absoluto. A alternativa (a) descreve
substitutos perfeitos; (b) descreve preferências côncavas (não
convexas); (d) não é uma forma padrão.

\end{infobox}

\begin{infobox}{5. Qual das seguintes funções de utilidade representa preferências por complementos perfeitos?}

\begin{itemize}
\tightlist
\item
  \begin{enumerate}
  \def\labelenumi{(\alph{enumi})}
  \tightlist
  \item
    \(u(x_1, x_2) = x_1 + x_2\)
  \end{enumerate}
\item
  \begin{enumerate}
  \def\labelenumi{(\alph{enumi})}
  \setcounter{enumi}{1}
  \tightlist
  \item
    \(u(x_1, x_2) = \min(x_1, x_2)\)
  \end{enumerate}
\item
  \begin{enumerate}
  \def\labelenumi{(\alph{enumi})}
  \setcounter{enumi}{2}
  \tightlist
  \item
    \(u(x_1, x_2) = x_1 \cdot x_2\)
  \end{enumerate}
\item
  \begin{enumerate}
  \def\labelenumi{(\alph{enumi})}
  \setcounter{enumi}{3}
  \tightlist
  \item
    \(u(x_1, x_2) = x_1^2 + x_2^2\)
  \end{enumerate}
\end{itemize}

??? success ``Resposta'' \textbf{(b)} Complementos perfeitos são
consumidos em proporções fixas; a utilidade é determinada pelo bem em
menor quantidade relativa, representada pela função \(\min\). As curvas
de indiferença são em L. A alternativa (a) representa substitutos
perfeitos; (c) representa preferências Cobb-Douglas; (d) representa
preferências côncavas (não convexas), onde o consumidor prefere extremos
a médias.

\end{infobox}

\begin{infobox}{6. Preferências homotéticas implicam que:}

\begin{itemize}
\tightlist
\item
  \begin{enumerate}
  \def\labelenumi{(\alph{enumi})}
  \tightlist
  \item
    A demanda pelo bem 1 é independente da renda
  \end{enumerate}
\item
  \begin{enumerate}
  \def\labelenumi{(\alph{enumi})}
  \setcounter{enumi}{1}
  \tightlist
  \item
    A participação de cada bem na despesa total é constante quando a
    renda varia
  \end{enumerate}
\item
  \begin{enumerate}
  \def\labelenumi{(\alph{enumi})}
  \setcounter{enumi}{2}
  \tightlist
  \item
    A TMS depende apenas da quantidade absoluta de \(x_1\)
  \end{enumerate}
\item
  \begin{enumerate}
  \def\labelenumi{(\alph{enumi})}
  \setcounter{enumi}{3}
  \tightlist
  \item
    As curvas de indiferença são translações verticais
  \end{enumerate}
\end{itemize}

??? success ``Resposta'' \textbf{(b)} Preferências homotéticas geram
curvas de Engel lineares e elasticidade-renda unitária para todos os
bens, o que implica participação constante na despesa (Proposição
§3.6.6). A alternativa (a) descreve preferências quase-lineares, não
homotéticas; (c) também é propriedade da quase-linear (TMS depende só de
\(x_1\)); (d) descreve curvas de indiferença quase-lineares. A
assinatura das homotéticas é que a TMS depende da \emph{razão}
\(x_1/x_2\), gerando expansão radial --- não vertical --- das curvas.

\end{infobox}

\begin{center}\rule{0.5\linewidth}{0.5pt}\end{center}

\section{Resumo do Capítulo}\label{resumo-do-capuxedtulo}

\begin{itemize}
\tightlist
\item
  A teoria do consumidor parte de \textbf{axiomas sobre preferências}
  --- completude, transitividade, continuidade e monotonicidade --- que
  estabelecem regras mínimas de coerência para ordenar cestas de
  consumo. Esses axiomas são o alicerce sobre o qual todo o resto se
  constrói: sem eles, não há função de utilidade, não há curvas de
  indiferença e não há otimização.
\item
  Sob esses axiomas, o \textbf{Teorema de Debreu} (1954) garante a
  existência de uma \textbf{função de utilidade} contínua que representa
  as preferências. Essa função é \textbf{ordinal}: apenas o ordenamento
  das cestas importa, não os valores numéricos em si. A ordinalidade
  implica que qualquer transformação monotônica crescente da função de
  utilidade representa as mesmas preferências (Seção 3.7).
\item
  As \textbf{curvas de indiferença} são curvas de nível da função de
  utilidade: cobrem todo o espaço de consumo, não se cruzam (pela
  transitividade), têm inclinação negativa (pela monotonicidade) e, sob
  convexidade estrita, são abauladas em direção à origem --- refletindo
  a ideia de que o consumidor valoriza a diversidade na composição de
  sua cesta.
\item
  A \textbf{taxa marginal de substituição (TMS)} mede a taxa de troca
  subjetiva entre bens ao longo da curva de indiferença e equivale à
  razão das utilidades marginais:
  \(\text{TMS}_{12} = \text{UMg}_1 / \text{UMg}_2\). A TMS decrescente
  --- que reflete a disposição cada vez menor do consumidor a abrir mão
  de um bem à medida que ele se torna mais escasso na cesta --- é
  matematicamente equivalente à convexidade estrita das preferências.
\item
  O capítulo apresenta as principais famílias de funções de utilidade
  --- \textbf{Cobb-Douglas}, \textbf{substitutos perfeitos},
  \textbf{complementos perfeitos}, \textbf{CES} e \textbf{quase-linear}
  --- cada uma com formato de curvas de indiferença e elasticidade de
  substituição distintos. A função CES unifica as três primeiras como
  casos especiais de um único parâmetro \(\rho\), enquanto a
  quase-linear se distingue por eliminar o efeito renda sobre um dos
  bens.
\item
  \textbf{Preferências homotéticas} (TMS depende apenas da razão
  \(x_1/x_2\)) geram curvas de Engel lineares, elasticidade-renda
  unitária e participação constante na despesa --- propriedades que
  permitem a agregação em um consumidor representativo. Preferências
  \textbf{quase-lineares}, por contraste, concentram todo o efeito renda
  em um único bem e garantem que as medidas de bem-estar (VC, VE,
  \(\Delta\)EC) coincidam.
\end{itemize}

\section{Conceitos-Chave}\label{conceitos-chave}

{\def\LTcaptype{none} % do not increment counter
\begin{longtable}[]{@{}
  >{\raggedright\arraybackslash}p{(\linewidth - 2\tabcolsep) * \real{0.4762}}
  >{\raggedright\arraybackslash}p{(\linewidth - 2\tabcolsep) * \real{0.5238}}@{}}
\toprule\noalign{}
\begin{minipage}[b]{\linewidth}\raggedright
Conceito
\end{minipage} & \begin{minipage}[b]{\linewidth}\raggedright
Definição
\end{minipage} \\
\midrule\noalign{}
\endhead
\bottomrule\noalign{}
\endlastfoot
Relação de preferência (\(\succsim\)) & Ordenamento sobre cestas de
consumo indicando que uma cesta é ``pelo menos tão boa quanto''
outra. \\
Completude & Axioma que exige que o consumidor consiga comparar
quaisquer duas cestas. \\
Transitividade & Axioma de consistência: se
\(\mathbf{x} \succsim \mathbf{y}\) e \(\mathbf{y} \succsim \mathbf{z}\),
então \(\mathbf{x} \succsim \mathbf{z}\). \\
Monotonicidade & ``Mais é melhor'': cestas com mais de pelo menos um bem
são estritamente preferidas. \\
Função de utilidade & Função \(u: X \to \mathbb{R}\) que representa
preferências:
\(\mathbf{x} \succsim \mathbf{y} \iff u(\mathbf{x}) \geq u(\mathbf{y})\).
É ordinal. \\
Curva de indiferença & Conjunto de cestas que proporcionam o mesmo nível
de utilidade; curva de nível de \(u\). \\
Taxa marginal de substituição (TMS) & Quantidade do bem 2 que o
consumidor abre mão por uma unidade adicional do bem 1, mantendo a
utilidade constante; igual a \(\text{UMg}_1/\text{UMg}_2\). \\
Elasticidade de substituição (\(\sigma\)) & Mede a facilidade com que o
consumidor substitui entre bens; varia de 0 (complementos perfeitos) a
\(\infty\) (substitutos perfeitos). \\
Preferências homotéticas & Preferências cuja TMS depende apenas da razão
\(x_1/x_2\); geram elasticidade-renda unitária e participação constante
na despesa. \\
Utilidade quase-linear & Função \(u = v(x_1) + x_2\) em que todo aumento
de renda vai para \(x_2\); elimina o efeito renda sobre \(x_1\) e iguala
as medidas de bem-estar (VC = VE = \(\Delta\)EC). \\
\end{longtable}
}

\textbf{Tabela 3.2 --- Conceitos-chave.}

\begin{center}\rule{0.5\linewidth}{0.5pt}\end{center}

\section{Exercícios}\label{exercuxedcios}

Os exercícios a seguir cobrem os principais tópicos do capítulo ---
axiomas, funções de utilidade, TMS, formas funcionais e transformações
monotônicas. Eles progridem em dificuldade: os primeiros requerem
cálculos diretos de TMS para funções específicas, enquanto os últimos
exigem demonstrações e raciocínio mais abstrato sobre propriedades das
preferências. As soluções detalhadas estão disponíveis na seção de
soluções.

\textbf{Exercício 3.1.} Considere um consumidor com preferências sobre
dois bens (\(x_1, x_2\)) representadas pela função de utilidade
\(u(x_1, x_2) = x_1^{1/3} x_2^{2/3}\).

\begin{enumerate}
\def\labelenumi{(\alph{enumi})}
\item
  Calcule a TMS\(_{12}\).
\item
  Qual o valor da TMS no ponto \((x_1, x_2) = (4, 8)\)? Interprete
  economicamente.
\item
  Mostre que a TMS é decrescente ao longo de uma curva de indiferença
  (ou seja, \(\partial \text{TMS}_{12} / \partial x_1 < 0\) ao longo de
  \(u = \bar{u}\)).
\end{enumerate}

→ Ver solução (Cap. 3)

\begin{center}\rule{0.5\linewidth}{0.5pt}\end{center}

\textbf{Exercício 3.2.} Mostre que a função de utilidade
\(u(x_1, x_2) = \ln x_1 + 2 \ln x_2\) representa as mesmas preferências
que \(v(x_1, x_2) = x_1 \cdot x_2^2\). Verifique que ambas produzem a
mesma TMS.

→ Ver solução (Cap. 3)

\begin{center}\rule{0.5\linewidth}{0.5pt}\end{center}

\textbf{Exercício 3.3.} Um consumidor tem preferências do tipo CES com
\(\rho = -1\):

{[} u(x\_1, x\_2) = \left(x\_1\^{}\{-1\} +
x\_2\textsuperscript{\{-1\}\right)}\{-1\}. {]}

\begin{enumerate}
\def\labelenumi{(\alph{enumi})}
\item
  Calcule a elasticidade de substituição.
\item
  Derive a TMS\(_{12}\).
\item
  Esboce as curvas de indiferença. Elas estão mais próximas do caso
  Cobb-Douglas ou do caso de complementos perfeitos? Justifique.
\end{enumerate}

→ Ver solução (Cap. 3)

\begin{center}\rule{0.5\linewidth}{0.5pt}\end{center}

\textbf{Exercício 3.4.} Considere preferências quase-lineares
\(u(x_1, x_2) = \sqrt{x_1} + x_2\).

\begin{enumerate}
\def\labelenumi{(\alph{enumi})}
\item
  Calcule a TMS\(_{12}\) e mostre que ela depende apenas de \(x_1\).
\item
  Desenhe duas curvas de indiferença e mostre que elas são translações
  verticais uma da outra.
\item
  Explique por que, nesse caso, a demanda pelo bem 1 é independente da
  renda (para soluções interiores).
\end{enumerate}

→ Ver solução (Cap. 3)

\begin{center}\rule{0.5\linewidth}{0.5pt}\end{center}

\textbf{Exercício 3.5.} Um economista propõe representar as preferências
de um consumidor pela função \(u(x_1, x_2) = x_1^2 + x_2^2\).

\begin{enumerate}
\def\labelenumi{(\alph{enumi})}
\item
  As curvas de indiferença dessa função são convexas em relação à
  origem? Justifique.
\item
  A TMS é decrescente ao longo de uma curva de indiferença?
\item
  Essa função satisfaz o axioma de preferências estritamente convexas?
  Que problema isso gera para a existência de soluções interiores no
  problema de otimização do consumidor?
\end{enumerate}

→ Ver solução (Cap. 3)

\begin{center}\rule{0.5\linewidth}{0.5pt}\end{center}

\textbf{Exercício 3.6.} Considere a função de utilidade CES generalizada
com pesos:

{[} u(x\_1, x\_2) = \left(\alpha , x\_1\^{}\{\rho\} + (1-\alpha) ,
x\_2\textsuperscript{\{\rho\}\right)}\{1/\rho\}, \quad 0 \textless{}
\alpha \textless{} 1, ; \rho \textless{} 1, ; \rho \neq 0. {]}

\begin{enumerate}
\def\labelenumi{(\alph{enumi})}
\item
  Derive a TMS\(_{12}\) e mostre que ela depende de \(\alpha\), \(\rho\)
  e da razão \(x_1/x_2\).
\item
  Mostre que, no caso \(\alpha = 1/2\), a TMS no ponto \(x_1 = x_2\) é
  igual a 1, independentemente de \(\rho\).
\item
  Interprete economicamente o papel do parâmetro \(\alpha\): o que
  acontece com as curvas de indiferença quando \(\alpha\) aumenta?
\end{enumerate}

→ Ver solução (Cap. 3)

\begin{center}\rule{0.5\linewidth}{0.5pt}\end{center}

\textbf{Exercício 3.7.} Sejam duas funções de utilidade:

{[} u(x\_1, x\_2) = x\_1\^{}\{1/2\} , x\_2\^{}\{1/2\}, \qquad v(x\_1,
x\_2) = -\frac{1}{x_1 \, x_2}. {]}

\begin{enumerate}
\def\labelenumi{(\alph{enumi})}
\item
  Mostre que \(v\) é uma transformação monotônica de \(u\) e identifique
  a função \(f\) tal que \(v = f(u)\).
\item
  Calcule a TMS\(_{12}\) usando \(u\) e usando \(v\). Verifique que
  ambas coincidem.
\item
  Compare as utilidades marginais de \(u\) e de \(v\). Embora ambas
  representem as mesmas preferências, os valores das utilidades
  marginais coincidem? Explique por que a utilidade marginal não possui
  significado econômico isolado.
\end{enumerate}

→ Ver solução (Cap. 3)

\begin{center}\rule{0.5\linewidth}{0.5pt}\end{center}

\textbf{Exercício 3.8.} Um consumidor tem preferências representadas por
\(u(x_1, x_2) = \min\{2x_1, x_2\}\).

\begin{enumerate}
\def\labelenumi{(\alph{enumi})}
\item
  Em que proporção o consumidor deseja consumir os dois bens? Justifique
  a partir da função de utilidade.
\item
  Suponha que o consumidor possui renda \(m = 120\) e enfrenta preços
  \(p_1 = 10\) e \(p_2 = 5\). Encontre a cesta ótima.
\item
  Se o preço do bem 1 dobrar para \(p_1 = 20\), qual será a nova cesta
  ótima? Calcule a variação percentual no consumo de cada bem e
  interprete.
\end{enumerate}

→ Ver solução (Cap. 3)

\begin{center}\rule{0.5\linewidth}{0.5pt}\end{center}

\textbf{Exercício 3.9.} Demonstre que, para preferências homotéticas, a
TMS depende apenas da razão \(x_1/x_2\).

\emph{Dica:} Seja \(u(x_1, x_2) = g(h(x_1, x_2))\), onde \(g' > 0\) e
\(h\) é homogênea de grau 1. Use o Teorema de Euler para mostrar que
\(\text{TMS}_{12} = h_1/h_2\), e depois explore a homogeneidade de grau
zero das derivadas de uma função homogênea de grau 1.

→ Ver solução (Cap. 3)

\begin{center}\rule{0.5\linewidth}{0.5pt}\end{center}

\textbf{Exercício 3.10.} A Pesquisa de Orçamentos Familiares (POF
2017-2018) mostra que famílias com renda mensal de R\$ 1.908 gastam 22\%
da renda em alimentação, enquanto famílias com renda de R\$ 23.850
gastam 7,6\%.

\begin{enumerate}
\def\labelenumi{(\alph{enumi})}
\tightlist
\item
  Calcule a elasticidade-renda aproximada da alimentação entre essas
  duas faixas de renda, usando a fórmula de arco-elasticidade:
\end{enumerate}

{[} \varepsilon \approx \frac{\Delta Q / \bar{Q}}{\Delta I / \bar{I}}
{]}

onde \(Q = s \cdot I\) é o gasto com alimentação (\(s\) = participação),
e \(\bar{Q}\) e \(\bar{I}\) são as médias dos dois pontos.

\begin{enumerate}
\def\labelenumi{(\alph{enumi})}
\setcounter{enumi}{1}
\item
  A alimentação é um bem normal ou inferior? É um bem de luxo ou de
  necessidade? Justifique com base na elasticidade calculada.
\item
  Esse padrão é compatível com preferências homotéticas? E com
  preferências Cobb-Douglas? Explique, relacionando com as propriedades
  da Seção 3.6.6.
\end{enumerate}

→ Ver solução (Cap. 3)

\begin{center}\rule{0.5\linewidth}{0.5pt}\end{center}

\section{Vem, ANPEC!}\label{vem-anpec}

\begin{infobox}{ANPEC 2019 — Questão 01}

Com relação às preferências do consumidor, indique quais das afirmações
a seguir são verdadeiras e quais são falsas:

{\def\LTcaptype{none} % do not increment counter
\begin{longtable}[]{@{}
  >{\raggedright\arraybackslash}p{(\linewidth - 2\tabcolsep) * \real{0.3529}}
  >{\raggedright\arraybackslash}p{(\linewidth - 2\tabcolsep) * \real{0.6471}}@{}}
\toprule\noalign{}
\begin{minipage}[b]{\linewidth}\raggedright
Item
\end{minipage} & \begin{minipage}[b]{\linewidth}\raggedright
Afirmação
\end{minipage} \\
\midrule\noalign{}
\endhead
\bottomrule\noalign{}
\endlastfoot
0 & Sendo \(U(x, y)\) a função de utilidade em dois bens \(x\) e \(y\),
\(U(x, y) = \ln x \cdot \ln y\) representa uma função de utilidade
quase-linear. \\
1 & Podemos sempre extrair a transformação monotônica da função de
utilidade do tipo Cobb-Douglas. \\
2 & Uma função de utilidade do tipo \(U(x, y) = (x + y)^{0,5}\) implica
que \(x\) e \(y\) são bens substitutos perfeitos. \\
3 & Uma função de utilidade do tipo \(U(x, y) = x + y\) implica que
\(x\) e \(y\) são bens complementares perfeitos. \\
4 & \(f(U) = U^2\) é uma transformação monotônica apenas para \(U\)
positivo. \\
\end{longtable}
}

??? success ``Gabarito'' \textbf{Respostas: F V V F V}

\begin{verbatim}
**Justificativa por item:**

- **Item 0 — F:** A função $U(x,y) = \ln x \cdot \ln y$ não é quase-linear. Uma função quase-linear tem a forma $v(x_1) + x_2$, isto é, é linear em um dos bens e não linear no outro. O produto de logaritmos não se enquadra nessa estrutura.
- **Item 1 — V:** Da Cobb-Douglas $u = x^a y^b$, podemos aplicar a transformação monotônica $f(u) = \ln(u)$, obtendo $\hat{u} = a\ln x + b\ln y$. Como $f$ é estritamente crescente para $u > 0$ (e a Cobb-Douglas é positiva no interior do consumo), a transformação é sempre válida.
- **Item 2 — V:** A função $U(x,y) = (x+y)^{0,5}$ é uma transformação monotônica de $V(x,y) = x + y$, via $f(V) = V^{0,5}$, que é estritamente crescente para $V > 0$. Como $V = x + y$ representa substitutos perfeitos (TMS constante igual a 1), $U$ também representa substitutos perfeitos (Proposição 3.3).
- **Item 3 — F:** A função $U(x,y) = x + y$ é a representação clássica de substitutos perfeitos, com curvas de indiferença retas e TMS constante igual a 1. Complementares perfeitos seriam representados por $\min\{x, y\}$.
- **Item 4 — V:** A função $f(U) = U^2$ tem derivada $f'(U) = 2U$. Para $U > 0$, $f'(U) > 0$ e a transformação é estritamente crescente (monotônica). Para $U < 0$, $f'(U) < 0$ e a função é decrescente, invertendo o ordenamento — logo, não é uma transformação monotônica nesse domínio.
\end{verbatim}

\end{infobox}

\begin{infobox}{ANPEC 2021 — Questão 01}

Seja um consumidor com função de utilidade dada por \(U = X^2 + Y^2\),
em que \(X\) é a quantidade consumida de entradas de cinema e \(Y\) é a
quantidade consumida de pizzas. Com relação a este consumidor, verifique
quais das seguintes afirmações são verdadeiras e quais são falsas:

{\def\LTcaptype{none} % do not increment counter
\begin{longtable}[]{@{}
  >{\raggedright\arraybackslash}p{(\linewidth - 2\tabcolsep) * \real{0.3529}}
  >{\raggedright\arraybackslash}p{(\linewidth - 2\tabcolsep) * \real{0.6471}}@{}}
\toprule\noalign{}
\begin{minipage}[b]{\linewidth}\raggedright
Item
\end{minipage} & \begin{minipage}[b]{\linewidth}\raggedright
Afirmação
\end{minipage} \\
\midrule\noalign{}
\endhead
\bottomrule\noalign{}
\endlastfoot
0 & A taxa marginal de substituição deste consumidor é \(X/Y\). \\
1 & A cesta \((X = 2, Y = 1)\) e a cesta \((X = 1, Y = 2)\) se encontram
sobre a mesma curva de indiferença. \\
2 & As curvas de indiferença do consumidor são estritamente convexas
entre as cestas \((X = 2, Y = 1)\) e \((X = 1, Y = 2)\). \\
3 & \(X\) e \(Y\) são substitutos perfeitos. \\
4 & O bem \(Y\) é um mal. \\
\end{longtable}
}

??? success ``Gabarito'' \textbf{Respostas: V V F F F}

\begin{verbatim}
**Justificativa por item:**

- **Item 0 — V:** $\text{TMS} = \frac{\text{UMg}_X}{\text{UMg}_Y} = \frac{2X}{2Y} = \frac{X}{Y}$. Correto.
- **Item 1 — V:** $U(2,1) = 4 + 1 = 5$ e $U(1,2) = 1 + 4 = 5$. Ambas geram o mesmo nível de utilidade, logo pertencem à mesma curva de indiferença.
- **Item 2 — F:** As curvas de indiferença de $U = X^2 + Y^2 = c$ são arcos de circunferência — côncavas em relação à origem, não convexas. O conjunto superior $\{(X,Y): X^2+Y^2 \geq c\}$ não é convexo. Este é exatamente o caso discutido no Exercício 3.5 do capítulo.
- **Item 3 — F:** Substitutos perfeitos têm curvas de indiferença retas (TMS constante). Aqui, a TMS $= X/Y$ varia com a composição da cesta e as curvas são circulares, não retas.
- **Item 4 — F:** $\text{UMg}_Y = 2Y > 0$ para $Y > 0$, logo mais pizza aumenta a utilidade. O bem $Y$ não é um mal.
\end{verbatim}

\end{infobox}

\begin{infobox}{ANPEC 2021 — Questão 02}

Considere a Teoria da Utilidade para responder quais das afirmações a
seguir são verdadeiras e quais são falsas:

{\def\LTcaptype{none} % do not increment counter
\begin{longtable}[]{@{}
  >{\raggedright\arraybackslash}p{(\linewidth - 2\tabcolsep) * \real{0.3529}}
  >{\raggedright\arraybackslash}p{(\linewidth - 2\tabcolsep) * \real{0.6471}}@{}}
\toprule\noalign{}
\begin{minipage}[b]{\linewidth}\raggedright
Item
\end{minipage} & \begin{minipage}[b]{\linewidth}\raggedright
Afirmação
\end{minipage} \\
\midrule\noalign{}
\endhead
\bottomrule\noalign{}
\endlastfoot
0 & Todos os tipos de preferências podem ser representados pela função
de utilidade. \\
1 & Sejam dois bens, \(X\) e \(Y\). A uma função de utilidade dada por
\(U(X, Y) = XY\) corresponde uma curva de indiferença típica dada por
\(Y = cX\), em que \(c\) é uma constante. \\
2 & Se dois bens (\(A\) e \(B\)) forem substitutos perfeitos, pode-se,
em geral, representar sua função de utilidade na forma
\(U(A, B) = c_1 A + c_2 B\), em que \(c_1\) e \(c_2\) são constantes
positivas. \\
3 & A inclinação de uma curva de indiferença típica da função de
utilidade \(U(A, B) = c_1 A + c_2 B\), em que \(c_1\) e \(c_2\) são
constantes positivas, é \(-c_1/c_2\). \\
4 & A transformação monotônica de uma função de utilidade não altera a
taxa marginal de substituição (TMS), porque a TMS é medida ao longo de
uma curva de indiferença, e a utilidade permanece constante ao longo da
curva de indiferença. \\
\end{longtable}
}

??? success ``Gabarito'' \textbf{Respostas: F F V V V}

\begin{verbatim}
**Justificativa por item:**

- **Item 0 — F:** Nem todos os tipos de preferências admitem representação por função de utilidade. As preferências lexicográficas, por exemplo, são completas e transitivas mas violam o axioma de continuidade — e, conforme discutido na Seção 3.1, sem continuidade não se garante a existência de uma função de utilidade contínua (Teorema 3.1).
- **Item 1 — F:** Para $U(X,Y) = XY$, as curvas de indiferença são $XY = k$, ou seja, $Y = k/X$ — hipérboles retangulares, não retas pela origem. A equação $Y = cX$ descreveria retas, o que é incorreto.
- **Item 2 — V:** A função $U(A,B) = c_1 A + c_2 B$, com $c_1, c_2 > 0$, gera curvas de indiferença retas com TMS constante $= c_1/c_2$, que é a definição de substitutos perfeitos.
- **Item 3 — V:** A inclinação da curva de indiferença é $dB/dA|_{U=\bar{u}} = -c_1/c_2$. Correto, conforme a Seção 3.6.2.
- **Item 4 — V:** Uma transformação monotônica $\hat{u} = f(u)$ preserva as curvas de indiferença (mesmos conjuntos de nível, apenas com rótulos diferentes). Como a TMS depende apenas da inclinação da curva de indiferença, ela é invariante (Proposição 3.3 e Seção 3.7).
\end{verbatim}

\end{infobox}

\begin{infobox}{ANPEC 2023 — Questão 03}

Com relação à Teoria do Consumidor, julgue as afirmações abaixo:

{\def\LTcaptype{none} % do not increment counter
\begin{longtable}[]{@{}
  >{\raggedright\arraybackslash}p{(\linewidth - 2\tabcolsep) * \real{0.3529}}
  >{\raggedright\arraybackslash}p{(\linewidth - 2\tabcolsep) * \real{0.6471}}@{}}
\toprule\noalign{}
\begin{minipage}[b]{\linewidth}\raggedright
Item
\end{minipage} & \begin{minipage}[b]{\linewidth}\raggedright
Afirmação
\end{minipage} \\
\midrule\noalign{}
\endhead
\bottomrule\noalign{}
\endlastfoot
0 & Se \(U(X, Y) = \sqrt{X} + \sqrt{Y}\), então a elasticidade de
substituição é igual a 2. \\
1 & Se \(U(X, Y)\) é uma função de utilidade diferenciável homogênea de
grau \(k\), com \(U(X, Y) \neq 0\), e se \(\eta_X(X, Y)\) e
\(\eta_Y(X, Y)\) denotam a elasticidade de \(U\) relativamente a \(X\) e
\(Y\), respectivamente, então \(\eta_X(X, Y) + \eta_Y(X, Y) = k\). \\
2 & Considere o conjunto \(X = \bigcup_{n=1}^{\infty}\{1 - 1/n\}\) e
defina \(\succsim\) sobre \(X\) por: \(x \succsim y\) se, e somente se,
\(\lvert x - 1/2 \rvert \geq \lvert y - 1/2 \rvert\). Então \(\succsim\)
é contínua. \\
3 & Considere o conjunto \(X = [0,2] \times [0,2]\) e defina
\(\succsim\) sobre \(X\) por: \((x,y) \succsim (z,w)\) se, e somente se,
\(\lVert (x,y) - (1,1) \rVert \geq \lVert (z,w) - (1,1) \rVert\). Então
\((0,0)\) é um elemento maximal. \\
4 & A ordem lexicográfica é contínua. \\
\end{longtable}
}

??? success ``Gabarito'' \textbf{Respostas: V V F V F}

\begin{verbatim}
**Justificativa por item:**

- **Item 0 — V:** A função $U = X^{1/2} + Y^{1/2}$ é uma CES com $\rho = 1/2$. A elasticidade de substituição é $\sigma = \frac{1}{1-\rho} = \frac{1}{1-1/2} = 2$. Correto — ver Seção 3.6.4.
- **Item 1 — V:** Pelo teorema de Euler para funções homogêneas de grau $k$: $X \frac{\partial U}{\partial X} + Y \frac{\partial U}{\partial Y} = kU$. Dividindo ambos os lados por $U \neq 0$: $\eta_X + \eta_Y = k$, onde $\eta_i = \frac{x_i}{U}\frac{\partial U}{\partial x_i}$ é a elasticidade de $U$ em relação ao bem $i$.
- **Item 2 — F:** A relação $\succsim$ definida pela distância a $1/2$ não satisfaz a propriedade de continuidade nesse conjunto. O conjunto $X = \{0, 1/2, 2/3, 3/4, \ldots\}$ acumula em $1 \notin X$, e a "utilidade" $|x - 1/2|$ converge para $1/2$ ao longo da sequência $x_n = 1 - 1/n$ sem nunca atingir esse supremo em $X$, o que compromete o fechamento dos conjuntos de contorno superiores quando se considera a topologia induzida pela reta real.
- **Item 3 — V:** A relação ordena cestas pela distância à cesta $(1,1)$: quanto **maior** a distância, melhor. A cesta $(0,0)$ tem distância $\lVert(0,0)-(1,1)\rVert = \sqrt{2}$, que é a distância máxima possível em $[0,2] \times [0,2]$ a partir de $(1,1)$ (empatando com $(2,2)$, $(0,2)$ e $(2,0)$). Logo, $(0,0)$ é maximal.
- **Item 4 — F:** A ordem lexicográfica não é contínua — este é o exemplo clássico discutido na Seção 3.1. Os conjuntos de contorno superior não são fechados, o que impede a representação por função de utilidade contínua (Teorema 3.1).
\end{verbatim}

\end{infobox}

\chapter{Pesquisa em Ação}\label{pesquisa-em-auxe7uxe3o}

\begin{pesquisabox}{Falk, A.; Becker, A.; Dohmen, T.; Enke, B.; Huffman, D.; Sunde, U. (2018). [Global Evidence on Economic Preferences](https://doi.org/10.1093/qje/qjy013). *Quarterly Journal of Economics*, 133(4), 1645–1692.}

\textbf{Pergunta central:} Os axiomas de preferência apresentados na
Seção 3.1 são abstrações teóricas --- mas como as preferências reais dos
indivíduos variam entre países e dentro de cada sociedade? Existem
padrões sistemáticos que conectem preferências a características
demográficas, culturais e econômicas? Este artigo é o maior esforço
empírico já realizado para responder a essas perguntas.

\textbf{Método:} Falk e coautores construíram o \emph{Global Preference
Survey} (GPS), um instrumento de pesquisa experimentalmente validado,
aplicado a amostras representativas de 80.000 pessoas em 76 países ---
incluindo o Brasil. O GPS mede seis dimensões de preferências: paciência
(preferência temporal), disposição a assumir riscos, reciprocidade
positiva e negativa, altruísmo e confiança. A validação experimental foi
realizada comparando as respostas da pesquisa com escolhas em
experimentos com incentivos monetários reais.

\textbf{Resultado principal:} Há substancial heterogeneidade de
preferências entre países, mas a variação \emph{dentro} de cada país é
ainda maior. Globalmente, preferências variam sistematicamente com
idade, gênero e habilidade cognitiva --- porém essas relações são
parcialmente específicas de cada país. O Brasil apresenta níveis
intermediários de paciência e disposição ao risco na comparação
internacional, com heterogeneidade interna significativa --- consistente
com a elevada desigualdade socioeconômica do país.

\textbf{Por que isso importa:} O estudo demonstra que as preferências
--- fundamento de toda a teoria do consumidor --- não são uniformes nem
entre nem dentro de países. Diferenças de preferência ajudam a explicar
variações em poupança, investimento em educação e comportamento de
consumo, complementando os modelos baseados apenas em diferenças de
renda e preços.

\textbf{Relevância para o capítulo:} O artigo conecta diretamente os
axiomas da Seção 3.1 com evidências empíricas: se as preferências variam
sistematicamente entre indivíduos, as funções de utilidade que as
representam também diferem --- o que justifica a diversidade de formas
funcionais apresentadas na Seção 3.6. Além disso, a heterogeneidade de
preferências dentro do Brasil reforça a importância de modelos que
permitam diferenças individuais, como a análise por faixa de renda da
POF discutida no Box Brasil sobre Cobb-Douglas.

\end{pesquisabox}

\begin{pesquisabox}{Choi, S.; Kariv, S.; Müller, W.; Silverman, D. (2014). [Who Is (More) Rational?](https://doi.org/10.1257/aer.104.6.1518) *American Economic Review*, 104(6), 1518–1550.}

\textbf{Pergunta central:} Os axiomas de preferência --- especialmente a
transitividade e a completude --- são \emph{de fato} satisfeitos pelas
escolhas dos consumidores reais? E se houver variação na
``racionalidade'' das decisões, ela está correlacionada com resultados
econômicos importantes, como a acumulação de riqueza?

\textbf{Método:} Choi e coautores conduziram um experimento em larga
escala com uma amostra representativa da população holandesa (painel
CentERpanel). Cada participante tomou 25 decisões de alocação entre dois
bens sob restrições orçamentárias variadas. Os autores testaram se as
escolhas observadas satisfazem o GARP (\emph{Generalized Axiom of
Revealed Preference}) --- a condição necessária e suficiente para que os
dados sejam consistentes com a maximização de alguma função de utilidade
bem comportada (Teorema 3.1).

\textbf{Resultado principal:} A consistência com GARP varia
substancialmente entre indivíduos: enquanto muitos participantes fazem
escolhas quase perfeitamente racionais, outros violam sistematicamente
os axiomas. Crucialmente, a consistência com a maximização de utilidade
está fortemente correlacionada com a riqueza: um aumento de um
desvio-padrão na consistência está associado a 15-19\% a mais de riqueza
acumulada. Essa associação é robusta mesmo controlando para renda,
educação e outras variáveis.

\textbf{Por que isso importa:} O resultado sugere que a capacidade de
tomar decisões consistentes com uma função de utilidade --- ou seja, de
satisfazer os axiomas da Seção 3.1 na prática --- não é apenas uma
abstração teórica, mas uma habilidade com consequências econômicas reais
e mensuráveis.

\textbf{Relevância para o capítulo:} Este artigo testa empiricamente os
fundamentos do Capítulo 3. O GARP é a tradução operacional dos axiomas
de completude e transitividade para dados de consumo observados. Os
resultados mostram que, embora a maioria dos consumidores se comporte de
forma aproximadamente consistente com os axiomas, há variação
significativa --- o que justifica tanto o uso do arcabouço axiomático
como ponto de partida quanto a atenção a seus limites, discutidos na
observação sobre preferências lexicográficas (Seção 3.1).

\end{pesquisabox}

\begin{pesquisabox}{Atkin, D. (2013). [Trade, Tastes, and Nutrition in India](https://doi.org/10.1257/aer.103.5.1629). *American Economic Review*, 103(5), 1629–1663.}

\textbf{Pergunta central:} As preferências dos consumidores --- que a
Seção 3.1 trata como dadas e estáveis --- podem ser moldadas por forças
econômicas como a abertura comercial? E, se as preferências mudam, quais
são as consequências para o bem-estar, especialmente em dimensões como a
nutrição?

\textbf{Método:} Atkin explora a liberalização comercial na Índia a
partir de 1991 para investigar como a queda nos preços relativos de
certos alimentos afetou não apenas as quantidades consumidas, mas os
\emph{hábitos alimentares} de longo prazo. Utilizando dados de consumo
domiciliar de mais de 500.000 observações ao longo de duas décadas, o
autor estima um modelo estrutural de demanda que permite distinguir
entre movimentos \emph{ao longo} de curvas de indiferença (substituição
padrão) e \emph{deslocamentos} das próprias curvas (mudança de
preferências --- formação de hábitos).

\textbf{Resultado principal:} A abertura comercial reduziu os preços de
calorias, mas parte dos ganhos calóricos potenciais foi dissipada por
mudanças nos gostos: consumidores expostos a preços mais baixos de
certos alimentos desenvolveram preferências por esses alimentos que
persistiram mesmo após os preços retornarem. O efeito-hábito reduziu em
até 10\% os ganhos nutricionais da liberalização, especialmente entre os
mais pobres.

\textbf{Relevância para o capítulo:} O artigo questiona a hipótese de
preferências estáveis que sustenta o arcabouço axiomático da Seção 3.1.
Se as preferências são endógenas aos preços (via formação de hábitos), a
função de utilidade \(u(\mathbf{x})\) não é fixa --- ela se desloca com
o histórico de consumo. Isso tem implicações diretas para a
interpretação das curvas de indiferença e da TMS: a ``taxa de troca
subjetiva'' de hoje pode diferir da de amanhã, não por mudança de
preços, mas por mudança nos próprios gostos. O resultado é especialmente
relevante para países em desenvolvimento como o Brasil, onde mudanças
estruturais rápidas (como o Plano Real, discutido no Box Brasil §3.7)
alteram simultaneamente preços e hábitos.

\end{pesquisabox}

\begin{pesquisabox}{Aguiar, M.; Bils, M. (2015). [Has Consumption Inequality Mirrored Income Inequality?](https://doi.org/10.1257/aer.20120599) *American Economic Review*, 105(9), 2725–2756.}

\textbf{Pergunta central:} A desigualdade de consumo acompanhou o
aumento da desigualdade de renda nos Estados Unidos nas últimas décadas?
Se as preferências são homotéticas (Seção 3.6.6), a resposta deveria ser
sim --- mas os dados contam uma história mais complexa.

\textbf{Método:} Aguiar e Bils utilizam dados da \emph{Consumer
Expenditure Survey} (CEX) e da \emph{Nielsen Homescan} para construir
medidas de desigualdade de consumo que corrigem problemas de mensuração
nos dados de pesquisas domiciliares. A estratégia empírica central
explora a Lei de Engel: bens com alta elasticidade-renda (bens de luxo)
servem como ``indicadores'' de renda permanente, permitindo inferir
desigualdade de consumo a partir da composição das cestas --- e não
apenas dos valores totais declarados, que sofrem viés de subdeclaração.

\textbf{Resultado principal:} Contrariamente a estudos anteriores que
sugeriam estabilidade, a desigualdade de consumo \emph{cresceu
significativamente} entre 1980 e 2010, acompanhando de perto a
desigualdade de renda. A chave metodológica foi usar a composição das
cestas (participações engelianas) em vez de valores totais. As famílias
mais ricas deslocaram seu consumo fortemente para bens de luxo
(educação, lazer, serviços), enquanto as mais pobres mantiveram alta
participação em necessidades (alimentação, moradia).

\textbf{Relevância para o capítulo:} O artigo mobiliza diretamente dois
conceitos centrais do capítulo: a \textbf{Lei de Engel} (discutida no
Box Brasil sobre homoteticidade, §3.6.6) e a distinção entre
\textbf{preferências homotéticas e não homotéticas}. Se as preferências
fossem homotéticas, a composição das cestas seria idêntica entre ricos e
pobres --- e a estratégia de identificação dos autores não funcionaria.
O fato de que a estratégia é bem-sucedida é, em si, evidência empírica
contra a homoteticidade. O resultado reforça a importância de modelos
não homotéticos (como o AIDS) para análises distributivas --- inclusive
no contexto brasileiro, onde a POF revela padrões engelianos ainda mais
pronunciados.

\end{pesquisabox}

\begin{pesquisabox}{Thomas, D.; Strauss, J.; Henriques, M. H. (1991). [How Does Mother's Education Affect Child Height?](https://doi.org/10.2307/145920) *Journal of Human Resources*, 26(2), 183–211.}

\textbf{Pergunta central:} A educação materna afeta a nutrição infantil
por meio de mudanças nas \emph{preferências} (função de utilidade) ou
apenas por mudanças na \emph{restrição orçamentária} (renda e preços)?
Utilizando dados brasileiros, este artigo seminal testa se mães mais
educadas fazem escolhas alimentares qualitativamente diferentes para
seus filhos --- e não apenas quantitativamente maiores.

\textbf{Método:} Thomas, Strauss e Henriques utilizam microdados do
\textbf{ENDEF} (\emph{Estudo Nacional da Despesa Familiar}) de 1974-75,
com informações antropométricas de crianças e características
socioeconômicas de mais de 50.000 domicílios brasileiros. O modelo
estima o efeito da educação materna sobre a estatura infantil
(\emph{height-for-age}), controlando para renda familiar, preços locais,
acesso a serviços de saúde e região geográfica. A estratégia de
identificação compara crianças dentro de faixas de renda similares,
isolando o efeito da educação que opera \emph{via preferências} (escolha
de alimentos mais nutritivos, práticas de higiene) do efeito que opera
\emph{via renda}.

\textbf{Resultado principal:} A educação materna tem um efeito forte e
robusto sobre a nutrição infantil, mesmo após controlar para renda: um
ano adicional de escolaridade da mãe aumenta a estatura da criança em
0,3-0,5 cm. Crucialmente, o efeito é muito maior para a educação da
\emph{mãe} do que para a do \emph{pai}, sugerindo que o canal relevante
não é apenas renda (ambos os pais contribuem para a renda), mas a forma
como a mãe \emph{aloca} os recursos do domicílio entre as diferentes
categorias de consumo --- alimentação nutritiva, cuidados de saúde,
saneamento.

\textbf{Relevância para o capítulo:} O artigo ilustra, com dados
brasileiros, que a \emph{forma} da função de utilidade --- não apenas a
restrição orçamentária --- é determinante para os resultados de consumo.
Nos termos do capítulo, mães com diferentes níveis de educação possuem
\emph{curvas de indiferença} com formatos distintos: mães mais educadas
atribuem TMS mais alta a alimentos nutritivos em relação a outros bens,
mesmo quando enfrentam os mesmos preços e a mesma renda. Isso reforça a
mensagem do artigo de Falk et al.~(2018) sobre a heterogeneidade de
preferências e conecta-se à discussão sobre a Lei de Engel (Box Brasil
§3.6.6): a composição da cesta depende não apenas de quanto se ganha,
mas de \emph{quem} decide como gastar.

\end{pesquisabox}

\section{Referências do Capítulo}\label{referuxeancias-do-capuxedtulo}

\begin{itemize}
\tightlist
\item
  Aguiar, Mark, e Mark Bils. 2015.
  ``\href{https://doi.org/10.1257/aer.20120599}{Has Consumption
  Inequality Mirrored Income Inequality?}'' \emph{American Economic
  Review} 105 (9): 2725--2756.
\item
  Arrow, Kenneth J., Hollis B. Chenery, Bagicha S. Minhas, e Robert M.
  Solow. 1961. ``\href{https://doi.org/10.2307/1927286}{Capital-Labor
  Substitution and Economic Efficiency.}'' \emph{Review of Economics and
  Statistics} 43 (3): 225--250.
\item
  Atkin, David. 2013.
  ``\href{https://doi.org/10.1257/aer.103.5.1629}{Trade, Tastes, and
  Nutrition in India.}'' \emph{American Economic Review} 103 (5):
  1629--1663.
\item
  Barros, Ricardo Paes de, Miguel Nathan Foguel, e Gabriel Ulyssea,
  orgs. 2007.
  \href{https://repositorio.ipea.gov.br/handle/11058/3249}{\emph{Desigualdade
  de renda no Brasil: uma análise da queda recente}}. Brasília: IPEA.
\item
  Besanko, David, e Ronald R. Braeutigam. 2014.
  \href{https://books.google.com.br/books?id=BeoengEACAAJ}{\emph{Microeconomics}}.
  5ª ed.~Hoboken: Wiley. Capítulos 3--4.
\item
  Choi, Syngjoo, Shachar Kariv, Wieland Müller, e Dan Silverman. 2014.
  ``\href{https://doi.org/10.1257/aer.104.6.1518}{Who Is (More)
  Rational?}'' \emph{American Economic Review} 104 (6): 1518--1550.
\item
  Debreu, Gerard. 1954.
  ``\href{https://www.cambridge.org/core/books/abs/mathematical-economics/representation-of-a-preference-ordering-by-a-numerical-function/009AEFFF2E07235C5BBFFEF226353FE8}{Representation
  of a preference ordering by a numerical function}.'' In \emph{Decision
  Processes}, editado por Robert M. Thrall, Clyde H. Coombs, e Robert L.
  Davis, 159--165. New York: Wiley.
\item
  Falk, Armin, Anke Becker, Thomas Dohmen, Benjamin Enke, David Huffman,
  e Uwe Sunde. 2018. ``\href{https://doi.org/10.1093/qje/qjy013}{Global
  Evidence on Economic Preferences.}'' \emph{Quarterly Journal of
  Economics} 133 (4): 1645--1692.
\item
  IBGE --- Instituto Brasileiro de Geografia e Estatística. 2019.
  \href{https://www.ibge.gov.br/estatisticas/sociais/rendimento-despesa-e-consumo/9050-pesquisa-de-orcamentos-familiares.html}{\emph{Pesquisa
  de Orçamentos Familiares (POF) 2017--2018: Primeiros Resultados}}. Rio
  de Janeiro: IBGE.
\item
  Mas-Colell, Andreu, Michael D. Whinston, e Jerry R. Green. 1995.
  \href{https://books.google.com/books/about/Microeconomic_Theory.html?id=KGtegVXqD8wC}{\emph{Microeconomic
  Theory}}. New York: Oxford University Press. Capítulo 3.
\item
  Thomas, Duncan, John Strauss, e Maria-Helena Henriques. 1991.
  ``\href{https://doi.org/10.2307/145920}{How Does Mother's Education
  Affect Child Height?}'' \emph{Journal of Human Resources} 26 (2):
  183--211.
\item
  Nicholson, Walter, e Christopher M. Snyder. 2017.
  \href{https://books.google.com/books/about/Microeconomic_Theory_Basic_Principles_an.html?id=YdkhCwAAQBAJ}{\emph{Microeconomic
  Theory}}. 12ª ed.~Boston: Cengage Learning. Capítulo 3.
\item
  Vasconcellos, Marco Antonio Sandoval de, e Manuel Enriquez Garcia.
  2014.
  \href{https://books.google.com.br/books?id=LSb8oAEACAAJ}{\emph{Fundamentos
  de economia}}. 5ª ed.~São Paulo: Saraiva.
\item
  Varian, Hal R. 2015.
  \href{https://books.google.com/books/about/Intermediate_Microeconomics_with_Calculu.html?id=9mabDwAAQBAJ}{\emph{Microeconomia:
  uma abordagem moderna}}. 9ª ed.~Rio de Janeiro: Elsevier. Capítulos
  3--5.
\end{itemize}

\newpage

\chapter{Capítulo 15 --- Sozinho, Feliz e Cobrando
Caro}\label{capuxedtulo-15-sozinho-feliz-e-cobrando-caro}

Se no Capítulo 14 vivíamos no paraíso da concorrência perfeita ---
milhares de firmas, ninguém mandando em nada ---, agora entramos no lado
sombrio do mercado: um único vendedor, zero concorrência e um sorriso no
rosto de quem define o preço. O monopólio representa a antítese da
concorrência perfeita. Enquanto no modelo competitivo cada firma é uma
tomadora de preço, incapaz de influenciar individualmente as condições
de mercado, o monopolista é o único ofertante e, portanto, enfrenta toda
a curva de demanda do mercado. Essa posição privilegiada lhe confere
\textbf{poder de mercado} --- a capacidade de fixar preços acima do
custo marginal e obter lucros econômicos persistentes.\footnote{O
  monopolista é como o Black Knight de \emph{Monty Python and the Holy
  Grail}: ``'Tis but a scratch!'' --- ele insiste que está tudo bem
  enquanto o CADE corta suas margens. A diferença é que, no caso do
  monopólio, o cavaleiro negro realmente tem uma espada grande o
  suficiente para fazer estrago. E ao contrário do sketch, ele não perde
  os braços: a barreira à entrada os protege.}

A análise do monopólio não é meramente teórica: monopólios naturais
regulados dominam setores fundamentais da economia brasileira ---
energia elétrica, saneamento, telecomunicações fixas ---, e o exercício
de poder de mercado é uma preocupação central da política antitruste em
todo o mundo. A fusão que criou a Ambev, os contratos de concessão da
ANEEL, o teto de preços da ANATEL --- todos esses casos envolvem, em sua
essência, a economia do monopólio apresentada neste capítulo.
Compreender como o monopolista fixa preços, por que essa fixação gera
ineficiência e quais instrumentos regulatórios podem mitigá-la é,
portanto, indispensável tanto para o economista teórico quanto para o
formulador de políticas públicas.

Este capítulo examina as causas do monopólio, a lógica de sua
maximização de lucro, as perdas de eficiência decorrentes do poder de
mercado, as estratégias de discriminação de preços e os mecanismos de
regulação. Ao longo da exposição, utilizaremos intensivamente os
conceitos de receita marginal, elasticidade-preço e excedente
desenvolvidos nos capítulos anteriores --- em particular, a relação
entre preço, custo marginal e elasticidade que constituirá o fio
condutor de toda a análise. Nos Capítulos 16a e 16b, generalizaremos
essa estrutura para mercados com poucos competidores --- os oligopólios
---, onde a interação estratégica entre firmas substitui o isolamento do
monopolista.\footnote{``Nobody expects the Spanish Inquisition!'' E
  ninguém espera que um capítulo sobre monopólio termine com a
  preparação para o oligopólio. Mas como os Pythons, nós vamos aos
  próximos capítulos (16a--16b) com três armas: medo, surpresa e uma
  devoção quase fanática pela interação estratégica.}

\begin{center}\rule{0.5\linewidth}{0.5pt}\end{center}

\section{Roteiro do Capítulo}\label{roteiro-do-capuxedtulo}

{\def\LTcaptype{none} % do not increment counter
\begin{longtable}[]{@{}
  >{\centering\arraybackslash}p{(\linewidth - 6\tabcolsep) * \real{0.3125}}
  >{\raggedright\arraybackslash}p{(\linewidth - 6\tabcolsep) * \real{0.1875}}
  >{\raggedright\arraybackslash}p{(\linewidth - 6\tabcolsep) * \real{0.1875}}
  >{\centering\arraybackslash}p{(\linewidth - 6\tabcolsep) * \real{0.3125}}@{}}
\toprule\noalign{}
\begin{minipage}[b]{\linewidth}\centering
Seção
\end{minipage} & \begin{minipage}[b]{\linewidth}\raggedright
Pergunta-guia
\end{minipage} & \begin{minipage}[b]{\linewidth}\raggedright
O que você vai aprender
\end{minipage} & \begin{minipage}[b]{\linewidth}\centering
Página
\end{minipage} \\
\midrule\noalign{}
\endhead
\bottomrule\noalign{}
\endlastfoot
15.1 & O que impede a concorrência de derrubar o monopolista? &
Barreiras à entrada: legais, tecnológicas, estratégicas & Barreiras \\
15.2 & Como o monopolista escolhe preço e quantidade? & Maximização de
lucro, RMg = CMg & Maximização \\
15.3 & Como medir o poder de mercado com uma fórmula? & Índice de Lerner
e elasticidade & Lerner \\
15.4 & Quanto a sociedade perde com o monopólio? & Peso morto,
ineficiência alocativa & Ineficiência \\
15.5 & O que acontece com o monopólio quando a demanda ou os custos
mudam? & Estática comparativa, pass-through & Estática comp. \\
15.6 & O monopolista escolhe qualidade alta ou baixa? & Qualidade
endógena sob monopólio & Qualidade \\
15.7 & Por que a Netflix cobra preços diferentes em países diferentes? &
Discriminação de preços: 1º, 2º e 3º graus & Discriminação \\
15.8 & Por que o parque cobra entrada e depois cobra por brinquedo? &
Tarifa em duas partes, bundling & Tarifas \\
15.9 & Como o regulador doma o monopolista sem destruir o incentivo a
investir? & Regulação: preço teto, taxa de retorno, price cap &
Regulação \\
15.10 & O monopólio pode ser bom para a inovação? & Schumpeter,
destruição criativa, patentes & Visão dinâmica \\
\textbf{Exercícios} & Teste, pratique, resolva & Revisão, exercícios,
ANPEC & Exercícios \\
\textbf{Pesquisa} & O que a pesquisa recente diz? & Artigos seminais e
fronteira empírica & Pesquisa \\
\end{longtable}
}

\begin{center}\rule{0.5\linewidth}{0.5pt}\end{center}

\begin{classroombox}{Atividade de Sala — O Leilão do Monopólio: Poder de Mercado, Discriminação e Regulação}

\textbf{Formato:} Simulação de mercado + debrief analítico (45--55 min)

\textbf{Objetivo:} Vivenciar como o poder de mercado afeta preços,
quantidades e bem-estar; comparar preço único com discriminação de
preços; experimentar regulação como resposta.

\textbf{Preparação (professor):}

\begin{itemize}
\tightlist
\item
  Divida a turma em grupos de 5--6 alunos. Em cada grupo: 1 aluno é o
  monopolista, os demais são consumidores.
\item
  Distribua fichas de \textbf{valoração} aos consumidores: cada
  consumidor recebe uma carta com sua disposição a pagar (valores
  diferentes: R\$ 10, R\$ 8, R\$ 6, R\$ 4, R\$ 2).
\item
  O monopolista tem custo marginal constante de R\$ 2 por unidade.
\item
  Prepare uma planilha de registro no quadro (preço, quantidade, lucro,
  EC, PPM).
\end{itemize}

\textbf{Rodada 1 --- Preço único (10 min):}

\begin{enumerate}
\def\labelenumi{\arabic{enumi}.}
\tightlist
\item
  O monopolista anuncia um preço único. Consumidores cuja valoração ≥
  preço compram; os demais ficam de fora.
\item
  O monopolista pode mudar o preço em até 3 tentativas (tâtonnement),
  observando quantos compram.
\item
  Registre o preço final, a quantidade vendida, o lucro e o excedente do
  consumidor.
\item
  Calcule a perda de peso morto (consumidores excluídos cujo valor
  \textgreater{} CMg).
\end{enumerate}

\textbf{Rodada 2 --- Discriminação perfeita (10 min):}

\begin{enumerate}
\def\labelenumi{\arabic{enumi}.}
\tightlist
\item
  O monopolista agora pode fazer ofertas individuais secretas a cada
  consumidor (bilhetes escritos).
\item
  Cada consumidor aceita ou rejeita a oferta.
\item
  Registre os resultados: quantos compraram? A que preço? Qual o lucro?
  Qual o EC?
\end{enumerate}

\textbf{Rodada 3 --- Regulação (10 min):}

\begin{enumerate}
\def\labelenumi{\arabic{enumi}.}
\tightlist
\item
  Um aluno assume o papel de regulador (ANEEL/ANATEL) e impõe
  \(p = CMg = 2\).
\item
  O monopolista reclama: ``Não cubro meus custos fixos!'' (professor
  adiciona CF = R\$ 8).
\item
  O regulador tenta a alternativa \(p = CMe\). Calcule e compare.
\end{enumerate}

\textbf{Debrief (15--20 min):}

\begin{itemize}
\tightlist
\item
  Compare os três regimes: preço único, discriminação perfeita e
  regulação. Quem ganhou e quem perdeu em cada caso?
\item
  O monopolista descobriu sozinho o preço ótimo (RMg = CMg) ou precisou
  de tentativa e erro?
\item
  Na Rodada 2, o EC caiu a zero --- isso é ``justo''? Discuta eficiência
  vs.~equidade.
\item
  Na Rodada 3, como o regulador resolveu o dilema do prejuízo? Conecte
  com o modelo da ANEEL.
\item
  Referência Monty Python: ``Bring me a shrubbery!'' --- os cavaleiros
  Ni são monopolistas perfeitos: demanda inelástica (sem shrubbery, não
  te deixo passar) e preço acima do custo marginal. A regulação equivale
  a chamar o Rei Arthur para negociar.
\end{itemize}

\textbf{Variante avançada:} Na Rodada 2, em vez de discriminação
perfeita, use discriminação de 3º grau --- divida os consumidores em
dois grupos (estudantes e executivos) e peça ao monopolista para definir
preços diferentes para cada grupo.

\textbf{Referência:} Holt, C.A. (2007). \emph{Markets, Games, and
Strategic Behavior}, Cap. 24 (Monopoly).

\end{classroombox}

\chapter{15.1--15.2 Barreiras à Entrada e Maximização de
Lucro}\label{barreiras-uxe0-entrada-e-maximizauxe7uxe3o-de-lucro}

\section{15.1 Muralhas, Fossos e Portões Levadiços: Barreiras à
Entrada}\label{151}

Se o monopólio é tão lucrativo, por que todo mundo não entra nesse
mercado e acaba com a festa? Imagine uma fortaleza medieval: o castelo
só sobrevive porque tem muralhas, fossos e portões levadiços. O
monopólio funciona da mesma forma --- ele só se sustenta porque existem
\textbf{barreiras à entrada}. Sem elas, o monopólio seria uma posição
transitória, rapidamente contestada por rivais ávidos por lucro --- como
veremos na Seção 15.10 ao discutir mercados contestáveis. É precisamente
a existência de obstáculos à entrada que permite ao monopolista
sustentar sua posição privilegiada no longo prazo.

Historicamente, a reflexão sobre barreiras à entrada remonta pelo menos
a Joe Bain (1956), que as classificou em três categorias --- vantagens
absolutas de custo, economias de escala e diferenciação de produto. A
taxonomia moderna, apresentada a seguir, preserva a essência dessa
classificação enquanto incorpora contribuições posteriores de Stigler
(1968) e da teoria dos jogos.

\subsection{Barreiras legais}\label{barreiras-legais}

O Estado pode conceder a uma firma o direito exclusivo de operar em um
mercado. Exemplos incluem:

\begin{itemize}
\tightlist
\item
  \textbf{Patentes}: conferem ao inventor o monopólio temporário sobre a
  exploração de uma inovação. No Brasil, a Lei de Propriedade Industrial
  (Lei 9.279/1996) estabelece prazo de 20 anos para patentes de
  invenção. Esse prazo resulta de um trade-off deliberado de política
  pública: um período mais longo aumenta o incentivo à inovação
  (recompensando o inventor com lucros de monopólio), mas prolonga a
  ineficiência alocativa associada ao poder de mercado.
\item
  \textbf{Concessões e licenças}: o poder público pode restringir o
  número de operadores em um mercado, como ocorre nos serviços de
  distribuição de energia elétrica e nas concessões rodoviárias.
\item
  \textbf{Direitos autorais}: protegem obras intelectuais contra
  reprodução não autorizada, garantindo aos criadores o controle sobre a
  exploração comercial de suas obras.
\end{itemize}

\subsection{Barreiras naturais}\label{barreiras-naturais}

Em alguns mercados, as condições tecnológicas tornam ineficiente a
presença de mais de uma firma. Isso ocorre quando há \textbf{economias
de escala} significativas ao longo de toda a faixa relevante de
produção. Note que a barreira é ``natural'' no sentido de que decorre da
estrutura de custos, não de uma decisão governamental ou de uma ação
estratégica --- nenhuma lei proíbe a entrada, mas a tecnologia a torna
economicamente inviável.

\begin{definitionbox}{Monopólio Natural}

Um mercado constitui um \textbf{monopólio natural} quando a função de
custo é \textbf{subaditiva}, isto é, quando uma única firma pode
produzir qualquer quantidade a um custo total menor do que duas ou mais
firmas produzindo conjuntamente a mesma quantidade:

{[} C(q) \textless{} C(q\_1) + C(q\_2), \quad \text{para todo } q\_1,
q\_2 \textgreater{} 0 \text{ com } q\_1 + q\_2 = q {]}

No caso de um único produto, a subaditividade é implicada por economias
de escala ao longo de toda a faixa relevante de produção --- ou seja, o
custo médio é decrescente.

\end{definitionbox}

Exemplos clássicos incluem redes de distribuição de água, gás,
eletricidade e telecomunicações fixas, nas quais a duplicação da
infraestrutura seria socialmente custosa. É importante distinguir
monopólio natural de monopólio \emph{de fato}: uma firma pode ser a
única no mercado por razões históricas ou estratégicas sem que o mercado
seja um monopólio natural. A subaditividade dos custos é condição
necessária para caracterizar o monopólio natural.

\begin{intuitionbox}{Intuição Econômica}

\textbf{Em uma frase:} Em um monopólio natural, uma única empresa atende
todo o mercado a um custo menor do que duas ou mais fariam --- duplicar
a infraestrutura seria desperdício.

\textbf{Pense assim:} Pense na rede de água encanada do seu bairro. Faz
sentido ter uma empresa com uma rede de canos, não duas redes paralelas
competindo na mesma rua. O custo fixo de enterrar a tubulação é tão alto
que, quanto mais casas uma única rede atende, menor o custo por casa.
Construir uma segunda rede seria jogar dinheiro fora --- literalmente
cavar buracos à toa.

\textbf{Por que isso importa:} É por isso que setores como saneamento,
distribuição de energia e ferrovias são regulados por agências como
\href{https://www.aneel.gov.br}{ANEEL} e
\href{https://www.gov.br/ana}{ANA}, em vez de simplesmente liberados à
concorrência.

\end{intuitionbox}

\subsection{Barreiras estratégicas}\label{barreiras-estratuxe9gicas}

Além das barreiras criadas pela natureza da tecnologia ou pela ação do
Estado, as próprias firmas incumbentes podem erguer obstáculos
deliberados à entrada de concorrentes. Essas barreiras estratégicas
refletem o comportamento racional de uma firma que antecipa a
possibilidade de entrada e age preventivamente para torná-la não
lucrativa. A análise formal dessas estratégias envolve teoria dos jogos
e será aprofundada nos Capítulos 16a--16b; aqui, destacamos os
principais mecanismos.

Firmas incumbentes podem adotar comportamentos deliberados para
dificultar a entrada de rivais:

\begin{itemize}
\tightlist
\item
  \textbf{Excesso de capacidade instalada}: sinaliza que o incumbente
  pode expandir a produção rapidamente em resposta à entrada, tornando-a
  não lucrativa. Spence (1977) formalizou esse argumento mostrando que o
  investimento em capacidade funciona como um compromisso crível
  (\emph{commitment}).
\item
  \textbf{Preços predatórios}: fixar preços temporariamente abaixo do
  custo para expulsar ou dissuadir concorrentes (prática ilícita no
  direito concorrencial brasileiro, conforme a Lei 12.529/2011). A
  eficácia dessa estratégia depende da capacidade financeira do
  incumbente de sustentar prejuízos temporários.
\item
  \textbf{Proliferação de marcas}: ocupar nichos de mercado para reduzir
  o espaço disponível para entrantes --- estratégia documentada
  empiricamente no mercado americano de cereais matinais por Schmalensee
  (1978).
\item
  \textbf{Controle de insumos essenciais}: adquirir ou controlar o
  acesso a recursos sem os quais rivais não podem operar.
\end{itemize}

\begin{boxmundobox}{Box Mundo 15.1 — Patentes farmacêuticas nos EUA: o preço do monopólio legal}

\textbf{Contexto:} A indústria farmacêutica dos Estados Unidos é o
exemplo contemporâneo mais saliente de monopólio criado por barreiras
legais. Patentes de 20 anos conferem às empresas farmacêuticas o direito
exclusivo de comercializar novas moléculas, gerando margens de lucro que
figuram entre as maiores de qualquer indústria. Além disso, o
\emph{Hatch-Waxman Act} (1984) criou extensões de patente que podem
estender a exclusividade por até 5 anos adicionais, além de períodos de
exclusividade de dados de 5 a 12 anos para biológicos.

\textbf{Dados:} Segundo o \emph{Congressional Budget Office} (CBO), o
preço médio dos medicamentos de marca nos EUA é 3,4 vezes maior do que
em outros países do G7 (CBO, 2022). Quando uma patente expira e
genéricos entram no mercado, o preço cai tipicamente 80--90\% dentro de
dois anos --- uma demonstração direta do efeito do poder de monopólio
sobre os preços. Em 2022, os americanos gastaram US\$ 405 bilhões em
medicamentos prescritos (CMS National Health Expenditure Data).
Estratégias como \emph{pay-for-delay} --- acordos nos quais o detentor
da patente paga ao fabricante de genéricos para atrasar a entrada ---
custaram aos consumidores americanos US\$ 3,5 bilhões por ano, segundo
estimativa da FTC (2010).

\textbf{Análise:} A patente farmacêutica ilustra o dilema central do
monopólio legal: sem a promessa de lucros temporários de monopólio,
empresas teriam menos incentivo para investir os bilhões necessários ao
desenvolvimento de novos medicamentos (o custo médio de levar um fármaco
ao mercado é estimado em US\$ 1,3 bilhão pelo \emph{Tufts Center for the
Study of Drug Development}). Mas a duração e amplitude da proteção
patentária determinam o tamanho da perda de peso morto imposta à
sociedade durante a vigência do monopólio. O debate sobre a duração
ótima de patentes é, essencialmente, uma análise de custo-benefício
entre incentivo à inovação (benefício dinâmico) e ineficiência alocativa
(custo estático) --- tema que retomaremos na Seção 15.10.

\textbf{Fonte:} Congressional Budget Office (2022). \emph{Prescription
Drugs: Spending, Use, and Prices}. FTC (2010). \emph{Pay-for-Delay: How
Drug Company Pay-Offs Cost Consumers Billions}. CMS National Health
Expenditure Data, 2022.

\end{boxmundobox}

\begin{center}\rule{0.5\linewidth}{0.5pt}\end{center}

\section{15.2 O Sorriso de Quem Define o Preço: Maximização de Lucro do
Monopolista}\label{152}

Agora que sabemos \emph{por que} o monopolista reina sozinho, vem a
pergunta que realmente interessa: o que ele faz com todo esse poder?
Imagine que você é o único vendedor de guarda-chuvas em uma cidade onde
chove todos os dias. Você cobra quanto quiser? Não exatamente --- se
exagerar no preço, ninguém compra. O monopolista vive esse dilema: ele
tem poder, mas não é onipotente. A resposta exige uma análise cuidadosa
da relação entre receita marginal, custo marginal e elasticidade da
demanda --- relação esta que se tornará um instrumento analítico
recorrente ao longo do restante do livro.

Diferentemente da firma competitiva, que toma o preço como dado e
escolhe apenas a quantidade, o monopolista reconhece que suas decisões
de quantidade afetam o preço de mercado. Ao expandir a produção, ele
enfrenta dois efeitos opostos sobre a receita: o \emph{efeito
quantidade} (vende mais unidades) e o \emph{efeito preço} (precisa
reduzir o preço de \emph{todas} as unidades, inclusive as que já
venderia ao preço mais alto). Essa tensão entre os dois efeitos é a
essência do problema de monopólio e explica por que a receita marginal é
sempre inferior ao preço.

\subsection{O problema do monopolista}\label{o-problema-do-monopolista}

O monopolista enfrenta toda a curva de demanda do mercado. Se a demanda
inversa é \(p(q)\), o problema de maximização é:

{[} \max\_q ; \pi(q) = p(q) \cdot q - C(q) \label{eq:15.1} \tag{15.1}
{]}

A condição de primeira ordem é:

{[} \frac{d\pi}{dq} =
\underbrace{p(q) + q \cdot p'(q)}\emph{\{\text{Receita Marginal (RMg)}\}
- \underbrace{C'(q)}}\{\text{Custo Marginal (CMg)}\} = 0 \label{eq:15.2}
\tag{15.2} {]}

Portanto, a regra de maximização é:

{[} \boxed{RMg(q^m) = CMg(q^m)} \label{eq:15.3} \tag{15.3} {]}

A condição de segunda ordem exige que \(\frac{d^2\pi}{dq^2} < 0\), ou
seja, que a receita marginal corte o custo marginal ``de cima para
baixo''. Geometricamente, essa condição garante que o ponto encontrado é
um máximo (e não um mínimo) de lucro.

Note que a expressão \(RMg = p + q \cdot p'(q)\) tem uma interpretação
econômica direta: o primeiro termo é a receita obtida pela unidade
adicional vendida; o segundo (negativo, pois \(p'(q) < 0\) para demanda
decrescente) é a perda de receita nas unidades que já seriam vendidas ao
preço mais alto. Para uma firma competitiva, \(p'(q) = 0\) (a firma não
afeta o preço) e \(RMg = p\). Para o monopolista, \(RMg < p\): cada
unidade adicional vale menos que seu preço, porque deprime o preço de
todas as unidades.

No gráfico interativo abaixo, ajuste os parâmetros da demanda e do custo
marginal para visualizar como o equilíbrio de monopólio se modifica.
Observe que o monopolista sempre escolhe a quantidade onde a curva de
\(RMg\) cruza a de \(CMg\), e o preço é lido sobre a curva de demanda
--- acima do \(CMg\).

\figurePlaceholder{/micro-book/graficos/cap15/monopolio.html}

\textbf{Figura 15.1 --- Equilíbrio de monopólio.} Ajuste os parâmetros
da demanda (\(a\), \(b\)) e do custo marginal (\(c\)) para visualizar o
equilíbrio de monopólio, o lucro, o excedente do consumidor, a perda de
peso morto e o índice de Lerner. Compare com o resultado competitivo.

\subsection{Receita marginal e
elasticidade}\label{receita-marginal-e-elasticidade}

A receita marginal pode ser expressa em termos da elasticidade-preço da
demanda \(\varepsilon_{p}\) (definida como valor negativo,
\(\varepsilon_p < 0\)):

{[} RMg = p\left(1 + \frac{1}{\varepsilon_p}\right) = p\left(1 -
\frac{1}{|\varepsilon_p|}\right) \label{eq:15.4} \tag{15.4} {]}

Pela equação (\ref{eq:15.4}), como o monopolista opera onde \(RMg > 0\),
ele necessariamente produz na faixa \textbf{elástica} da demanda
(\(|\varepsilon_p| > 1\)). Na faixa inelástica, a receita marginal seria
negativa, e o monopolista poderia aumentar a receita (e o lucro)
simplesmente reduzindo a produção --- uma situação que contradiz a
otimização. Esse resultado, embora simples, é frequentemente esquecido
por alunos e tem implicações práticas importantes: se um governo observa
que a demanda por um bem fornecido por um monopolista é inelástica, pode
concluir que o monopolista não está maximizando lucro --- ou que há
restrições regulatórias limitando seu preço.

\subsection{A regra de markup}\label{a-regra-de-markup}

A relação entre receita marginal e elasticidade nos conduz a um
resultado fundamental: o monopolista não fixa preços arbitrariamente,
mas obedece a uma regra de markup que conecta o preço cobrado ao custo
marginal e à sensibilidade dos consumidores a variações de preço. Quanto
menos sensíveis forem os consumidores --- isto é, quanto mais inelástica
a demanda ---, maior a margem que o monopolista consegue extrair.

Da condição \(RMg = CMg\), obtemos:

{[} p\left(1 - \frac{1}{|\varepsilon_p|}\right) = CMg \implies p =
\frac{CMg}{1 - \frac{1}{|\varepsilon_p|}} \label{eq:15.5} \tag{15.5} {]}

O monopolista aplica um \textbf{markup} sobre o custo marginal que é
inversamente proporcional à elasticidade-preço da demanda. Quando
\(|\varepsilon_p| \to \infty\) (demanda perfeitamente elástica), o
markup desaparece e \(p \to CMg\) --- recuperamos o resultado
competitivo. Quando \(|\varepsilon_p|\) é baixo (demanda inelástica), o
markup explode --- o monopolista extrai margens cada vez maiores. É essa
relação que fundamenta toda a análise empírica do poder de mercado.

\begin{definitionbox}{Índice de Lerner}

O \textbf{índice de Lerner} mede o poder de mercado de uma firma como a
diferença percentual entre preço e custo marginal:

{[} L = \frac{p - CMg}{p} \label{eq:15.6} \tag{15.6} {]}

Para uma firma em concorrência perfeita, \(L = 0\). Para um monopolista,
\(L = \frac{1}{|\varepsilon_p|}\). O índice varia entre 0 e 1.

\end{definitionbox}

\begin{intuitionbox}{Intuição Econômica}

\textbf{Em uma frase:} O markup do monopolista é maior quando os
consumidores têm poucas alternativas --- quanto mais ``preso'' o
cliente, mais caro o produto.

\textbf{Pense assim:} Compare o preço de uma garrafa d'água no
supermercado (R\$ 2) com o preço no estádio de futebol (R\$ 10). No
estádio, você não pode sair para comprar em outro lugar --- sua demanda
é inelástica --- e o vendedor explora isso cobrando um markup enorme. No
supermercado, com dez marcas na prateleira, ninguém consegue cobrar
muito acima do custo. O índice de Lerner mede exatamente esse ``poder de
espremer'' o consumidor.

\textbf{Por que isso importa:} O \href{https://www.gov.br/cade}{CADE}
usa medidas como o índice de Lerner para avaliar se fusões e aquisições
criarão poder de mercado excessivo --- como na análise da fusão que
criou a Ambev.

\end{intuitionbox}

\begin{warningbox}{Erro Comum}

\textbf{Achar que o monopolista pode operar na faixa inelástica da
demanda.}

Um erro frequente é esquecer que o monopolista maximizador de lucro
\textbf{nunca} produz na região onde \(|\varepsilon_p| < 1\). Na faixa
inelástica, a receita marginal é negativa: ao reduzir a produção, a
firma aumenta a receita \emph{e} reduz o custo --- logo, o lucro sobe. O
monopolista sempre opera onde \(|\varepsilon_p| > 1\). Se uma questão de
prova apresenta uma demanda com \(|\varepsilon_p| < 1\) no ponto
calculado, revise o cálculo: algo está errado, ou trata-se de uma
solução de canto.

\end{warningbox}

O exercício a seguir aplica as ferramentas recém-desenvolvidas a um
problema numérico completo, ilustrando cada etapa do cálculo do
equilíbrio de monopólio.

\begin{exresolvidobox}{Exercício Resolvido 15.1 — Monopólio com demanda linear: preço, lucro e PPM}

\textbf{Enunciado.} Um monopolista enfrenta demanda \(p = 80 - 2q\) e
tem custo total \(C(q) = 200 + 8q\). (a) Encontre preço, quantidade e
lucro de monopólio. (b) Calcule o índice de Lerner. (c) Calcule a perda
de peso morto e o excedente total. (d) Compare com o resultado
competitivo.

\begin{center}\rule{0.5\linewidth}{0.5pt}\end{center}

\textbf{(a)} A receita marginal é \(RMg = 80 - 4q\). Igualando ao custo
marginal \(CMg = 8\):

{[} 80 - 4q = 8 \implies q\^{}m = 18, \quad p\^{}m = 80 - 36 = 44 {]}

Lucro: \(\pi = 44 \times 18 - (200 + 8 \times 18) = 792 - 344 = 448\).

\textbf{(b)} Índice de Lerner:
\(L = (p - CMg)/p = (44 - 8)/44 = 36/44 \approx 0{,}818\).

Verificação pela elasticidade: no ponto de monopólio,
\(\varepsilon_p = (dq/dp)(p/q) = (-1/2)(44/18) \approx -1{,}222\). Logo
\(L = 1/|\varepsilon_p| = 1/1{,}222 \approx 0{,}818\). ✓

\textbf{(c)} Resultado competitivo: \(p^c = CMg = 8\),
\(q^c = (80-8)/2 = 36\).

{[} PPM = \frac{1}{2}(p\^{}m - CMg)(q\^{}c - q\^{}m) = \frac{1}{2}(44 -
8)(36 - 18) = \frac{1}{2} \times 36 \times 18 = 324 {]}

Excedente total sob monopólio: \(W^m = EC^m + EP^m\).

\(EC^m = \frac{1}{2}(80 - 44)(18) = \frac{1}{2} \times 36 \times 18 = 324\).

O excedente do produtor é a área acima do CMg:
\(EP^m = (p^m - CMg) \times q^m = (44 - 8)(18) = 648\).

\(W^m = 324 + 648 = 972\).

\textbf{(d)} Sob concorrência perfeita:
\(EC^c = \frac{1}{2}(80-8)(36) = 1.296\), \(EP^c = 0\) (com CMg
constante). \(W^c = 1.296\).

\(PPM = W^c - W^m = 1.296 - 972 = 324\). ✓

\textbf{Interpretação:} O monopolista captura R\$ 648 de excedente, mas
destrói R\$ 324 de bem-estar social que ninguém captura --- o triângulo
de Harberger.

\end{exresolvidobox}

\figurePlaceholder{/micro-book/graficos/cap15/webr-monopolio-basico.html}

\textbf{WebR 15.1 --- Monopólio de ponta a ponta: do cálculo ao
gráfico.} O código resolve o ER 15.1 completo em R: calcula preço,
quantidade, lucro, índice de Lerner, excedentes e perda de peso morto,
comparando monopólio com concorrência perfeita. Altere os parâmetros
\texttt{a}, \texttt{b} e \texttt{c} e observe como o triângulo de
Harberger se expande ou encolhe.

\chapter{15.3--15.5 Índice de Lerner, Ineficiência e Estática
Comparativa}\label{uxedndice-de-lerner-ineficiuxeancia-e-estuxe1tica-comparativa}

\section{15.3 A Régua do Poder: Demonstração do Índice de
Lerner}\label{153}

A derivação formal a seguir explicita a conexão algébrica entre o poder
de mercado do monopolista --- medido pelo índice de Lerner --- e a
elasticidade-preço da demanda que ele enfrenta. Embora o resultado já
tenha sido antecipado na seção anterior, a demonstração rigorosa é
instrutiva porque revela a estrutura lógica subjacente e prepara o
terreno para generalizações em contextos de oligopólio (Capítulos
16a--16b), onde o índice de Lerner de cada firma dependerá não apenas da
elasticidade de mercado, mas também de sua participação de mercado e do
grau de colusão.

\begin{proofbox}{Demonstração}

Seja \(p(q)\) a demanda inversa do mercado e \(C(q)\) a função de custo
do monopolista. O problema de maximização de lucro é:

{[} \max\_q ; \pi(q) = p(q) \cdot q - C(q) {]}

A condição de primeira ordem é:

{[} \frac{d\pi}{dq} = p(q) + q \cdot p'(q) - C'(q) = 0 {]}

Reescrevendo:

{[} p + q \cdot p'(q) = CMg {]}

Multiplicamos e dividimos o segundo termo por \(p\):

{[} p + p \cdot \frac{q \cdot p'(q)}{p} = CMg {]}

Observe que a elasticidade-preço da demanda é definida como:

{[} \varepsilon\_p = \frac{dq}{dp} \cdot \frac{p}{q} {]}

Portanto:

{[} \frac{q \cdot p'(q)}{p} = \frac{q}{p} \cdot \frac{dp}{dq} =
\frac{1}{\varepsilon_p} {]}

Substituindo:

{[} p\left(1 + \frac{1}{\varepsilon_p}\right) = CMg {]}

Rearranjando:

{[} p - CMg = -\frac{p}{\varepsilon_p} {]}

Dividindo ambos os lados por \(p\):

{[} \frac{p - CMg}{p} = -\frac{1}{\varepsilon_p} =
\frac{1}{|\varepsilon_p|} {]}

Portanto:

{[} \boxed{L = \frac{p - CMg}{p} = \frac{1}{|\varepsilon_p|}}
\label{eq:15.7} \tag{15.7} {]}

O índice de Lerner é igual ao inverso do valor absoluto da
elasticidade-preço da demanda. Quanto mais inelástica a demanda (menor
\(|\varepsilon_p|\)), maior o poder de mercado. Em concorrência
perfeita, a firma enfrenta demanda perfeitamente elástica
(\(|\varepsilon_p| \to \infty\)), de modo que \(L = 0\).
\(\blacksquare\)

\end{proofbox}

\begin{center}\rule{0.5\linewidth}{0.5pt}\end{center}

\section{15.4 Quem Paga a Conta: Ineficiência do Monopólio}\label{154}

Até aqui, contamos a história do monopólio pelo ângulo do vilão --- e,
convenhamos, para o monopolista a vida é boa: ele escolhe preço,
quantidade e vai dormir tranquilo. Mas e o resto da sociedade? Alguém
está pagando a conta dessa tranquilidade. Em que medida o poder de
mercado do monopolista compromete a eficiência alocativa que os Teoremas
do Bem-Estar (Capítulo 14) garantem para mercados competitivos? A
resposta a essas perguntas fundamenta a justificativa econômica para a
regulação e a política antitruste --- e, como veremos, envolve não
apenas o conhecido ``triângulo de Harberger'', mas também custos mais
sutis associados ao \emph{rent-seeking} e à ineficiência organizacional.

\subsection{Perda de peso morto}\label{perda-de-peso-morto}

O monopolista produz menos e cobra mais do que o resultado competitivo.
Essa distorção gera uma \textbf{perda de peso morto} (PPM) --- uma
redução no excedente total que não é capturada por nenhum agente. A PPM
representa transações mutuamente benéficas que deixam de ocorrer:
consumidores que estariam dispostos a pagar mais do que o custo marginal
de produção, mas são excluídos pelo preço de monopólio.

Considere um mercado com demanda linear \(p = a - bq\) e custo marginal
constante \(CMg = c\).

\textbf{Resultado competitivo} (\(p = CMg\)):

{[} q\^{}c = \frac{a - c}{b}, \quad p\^{}c = c {]}

\textbf{Resultado monopolístico} (\(RMg = CMg\)):

{[} RMg = a - 2bq = c \implies q\^{}m = \frac{a - c}{2b}, \quad p\^{}m =
\frac{a + c}{2} \label{eq:15.8} \tag{15.8} {]}

Pela equação (\ref{eq:15.8}), o monopolista produz \textbf{metade} da
quantidade competitiva e cobra um preço \textbf{acima} do custo
marginal. Essa relação --- com demanda linear, a produção de monopólio é
exatamente metade da competitiva --- é um resultado útil que merece ser
lembrado; ele decorre do fato de que, com demanda linear, a curva de
\(RMg\) tem o mesmo intercepto e o dobro da inclinação da curva de
demanda.

\subsection{O triângulo de Harberger}\label{o-triuxe2ngulo-de-harberger}

A perda de peso morto corresponde à área do triângulo entre a curva de
demanda e o custo marginal, entre \(q^m\) e \(q^c\):

{[} PPM = \frac{1}{2}(p\^{}m - c)(q\^{}c - q\^{}m) = \frac{(a-c)^2}{8b}
\label{eq:15.9} \tag{15.9} {]}

Arnold Harberger (1954) pioneirou a mensuração empírica dessa perda e
estimou que a PPM nos Estados Unidos seria da ordem de 0,1\% do PIB ---
um valor surpreendentemente pequeno que provocou intenso debate na
profissão. Estimativas posteriores, incorporando \emph{rent-seeking} (os
recursos desperdiçados na tentativa de obter ou manter o monopólio, como
atividades de lobby) e ineficiências organizacionais
(\emph{X-inefficiency} de Leibenstein, 1966 --- a ideia de que a
ausência de pressão competitiva permite à firma operar aquém de sua
fronteira de custos), sugerem valores significativamente maiores.
Cowling e Mueller (1978) estimaram que o custo social do monopólio
poderia chegar a 4--13\% do PIB quando o lucro do monopolista é incluído
como custo social (na hipótese de que é inteiramente dissipado em
\emph{rent-seeking}).

Richard Posner (1975) formalizou esse argumento: se firmas gastam
recursos para obter e manter poder de mercado (por exemplo, em lobby,
litígios e barreiras estratégicas), o custo social total pode ser muito
maior que o triângulo de Harberger --- podendo, no limite, ser tão
grande quanto o retângulo de lucro do monopolista, como veremos no
Exercício Resolvido 15.2.

\begin{examplebox}{Poder de mercado e o CADE}

O Conselho Administrativo de Defesa Econômica (CADE) utiliza o índice de
Lerner e medidas correlatas para avaliar poder de mercado em atos de
concentração e investigações de conduta. No caso da fusão Ambev (1999),
o CADE estimou que a empresa resultante deteria cerca de 70\% do mercado
de cervejas, com significativo poder de precificação. Mais recentemente,
em 2022, o CADE aplicou multa de R\$ 539 milhões ao Google por abuso de
posição dominante no mercado de \emph{ad tech} --- um caso em que o
poder de mercado decorre não de barreiras tecnológicas clássicas, mas de
efeitos de rede e controle de dados. O índice de Lerner é uma ferramenta
analítica central nessas investigações: um \(L\) elevado e persistente
sinaliza poder de mercado que justifica intervenção regulatória. A
\hyperref[tabela-15-1]{Tabela 15.1} compara os resultados dos dois
regimes.

\end{examplebox}

{\def\LTcaptype{none} % do not increment counter
\begin{longtable}[]{@{}
  >{\raggedright\arraybackslash}p{(\linewidth - 6\tabcolsep) * \real{0.2500}}
  >{\raggedright\arraybackslash}p{(\linewidth - 6\tabcolsep) * \real{0.2500}}
  >{\raggedright\arraybackslash}p{(\linewidth - 6\tabcolsep) * \real{0.2500}}
  >{\raggedright\arraybackslash}p{(\linewidth - 6\tabcolsep) * \real{0.2500}}@{}}
\toprule\noalign{}
\begin{minipage}[b]{\linewidth}\raggedright
Variável
\end{minipage} & \begin{minipage}[b]{\linewidth}\raggedright
Concorrência Perfeita
\end{minipage} & \begin{minipage}[b]{\linewidth}\raggedright
Monopólio
\end{minipage} & \begin{minipage}[b]{\linewidth}\raggedright
Variação
\end{minipage} \\
\midrule\noalign{}
\endhead
\bottomrule\noalign{}
\endlastfoot
\textbf{Preço} & \(p^c = c\) & \(p^m = \frac{a+c}{2}\) & \(\uparrow\) \\
\textbf{Quantidade} & \(q^c = \frac{a-c}{b}\) & \(q^m = \frac{a-c}{2b}\)
& \(\downarrow\) \\
\textbf{Excedente do Consumidor} & \(\frac{(a-c)^2}{2b}\) &
\(\frac{(a-c)^2}{8b}\) & \(\downarrow\) \\
\textbf{Excedente do Produtor} & \(0\) & \(\frac{(a-c)^2}{4b}\) &
\(\uparrow\) \\
\textbf{Excedente Total} & \(\frac{(a-c)^2}{2b}\) &
\(\frac{3(a-c)^2}{8b}\) & \(\downarrow\) \\
\textbf{Perda de Peso Morto} & \(0\) & \(\frac{(a-c)^2}{8b}\) & --- \\
\end{longtable}
}

\textbf{Tabela 15.1 --- Comparação concorrência perfeita versus
monopólio.}

\begin{tipbox}{Interpretação}

A passagem de concorrência perfeita para monopólio envolve duas
transferências: (i) parte do excedente do consumidor é capturada pelo
monopolista (retângulo de transferência); (ii) parte do excedente total
simplesmente desaparece (triângulo de Harberger). A PPM representa uma
perda \textbf{líquida} de bem-estar social.

\end{tipbox}

\begin{intuitionbox}{Intuição Econômica}

\textbf{Em uma frase:} O monopólio produz menos e cobra mais do que um
mercado competitivo, e a diferença de bem-estar que se perde no caminho
não vai para ninguém --- simplesmente evapora.

\textbf{Pense assim:} Imagine que, em vez de vários postos de gasolina
no bairro, houvesse apenas um. Ele reduziria a quantidade vendida e
aumentaria o preço. Parte do que os consumidores perdem vai para o bolso
do dono do posto (lucro de monopólio), mas outra parte --- os litros que
deixaram de ser vendidos porque o preço ficou alto demais --- é pura
perda. É como se alguém queimasse dinheiro: nem o consumidor nem o
produtor ficam com ele.

\textbf{Por que isso importa:} A perda de peso morto do monopólio é o
principal argumento econômico para a existência de órgãos como o CADE,
que combatem concentração excessiva de mercado para proteger o bem-estar
dos consumidores brasileiros.

\end{intuitionbox}

\begin{boxmundobox}{Box Mundo 15.2 — O custo social do monopólio: de Harberger (EUA) à OCDE}

\textbf{Contexto:} Quanto custa o monopólio para a sociedade? A resposta
depende criticamente do que se inclui na conta. Arnold Harberger (1954),
em estudo seminal com dados de 73 indústrias manufatureiras dos EUA no
período 1924--1928, inaugurou a tradição de mensuração empírica das
perdas de eficiência do poder de mercado --- e obteve um resultado que
surpreendeu a profissão por sua modéstia.

\textbf{Dados:} A estimativa original de Harberger foi de apenas 0,1\%
do PIB americano --- tão pequena que levantou a questão de se o
monopólio era realmente um problema econômico relevante. Entretanto,
estimativas posteriores usando metodologias mais abrangentes chegaram a
valores muito maiores. Cowling e Mueller (1978), utilizando dados de
firmas individuais nos EUA (734 firmas) e no Reino Unido, estimaram
perdas de 4--13\% do PIB americano e 3--7\% do PIB britânico, ao incluir
o lucro de monopólio como custo social (seguindo a lógica de Posner de
que os lucros de monopólio são dissipados em \emph{rent-seeking}). Mais
recentemente, o \emph{OECD Competition Assessment Toolkit} (2019) estima
que reformas pró-concorrência em mercados com poder de mercado
significativo podem elevar o PIB em 0,5\% a 1,5\% nos países membros ---
cifras que indicam perdas correntes nessa magnitude.

\textbf{Análise:} A discrepância entre as estimativas reflete diferenças
conceituais profundas sobre o que constitui o ``custo social'' do
monopólio. O triângulo de Harberger captura apenas a ineficiência
alocativa estática (transações não realizadas). A abordagem de
Posner/Cowling-Mueller adiciona os custos de \emph{rent-seeking} ---
recursos reais desviados para obter ou manter poder de mercado, como
lobby, litígios e publicidade persuasiva. A teoria da
\emph{X-inefficiency} de Leibenstein (1966) sugere ainda que, na
ausência de pressão competitiva, as firmas operam acima de sua fronteira
de custos mínimos, gerando desperdício interno. Essas três fontes de
custo --- alocativo, \emph{rent-seeking} e organizacional --- se somam e
podem tornar o monopólio um problema econômico substancial.

\textbf{Fonte:} Harberger, A. C. (1954). Monopoly and resource
allocation. \emph{AER}, 44(2), 77--87. Cowling, K.; Mueller, D. C.
(1978). The social costs of monopoly power. \emph{Economic Journal},
88(352), 727--748. OECD (2019). \emph{Competition Assessment Toolkit}.

\end{boxmundobox}

O exercício a seguir formaliza o argumento de Posner sobre os custos de
\emph{rent-seeking}.

\begin{exresolvidobox}{Exercício Resolvido 15.2 — Custo social do monopólio com rent-seeking}

\textbf{Enunciado.} Em um mercado com demanda \(p = 120 - 3q\), um
monopolista tem \(CMg = 30\). (a) Calcule o equilíbrio de monopólio e a
PPM (triângulo de Harberger). (b) Se o monopolista gasta até a
totalidade do lucro em atividades de \emph{rent-seeking} para manter o
monopólio, qual é o custo social máximo? (c) Compare com o excedente
total sob concorrência perfeita.

\begin{center}\rule{0.5\linewidth}{0.5pt}\end{center}

\textbf{(a)} \(RMg = 120 - 6q = 30 \implies q^m = 15\),
\(p^m = 120 - 45 = 75\).

Lucro: \(\pi^m = (75 - 30) \times 15 = 675\).

Resultado competitivo: \(p^c = 30\), \(q^c = (120-30)/3 = 30\).

{[} PPM = \frac{1}{2}(75 - 30)(30 - 15) = \frac{1}{2} \times 45
\times 15 = 337\{,\}5 {]}

\textbf{(b)} Segundo Posner (1975), se o monopolista dissipa todo o
lucro em \emph{rent-seeking}, o custo social total é:

{[} CS\_\{\text{total}\} = PPM + \pi\^{}m = 337\{,\}5 + 675 =
1.012\{,\}5 {]}

\textbf{(c)} Excedente total sob concorrência perfeita:

{[} W\^{}c = EC\^{}c = \frac{1}{2}(120 - 30)(30) = 1.350 {]}

O custo social máximo (1.012,5) representa 75\% do excedente
competitivo. Sob monopólio com \emph{rent-seeking} total, o bem-estar
social restante seria apenas
\(EC^m = \frac{1}{2}(120-75)(15) = 337{,}5\) --- exatamente 25\% do
excedente competitivo.

\textbf{Interpretação:} O triângulo de Harberger sozinho subestima o
custo social do monopólio. Quando incluímos o desperdício de recursos em
\emph{rent-seeking}, o custo pode ser até três vezes maior. Este é o
argumento central de Posner para uma política antitruste vigorosa.

\end{exresolvidobox}

\begin{center}\rule{0.5\linewidth}{0.5pt}\end{center}

\section{15.5 Mexeu no Tabuleiro: Estática Comparativa do
Monopólio}\label{155}

Estabelecido o equilíbrio de monopólio, podemos agora investigar como
ele se altera diante de mudanças nos parâmetros exógenos. Que efeito tem
um imposto sobre o preço e a quantidade de monopólio? Como o monopolista
responde a um deslocamento da demanda? Essas perguntas de estática
comparativa são particularmente relevantes para a formulação de
políticas públicas, pois revelam como instrumentos tributários e
regulatórios interagem com o poder de mercado. Os resultados, como
veremos, diferem qualitativamente do caso competitivo --- o que implica
que intuições desenvolvidas para mercados competitivos podem ser
enganosas quando aplicadas a mercados com poder de mercado.

\subsection{Efeito de um imposto
unitário}\label{efeito-de-um-imposto-unituxe1rio}

Suponha que o governo impõe um imposto unitário \(t\) sobre cada unidade
produzida. O novo custo marginal efetivo é \(CMg + t\). Com demanda
linear:

{[} q\^{}m(t) = \frac{a - c - t}{2b}, \quad p\^{}m(t) =
\frac{a + c + t}{2} \label{eq:15.10} \tag{15.10} {]}

Observe que \(\frac{dp^m}{dt} = \frac{1}{2}\): o monopolista repassa
apenas \textbf{metade} do imposto ao consumidor. Isso contrasta com a
concorrência perfeita, na qual, com oferta perfeitamente elástica, o
repasse é integral. A diferença tem implicações diretas para a política
tributária: um imposto sobre um monopolista é parcialmente absorvido
pelo lucro da firma, o que pode ser desejável do ponto de vista
distributivo --- o monopolista ``paga'' metade do imposto.

\begin{notebox}{Repasse imperfeito}

O fato de o monopolista absorver parte do imposto decorre da curvatura
da curva de demanda que ele enfrenta. A incidência do imposto depende
crucialmente da forma funcional da demanda: com demanda isoelástica
(\(q = Ap^{\varepsilon}\)), por exemplo, o repasse pode \emph{exceder}
100\% --- o preço sobe mais do que o valor do imposto. Esse resultado
aparentemente paradoxal ocorre porque a demanda isoelástica é mais
convexa que a linear, fazendo com que a \(RMg\) se desloque mais do que
proporcionalmente. A forma funcional da demanda, portanto, é crucial
para a análise de incidência tributária em mercados com poder de
mercado.

\end{notebox}

\subsection{Deslocamento da demanda}\label{deslocamento-da-demanda}

Um aumento na demanda (deslocamento paralelo de \(a\) para
\(a + \Delta\)) eleva tanto o preço quanto a quantidade de monopólio:

{[} \Delta p\^{}m = \frac{\Delta}{2}, \quad \Delta q\^{}m =
\frac{\Delta}{2b} {]}

O lucro do monopolista é crescente e convexo na demanda, o que implica
que o monopolista tem incentivos fortes para investir em publicidade ou
atividades que expandam a demanda. Essa convexidade do lucro na demanda
é uma propriedade geral do monopólio (não restrita à demanda linear) e
tem uma consequência importante: o monopolista valora aumentos de
demanda mais do que proporcionalmente, o que explica os elevados gastos
em publicidade observados em mercados monopolísticos ou altamente
concentrados.

No gráfico interativo abaixo, experimente ajustar o slider do imposto
\(t\) e observe como o equilíbrio de monopólio se desloca. Note que o
preço sobe apenas metade do valor do imposto, enquanto a quantidade cai.
Observe também o retângulo de receita fiscal e como a perda de peso
morto se altera.

\figurePlaceholder{/micro-book/graficos/cap15/estatica-comparativa.html}

\textbf{Figura 15.2 --- Estática comparativa: efeito de um imposto sobre
o monopólio.} Ajuste o imposto unitário \(t\) e os parâmetros da demanda
para visualizar como o equilíbrio se desloca. Observe que o monopolista
repassa apenas metade do imposto (com demanda linear) e compare com o
resultado competitivo.

\begin{exresolvidobox}{Exercício Resolvido 15.3 — Imposto sobre o monopolista}

\textbf{Enunciado.} Um monopolista enfrenta demanda \(p = 60 - q\) e tem
\(CMg = 12\). O governo impõe imposto unitário \(t = 8\). (a) Encontre
\((p, q, \pi)\) antes e depois do imposto. (b) Que fração do imposto é
repassada ao consumidor? (c) Calcule a variação no excedente do
consumidor.

\begin{center}\rule{0.5\linewidth}{0.5pt}\end{center}

\textbf{(a) Antes do imposto:}

\(RMg = 60 - 2q = 12 \implies q^m = 24\), \(p^m = 36\).
\(\pi = (36-12)(24) = 576\).

\textbf{Depois do imposto:} CMg efetivo = \(12 + 8 = 20\).

\(RMg = 60 - 2q = 20 \implies q^m(t) = 20\), \(p^m(t) = 40\).
\(\pi = (40 - 20)(20) = 400\).

\textbf{(b)} Repasse: \(\Delta p = 40 - 36 = 4\). Fração repassada:
\(4/8 = 50\%\). ✓ (Resultado geral para demanda linear.)

\textbf{(c)} \(EC_{\text{antes}} = \frac{1}{2}(60-36)(24) = 288\).

\(EC_{\text{depois}} = \frac{1}{2}(60-40)(20) = 200\).

\(\Delta EC = 200 - 288 = -88\).

\textbf{Interpretação:} O consumidor perde R\$ 88 de excedente. O
monopolista absorve R\$ 176 de perda de lucro (\(576 - 400\)). A receita
fiscal é \(8 \times 20 = 160\). A soma \(88 + 176 = 264\) excede a
receita fiscal de 160 em \(264 - 160 = 104\) --- esta diferença é o
aumento na perda de peso morto causado pelo imposto.

\end{exresolvidobox}

\figurePlaceholder{/micro-book/graficos/cap15/webr-imposto-monopolio.html}

\textbf{WebR 15.2 --- O monopolista absorve metade do imposto.} O código
implementa o ER 15.3: calcula o equilíbrio antes e depois do imposto,
mostrando que o repasse é exatamente 50\% com demanda linear. No gráfico
da direita, observe que o lucro é uma parábola invertida no imposto ---
existe um \(t\) que maximiza a receita fiscal. Experimente alterar
\texttt{t} e compare o efeito sobre consumidor, produtor e governo.

\chapter{15.6--15.8 Qualidade, Discriminação de Preços e
Tarifas}\label{qualidade-discriminauxe7uxe3o-de-preuxe7os-e-tarifas}

\section{15.6 Bom o Bastante para o Último Cliente: Qualidade sob
Monopólio}\label{156}

Até este ponto, tratamos o produto do monopolista como dado --- variando
apenas preço e quantidade. Na prática, porém, o monopolista também
decide sobre a \textbf{qualidade} do que oferece. Será que um
monopolista investe adequadamente em qualidade, ou a ausência de pressão
competitiva o leva a oferecer produtos inferiores ao que seria
socialmente desejável? A resposta, como veremos, é surpreendentemente
ambígua --- e depende de detalhes da distribuição de preferências entre
os consumidores.

O monopolista também escolhe a \textbf{qualidade} \(s\) do produto. Se a
demanda é \(p(q, s)\) e o custo é \(C(q, s)\), o problema torna-se:

{[} \max\_\{q, s\} ; \pi(q, s) = p(q, s) \cdot q - C(q, s) {]}

A condição para a escolha ótima de qualidade é:

{[} \frac{\partial p}{\partial s} \cdot q =
\frac{\partial C}{\partial s} \label{eq:15.11} \tag{15.11} {]}

O monopolista iguala a receita marginal da qualidade ao custo marginal
da qualidade. A receita marginal da qualidade é avaliada no
\textbf{consumidor marginal} (o último a adquirir o bem), não no
consumidor médio. Se o consumidor marginal valora a qualidade menos que
o consumidor médio --- o que é plausível quando os consumidores com
maior disposição a pagar também valorizam mais a qualidade ---, o
monopolista subprovê qualidade em relação ao ótimo social.
Intuitivamente, o monopolista ajusta a qualidade para agradar ao
consumidor que está ``na margem'' de comprar ou não, não ao consumidor
que já está disposto a pagar caro.

\begin{notebox}{Qualidade e bem-estar}

Spence (1975) demonstrou que o monopolista provê qualidade eficiente
quando o consumidor marginal tem a mesma valoração de qualidade que o
consumidor médio. Em geral, não há resultado inequívoco: o monopolista
pode oferecer qualidade excessiva ou insuficiente. Esse resultado
contrasta com a intuição popular de que ``monopólio implica qualidade
ruim'' --- na verdade, a direção da distorção depende de como as
preferências por qualidade se distribuem entre os consumidores.

\end{notebox}

\begin{intuitionbox}{Intuição Econômica}

\textbf{Em uma frase:} O monopolista ajusta a qualidade pensando no
cliente que está quase desistindo de comprar --- não no cliente mais
entusiasmado.

\textbf{Pense assim:} Imagine uma operadora de TV a cabo monopolista
decidindo quantos canais oferecer. Ela não pergunta ``quantos canais o
cinéfilo fanático quer?'', mas ``quantos canais preciso para convencer a
última família a assinar?''. Se as famílias que assinam mais facilmente
(alta disposição a pagar) também são as que mais valorizam variedade de
canais, a operadora oferecerá \emph{menos} canais do que seria
socialmente ótimo --- porque está calibrando pelo consumidor marginal,
menos exigente.

\textbf{Por que isso importa:} Esse resultado é relevante para a
regulação de serviços monopolísticos no Brasil: a ANATEL, ao regular
planos de telecomunicações, precisa considerar não apenas preços, mas
também padrões mínimos de qualidade --- porque o monopolista não tem
incentivos automáticos para fornecer a qualidade socialmente desejável.

\end{intuitionbox}

\begin{center}\rule{0.5\linewidth}{0.5pt}\end{center}

\section{15.7 Um Preço para Cada Bolso: Discriminação de
Preços}\label{157}

Você já se perguntou por que a mesma passagem de avião custa R\$ 400
para quem compra com antecedência e R\$ 2.000 para quem compra na
véspera --- sendo que os dois sentam no mesmo avião, comem o mesmo
amendoim e chegam ao mesmo destino? A resposta é que o monopolista (ou
quase-monopolista) aprendeu um truque poderoso: abandonar a prática de
preço único. Se o monopolista consegue identificar --- ou induzir a
autorevelação de --- diferentes disposições a pagar entre consumidores,
ele pode explorar essa heterogeneidade para extrair mais excedente do
que seria possível com um preço uniforme. Essa prática, conhecida como
discriminação de preços, é onipresente no mundo real e assume formas
sofisticadas que variam de passagens aéreas a planos de software, de
ingressos de cinema a tarifas de eletricidade.

A discriminação de preços exige duas condições: (i) algum grau de poder
de mercado e (ii) impossibilidade de revenda (arbitragem) entre
consumidores. A segunda condição é crucial: se consumidores que pagam
preço baixo pudessem revender para os que pagam preço alto, a
discriminação seria impossível. A taxonomia clássica, devida a Pigou
(1920), distingue três graus de discriminação conforme a informação
disponível ao monopolista.

\subsection{15.7.1 Discriminação de primeiro grau
(perfeita)}\label{discriminauxe7uxe3o-de-primeiro-grau-perfeita}

Na discriminação perfeita, o monopolista cobra de cada consumidor
exatamente sua \textbf{disposição a pagar}. Cada unidade é vendida a um
preço diferente --- é o caso limite em que o monopolista possui
informação perfeita sobre as preferências de cada consumidor individual.

\begin{definitionbox}{Discriminação de Primeiro Grau}

Na \textbf{discriminação de preços de primeiro grau}, o monopolista
extrai todo o excedente do consumidor, cobrando o preço máximo que cada
consumidor está disposto a pagar por cada unidade.

\end{definitionbox}

Resultado:

\begin{itemize}
\tightlist
\item
  A quantidade produzida é \textbf{eficiente}: \(q = q^c\) (idêntica à
  concorrência perfeita).
\item
  Não há perda de peso morto.
\item
  Todo o excedente é capturado pelo produtor: \(EC = 0\),
  \(EP = \frac{(a-c)^2}{2b}\).
\end{itemize}

A discriminação perfeita é um caso limite teórico --- na prática, o
monopolista raramente conhece as disposições a pagar individuais. Sua
importância analítica reside em dois pontos: primeiro, ela serve como
benchmark contra o qual comparar as demais formas de discriminação;
segundo, ela mostra que o poder de mercado \emph{per se} não gera
ineficiência alocativa --- é a \emph{uniformidade} do preço que gera a
PPM. Se o monopolista pudesse cobrar preços diferentes de cada
consumidor, a eficiência seria restaurada (embora a distribuição do
excedente fosse radicalmente diferente).

\subsection{15.7.2 Discriminação de segundo grau
(não-linear)}\label{discriminauxe7uxe3o-de-segundo-grau-nuxe3o-linear}

A discriminação de primeiro grau é um caso limite que pressupõe
informação perfeita sobre cada consumidor --- uma hipótese heroica na
maioria dos contextos reais. O que acontece quando o monopolista sabe
que os consumidores diferem em suas preferências, mas não consegue
observar diretamente o ``tipo'' de cada um? Nesse cenário, o monopolista
recorre a uma estratégia engenhosa: em vez de atribuir preços a
consumidores, ele oferece um menu de opções e deixa que os próprios
consumidores se revelem por meio de suas escolhas. Esse é um caso
particular do problema de desenho de mecanismos sob informação
assimétrica, formalizado por Mirrlees (Nobel 1996) e Maskin (Nobel
2007).

Quando o monopolista não observa as características dos consumidores,
ele pode oferecer um \textbf{menu de contratos} (combinações
preço-quantidade) e permitir que os consumidores se
\textbf{autosselecionem}. Consumidores com alta disposição a pagar
escolhem o pacote premium; consumidores com baixa disposição escolhem o
pacote básico.

A estrutura ótima envolve:

\begin{itemize}
\tightlist
\item
  O tipo de \textbf{alta valoração} recebe a quantidade eficiente, mas
  paga um preço que lhe deixa algum excedente (renda informacional).
\item
  O tipo de \textbf{baixa valoração} recebe uma quantidade
  \textbf{distorcida para baixo} (abaixo do eficiente) e extrai
  exatamente zero de excedente.
\end{itemize}

Essa distorção no consumo do tipo baixo é o custo da
\textbf{compatibilidade de incentivos}: é necessário tornar o pacote
básico suficientemente pouco atrativo para que o tipo alto não queira
``imitar'' o tipo baixo e capturar renda informacional excessiva.

\begin{tipbox}{Exemplos práticos}

\begin{itemize}
\tightlist
\item
  Passagens aéreas com classes tarifárias (econômica, executiva,
  primeira classe).
\item
  Versões de software (básica, profissional, enterprise).
\item
  Descontos por quantidade em supermercados.
\item
  Planos de streaming com e sem anúncios.
\end{itemize}

\end{tipbox}

\begin{boxbrasilbox}{Box Brasil 15.1 — Discriminação de preços nas passagens aéreas}

O mercado brasileiro de aviação civil, dominado por três companhias
(Latam, Gol e Azul, que juntas detêm mais de 98\% do mercado doméstico
segundo dados da \href{https://www.anac.gov.br}{ANAC}), é um exemplo
rico de discriminação de preços de segundo e terceiro grau.

\textbf{Discriminação de segundo grau (autosseleção):}

As companhias oferecem múltiplas classes tarifárias para o mesmo voo ---
desde tarifas promocionais sem direito a bagagem ou remarcação até
tarifas flexíveis com reembolso integral. Passageiros corporativos, com
alta disposição a pagar e demanda inelástica (viagens de última hora),
autosselecionam-se para tarifas mais caras. Passageiros a lazer, com
demanda elástica e flexibilidade de datas, escolhem as tarifas
promocionais. Essa estrutura replica o menu de contratos descrito na
teoria: o tipo de alta valoração paga mais, mas recebe serviços
adicionais.

\textbf{Discriminação de terceiro grau (segmentação observável):}

Companhias oferecem tarifas diferenciadas para idosos (desconto de 5\%
obrigatório pela ANAC), estudantes e militares. A segmentação por
antecedência de compra (preços mais baixos para compras com 30-60 dias
de antecedência) também funciona como discriminação de terceiro grau,
separando viajantes a lazer (antecedência) de viajantes corporativos
(última hora).

Segundo dados da ANAC, a tarifa aérea média doméstica por km variou
entre R\$ 0,35 e R\$ 0,55 em 2023, mas a dispersão de preços dentro de
um mesmo voo pode facilmente superar 300\% --- evidência direta de
discriminação de preços.

\end{boxbrasilbox}

\subsection{15.7.3 Discriminação de terceiro grau (segmentação de
mercados)}\label{discriminauxe7uxe3o-de-terceiro-grau-segmentauxe7uxe3o-de-mercados}

Na discriminação de terceiro grau, o monopolista divide os consumidores
em \textbf{grupos observáveis} (por exemplo, por idade, localização,
status profissional) e cobra preços diferentes de cada grupo.
Diferentemente da discriminação de segundo grau, aqui a firma observa
diretamente a que grupo cada consumidor pertence e condiciona o preço a
essa observação.

Se há dois mercados com demandas \(p_1(q_1)\) e \(p_2(q_2)\), o
monopolista resolve:

{[} \max\_\{q\_1, q\_2\} ; p\_1(q\_1) \cdot q\_1 + p\_2(q\_2) \cdot q\_2
- C(q\_1 + q\_2) {]}

As condições de primeira ordem são:

{[} RMg\_1(q\_1) = RMg\_2(q\_2) = CMg(q\_1 + q\_2) \label{eq:15.12}
\tag{15.12} {]}

A equação (\ref{eq:15.12}) mostra que o monopolista iguala as receitas
marginais em todos os mercados ao custo marginal. O mercado com demanda
\textbf{mais inelástica} paga o preço \textbf{mais alto}:

{[} \frac{p_1}{p_2} =
\frac{1 - \frac{1}{|\varepsilon_2|}}{1 - \frac{1}{|\varepsilon_1|}}
\label{eq:15.13} \tag{15.13} {]}

Essa equação formaliza a intuição de que o monopolista ``espreme'' mais
os consumidores com poucas alternativas. A razão de preços entre
mercados depende exclusivamente da razão de elasticidades --- e não dos
custos, que são comuns. É por isso que passagens aéreas para executivos
custam mais do que para turistas: não porque custe mais transportar um
executivo, mas porque ele é menos sensível ao preço.

\begin{intuitionbox}{Intuição Econômica}

\textbf{Em uma frase:} Discriminar preços significa cobrar mais de quem
está mais disposto a pagar e menos de quem é mais sensível ao preço.

\textbf{Pense assim:} Pense nas passagens aéreas da Latam ou Gol. O
executivo que compra na véspera paga R\$ 2.000 porque \emph{precisa}
viajar (demanda inelástica). O estudante que compra com dois meses de
antecedência paga R\$ 400 porque tem flexibilidade (demanda elástica). A
companhia aérea não está sendo ``generosa'' com o estudante --- está
maximizando lucro ao extrair mais de quem pode pagar mais e ainda assim
vender para quem pagaria menos.

\textbf{Por que isso importa:} A discriminação de preços pode, em alguns
casos, aumentar a quantidade total vendida e até reduzir a perda de peso
morto --- o que complica o juízo de política concorrencial e exige
análise caso a caso pelo CADE.

\end{intuitionbox}

\begin{boxmundobox}{Box Mundo 15.3 — Discriminação de preços nas ferrovias europeias}

\textbf{Contexto:} O transporte ferroviário europeu oferece um exemplo
notável de discriminação de preços sofisticada, combinando técnicas de
segundo e terceiro grau em escala continental. Operadoras como a
Deutsche Bahn (Alemanha), a SNCF (França) e a Eurostar
(Londres--Paris--Bruxelas) praticam estratégias de precificação que
espelham as do setor aéreo, com uma particularidade: as ferrovias
europeias operam em um regime de concorrência limitada, protegidas pela
infraestrutura fixa das linhas férreas (uma barreira natural clássica).

\textbf{Dados:} Uma passagem Eurostar Londres--Paris pode variar de £39
(compra com 3 meses de antecedência, horário fora de pico) a £400+
(compra de última hora, horário de pico, classe business) --- uma razão
de preços superior a 10:1 para o mesmo trajeto. A Deutsche Bahn oferece
o sistema \emph{BahnCard}: uma tarifa em duas partes (Seção 15.8) na
qual o passageiro paga uma anuidade (€59,90 para BahnCard 25; €244,90
para BahnCard 50) e recebe, em troca, desconto de 25\% ou 50\% em todas
as viagens. Segundo dados da Deutsche Bahn (Relatório Anual 2023), mais
de 5 milhões de passageiros possuem algum tipo de BahnCard --- evidência
de que a tarifa em duas partes é eficaz na extração de excedente de
viajantes frequentes.

\textbf{Análise:} A estrutura tarifária das ferrovias europeias combina
discriminação de terceiro grau (descontos para jovens, idosos e grupos)
com discriminação de segundo grau (autosseleção por antecedência de
compra e classe de serviço). O sistema BahnCard é particularmente
elegante do ponto de vista teórico: funciona exatamente como a tarifa em
duas partes da Seção 15.8, onde a anuidade é a taxa de entrada \(T\) e o
desconto reduz o preço marginal por viagem \(p\). Viajantes frequentes
(alta demanda, tipo A) se autosselecionam para BahnCards mais caras,
revelando sua alta disposição a pagar.

\textbf{Fonte:} Eurostar Group Limited, \emph{Annual Report 2023}.
Deutsche Bahn AG, \emph{Integrated Report 2023}. European Commission, DG
MOVE, \emph{Statistical Pocketbook: EU Transport in Figures 2023}.

\end{boxmundobox}

No gráfico interativo abaixo, alterne entre os três graus de
discriminação e observe como cada tipo afeta a extração de excedente e a
eficiência alocativa.

\figurePlaceholder{/micro-book/graficos/cap15/discriminacao-precos.html}

\textbf{Figura 15.3 --- Discriminação de preços.} Alterne entre
discriminação de 1º grau (perfeita), 2º grau (tarifa em duas partes) e
3º grau (dois mercados). Observe como cada tipo afeta a extração de
excedente e a eficiência alocativa.

O exercício a seguir aplica a discriminação de terceiro grau a um
problema com dois mercados.

\begin{exresolvidobox}{Exercício Resolvido 15.4 — Discriminação de terceiro grau com dois mercados}

\textbf{Enunciado.} Um monopolista vende em dois mercados:
\(p_1 = 120 - 2q_1\) e \(p_2 = 80 - q_2\). O custo marginal é constante
\(CMg = 20\). (a) Com discriminação de 3º grau, encontre preços e
quantidades em cada mercado. (b) Calcule os índices de Lerner e
relacione com as elasticidades. (c) Calcule o lucro total.

\begin{center}\rule{0.5\linewidth}{0.5pt}\end{center}

\textbf{(a)} Receitas marginais: \(RMg_1 = 120 - 4q_1\) e
\(RMg_2 = 80 - 2q_2\).

Igualando ao CMg:

{[} 120 - 4q\_1 = 20 \implies q\_1 = 25, \quad p\_1 = 120 - 50 = 70 {]}

{[} 80 - 2q\_2 = 20 \implies q\_2 = 30, \quad p\_2 = 80 - 30 = 50 {]}

\textbf{(b)} Lerner no mercado 1:
\(L_1 = (70-20)/70 = 50/70 \approx 0{,}714\).

Elasticidade: \(\varepsilon_1 = (-1/2)(70/25) = -1{,}4\), logo
\(L_1 = 1/1{,}4 \approx 0{,}714\). ✓

Lerner no mercado 2: \(L_2 = (50-20)/50 = 30/50 = 0{,}60\).

Elasticidade: \(\varepsilon_2 = (-1)(50/30) = -1{,}667\), logo
\(L_2 = 1/1{,}667 = 0{,}60\). ✓

O mercado 1 tem demanda mais inelástica
(\(|\varepsilon_1| = 1{,}4 < |\varepsilon_2| = 1{,}667\)) e paga o preço
mais alto (\(p_1 = 70 > p_2 = 50\)). Isso confirma a regra: o
monopolista cobra mais do segmento com demanda mais inelástica.

\textbf{(c)} Lucro total:

{[} \pi = (70 - 20)(25) + (50 - 20)(30) = 1.250 + 900 = 2.150 {]}

\end{exresolvidobox}

\figurePlaceholder{/micro-book/graficos/cap15/webr-discriminacao-3grau.html}

\textbf{WebR 15.3 --- Um preço para cada bolso: discriminação de 3º
grau.} O código implementa o ER 15.4 com dois mercados segmentados.
Observe que o mercado com demanda mais inelástica paga o preço mais alto
--- e verifique pela fórmula das elasticidades. Altere as inclinações
\texttt{b1} e \texttt{b2} para ver como a diferença de preços responde à
diferença de elasticidades.

\begin{exresolvidobox}{Exercício Resolvido 15.5 — Comparação dos três regimes: concorrência, monopólio e discriminação perfeita}

\textbf{Enunciado.} Considere demanda \(p = 50 - q\) e \(CMg = 10\)
(constante, sem custos fixos). (a) Encontre \((p, q, \pi)\) sob: (i)
concorrência perfeita; (ii) monopólio com preço único; (iii)
discriminação de primeiro grau. (b) Para cada caso, calcule EC, EP e
PPM.

\begin{center}\rule{0.5\linewidth}{0.5pt}\end{center}

\textbf{(a)}

\textbf{(i) Concorrência perfeita:} \(p = CMg = 10\), \(q^c = 40\).
\(\pi = 0\).

\textbf{(ii) Monopólio:} \(RMg = 50 - 2q = 10 \implies q^m = 20\),
\(p^m = 30\). \(\pi = (30-10)(20) = 400\).

\textbf{(iii) Discriminação perfeita:} \(q^{dp} = q^c = 40\) (quantidade
eficiente). O monopolista cobra a disposição a pagar de cada consumidor.
\(\pi = \frac{1}{2}(50-10)(40) = 800\).

\textbf{(b)}

{\def\LTcaptype{none} % do not increment counter
\begin{longtable}[]{@{}lllll@{}}
\toprule\noalign{}
Regime & EC & EP & PPM & W = EC + EP \\
\midrule\noalign{}
\endhead
\bottomrule\noalign{}
\endlastfoot
Concorrência & 800 & 0 & 0 & 800 \\
Monopólio & 200 & 400 & 200 & 600 \\
Disc. perfeita & 0 & 800 & 0 & 800 \\
\end{longtable}
}

\textbf{Interpretação:} A discriminação perfeita restaura a eficiência
alocativa (\(PPM = 0\)), mas redistribui todo o excedente para o
produtor. O excedente total é o mesmo que em concorrência perfeita
(800), mas o consumidor não fica com nada. Isso mostra que eficiência e
equidade são questões distintas --- um tema central do Capítulo 14
(Teoremas do Bem-Estar).

\end{exresolvidobox}

\begin{center}\rule{0.5\linewidth}{0.5pt}\end{center}

\section{15.8 Pague para Entrar, Pague para Brincar: Tarifas em Duas
Partes}\label{158}

Se você já pagou mensalidade de academia e \emph{depois} pagou por aula
extra, ou comprou o cartão do Sam's Club para ter o ``privilégio'' de
comprar lá dentro, você já foi vítima (voluntária) de uma tarifa em duas
partes. Essa é uma das formas mais elegantes e difundidas de
precificação não-linear. Diferentemente dos esquemas de discriminação
discutidos na seção anterior, que exigem observar ou induzir a revelação
de tipos, a tarifa em duas partes combina simplicidade operacional com
capacidade de extração de excedente, razão pela qual é amplamente
utilizada em setores que vão de telecomunicações a academias de
ginástica, de parques de diversões a clubes de compras.

Uma tarifa em duas partes consiste em uma \textbf{taxa de entrada}
(tarifa fixa \(T\)) e um \textbf{preço por unidade} (\(p\)). O gasto
total do consumidor é \(G = T + p \cdot q\). O papel de cada componente
é distinto: a taxa fixa extrai excedente inframarginal (que o consumidor
obteria ao preço \(p\)); o preço por unidade governa a quantidade
demandada e, portanto, a eficiência alocativa.

\subsection{Consumidores homogêneos}\label{consumidores-homoguxeaneos}

Se todos os consumidores são idênticos, a estratégia ótima é simples e
elegante:

\begin{enumerate}
\def\labelenumi{\arabic{enumi}.}
\tightlist
\item
  Fixar \(p = CMg\) (preço eficiente).
\item
  Fixar \(T = EC\) (capturar todo o excedente do consumidor).
\end{enumerate}

Isso replica o resultado da discriminação de primeiro grau: quantidade
eficiente e extração completa do excedente. A intuição é direta: o preço
unitário garante eficiência (elimina a PPM) e a taxa fixa transfere todo
o benefício para o produtor.

\subsection{Consumidores
heterogêneos}\label{consumidores-heteroguxeaneos}

Com consumidores heterogêneos, o problema é mais complexo e envolve o
trade-off fundamental da discriminação de segundo grau: extrair mais de
cada consumidor versus manter consumidores de baixa demanda no mercado.
Se a taxa de entrada for muito alta, consumidores com baixa demanda
deixam o mercado. O monopolista enfrenta um trade-off entre:

\begin{itemize}
\tightlist
\item
  Cobrar uma taxa alta e perder consumidores marginais.
\item
  Cobrar uma taxa baixa e atender mais consumidores, mas extrair menos
  excedente de cada um.
\end{itemize}

A solução ótima tipicamente envolve \(p > CMg\) e \(T < EC\) do
consumidor com menor demanda. O preço acima do custo marginal gera uma
PPM residual --- o custo da heterogeneidade informacional.

\begin{tipbox}{Exemplos de tarifas em duas partes}

\begin{itemize}
\tightlist
\item
  Clubes de compras como o Sam's Club (anuidade + preço por produto).
\item
  Parques de diversões como a Disney (entrada + valor por brinquedo ---
  ou entrada com tudo incluso, que é um caso limite).
\item
  Planos telefônicos (assinatura mensal + tarifa por minuto/dados).
\item
  Serviços de streaming (assinatura fixa com acesso ao catálogo).
\item
  Academias de ginástica (mensalidade + taxa por aula ou mensalidade
  fixa com acesso livre).
\end{itemize}

\end{tipbox}

\begin{exresolvidobox}{Exercício Resolvido 15.6 — Tarifa em duas partes com dois tipos de consumidores}

\textbf{Enunciado.} Uma academia de ginástica (monopolista local) cobra
tarifa em duas partes \((T, p)\) por aula. Tipo A: \(q_A = 30 - p\), 50
consumidores. Tipo B: \(q_B = 20 - p\), 50 consumidores. \(CMg = 2\) por
aula. (a) Se a academia atende ambos os tipos, encontre \((T^*, p^*)\).
(b) Se atende apenas o Tipo A, encontre \((T^*, p^*)\) e compare os
lucros.

\begin{center}\rule{0.5\linewidth}{0.5pt}\end{center}

\textbf{(a) Atendendo ambos os tipos:}

A taxa fixa é limitada pelo excedente do Tipo B (o mais restrito):
\(T = EC_B(p) = \frac{1}{2}(20-p)^2\).

Lucro total:
\(\Pi = 50 \times [\text{lucro por A}] + 50 \times [\text{lucro por B}]\),
onde o lucro por consumidor é \(T + (p-2) \cdot q_i\).

Após simplificação e otimização (derivando em relação a \(p\)):

{[} \frac{d\Pi}{dp} = 0 \implies p\^{}* = 7 {]}

Logo \(T^* = \frac{1}{2}(20-7)^2 = 84{,}50\). Quantidades: \(q_A = 23\),
\(q_B = 13\).

Lucro por A: \(84{,}50 + (7-2)(23) = 199{,}50\). Lucro por B:
\(84{,}50 + (7-2)(13) = 149{,}50\).

\(\Pi = 50 \times (199{,}50 + 149{,}50) = 17.450\).

\textbf{(b) Atendendo apenas o Tipo A:}

Maximizar extração sobre A: \(p^* = CMg = 2\),
\(T^* = EC_A(2) = \frac{1}{2}(30-2)^2 = 392\).

Verificação: Tipo B não participa, pois
\(EC_B(2) = \frac{1}{2}(18)^2 = 162 < T = 392\). ✓

\(\Pi = 50 \times 392 = 19.600\).

\textbf{Resultado:} Atender apenas o Tipo A (\(\Pi = 19.600\)) é mais
lucrativo do que atender ambos (\(\Pi = 17.450\)). O monopolista prefere
excluir o Tipo B e extrair o máximo do Tipo A --- ilustrando como a
heterogeneidade de demanda pode levar à exclusão ineficiente de
consumidores.

\end{exresolvidobox}

\chapter{15.9--15.10 Regulação e Visão Dinâmica do
Monopólio}\label{regulauxe7uxe3o-e-visuxe3o-dinuxe2mica-do-monopuxf3lio}

\section{15.9 Domar a Fera: Regulação de Monopólios}\label{159}

As seções anteriores documentaram os custos sociais do monopólio ---
perda de peso morto, distorção de preços e quantidades --- e as
estratégias que o monopolista utiliza para maximizar a extração de
excedente. Diante dessas ineficiências, surge naturalmente a questão: o
que a sociedade pode fazer? Em mercados onde o monopólio pode ser
eliminado (por exemplo, removendo barreiras legais desnecessárias), a
promoção da concorrência é o caminho mais direto. Mas em monopólios
naturais --- nos quais a tecnologia torna ineficiente a presença de
múltiplas firmas --- a regulação torna-se o instrumento central de
política pública.

O dilema central da regulação é: como induzir a firma a produzir mais e
cobrar menos, sem eliminar seus incentivos ao investimento e à
eficiência? Como veremos, os diferentes regimes regulatórios representam
soluções distintas para esse dilema, cada uma com seus trade-offs.

\subsection{Regulação pelo custo marginal
(first-best)}\label{regulauxe7uxe3o-pelo-custo-marginal-first-best}

A solução ideal seria impor \(p = CMg\). Entretanto, em um monopólio
natural com custos médios decrescentes, \(CMg < CMe\), de modo que
\(p = CMg\) gera \textbf{prejuízo}. A firma não cobre seus custos fixos
e, a menos que receba um subsídio, sairá do mercado. O governo
precisaria subsidiar a firma, o que acarreta distorções em outros
mercados (custo dos fundos públicos). Estimativas típicas do custo
marginal dos fundos públicos no Brasil variam de 0,20 a 0,50 --- ou
seja, cada R\$ 1 de subsídio impõe um custo adicional de R\$ 0,20 a R\$
0,50 em distorções tributárias em outros mercados.

\subsection{Regulação pelo custo médio
(second-best)}\label{regulauxe7uxe3o-pelo-custo-muxe9dio-second-best}

A alternativa mais comum é impor \(p = CMe\), garantindo lucro zero. A
firma produz mais do que o monopólio sem regulação, mas menos do que o
first-best. Há uma PPM residual, menor que a do monopólio não regulado.
Esse é o modelo conceitual por trás da regulação das distribuidoras de
energia no Brasil pela ANEEL: nas revisões tarifárias periódicas, a
agência calcula uma ``receita requerida'' que cobre os custos eficientes
(não os custos reais, para preservar incentivos) e divide pelo volume
esperado de energia, chegando a uma tarifa que, em princípio, garante
lucro econômico zero.

\subsection{Regulação por teto de preços (price
cap)}\label{regulauxe7uxe3o-por-teto-de-preuxe7os-price-cap}

O regulador fixa um teto de preço que é reajustado periodicamente pela
inflação menos um fator de produtividade:

{[} \Delta p \leq \text{inflação} - X {]}

O fator \(X\) captura ganhos de eficiência esperados. A firma retém
ganhos de produtividade acima de \(X\) até a próxima revisão, o que
incentiva a redução de custos. Entre as revisões, a firma é uma
``maximizadora residual'' --- quanto mais cortar custos, maior seu
lucro. Esse incentivo à eficiência é a principal vantagem do price cap
sobre a regulação por taxa de retorno.

\begin{notebox}{Vantagem do price cap}

O sistema de price cap, introduzido no Reino Unido por Stephen
Littlechild (1983) para a privatização da British Telecom, possui
vantagens informacionais sobre a regulação por taxa de retorno: o
regulador não precisa conhecer detalhadamente a estrutura de custos da
firma. Em contrapartida, pode induzir redução de qualidade se os padrões
de serviço não forem adequadamente monitorados --- razão pela qual
reguladores como a ANATEL complementam o price cap com metas de
qualidade.

\end{notebox}

\subsection{Regulação por taxa de retorno (rate of
return)}\label{regulauxe7uxe3o-por-taxa-de-retorno-rate-of-return}

O regulador permite que a firma obtenha uma taxa de retorno ``justa''
\(s\) sobre o capital investido \(K\):

{[} pq - wL - rK \leq (s - r)K \label{eq:15.14} \tag{15.14} {]}

Averch e Johnson (1962) demonstraram que esse esquema induz a firma a
\textbf{sobreinvestir em capital} para inflar a base de remuneração ---
o chamado \textbf{efeito Averch-Johnson}. A firma substitui trabalho por
capital além do socialmente ótimo, resultando em ineficiência produtiva.
Intuitivamente, se a firma lucra proporcionalmente ao capital empregado,
ela tem incentivo para empregar mais capital do que o necessário ---
como um taxista que compra um carro de luxo desnecessário se pudesse
repassar o custo ao regulador. Esse resultado é um exemplo clássico de
consequências não intencionais da regulação.

\begin{infobox}{Prêmio Nobel — Jean Tirole (2014)}

\textbf{Jean Tirole} (1953--presente) é um economista francês. Obteve o
PhD no MIT sob orientação de Eric Maskin e é professor na Toulouse
School of Economics (TSE), que ajudou a transformar em um dos principais
centros de pesquisa econômica do mundo.

\textbf{Por que ganhou o Nobel:} Premiado por sua análise do poder de
mercado e da regulação. Tirole desenvolveu a teoria moderna de regulação
de monopólios sob informação assimétrica (com Jean-Jacques Laffont),
mostrando como desenhar contratos regulatórios quando o regulador não
conhece os custos da firma. Suas contribuições à organização industrial
unificaram a teoria do monopólio, oligopólio e barreiras à entrada em um
framework analítico coerente baseado na teoria dos jogos.

\textbf{Conexão com este capítulo:} A análise de regulação de monopólio
natural --- regulação por custo marginal, custo médio e tarifas em duas
partes --- apresentada neste capítulo é o ponto de partida da agenda de
pesquisa de Tirole. Sua contribuição foi mostrar que, na prática, o
regulador enfrenta informação assimétrica sobre os custos da firma, o
que torna o problema de regulação fundamentalmente um problema de
desenho de mecanismos --- conectando este capítulo diretamente ao
Capítulo 19 (Informação Assimétrica).

\end{infobox}

\begin{boxbrasilbox}{Box Brasil 15.2 — Monopólios naturais regulados: o setor elétrico e telecomunicações pós-privatização}

O Brasil passou por extenso processo de privatização e regulação de
monopólios naturais nos anos 1990. Dois casos emblemáticos ilustram os
desafios da regulação:

\textbf{Setor Elétrico}

A reestruturação do setor elétrico brasileiro, iniciada em 1995 (Lei
8.987/1995 e Lei 9.074/1995), separou as atividades de geração,
transmissão e distribuição. A distribuição de energia --- um monopólio
natural em cada área de concessão --- é regulada pela ANEEL (Agência
Nacional de Energia Elétrica), criada em 1996. A ANEEL realiza revisões
tarifárias periódicas (a cada 4 ou 5 anos) utilizando o modelo de
\textbf{empresa de referência} e aplica reajustes anuais baseados no
IGP-M. Segundo dados da ANEEL, o Brasil possui 53 distribuidoras
reguladas, atendendo mais de 90 milhões de unidades consumidoras. A
tarifa média residencial brasileira, em 2024, situava-se entre R\$ 0,60
e R\$ 0,90 por kWh (com tributos), uma das mais altas do mundo em
proporção à renda.

\textbf{Telecomunicações}

A privatização do Sistema Telebrás em 1998 transferiu ao setor privado o
monopólio estatal de telecomunicações. A
\href{https://www.anatel.gov.br}{ANATEL} (Agência Nacional de
Telecomunicações), criada pela Lei 9.472/1997, regula concessões e
autorizações. Na telefonia fixa, as concessionárias originais (Oi,
Telefônica/Vivo, Embratel) enfrentaram obrigações de universalização e
controle tarifário (price cap com fator X de produtividade). A telefonia
móvel, operando em regime de autorização, desenvolveu-se em um
oligopólio com quatro operadoras principais (Vivo, Claro, TIM, Oi). A
venda dos ativos móveis da Oi em 2022 para Vivo, Claro e TIM reduziu o
mercado a três grandes operadoras, levantando preocupações
concorrenciais analisadas pelo CADE.

{\def\LTcaptype{none} % do not increment counter
\begin{longtable}[]{@{}
  >{\raggedright\arraybackslash}p{(\linewidth - 4\tabcolsep) * \real{0.3333}}
  >{\raggedright\arraybackslash}p{(\linewidth - 4\tabcolsep) * \real{0.3333}}
  >{\raggedright\arraybackslash}p{(\linewidth - 4\tabcolsep) * \real{0.3333}}@{}}
\toprule\noalign{}
\begin{minipage}[b]{\linewidth}\raggedright
Indicador
\end{minipage} & \begin{minipage}[b]{\linewidth}\raggedright
Setor Elétrico
\end{minipage} & \begin{minipage}[b]{\linewidth}\raggedright
Telecomunicações
\end{minipage} \\
\midrule\noalign{}
\endhead
\bottomrule\noalign{}
\endlastfoot
Agência reguladora & ANEEL (1996) & ANATEL (1997) \\
Modelo regulatório & Empresa de referência + price cap & Price cap com
fator X \\
Nº de distribuidoras/operadoras & 53 distribuidoras & 3 grandes
operadoras (móvel) \\
Universalização & \textasciitilde99,8\% de acesso & \textasciitilde98\%
de cobertura 4G \\
Principal desafio atual & Transição energética e tarifas & 5G e
competição em banda larga \\
\end{longtable}
}

Ambos os setores ilustram o dilema fundamental da regulação de
monopólios naturais: garantir eficiência produtiva e preços acessíveis
sem eliminar os incentivos ao investimento. Para o contexto
institucional e fiscal da regulação de monopólios naturais no Brasil,
ver Giambiagi \& Além (2016).

\end{boxbrasilbox}

\begin{boxbrasilbox}{Box Brasil 15.3 — ANS e a regulação dos planos de saúde como controle de monopólio}

\textbf{Contexto:} A \textbf{Agência Nacional de Saúde Suplementar
(ANS)}, criada pela Lei n.º 9.961/2000, regula o mercado de planos de
saúde no Brasil --- um setor com características de oligopólio
concentrado e fortes assimetrias de informação. Em 2024, o setor atendia
cerca de 51 milhões de beneficiários (aproximadamente 25\% da
população), movimentando mais de R\$ 300 bilhões em receitas anuais. As
quatro maiores operadoras (Hapvida-NotreDame Intermédica, Bradesco
Saúde, SulAmérica e Amil) concentram mais de 35\% do mercado, e em
muitas regiões uma ou duas operadoras detêm posição dominante ---
configurando monopólios ou duopólios locais.

\textbf{Dados:} A ANS utiliza diversos instrumentos regulatórios que
espelham os mecanismos discutidos neste capítulo: (i) \textbf{teto de
reajuste} para planos individuais --- análogo ao \emph{price cap} (Seção
15.9), fixado anualmente pela ANS com base na variação de custos do
setor (reajuste máximo de 6,91\% em 2024); (ii) \textbf{rol de
procedimentos obrigatórios} --- que funciona como regulação de qualidade
mínima, impedindo que operadoras reduzam custos eliminando coberturas
essenciais; (iii) \textbf{ressarcimento ao SUS} --- quando beneficiários
de planos usam o sistema público, a operadora deve reembolsar o SUS,
internalizando parte do custo social. A sinistralidade média do setor
(razão entre despesas assistenciais e receitas de contraprestações)
situou-se em torno de 85\% em 2023, indicando margens operacionais
comprimidas.

\textbf{Análise:} O mercado de planos de saúde combina duas falhas de
mercado estudadas neste livro: \textbf{poder de mercado} (Capítulo 15) e
\textbf{informação assimétrica} (Capítulo 19). A concentração regional
confere às operadoras poder de precificação, enquanto a complexidade dos
contratos gera assimetria entre operadora e beneficiário. A regulação da
ANS enfrenta o dilema clássico de Tirole: sem informação perfeita sobre
os custos das operadoras, o regulador não consegue distinguir
ineficiência operacional de custos genuinamente elevados. O teto de
reajuste para planos individuais protege consumidores cativos, mas pode
desestimular a entrada de novos competidores e levar à deterioração da
qualidade --- exatamente o trade-off previsto pela teoria de regulação
por price cap.

\textbf{Para refletir:} Os planos coletivos (empresariais e por adesão),
que representam mais de 80\% do mercado, não estão sujeitos ao teto de
reajuste da ANS --- seus preços são negociados livremente. Usando a
teoria de discriminação de preços do Capítulo 15, analise: por que a ANS
regula os reajustes individuais mas não os coletivos? A resposta envolve
a diferença de poder de barganha (elasticidade da demanda) entre os dois
segmentos.

\end{boxbrasilbox}

No gráfico interativo abaixo, visualize como diferentes regimes
regulatórios afetam preço, quantidade e bem-estar em um monopólio
natural com custo médio decrescente. Alterne entre os três regimes e
observe os trade-offs.

\figurePlaceholder{/micro-book/graficos/cap15/regulacao-monopolio.html}

\textbf{Figura 15.4 --- Regulação de monopólio natural.} Alterne entre
monopólio não regulado, regulação pelo custo marginal (first-best, com
subsídio) e regulação pelo custo médio (second-best, lucro zero).
Compare preços, quantidades e perda de peso morto em cada regime.

O exercício a seguir aplica os três regimes regulatórios a um caso
numérico de monopólio natural.

\begin{exresolvidobox}{Exercício Resolvido 15.7 — Regulação de monopólio natural}

\textbf{Enunciado.} Uma distribuidora de água (monopólio natural) tem
custo total \(C(q) = 500 + 10q\) e enfrenta demanda \(p = 60 - 0{,}5q\).
(a) Calcule o equilíbrio sem regulação. (b) Calcule o resultado com
regulação \emph{first-best} (\(p = CMg\)) e mostre o prejuízo. (c)
Calcule o resultado \emph{second-best} (\(p = CMe\)).

\begin{center}\rule{0.5\linewidth}{0.5pt}\end{center}

\textbf{(a) Monopólio não regulado:}

\(RMg = 60 - q\), \(CMg = 10\).

{[} 60 - q = 10 \implies q\^{}m = 50, \quad p\^{}m = 60 - 25 = 35 {]}

\(\pi^m = (p^m - CMg) \times q^m - CF = (35-10)(50) - 500 = 1.250 - 500 = 750\).

{[} q\^{}m = 50, \quad p\^{}m = 35, \quad \pi\^{}m = 750 {]}

\textbf{(b) Regulação first-best (\(p = CMg = 10\)):}

{[} 10 = 60 - 0\{,\}5q \implies q\^{}\{fb\} = 100, \quad p\^{}\{fb\} =
10 {]}

\(\pi^{fb} = 10 \times 100 - (500 + 1.000) = 1.000 - 1.500 = -500\).

O prejuízo de R\$ 500 corresponde exatamente ao custo fixo: com
\(p = CMg\), a receita cobre apenas o custo variável. O governo
precisaria subsidiar R\$ 500 para manter a firma no mercado.

\textbf{(c) Regulação second-best (\(p = CMe\)):}

\(CMe = 500/q + 10\). Igualando à demanda inversa:

{[} 60 - 0\{,\}5q = \frac{500}{q} + 10 \implies 50 - 0\{,\}5q =
\frac{500}{q} {]}

{[} 50q - 0\{,\}5q\^{}2 = 500 \implies 0{,}5q\^{}2 - 50q + 500 = 0
\implies q\^{}2 - 100q + 1.000 = 0 {]}

{[} q = \frac{100 \pm \sqrt{10.000 - 4.000}}{2} =
\frac{100 \pm \sqrt{6.000}}{2} = \frac{100 \pm 77{,}46}{2} {]}

Tomando a raiz maior (mais produção): \(q^{sb} \approx 88{,}7\),
\(p^{sb} = 60 - 44{,}4 = 15{,}6\).

Lucro: \(\pi^{sb} = 0\) (por construção). A PPM residual é a área do
triângulo entre \(q^{sb}\) e \(q^{fb}\):

{[} PPM\^{}\{sb\} = \frac{1}{2}(p\^{}\{sb\} - CMg)(q\^{}\{fb\} -
q\^{}\{sb\}) = \frac{1}{2}(15\{,\}6 - 10)(100 - 88\{,\}7) =
\frac{1}{2}(5\{,\}6)(11\{,\}3) \approx 31{,}6 {]}

Compare: \(PPM^m = \frac{1}{2}(35-10)(100-50) = 625\). A regulação
\emph{second-best} reduz a PPM de 625 para apenas 31,6 --- uma melhoria
de 95\%.

\textbf{Interpretação:} A regulação por custo médio é o modelo adotado
pela ANEEL para distribuidoras de energia no Brasil. Ela garante lucro
zero para a concessionária enquanto expande significativamente o acesso,
a um custo de eficiência relativamente baixo.

\end{exresolvidobox}

\figurePlaceholder{/micro-book/graficos/cap15/webr-regulacao.html}

\textbf{WebR 15.4 --- Três regimes de regulação: do monopólio ao
first-best.} O código resolve o ER 15.7 completo: monopólio não
regulado, regulação por custo marginal (first-best, com prejuízo) e
regulação por custo médio (second-best, lucro zero). Compare os três
pontos no gráfico e observe como a PPM cai 95\% com a regulação
second-best. Altere \texttt{CF} para ver como o custo fixo afeta a
viabilidade do first-best.

\begin{center}\rule{0.5\linewidth}{0.5pt}\end{center}

\section{15.10 O Vilão Também Inova: Visões Dinâmicas do
Monopólio}\label{1510}

Hora de virar a mesa. Passamos o capítulo inteiro tratando o monopolista
como o grande vilão da economia --- e com razão: ele produz menos, cobra
mais e destrói excedente. Mas e se o monopólio, visto com outros olhos,
fosse também o motor que faz a economia avançar? A análise conduzida até
aqui é predominantemente estática: comparamos o monopólio com a
concorrência perfeita em um dado momento do tempo. Entretanto, quando
adotamos uma perspectiva dinâmica --- incorporando inovação, progresso
tecnológico e entrada potencial ---, o julgamento sobre os custos e
benefícios do monopólio torna-se mais nuançado. Será que o monopólio é
sempre prejudicial à sociedade, ou pode ele desempenhar um papel
positivo como motor de inovação e progresso técnico? Essa questão,
levantada originalmente por Joseph Schumpeter no contexto da Grande
Depressão, permanece no centro dos debates contemporâneos sobre política
antitruste --- especialmente em setores tecnológicos dominados por
grandes plataformas como Google, Amazon e Meta.

\subsection{Schumpeter e a destruição
criativa}\label{schumpeter-e-a-destruiuxe7uxe3o-criativa}

Joseph Schumpeter (1942) argumentou que o monopólio pode ser socialmente
benéfico quando considerado em perspectiva dinâmica. Segundo essa visão:

\begin{itemize}
\tightlist
\item
  Lucros de monopólio são a \textbf{recompensa pela inovação} e
  constituem o incentivo fundamental para o progresso tecnológico. Sem a
  perspectiva de lucros extraordinários, as firmas não assumiriam os
  riscos e custos do investimento em pesquisa e desenvolvimento.
\item
  A concorrência relevante não é a concorrência de preços estática, mas
  a \textbf{concorrência por inovação} --- a ``destruição criativa''
  pela qual novos produtos e processos tornam obsoletos os anteriores. O
  smartphone destruiu o mercado de câmeras fotográficas compactas; o
  streaming destruiu as videolocadoras; os aplicativos de transporte
  transformaram o mercado de táxis.
\item
  Firmas com poder de mercado têm maiores recursos para investir em P\&D
  e podem proteger os frutos da inovação por mais tempo, o que fortalece
  os incentivos para inovar.
\end{itemize}

A hipótese schumpeteriana gerou extensa literatura empírica, com
resultados ambíguos. A relação entre concentração de mercado e inovação
parece seguir uma curva em U invertido (Aghion et al., 2005): algum grau
de poder de mercado estimula a inovação, mas poder excessivo a inibe. Em
setores com poder de mercado moderado, as firmas inovam mais para
escapar da concorrência; em setores altamente concentrados, o incentivo
diminui porque o incumbente já é lucrativo.

\subsection{Mercados contestáveis}\label{mercados-contestuxe1veis}

Baumol, Panzar e Willig (1982) propuseram a teoria dos \textbf{mercados
contestáveis}: se a entrada e a saída do mercado são livres e sem custos
irrecuperáveis (sunk costs), a mera \textbf{ameaça} de entrada é
suficiente para disciplinar o monopolista. Mesmo um monopolista cobraria
preços competitivos se a entrada fosse perfeitamente livre.

\begin{definitionbox}{Mercado Contestável}

Um mercado é \textbf{perfeitamente contestável} se: (i) não há custos
irrecuperáveis de entrada e saída; (ii) entrantes potenciais têm acesso
à mesma tecnologia do incumbente; (iii) consumidores respondem
instantaneamente a diferenças de preço. Nessas condições, a ameaça de
``hit-and-run'' --- entrada para capturar lucros seguida de saída ---
disciplina o incumbente.

\end{definitionbox}

Na prática, poucos mercados satisfazem essas condições rigorosas. Custos
irrecuperáveis são a norma, não a exceção --- fábricas, equipamentos
especializados, marcas construídas ao longo de anos não podem ser
recuperados à saída. Ainda assim, a teoria da contestabilidade fornece
um benchmark útil e lembra que barreiras à entrada --- não a estrutura
de mercado per se --- são o determinante fundamental do poder de
mercado. Um mercado com uma única firma mas baixas barreiras à entrada
pode ser mais competitivo do que um mercado com muitas firmas mas altas
barreiras.

\begin{intuitionbox}{Intuição Econômica}

\textbf{Em uma frase:} Para Schumpeter, os lucros de monopólio são o
combustível da inovação --- sem a promessa de lucros extraordinários,
ninguém assume o risco de inventar algo novo.

\textbf{Pense assim:} Imagine que desenvolver um novo medicamento custa
R\$ 5 bilhões e leva 10 anos. Se, ao final do processo, qualquer
concorrente pudesse copiar a fórmula e vender a preço de custo, nenhuma
empresa investiria na pesquisa. A patente --- que cria um monopólio
temporário --- é a ``cenoura'' que justifica o investimento. A perda de
peso morto durante a vigência da patente é o ``preço'' que a sociedade
paga pela inovação futura.

\textbf{Por que isso importa:} Esse argumento é central no debate sobre
a duração de patentes farmacêuticas e sobre a regulação de big techs:
punir poder de mercado de forma muito agressiva pode sufocar a inovação
que gera bem-estar no longo prazo. A política antitruste ótima precisa
equilibrar eficiência estática (eliminar PPM) e eficiência dinâmica
(preservar incentivos à inovação).

\end{intuitionbox}

\begin{exresolvidobox}{Exercício Resolvido 15.8 — Inovação sob monopólio vs. concorrência (hipótese schumpeteriana)}

\textbf{Enunciado.} Um monopolista farmacêutico pode investir \(I\) em
P\&D, reduzindo seu custo marginal de \(c_0 = 40\) para
\(c_1 = 40 - \sqrt{I}\) (com \(I \leq 1.600\)). A demanda é
\(p = 100 - q\). (a) Qual o investimento ótimo em P\&D? (b) Se o mercado
fosse competitivo e a inovação fosse imediatamente imitável (spillover
total), qual seria o investimento? (c) Comente sobre a hipótese
schumpeteriana.

\begin{center}\rule{0.5\linewidth}{0.5pt}\end{center}

\textbf{(a)} Com CMg = \(40 - \sqrt{I}\), o monopolista resolve:

\(RMg = 100 - 2q = 40 - \sqrt{I} \implies q^m = 30 + \frac{\sqrt{I}}{2}\).

\(p^m = 100 - q^m = 70 - \frac{\sqrt{I}}{2}\).

\(\pi = (p^m - CMg) \cdot q^m - I = \left(30 + \frac{\sqrt{I}}{2}\right)^2 - I\).

Seja \(s = \sqrt{I}\):

\(\pi(s) = \left(30 + \frac{s}{2}\right)^2 - s^2 = 900 + 30s + \frac{s^2}{4} - s^2 = 900 + 30s - \frac{3s^2}{4}\).

\(\frac{d\pi}{ds} = 30 - \frac{3s}{2} = 0 \implies s^* = 20 \implies I^* = 400\).

Resultado: CMg cai para 20, \(q^m = 40\), \(p^m = 60\), \(\pi = 1.200\)
(líquido de I).

\textbf{(b)} Sob concorrência com spillover total, qualquer redução de
custo é imediatamente copiada. Como \(p = CMg\) e \(\pi = 0\), nenhuma
firma recupera o investimento. Logo \(I^* = 0\).

\textbf{(c)} Este resultado ilustra a hipótese schumpeteriana: o
monopolista investe 400 em P\&D (reduzindo o custo marginal pela
metade), enquanto o mercado competitivo não investe nada. O poder de
mercado é a ``recompensa'' que viabiliza a inovação. Contudo, o
argumento depende criticamente da hipótese de spillover total; na
prática, patentes e segredos industriais permitem que firmas
competitivas também se apropriem de inovações, atenuando a vantagem
schumpeteriana do monopólio.

\end{exresolvidobox}

\figurePlaceholder{/micro-book/graficos/cap15/webr-inovacao-schumpeter.html}

\textbf{WebR 15.5 --- O vilão também inova: P\&D sob monopólio.} O
código implementa o ER 15.8: o monopolista escolhe quanto investir em
P\&D para reduzir seu custo marginal. Compare o investimento ótimo
(\(I^* = 400\)) com o investimento zero da concorrência com spillover
total. Altere \texttt{c0} e observe como o incentivo à inovação muda com
o nível inicial de custo.

\begin{center}\rule{0.5\linewidth}{0.5pt}\end{center}

Abrimos o capítulo com um único vendedor, zero concorrência e um sorriso
no rosto. Ao longo da análise, vimos que esse sorriso tem consequências:
preço acima do custo marginal, produção abaixo do socialmente desejável
e uma perda de peso morto que ninguém captura. Mas também vimos que o
monopólio não é um vilão unidimensional --- Schumpeter nos lembrou que,
sem a promessa de lucros extraordinários, talvez não houvesse iPhone,
Netflix nem vacina de mRNA.

\emph{O monopolista reina sozinho. Mas e quando há dois? Três? Poucos
--- e isso muda tudo. É o próximo capítulo.}

\chapter{Exercícios e ANPEC --- Capítulo
15}\label{exercuxedcios-e-anpec-capuxedtulo-15}

\begin{tipbox}{Dados para exercicios empiricos}

\begin{itemize}
\tightlist
\item
  \textbf{ANEEL (tarifas de energia):}
  \href{https://www.aneel.gov.br/dados-abertos}{aneel.gov.br/dados-abertos}
  --- Tarifas por distribuidora para calcular indices de Lerner e
  avaliar regulacao de monopolio natural (price cap vs.~taxa de
  retorno).
\item
  \textbf{CADE (atos de concentracao):}
  \href{https://www.gov.br/cade/pt-br}{cade.gov.br} --- Pareceres
  publicos sobre monopolio e abuso de posicao dominante. Analise o caso
  Ambev ou da Souza Cruz.
\item
  \textbf{ANS (planos de saude):}
  \href{https://dados.gov.br/dados/organizacoes/visualizar/agencia-nacional-de-saude-suplementar-ans}{ans.gov.br/dados-abertos}
  --- Dados de concentracao de mercado por municipio para identificar
  monopolios regionais em saude.
\end{itemize}

\end{tipbox}

\section{Revisão Rápida}\label{revisuxe3o-ruxe1pida-1}

Teste seu entendimento dos conceitos centrais deste capítulo.

\begin{infobox}{1. O monopolista maximiza lucro igualando receita marginal ao custo marginal. A receita marginal é menor que o preço porque:}

\begin{itemize}
\tightlist
\item
  \begin{enumerate}
  \def\labelenumi{(\alph{enumi})}
  \tightlist
  \item
    O monopolista opera com custos mais altos que firmas competitivas
  \end{enumerate}
\item
  \begin{enumerate}
  \def\labelenumi{(\alph{enumi})}
  \setcounter{enumi}{1}
  \tightlist
  \item
    Para vender uma unidade adicional, o monopolista deve reduzir o
    preço de todas as unidades, gerando um efeito negativo sobre a
    receita das unidades inframarginais
  \end{enumerate}
\item
  \begin{enumerate}
  \def\labelenumi{(\alph{enumi})}
  \setcounter{enumi}{2}
  \tightlist
  \item
    O governo taxa a receita do monopolista
  \end{enumerate}
\item
  \begin{enumerate}
  \def\labelenumi{(\alph{enumi})}
  \setcounter{enumi}{3}
  \tightlist
  \item
    A demanda pelo bem do monopolista é perfeitamente elástica
  \end{enumerate}
\end{itemize}

??? success ``Resposta'' \textbf{(b)} Diferentemente da firma
competitiva (que é tomadora de preço), o monopolista enfrenta toda a
curva de demanda. Vender mais exige reduzir o preço, o que diminui a
receita das unidades já vendidas. Assim,
\(\text{RMg} = p + q \cdot dp/dq = p(1 + 1/\varepsilon_d) < p\). A
alternativa (d) descreve concorrência perfeita, não monopólio.

\end{infobox}

\begin{infobox}{2. O índice de Lerner mede o poder de mercado como:}

\begin{itemize}
\tightlist
\item
  \begin{enumerate}
  \def\labelenumi{(\alph{enumi})}
  \tightlist
  \item
    \(L = (p - \text{CMe})/p\)
  \end{enumerate}
\item
  \begin{enumerate}
  \def\labelenumi{(\alph{enumi})}
  \setcounter{enumi}{1}
  \tightlist
  \item
    \(L = (p - \text{CMg})/p = -1/\varepsilon_d\)
  \end{enumerate}
\item
  \begin{enumerate}
  \def\labelenumi{(\alph{enumi})}
  \setcounter{enumi}{2}
  \tightlist
  \item
    \(L = \text{lucro total}/\text{receita total}\)
  \end{enumerate}
\item
  \begin{enumerate}
  \def\labelenumi{(\alph{enumi})}
  \setcounter{enumi}{3}
  \tightlist
  \item
    \(L = \text{CMg}/p\)
  \end{enumerate}
\end{itemize}

??? success ``Resposta'' \textbf{(b)} O índice de Lerner
\(L = (p - \text{CMg})/p\) mede o markup percentual sobre o custo
marginal. Em equilíbrio, \(L = -1/\varepsilon_d\): quanto menos elástica
a demanda, maior o poder de mercado. Em concorrência perfeita, \(L = 0\)
(\(p = \text{CMg}\)). A alternativa (a) usa custo médio em vez de
marginal; (c) é uma medida de lucratividade, não de poder de mercado.

\end{infobox}

\begin{infobox}{3. Na discriminação de preços de primeiro grau (perfeita), o monopolista:}

\begin{itemize}
\tightlist
\item
  \begin{enumerate}
  \def\labelenumi{(\alph{enumi})}
  \tightlist
  \item
    Cobra o mesmo preço de todos os consumidores
  \end{enumerate}
\item
  \begin{enumerate}
  \def\labelenumi{(\alph{enumi})}
  \setcounter{enumi}{1}
  \tightlist
  \item
    Cobra de cada consumidor exatamente sua disposição a pagar,
    extraindo todo o excedente do consumidor
  \end{enumerate}
\item
  \begin{enumerate}
  \def\labelenumi{(\alph{enumi})}
  \setcounter{enumi}{2}
  \tightlist
  \item
    Divide os consumidores em dois grupos e cobra preços diferentes
  \end{enumerate}
\item
  \begin{enumerate}
  \def\labelenumi{(\alph{enumi})}
  \setcounter{enumi}{3}
  \tightlist
  \item
    Reduz a quantidade abaixo do nível competitivo
  \end{enumerate}
\end{itemize}

??? success ``Resposta'' \textbf{(b)} Na discriminação perfeita, o
monopolista conhece a disposição a pagar de cada consumidor e cobra
preços personalizados. O excedente do consumidor é zero e todo o
excedente vai para o monopolista. Curiosamente, a quantidade produzida é
eficiente (igual à competitiva) --- não há peso morto, embora a
distribuição seja extremamente desigual. A alternativa (d) descreve o
monopólio sem discriminação.

\end{infobox}

\begin{infobox}{4. Um monopólio natural ocorre quando:}

\begin{itemize}
\tightlist
\item
  \begin{enumerate}
  \def\labelenumi{(\alph{enumi})}
  \tightlist
  \item
    A firma detém uma patente sobre o produto
  \end{enumerate}
\item
  \begin{enumerate}
  \def\labelenumi{(\alph{enumi})}
  \setcounter{enumi}{1}
  \tightlist
  \item
    O custo médio é decrescente em toda a faixa relevante de demanda, de
    modo que uma única firma atende o mercado a custo menor do que duas
    ou mais
  \end{enumerate}
\item
  \begin{enumerate}
  \def\labelenumi{(\alph{enumi})}
  \setcounter{enumi}{2}
  \tightlist
  \item
    O governo concede exclusividade legal à firma
  \end{enumerate}
\item
  \begin{enumerate}
  \def\labelenumi{(\alph{enumi})}
  \setcounter{enumi}{3}
  \tightlist
  \item
    A firma possui o único recurso natural necessário à produção
  \end{enumerate}
\end{itemize}

??? success ``Resposta'' \textbf{(b)} O monopólio natural surge de
economias de escala extremas: subaditividade de custos na faixa
relevante. Duplicar a infraestrutura (ex.: rede elétrica, saneamento)
seria socialmente custoso. A alternativa (a) descreve monopólio por
patente; (c) descreve monopólio legal; (d) descreve monopólio por
controle de recurso --- todos são fontes de monopólio, mas não
`natural'.

\end{infobox}

\begin{infobox}{5. A regulação por preço-teto (price cap) de um monopólio natural, fixando $p = \text{CMe}$, resulta em:}

\begin{itemize}
\tightlist
\item
  \begin{enumerate}
  \def\labelenumi{(\alph{enumi})}
  \tightlist
  \item
    Lucro econômico positivo para a firma
  \end{enumerate}
\item
  \begin{enumerate}
  \def\labelenumi{(\alph{enumi})}
  \setcounter{enumi}{1}
  \tightlist
  \item
    Produção eficiente ao nível de custo marginal
  \end{enumerate}
\item
  \begin{enumerate}
  \def\labelenumi{(\alph{enumi})}
  \setcounter{enumi}{2}
  \tightlist
  \item
    Lucro econômico zero (second-best), mas quantidade inferior ao ótimo
    de primeiro melhor (\(p = \text{CMg}\))
  \end{enumerate}
\item
  \begin{enumerate}
  \def\labelenumi{(\alph{enumi})}
  \setcounter{enumi}{3}
  \tightlist
  \item
    Prejuízo para a firma, exigindo subsídio
  \end{enumerate}
\end{itemize}

??? success ``Resposta'' \textbf{(c)} Fixar \(p = \text{CMe}\) garante
viabilidade financeira (lucro zero) mas a quantidade é menor que a de
\(p = \text{CMg}\). A alternativa (d) descreve o que ocorre quando se
tenta impor \(p = \text{CMg}\) em monopólio natural com custos médios
decrescentes --- nesse caso, \(\text{CMg} < \text{CMe}\) e a firma tem
prejuízo. A regulação por CMe é o second-best que evita esse problema.

\end{infobox}

\begin{center}\rule{0.5\linewidth}{0.5pt}\end{center}

\section{Resumo do Capítulo}\label{resumo-do-capuxedtulo-1}

\begin{itemize}
\tightlist
\item
  O monopólio surge quando barreiras à entrada (legais, naturais ou
  estratégicas) impedem a concorrência, permitindo que uma única firma
  fixe preços acima do custo marginal e obtenha lucros persistentes.
\item
  O monopolista maximiza lucro igualando receita marginal ao custo
  marginal (RMg = CMg), operando sempre na faixa elástica da demanda. O
  markup sobre o custo marginal é inversamente proporcional à
  elasticidade-preço, conforme medido pelo índice de Lerner.
\item
  O monopólio gera ineficiência alocativa: produz menos e cobra mais do
  que o resultado competitivo, criando uma perda de peso morto
  (triângulo de Harberger) que não é capturada por nenhum agente. O
  custo social total pode ser significativamente maior quando se incluem
  gastos com \emph{rent-seeking}.
\item
  A discriminação de preços (primeiro, segundo e terceiro graus) permite
  ao monopolista extrair mais excedente do consumidor, podendo, em
  alguns casos, aumentar a quantidade produzida e reduzir a perda de
  peso morto.
\item
  A regulação de monopólios naturais enfrenta um dilema fundamental: a
  precificação pelo custo marginal gera prejuízos (first-best
  impraticável), enquanto a precificação pelo custo médio ou por teto de
  preços (price cap) representa soluções de second-best com trade-offs
  entre eficiência e incentivos ao investimento.
\item
  As visões dinâmicas de Schumpeter (destruição criativa) e de Baumol
  (mercados contestáveis) relativizam os custos do monopólio, sugerindo
  que lucros de monopólio podem incentivar inovação e que a ameaça de
  entrada pode disciplinar o incumbente.
\end{itemize}

\section{Conceitos-Chave}\label{conceitos-chave-1}

{\def\LTcaptype{none} % do not increment counter
\begin{longtable}[]{@{}
  >{\raggedright\arraybackslash}p{(\linewidth - 2\tabcolsep) * \real{0.4762}}
  >{\raggedright\arraybackslash}p{(\linewidth - 2\tabcolsep) * \real{0.5238}}@{}}
\toprule\noalign{}
\begin{minipage}[b]{\linewidth}\raggedright
Conceito
\end{minipage} & \begin{minipage}[b]{\linewidth}\raggedright
Definição
\end{minipage} \\
\midrule\noalign{}
\endhead
\bottomrule\noalign{}
\endlastfoot
Poder de mercado & Capacidade de uma firma fixar preços acima do custo
marginal e obter lucros econômicos persistentes \\
Monopólio natural & Mercado em que a função de custo é subaditiva, de
modo que uma única firma produz a qualquer quantidade a custo menor do
que duas ou mais firmas \\
Barreiras à entrada & Fatores (legais, naturais ou estratégicos) que
impedem a entrada de concorrentes e protegem os lucros do monopolista \\
Índice de Lerner & Medida de poder de mercado dada por
\(L = (p - CMg)/p = 1/|\varepsilon_p|\); varia entre 0 (concorrência
perfeita) e 1 \\
Perda de peso morto (triângulo de Harberger) & Redução no excedente
total causada pela restrição de produção do monopolista, que não é
capturada por nenhum agente \\
Discriminação de preços & Prática de cobrar preços diferentes de
consumidores diferentes (ou por unidades diferentes), classificada em
primeiro, segundo e terceiro graus \\
Tarifa em duas partes & Esquema de precificação com taxa fixa de entrada
mais preço por unidade consumida \\
Regulação por price cap & Regime regulatório em que o preço é reajustado
pela inflação menos um fator de produtividade (X), incentivando redução
de custos \\
Efeito Averch-Johnson & Tendência de firmas reguladas por taxa de
retorno a sobreinvestir em capital para inflar a base de remuneração \\
Mercado contestável & Mercado em que a ameaça de entrada e saída sem
custos irrecuperáveis disciplina o incumbente, mesmo que haja apenas uma
firma \\
\end{longtable}
}

\textbf{Tabela 15.2 --- Conceitos-chave.}

\begin{center}\rule{0.5\linewidth}{0.5pt}\end{center}

\section{Exercícios}\label{exercuxedcios-1}

\begin{exresolvidobox}{Exercício 15.1}

Considere um monopolista com custo total \(C(q) = 100 + 10q\)
enfrentando demanda \(p = 50 - 2q\).

\begin{enumerate}
\def\labelenumi{\alph{enumi})}
\item
  Encontre o preço, a quantidade e o lucro de monopólio.
\item
  Calcule o índice de Lerner e a elasticidade-preço da demanda no ponto
  de equilíbrio.
\item
  Calcule a perda de peso morto e compare com o excedente total sob
  concorrência perfeita.
\end{enumerate}

→ Ver solução (Cap. 15)

\end{exresolvidobox}

\begin{exresolvidobox}{Exercício 15.2}

Um monopolista atende dois mercados segmentados com demandas
\(p_1 = 100 - q_1\) e \(p_2 = 60 - 2q_2\). O custo marginal é constante
e igual a 20.

\begin{enumerate}
\def\labelenumi{\alph{enumi})}
\item
  Encontre os preços e quantidades ótimos em cada mercado com
  discriminação de terceiro grau.
\item
  Calcule os índices de Lerner em cada mercado e relacione com as
  elasticidades.
\item
  Compare o lucro com discriminação ao lucro sem discriminação (preço
  uniforme ótimo).
\end{enumerate}

→ Ver solução (Cap. 15)

\end{exresolvidobox}

\begin{exresolvidobox}{Exercício 15.3}

O governo impõe um imposto unitário \(t = 4\) sobre um monopolista com
custo marginal constante \(c = 10\) e demanda \(p = 30 - q\).

\begin{enumerate}
\def\labelenumi{\alph{enumi})}
\item
  Calcule preço, quantidade e lucro antes e depois do imposto.
\item
  Qual fração do imposto é repassada ao consumidor?
\item
  Compare a receita fiscal com a variação na perda de peso morto.
\end{enumerate}

→ Ver solução (Cap. 15)

\end{exresolvidobox}

\begin{exresolvidobox}{Exercício 15.4}

Um monopolista natural tem custo total \(C(q) = 1000 + 5q\). A demanda é
\(p = 45 - q\).

\begin{enumerate}
\def\labelenumi{\alph{enumi})}
\item
  Calcule o resultado do monopólio não regulado.
\item
  Calcule o resultado sob regulação de custo marginal (\(p = CMg\)).
  Mostre que a firma tem prejuízo.
\item
  Calcule o resultado sob regulação de custo médio (\(p = CMe\)) e a PPM
  residual.
\end{enumerate}

→ Ver solução (Cap. 15)

\end{exresolvidobox}

\begin{exresolvidobox}{Exercício 15.5}

Um monopolista pode adotar uma tarifa em duas partes \((T, p)\) para
atender dois tipos de consumidores. O tipo 1 tem demanda
\(q_1 = 20 - p\) e o tipo 2 tem demanda \(q_2 = 10 - p\). Há 100
consumidores de cada tipo. O custo marginal é \(c = 2\).

\begin{enumerate}
\def\labelenumi{\alph{enumi})}
\item
  Se o monopolista usa preço uniforme, qual é o preço ótimo?
\item
  Se usa tarifa em duas partes atendendo ambos os tipos, encontre
  \((T^*, p^*)\).
\item
  Se usa tarifa em duas partes atendendo apenas o tipo 1 (excluindo o
  tipo 2), encontre \((T^*, p^*)\) e compare os lucros.
\end{enumerate}

→ Ver solução (Cap. 15)

\end{exresolvidobox}

\begin{exresolvidobox}{Exercício 15.6}

Um monopolista enfrenta a demanda \(p = 80 - 0{,}5q\) e tem custo total
\(C(q) = 200 + 20q\).

\begin{enumerate}
\def\labelenumi{\alph{enumi})}
\item
  Encontre a quantidade, o preço e o lucro de monopólio.
\item
  Calcule o índice de Lerner e verifique a relação
  \(L = 1/|\varepsilon_d|\).
\item
  Se a demanda mudar para \(p = 80 - 2q\) (mantendo o mesmo custo),
  recalcule o equilíbrio de monopólio e compare o poder de mercado nos
  dois casos.
\end{enumerate}

→ Ver solução (Cap. 15)

\end{exresolvidobox}

\begin{exresolvidobox}{Exercício 15.7}

Um monopolista pode praticar discriminação de preços. A demanda
individual de cada consumidor é \(q = 10 - p\) e há 50 consumidores
idênticos. O custo marginal é constante \(c = 2\).

\begin{enumerate}
\def\labelenumi{\alph{enumi})}
\item
  Calcule o lucro sob preço uniforme de monopólio (sem discriminação).
\item
  Calcule o lucro sob discriminação de primeiro grau (preços
  personalizados perfeitos).
\item
  Compare os excedentes do consumidor e do produtor nos dois casos. Há
  perda de peso morto sob discriminação perfeita?
\end{enumerate}

→ Ver solução (Cap. 15)

\end{exresolvidobox}

\begin{exresolvidobox}{Exercício 15.8}

Um monopolista natural opera com custo total \(C(q) = 500 + 2q\) e
demanda \(p = 52 - q\).

\begin{enumerate}
\def\labelenumi{\alph{enumi})}
\item
  Calcule o resultado do monopólio não regulado (preço, quantidade,
  lucro).
\item
  Calcule o resultado sob regulação de custo marginal (\(p = CMg\)). A
  firma é viável?
\item
  Calcule o resultado sob regulação de custo médio (\(p = CMe\)) e
  determine o subsídio necessário (se houver) para viabilizar a firma em
  cada regime regulatório.
\end{enumerate}

→ Ver solução (Cap. 15)

\end{exresolvidobox}

\begin{exresolvidobox}{Exercício 15.9}

Considere um monopolista com custo marginal constante \(c = 10\) e
demanda \(p = 100 - 2q\).

\begin{enumerate}
\def\labelenumi{\alph{enumi})}
\item
  Calcule o equilíbrio de monopólio e a perda de peso morto (DWL).
\item
  Calcule a DWL como fração da receita do monopolista e como fração do
  excedente total competitivo.
\item
  Suponha que o custo marginal suba para \(c' = 30\). Recalcule a DWL e
  mostre como ela varia com o nível do custo marginal. Interprete.
\end{enumerate}

→ Ver solução (Cap. 15)

\end{exresolvidobox}

\begin{exresolvidobox}{Exercício 15.10}

Um monopolista de parque de diversões enfrenta consumidores idênticos,
cada um com demanda por brinquedos \(q = 24 - 2p\), onde \(p\) é o preço
por brinquedo. O custo marginal de cada brinquedo é \(c = 2\). O
monopolista usa uma tarifa em duas partes: cobra uma entrada fixa \(T\)
e um preço por brinquedo \(p\).

\begin{enumerate}
\def\labelenumi{\alph{enumi})}
\item
  Encontre a tarifa ótima em duas partes \((T^*, p^*)\) e o lucro por
  consumidor.
\item
  Compare com o lucro que o monopolista obteria usando apenas preço
  uniforme (sem entrada).
\item
  Suponha agora que haja dois tipos de consumidores: tipo A com demanda
  \(q_A = 24 - 2p\) e tipo B com demanda \(q_B = 12 - 2p\), em
  proporções iguais. Encontre a tarifa em duas partes ótima que atenda
  ambos os tipos.
\end{enumerate}

→ Ver solução (Cap. 15)

\end{exresolvidobox}

\begin{center}\rule{0.5\linewidth}{0.5pt}\end{center}

\section{Vem, ANPEC!}\label{vem-anpec-1}

Pratique com questões reais do Exame Nacional da ANPEC (Associação
Nacional dos Centros de Pós-Graduação em Economia). As questões seguem o
formato oficial: cinco itens (0 a 4) a serem julgados como Verdadeiro
(V) ou Falso (F).

\begin{infobox}{ANPEC 2016 — Questão 08}

Com relação a poder de mercado, monopólio e monopsônio, é correto
afirmar que:

\textbf{(0)} Se o custo marginal da empresa em monopólio for constante e
igual a \$ 10, e a elasticidade-preço da demanda for igual a \(-2\), o
preço do produto será \$ 20;

\textbf{(1)} Quanto menos elástica for a curva de demanda de uma
empresa, maior poder de monopólio ela terá;

\textbf{(2)} O poder de monopsônio permite à empresa compradora adquirir
o produto por um preço inferior ao competitivo;

\textbf{(3)} Quanto menos elástica for a curva de oferta, maior será a
diferença entre a despesa marginal e a despesa média, e maior será o
poder de monopsônio do comprador;

\textbf{(4)} Chama-se captura de renda ao processo pelo qual as empresas
aplicam recursos produtivos em atividade de lobby para adquirir, manter
ou exercer seu poder de monopólio.

??? success ``Gabarito'' \textbf{Gabarito oficial: V-V-V-V-V}

\begin{verbatim}
**(0) VERDADEIRO.** Pela regra de markup: $p = CMg/(1 - 1/|\varepsilon|) = 10/(1 - 1/2) = 10/0{,}5 = 20$.

**(1) VERDADEIRO.** O índice de Lerner é $L = 1/|\varepsilon|$. Quanto menos elástica (menor $|\varepsilon|$), maior o Lerner e maior o poder de monopólio.

**(2) VERDADEIRO.** O monopsonista iguala a despesa marginal (que excede o preço de oferta) ao valor do produto marginal do insumo, contratando menos unidades a um preço inferior ao que prevaleceria em concorrência.

**(3) VERDADEIRO.** Com oferta $w = w(L)$, a despesa marginal é $DM = w + L \cdot dw/dL$. A diferença $DM - w = L \cdot dw/dL$ é maior quanto mais inclinada (menos elástica) for a curva de oferta, pois $dw/dL$ é maior.

**(4) VERDADEIRO.** O comportamento descrito — aplicar recursos em lobby para obter ou manter poder de monopólio — corresponde ao conceito de *rent-seeking* (busca/captura de rendas), introduzido por Tullock (1967) e formalizado por Krueger (1974).
\end{verbatim}

\end{infobox}

\begin{infobox}{ANPEC 2015 — Questão 09}

Julgue as afirmações relativas à Teoria do Monopólio:

\textbf{(0)} Uma firma monopolista, que opera com várias fábricas, aloca
sua produção entre elas de forma a igualar o custo médio em cada uma das
fábricas.

\textbf{(1)} Uma firma capaz de discriminação de preços de terceiro grau
obtém lucro maior ou igual, em comparação com a situação na qual ela não
fosse capaz de discriminar.

\textbf{(2)} Uma firma monopolista, que se depara com curva de demanda
com elasticidade constante, é indiferente sobre a quantidade produzida.

\textbf{(3)} Para obter eficiência econômica, o regulador de um
monopólio natural deve escolher a alocação que minimiza o custo médio
unitário da firma.

\textbf{(4)} Se o monopolista for capaz de realizar discriminação de
preços de primeiro grau, a alocação de recursos será eficiente em termos
paretianos.

??? success ``Gabarito'' \textbf{Gabarito oficial: F-V-F-F-V}

\begin{verbatim}
**(0) FALSO.** A firma multi-planta minimiza custos igualando os **custos marginais** (não os custos médios) entre fábricas: $CMg_1(q_1) = CMg_2(q_2) = \ldots$. Essa é a condição de Lagrange para a minimização de custo total sujeita a uma meta de produção.

**(1) VERDADEIRO.** A capacidade de discriminar nunca reduz o lucro, pois o monopolista pode sempre replicar o preço uniforme como caso especial (cobrando o mesmo preço em todos os segmentos). Logo, $\pi_{\text{discriminação}} \geq \pi_{\text{uniforme}}$.

**(2) FALSO.** Com demanda isoelástica ($q = Ap^{\varepsilon}$) e custo marginal positivo, a regra de markup $p = CMg/(1 - 1/|\varepsilon|)$ determina univocamente o preço e a quantidade ótimos. O monopolista não é indiferente.

**(3) FALSO.** Eficiência econômica requer $p = CMg$ (*first-best*), que é a condição em que o benefício marginal do consumidor iguala o custo marginal da produção. Minimizar o custo médio gera uma alocação diferente e, em geral, ineficiente.

**(4) VERDADEIRO.** Com discriminação perfeita, o monopolista produz a quantidade eficiente ($q^c$, onde $p = CMg$), eliminando toda a perda de peso morto. A alocação é Pareto-eficiente, embora todo o excedente seja apropriado pelo produtor ($EC = 0$).
\end{verbatim}

\end{infobox}

\begin{infobox}{ANPEC 2018 — Questão 08}

Com relação à análise de discriminação de preços, indique quais das
afirmativas a seguir são verdadeiras e quais são falsas:

\textbf{(0)} Na discriminação de preços de terceiro grau, a receita
marginal deve ser igual para os diferentes grupos de consumidores e
igual ao custo marginal;

\textbf{(1)} Na discriminação de preços de terceiro grau, o preço mais
elevado será cobrado dos consumidores com demanda mais elástica;

\textbf{(2)} A discriminação de preços intertemporal cobra preços mais
elevados dos consumidores mais impacientes, reduzindo o preço mais tarde
para incentivar o consumo de massa;

\textbf{(3)} A tarifa em duas partes é eficaz quando as demandas dos
consumidores são relativamente homogêneas;

\textbf{(4)} Quando as demandas são heterogêneas e estão positivamente
correlacionadas, a prática do pacote pode ser uma técnica eficaz para a
fixação de preços.

??? success ``Gabarito'' \textbf{Gabarito oficial: V-F-V-V-F}

\begin{verbatim}
**(0) VERDADEIRO.** A condição de ótimo na discriminação de 3º grau é $RMg_1 = RMg_2 = \ldots = CMg$. O monopolista iguala as receitas marginais de todos os segmentos ao custo marginal.

**(1) FALSO.** O preço mais elevado é cobrado do grupo com demanda **mais inelástica** (menor $|\varepsilon|$), não mais elástica. Pela relação $p_i = CMg/(1 - 1/|\varepsilon_i|)$, menor elasticidade implica maior preço.

**(2) VERDADEIRO.** Na discriminação intertemporal, consumidores impacientes (alta disposição a pagar, demanda inelástica no tempo) compram imediatamente ao preço alto. O preço cai ao longo do tempo para capturar consumidores com menor disposição a pagar. Exemplos: livros em capa dura seguidos de edição de bolso, eletrônicos no lançamento.

**(3) VERDADEIRO.** Com consumidores homogêneos, a tarifa em duas partes é particularmente eficaz: fixa-se $p = CMg$ e $T = EC$ do consumidor representativo, extraindo todo o excedente com eficiência alocativa plena.

**(4) FALSO.** O *bundling* (venda em pacote) é mais eficaz quando as demandas dos bens estão **negativamente** correlacionadas, pois isso reduz a dispersão da disposição a pagar pelo pacote, permitindo melhor extração de excedente. Com correlação positiva, o pacote não reduz a heterogeneidade e o *bundling* é pouco eficaz.
\end{verbatim}

\end{infobox}

\begin{center}\rule{0.5\linewidth}{0.5pt}\end{center}

\begin{tipbox}{Exercício com IA}

\textbf{IA.3 --- Deadweight Loss do Monopólio}

Peça à IA para calcular o deadweight loss de um monopólio com demanda
\(P = 100 - 2Q\) e \(CMg = 20\). Depois peça para ela comparar com
discriminação de preços de 1º grau. A IA distinguiu corretamente os
excedentes? Desenhe os gráficos para confirmar.

→ Ver exercício completo

\end{tipbox}

\chapter{Pesquisa em Ação --- Capítulo
15}\label{pesquisa-em-auxe7uxe3o-capuxedtulo-15}

\section{Pesquisa em Ação}\label{pesquisa-em-auxe7uxe3o-1}

\begin{pesquisabox}{De Loecker, J., Eeckhout, J. & Unger, G. (2020). The Rise of Market Power and the Macroeconomic Implications. *Quarterly Journal of Economics*, 135(2), 561–644.}

\textbf{DOI:}
\href{https://doi.org/10.1093/qje/qjz041}{10.1093/qje/qjz041}

\textbf{Contexto.} O Capítulo 15 analisa o monopólio como estrutura de
mercado, mas até que ponto o poder de mercado é um fenômeno relevante
empiricamente? De Loecker, Eeckhout e Unger (2020) respondem a essa
pergunta com uma análise abrangente do \emph{markup} praticado pelas
firmas nos Estados Unidos ao longo de seis décadas.

\textbf{Metodologia.} Os autores estimam \emph{markups} firma a firma
usando dados contábeis (Compustat) e a abordagem de produção de Hall
(1988), que infere o \emph{markup} como a razão entre a
elasticidade-produto de um insumo variável e a participação desse insumo
na receita. A grande inovação é aplicar esse método a um painel de
milhares de firmas de 1955 a 2016, permitindo documentar a evolução do
poder de mercado agregado.

\textbf{Resultados.} O \emph{markup} médio na economia americana era
relativamente estável em torno de 1,21 (21\% acima do custo marginal)
até 1980, mas subiu para 1,61 em 2016 --- um aumento de 33\%. Esse
aumento é concentrado nas firmas do topo da distribuição: as 10\% mais
lucrativas elevaram seus \emph{markups} de 1,5 para mais de 2,5. Os
autores documentam que essa tendência está associada a menor
participação do trabalho na renda, menor dinamismo empresarial e maior
desigualdade.

\textbf{Conexão com o capítulo.} O artigo fornece evidência empírica
maciça sobre o índice de Lerner (Seção 15.3) em escala agregada. Os
\emph{markups} estimados podem ser diretamente interpretados como
\(1/(1 - L)\), conectando a teoria do monopólio à mensuração empírica do
poder de mercado. A tendência ascendente documentada sugere que o modelo
de concorrência perfeita se torna progressivamente menos adequado como
descrição da economia moderna.

\end{pesquisabox}

\begin{pesquisabox}{Bergemann, D., Brooks, B. & Morris, S. (2015). The Limits of Price Discrimination. *American Economic Review*, 105(3), 921–957.}

\textbf{DOI:}
\href{https://doi.org/10.1257/aer.20130848}{10.1257/aer.20130848}

\textbf{Contexto.} A Seção 15.7 classifica a discriminação de preços em
três graus, mas essa taxonomia assume que sabemos exatamente quanta
informação o monopolista possui sobre os consumidores. Bergemann, Brooks
e Morris (2015) fazem uma pergunta mais fundamental: dada uma demanda de
mercado, quais são os \textbf{limites} do que a discriminação de preços
pode alcançar, considerando toda estrutura de informação possível?

\textbf{Contribuição teórica.} Os autores mostram que, para qualquer
segmentação de mercado (qualquer partição informacional dos
consumidores), o resultado da discriminação de preços deve satisfazer
duas restrições: (i) o lucro do monopolista deve ser pelo menos tão
grande quanto o lucro de preço uniforme; (ii) o excedente do consumidor
agregado não pode ser negativo. O resultado central é que
\textbf{qualquer} par (lucro, excedente do consumidor) satisfazendo
essas duas restrições é alcançável por alguma segmentação.

\textbf{Resultados.} O artigo demonstra que os efeitos de bem-estar da
discriminação de preços dependem crucialmente da informação disponível
ao monopolista. A produção total pode subir, cair ou permanecer
constante conforme a segmentação. Isso contrasta com o resultado
clássico de Pigou (1920) para demandas lineares, em que a discriminação
de 3º grau não altera a quantidade total.

\textbf{Conexão com o capítulo.} O paper aprofunda a análise de
discriminação de preços das Seções 15.7.1--15.7.3, mostrando que a
distinção entre 1º, 2º e 3º grau é apenas uma parte de um espectro muito
mais rico de possibilidades informacionais.

\end{pesquisabox}

\begin{pesquisabox}{Posner, R. A. (1975). The Social Costs of Monopoly and Regulation. *Journal of Political Economy*, 83(4), 807–827.}

\textbf{DOI:} \href{https://doi.org/10.1086/260357}{10.1086/260357}

\textbf{Contexto.} A Seção 15.4 apresenta a perda de peso morto
(triângulo de Harberger) como a medida convencional do custo social do
monopólio. Posner (1975) argumenta que essa medida subestima
dramaticamente o custo real, ao ignorar os recursos desperdiçados na
\emph{obtenção e manutenção} do poder de mercado.

\textbf{Contribuição teórica.} Posner observa que, se firmas competem
para obter uma posição de monopólio lucrativa (por exemplo, disputando
uma concessão governamental, investindo em lobby ou em litígios para
bloquear concorrentes), os recursos gastos nessa competição representam
um custo social adicional. No limite, se a competição pelo monopólio
dissipa todos os lucros esperados, o custo social total é a soma do
triângulo de Harberger e do retângulo de lucro do monopolista ---
potencialmente muito maior que o triângulo sozinho.

\textbf{Resultados.} Posner estima que, para setores regulados nos EUA,
o custo social total (incluindo \emph{rent-seeking}) pode ser várias
vezes maior que o triângulo de Harberger. Ele argumenta,
provocativamente, que a regulação governamental --- ao criar monopólios
legais e oportunidades de \emph{rent-seeking} --- pode gerar custos
sociais tão grandes ou maiores que o monopólio não regulado.

\textbf{Conexão com o capítulo.} O artigo fundamenta o Exercício
Resolvido 15.2, que calcula o custo social total incluindo
\emph{rent-seeking}. A análise de Posner é essencial para entender por
que estimativas empíricas do custo do monopólio variam tão amplamente
(de 0,1\% a 13\% do PIB, conforme discutido no Box Mundo 15.2).

\end{pesquisabox}

\begin{pesquisabox}{Borenstein, S. & Rose, N. L. (1994). Competition and Price Dispersion in the U.S. Airline Industry. *Journal of Political Economy*, 102(4), 653–683.}

\textbf{DOI:} \href{https://doi.org/10.1086/261950}{10.1086/261950}

\textbf{Contexto.} A Seção 15.7 e o Box Brasil 15.1 discutem a
discriminação de preços no setor aéreo. Borenstein e Rose (1994)
fornecem a primeira análise empírica rigorosa da relação entre
concorrência e dispersão de preços no setor aéreo americano --- uma
relação que se revela mais complexa do que a teoria simples sugeriria.

\textbf{Metodologia.} Os autores utilizam dados do \emph{Databank 1A} do
Departamento de Transportes dos EUA (DOT) com informações sobre tarifas
individuais para milhares de rotas domésticas. A medida de dispersão de
preços é a diferença entre o percentil 80 e o percentil 20 das tarifas
em cada rota, normalizada pela tarifa mediana. A variável explicativa
central é a concentração de mercado (índice de Herfindahl-Hirschman) em
cada rota.

\textbf{Resultados.} O resultado central é surpreendente: rotas mais
competitivas apresentam \emph{maior} dispersão de preços, não menor. A
dispersão de preços aumenta cerca de 10\% quando a concentração cai de
monopólio para duopólio. Os autores interpretam isso como evidência de
que a concorrência intensifica a discriminação de preços: em rotas
competitivas, companhias aéreas oferecem tarifas promocionais agressivas
para preencher assentos, enquanto mantêm tarifas altas para viajantes de
última hora --- gerando ampla dispersão.

\textbf{Conexão com o capítulo.} O artigo desafia a intuição de que mais
concorrência sempre reduz diferenças de preço. Ele mostra que a
discriminação de preços (Seção 15.7) não é exclusividade do monopólio
--- firmas com algum poder de mercado em oligopólios também a praticam
intensamente. Esse resultado antecipa temas dos Capítulos 16a--16b
(Oligopólio) e conecta-se diretamente ao Box Brasil sobre passagens
aéreas.

\end{pesquisabox}

\begin{pesquisabox}{Wallsten, S. (2001). An Econometric Analysis of Telecom Competition, Privatization, and Regulation in Africa and Latin America. *Journal of Industrial Economics*, 49(1), 1–19.}

\textbf{DOI:}
\href{https://doi.org/10.1111/1468-2354.t01-1-00106}{10.1111/1468-2354.t01-1-00106}

\textbf{Contexto.} A Seção 15.9 discute a regulação de monopólios
naturais em teoria. Wallsten (2001) fornece evidência empírica crucial
sobre os efeitos reais da desregulamentação e privatização no setor de
telecomunicações --- um dos exemplos mais importantes de monopólio
natural regulado do século XX.

\textbf{Metodologia.} O autor utiliza dados em painel de 30 países da
África e da América Latina entre 1984 e 1997, período em que muitas
nações promoveram reformas regulatórias no setor de telecomunicações
(privatização, criação de agências reguladoras independentes e abertura
à competição). A estratégia empírica explora a variação temporal e entre
países nessas reformas para identificar seus efeitos sobre investimento,
acesso (linhas per capita), qualidade do serviço e tarifas.

\textbf{Resultados.} A privatização combinada com a entrada de
concorrentes aumentou significativamente o investimento em
telecomunicações e o número de linhas telefônicas per capita. No
entanto, a privatização \emph{sem} introdução de competição --- isto é,
a mera transferência do monopólio estatal para um monopólio privado ---
teve efeitos muito menores ou nulos sobre o bem-estar dos consumidores.
A criação de agências reguladoras independentes também se mostrou
relevante para a credibilidade das reformas e a atração de investimento
privado.

\textbf{Conexão com o capítulo.} O artigo conecta-se diretamente à
análise de regulação de monopólios naturais (Seção 15.9) e ao Box Brasil
15.2 sobre concessões de serviços públicos no Brasil. Ele ilustra
empiricamente o trade-off entre eficiência alocativa e viabilidade
financeira discutido nos Exercícios 15.4 e 15.8, mostrando que a
estrutura regulatória --- não apenas a propriedade --- determina os
resultados de bem-estar.

\end{pesquisabox}

\begin{center}\rule{0.5\linewidth}{0.5pt}\end{center}

\section{Referências do Capítulo}\label{referuxeancias-do-capuxedtulo-1}

\begin{itemize}
\tightlist
\item
  Aghion, Philippe, Nick Bloom, Richard Blundell, Rachel Griffith, e
  Peter Howitt. 2005.
  ``\href{https://doi.org/10.1093/qje/120.2.701}{Competition and
  Innovation: An Inverted-U Relationship.}'' \emph{Quarterly Journal of
  Economics} 120 (2): 701--728.
\item
  Averch, Harvey, e Leland L. Johnson. 1962.
  ``\href{https://doi.org/10.2307/1812181}{Behavior of the Firm Under
  Regulatory Constraint.}'' \emph{American Economic Review} 52 (5):
  1052--1069.
\item
  Bain, Joe S. 1956.
  \href{https://books.google.com/books/about/Barriers_to_New_Competition.html?id=KyFCAAAAIAAJ}{\emph{Barriers
  to New Competition}}. Cambridge: Harvard University Press.
\item
  Baumol, William J., John C. Panzar, e Robert D. Willig. 1982.
  \href{https://books.google.com/books/about/Contestable_Markets_and_the_Theory_of_In.html?id=N2c6AAAAIAAJ}{\emph{Contestable
  Markets and the Theory of Industry Structure}}. New York: Harcourt
  Brace Jovanovich.
\item
  Bergemann, Dirk, Benjamin Brooks, e Stephen Morris. 2015.
  ``\href{https://doi.org/10.1257/aer.20130848}{The Limits of Price
  Discrimination.}'' \emph{American Economic Review} 105 (3): 921--957.
\item
  Besanko, David, e Ronald Braeutigam. 2014.
  \href{https://books.google.com.br/books?id=BeoengEACAAJ}{\emph{Microeconomics}}.
  5ª ed.~New York: Wiley. Caps. 11--12.
\item
  Borenstein, Severin, e Nancy L. Rose. 1994.
  ``\href{https://doi.org/10.1086/261950}{Competition and Price
  Dispersion in the U.S. Airline Industry.}'' \emph{Journal of Political
  Economy} 102 (4): 653--683.
\item
  Cowling, Keith, e Dennis C. Mueller. 1978.
  ``\href{https://doi.org/10.2307/2231972}{The Social Costs of Monopoly
  Power.}'' \emph{Economic Journal} 88 (352): 727--748.
\item
  De Loecker, Jan, Jan Eeckhout, e Gabriel Unger. 2020.
  ``\href{https://doi.org/10.1093/qje/qjz041}{The Rise of Market Power
  and the Macroeconomic Implications.}'' \emph{Quarterly Journal of
  Economics} 135 (2): 561--644.
\item
  Giambiagi, Fabio, e Ana Cláudia Além. 2016.
  \href{https://books.google.com.br/books?id=Xqf-jwEACAAJ}{\emph{Finanças
  Públicas: Teoria e Prática no Brasil}}. 5ª ed.~Rio de Janeiro:
  Elsevier.
\item
  Harberger, Arnold C. 1954.
  ``\href{https://doi.org/10.2307/1923808}{Monopoly and Resource
  Allocation.}'' \emph{American Economic Review} 44 (2): 77--87.
\item
  Leibenstein, Harvey. 1966.
  ``\href{https://doi.org/10.2307/1823775}{Allocative Efficiency
  vs.~`X-Efficiency'.}'' \emph{American Economic Review} 56 (3):
  392--415.
\item
  Mas-Colell, Andreu, Michael D. Whinston, e Jerry R. Green. 1995.
  \href{https://books.google.com/books/about/Microeconomic_Theory.html?id=KGtegVXqD8wC}{\emph{Microeconomic
  Theory}}. New York: Oxford University Press. Cap. 12.
\item
  Nicholson, Walter, e Christopher M. Snyder. 2017.
  \href{https://books.google.com/books/about/Microeconomic_Theory_Basic_Principles_an.html?id=YdkhCwAAQBAJ}{\emph{Microeconomic
  Theory: Basic Principles and Extensions}}. 12ª ed.~Boston: Cengage.
  Caps. 14--15.
\item
  Posner, Richard A. 1975. ``\href{https://doi.org/10.1086/260357}{The
  Social Costs of Monopoly and Regulation.}'' \emph{Journal of Political
  Economy} 83 (4): 807--827.
\item
  Spence, Michael. 1975.
  ``\href{https://doi.org/10.2307/3003237}{Monopoly, Quality, and
  Regulation.}'' \emph{Bell Journal of Economics} 6 (2): 417--429.
\item
  Tirole, Jean. 1988.
  \href{https://books.google.com/books/about/The_Theory_of_Industrial_Organization.html?id=HIjsF0XONF8C}{\emph{The
  Theory of Industrial Organization}}. Cambridge: MIT Press. Caps. 1--3.
\item
  Varian, Hal R. 1992.
  \href{https://books.google.com/books/about/Microeconomic_Analysis.html?id=m20iQAAACAAJ}{\emph{Microeconomic
  Analysis}}. 3ª ed.~New York: W. W. Norton. Caps. 14, 24.
\item
  Wallsten, Scott. 2001.
  ``\href{https://doi.org/10.1111/1468-2354.t01-1-00106}{An Econometric
  Analysis of Telecom Competition, Privatization, and Regulation in
  Africa and Latin America.}'' \emph{Journal of Industrial Economics} 49
  (1): 1--19.
\end{itemize}

\newpage

\chapter{Capítulo 20 --- Fumaça, Caronas e
Tragédias}\label{capuxedtulo-20-fumauxe7a-caronas-e-traguxe9dias}

Seu vizinho resolve criar galos. Às 5h da manhã, o canto ecoa pelo
quarteirão --- e o preço do imóvel ao lado despenca. O vizinho não pagou
nada por isso. O mercado de galos funcionou perfeitamente; o mercado de
noites bem dormidas, não. Bem-vindo ao mundo das externalidades --- onde
o prejuízo de um não aparece na conta de ninguém.\footnote{\emph{``What
  have the Romans ever done for us?''} Em \emph{A Vida de Brian}, os
  rebeldes tentam enumerar motivos para odiar os romanos --- mas só
  conseguem listar bens públicos: aquedutos, estradas, saneamento,
  irrigação, saúde pública, educação, vinho, ordem pública e a paz. A
  cena é a melhor análise custo-benefício de bens públicos já filmada. O
  líder do grupo, frustrado, repete a pergunta --- e a cada repetição
  alguém lembra de mais um benefício. É a condição de Samuelson em ação:
  a soma dos benefícios marginais (aqueduto + estrada + saneamento +
  \ldots) excede claramente o custo marginal (a ocupação romana). Até os
  rebeldes, relutantemente, concordam. \emph{``All right, but apart from
  the sanitation, the medicine, education, wine, public order,
  irrigation, roads, a fresh water system, and public health, what have
  the Romans ever done for us?''}}

Nos capítulos anteriores, demonstramos que mercados competitivos com
informação completa conduzem a alocações Pareto-eficientes --- o
resultado central dos Teoremas do Bem-Estar. Entretanto, essa conclusão
depende criticamente da hipótese de que todas as interações entre
agentes econômicos ocorrem via mecanismo de preços. Quando as ações de
um agente afetam diretamente o bem-estar de outros \textbf{fora do
sistema de preços}, surgem as \textbf{externalidades} --- e a alocação
de mercado deixa de ser eficiente.

Da mesma forma, certos bens possuem características que tornam o
mecanismo de mercado inadequado para sua provisão: bens
\textbf{não-rivais} e \textbf{não-excludentes} --- os chamados
\textbf{bens públicos} --- tendem a ser subprovidos pelo mercado devido
ao problema do \textbf{carona} (free rider).

Este capítulo analisa essas duas importantes \textbf{falhas de mercado}
e discute as principais soluções propostas pela teoria econômica:
impostos pigouvianos, o Teorema de Coase, mercados de permissão de
emissão, preços de Lindahl e mecanismos de revelação de preferências.
Essas questões assumem importância crescente no debate sobre políticas
ambientais, saúde pública e infraestrutura.

A análise de externalidades e bens públicos tem raízes profundas na
história do pensamento econômico. Já Alfred Marshall, nos
\emph{Principles of Economics} (1890), reconhecia a existência de
``economias externas'' --- benefícios que a atividade de uma firma
gerava para outras firmas ou para a sociedade sem compensação pelo
mercado. Contudo, foi Arthur Cecil Pigou quem, em \emph{The Economics of
Welfare} (1920), sistematizou a análise das externalidades como
divergência entre custo privado e custo social, propondo a intervenção
fiscal como mecanismo de correção. Quatro décadas depois, Ronald Coase
(1960) desafiou a prescrição pigouviana, argumentando que a negociação
privada poderia, sob certas condições, alcançar a eficiência sem
intervenção estatal. Esse debate Pigou--Coase permanece no centro da
economia ambiental e da análise econômica do direito até hoje. Do lado
dos bens públicos, a contribuição seminal de Paul Samuelson (1954)
formalizou a condição de eficiência que distingue radicalmente bens
públicos de bens privados, inaugurando a moderna teoria da economia do
setor público.

\begin{center}\rule{0.5\linewidth}{0.5pt}\end{center}

\section{Roteiro do Capítulo}\label{roteiro-do-capuxedtulo-1}

{\def\LTcaptype{none} % do not increment counter
\begin{longtable}[]{@{}
  >{\centering\arraybackslash}p{(\linewidth - 6\tabcolsep) * \real{0.3125}}
  >{\raggedright\arraybackslash}p{(\linewidth - 6\tabcolsep) * \real{0.1875}}
  >{\raggedright\arraybackslash}p{(\linewidth - 6\tabcolsep) * \real{0.1875}}
  >{\centering\arraybackslash}p{(\linewidth - 6\tabcolsep) * \real{0.3125}}@{}}
\toprule\noalign{}
\begin{minipage}[b]{\linewidth}\centering
Seção
\end{minipage} & \begin{minipage}[b]{\linewidth}\raggedright
Pergunta-guia
\end{minipage} & \begin{minipage}[b]{\linewidth}\raggedright
O que você vai aprender
\end{minipage} & \begin{minipage}[b]{\linewidth}\centering
Página
\end{minipage} \\
\midrule\noalign{}
\endhead
\bottomrule\noalign{}
\endlastfoot
20.1 & Quando a ação de um afeta o bem-estar de outro sem passar pelo
mercado? & Definição e tipologia de externalidades & Externalidades \\
20.2 & Por que a poluição é excessiva se ninguém paga por ela? &
Externalidades e ineficiência alocativa & Ineficiência \\
20.3 & Como medir o custo social de uma externalidade? & Modelo de
equilíbrio parcial com externalidade & Modelo \\
20.4 & Imposto, negociação ou regulação --- o que resolve melhor? &
Pigou, Coase, cap-and-trade & Soluções \\
20.5 & O que torna um bem ``público'' --- e por que o mercado falha em
provê-lo? & Não rivalidade, não exclusibilidade & Bens públicos \\
20.6 & Quanto de bem público a sociedade deveria produzir? & Condição de
Samuelson, soma vertical & Alocação \\
20.7 & Existe um preço personalizado que faz todo mundo pagar pelo que
valoriza? & Preços de Lindahl & Lindahl \\
20.8 & Por que todo mundo quer pegar carona e ninguém quer pagar? &
Free-rider, provisão voluntária & Carona \\
20.9 & A maioria vota bem --- ou vota mal? & Votação, teorema do eleitor
mediano & Votação \\
20.10 & Como fazer as pessoas revelarem quanto realmente valorizam o bem
público? & Mecanismos de revelação, VCG & Revelação \\
\textbf{Exercícios} & Teste, pratique, resolva & Revisão, exercícios,
ANPEC & Exercícios \\
\textbf{Pesquisa} & O que a pesquisa recente diz? & Artigos seminais e
fronteira empírica & Pesquisa \\
\end{longtable}
}

\begin{center}\rule{0.5\linewidth}{0.5pt}\end{center}

\begin{classroombox}{Atividade de Sala — Jogo de Bens Públicos (Contribuição Voluntária)}

\textbf{Formato:} Experimento de contribuição voluntária com e sem
punição \textbf{Duração:} 40--45 minutos (20 rodadas + 15 debrief + 10
discussão) \textbf{Material:} Fichas ou moedas (10 por aluno), quadro
para registrar contribuições.

\textbf{Preparação (5 min):}

\begin{itemize}
\tightlist
\item
  Divida a turma em grupos de 4--5 alunos.
\item
  Cada aluno recebe 10 fichas por rodada.
\item
  Regra: cada ficha mantida rende 1 ponto ao aluno. Cada ficha
  contribuída ao ``fundo público'' do grupo rende 0,4 pontos
  \textbf{para cada membro} do grupo (incluindo o contribuidor).
\item
  Com 5 membros, o retorno social de 1 ficha contribuída = 5 × 0,4 = 2,0
  \textgreater{} 1,0 (custo privado). Contribuir é socialmente
  eficiente. Mas individualmente, contribuir 1 ficha retorna apenas 0,4
  para o contribuidor --- menos que manter (1,0). Free-riding é
  individualmente racional.
\end{itemize}

\textbf{Fase 1 --- Sem punição (rodadas 1--8, 10 min):}

\begin{itemize}
\tightlist
\item
  Cada rodada: alunos decidem secretamente quantas fichas contribuir
  (0--10).
\item
  O professor soma as contribuições, calcula o retorno e registra no
  quadro.
\item
  \textbf{Resultado esperado:} Contribuições começam em 40--60\% e caem
  para 10--20\% até a rodada 8. O free-riding emerge.
\end{itemize}

\textbf{Fase 2 --- Com punição (rodadas 9--16, 10 min):}

\begin{itemize}
\tightlist
\item
  Nova regra: após a revelação, cada aluno pode gastar 1 ficha para
  retirar 3 fichas de outro membro (punição).
\item
  \textbf{Resultado esperado:} Contribuições sobem para 60--90\%.
  Free-riders são punidos. Cooperação se sustenta (Fehr \& Gächter,
  2000).
\end{itemize}

\textbf{Fase 3 --- Com comunicação (rodadas 17--20, 5 min, opcional):}

\begin{itemize}
\tightlist
\item
  Permita 2 min de conversa antes de cada rodada.
\item
  \textbf{Resultado esperado:} Contribuições próximas de 100\%.
  Comunicação é poderosamente eficaz (Ostrom et al., 1992).
\end{itemize}

\textbf{Debrief (15 min):}

\begin{enumerate}
\def\labelenumi{\arabic{enumi}.}
\tightlist
\item
  Plote as contribuições médias por rodada no quadro. Identifique os
  padrões.
\item
  Pergunte: \emph{``Por que as contribuições caíram sem punição?''} →
  Problema do carona (Seção 20.8).
\item
  Pergunte: \emph{``Por que a punição funciona, mesmo sendo custosa para
  quem pune?''} → Punição altruísta, normas sociais.
\item
  Pergunte: \emph{``O que muda quando vocês podem conversar?''} →
  Compromissos, reputação (Ostrom).
\item
  Conecte: condição de Samuelson (retorno social \textgreater{} custo),
  subprovisão (Nash), Lindahl (preços personalizados).
\end{enumerate}

\textbf{Conexão com o capítulo:} Seções 20.5--20.8. O experimento segue
Fehr \& Gächter (2000) e Holt (2007, Cap. 23).

\end{classroombox}

\chapter{Externalidades: Definição e
Ineficiência}\label{externalidades-definiuxe7uxe3o-e-ineficiuxeancia}

\section{20.1 O Vizinho Que Cria Galos às 5h da Manhã: Definição de
Externalidades}\label{201}

\begin{definitionbox}{Externalidade}

Uma \textbf{externalidade} ocorre quando a ação de um agente econômico
afeta diretamente a utilidade ou a função de produção de outro agente,
sem que essa interação seja mediada pelo sistema de preços. A
externalidade representa um efeito \textbf{externo ao mercado} que não é
capturado nas decisões privadas dos agentes.

\end{definitionbox}

\subsection{Classificação das
externalidades}\label{classificauxe7uxe3o-das-externalidades}

As externalidades podem ser classificadas em duas dimensões:

\textbf{Pelo sinal (efeito sobre terceiros):}

\begin{itemize}
\tightlist
\item
  \textbf{Externalidade negativa}: a ação de um agente prejudica outros
  (poluição, congestionamento, ruído).
\item
  \textbf{Externalidade positiva}: a ação de um agente beneficia outros
  (vacinação, pesquisa básica, preservação de paisagem).
\end{itemize}

\textbf{Pela esfera (produção ou consumo):}

\begin{itemize}
\tightlist
\item
  \textbf{Externalidade de produção}: a função de produção de uma firma
  é afetada pela produção de outra. Exemplo: uma fábrica que polui o rio
  usado por um pesqueiro.
\item
  \textbf{Externalidade de consumo}: a utilidade de um consumidor é
  afetada pelo consumo de outro. Exemplo: o prazer (ou desprazer)
  causado pelo hábito de fumar de um vizinho.
\end{itemize}

Para além desses exemplos clássicos, externalidades permeiam a vida
cotidiana de maneiras muitas vezes sutis. A poluição sonora de um
aeroporto reduz o valor dos imóveis residenciais ao redor --- uma
externalidade negativa de produção sobre consumidores. As
\textbf{externalidades de rede} constituem uma categoria especialmente
relevante na economia digital: quando mais pessoas adotam uma plataforma
de comunicação (WhatsApp, por exemplo), o valor do serviço para cada
usuário existente aumenta. Trata-se de uma externalidade positiva de
consumo, pois a decisão de adesão de um indivíduo beneficia todos os
demais participantes da rede sem que haja compensação monetária por esse
benefício. De modo análogo, a polinização realizada por abelhas de um
apicultor beneficia lavouras vizinhas --- uma externalidade positiva de
produção que atravessa fronteiras de propriedade sem passar pelo sistema
de preços.

Formalmente, uma externalidade de produção existe quando a função de
produção da firma \(j\) depende do nível de produção da firma \(i\):

{[} q\_j = f\_j(L\_j, K\_j, q\_i) \label{eq:20.1} \tag{20.1} {]}

Uma externalidade de consumo existe quando a utilidade do indivíduo
\(B\) depende do consumo do indivíduo \(A\):

{[} U\_B = U\_B(x\_B\^{}1, x\_B\^{}2, \ldots, x\_A\^{}k) \label{eq:20.2}
\tag{20.2} {]}

onde \(x_A^k\) é o consumo do bem \(k\) pelo indivíduo \(A\).

\begin{notebox}{Externalidades Pecuniárias vs. Tecnológicas}

É importante distinguir \textbf{externalidades tecnológicas} (reais) de
\textbf{externalidades pecuniárias}. Externalidades pecuniárias operam
\emph{via preços}: quando uma firma expande sua produção e reduz o preço
do produto, prejudicando concorrentes. Essas não são verdadeiras
externalidades no sentido econômico, pois são mediadas pelo mercado e
não geram ineficiência. Apenas externalidades tecnológicas --- que
afetam diretamente funções de produção ou utilidade --- constituem
falhas de mercado.

\end{notebox}

A distinção entre externalidades pecuniárias e tecnológicas merece um
exemplo concreto. Suponha que uma grande rede varejista se instale em
uma cidade pequena, reduzindo os preços dos produtos e levando
comerciantes locais a perder clientes. Os comerciantes são
``prejudicados'', mas esse efeito opera \emph{via mercado} --- trata-se
de uma externalidade pecuniária que não gera ineficiência alocativa (na
verdade, reflete o funcionamento competitivo do mercado). Agora compare
com uma siderúrgica cujas emissões de material particulado causam
doenças respiratórias nos moradores vizinhos. Este efeito \emph{não}
passa pelo sistema de preços: a saúde dos moradores é afetada
diretamente, sem qualquer transação de mercado. Trata-se de uma
externalidade tecnológica --- e é esta, e somente esta, que constitui
uma falha de mercado requerendo correção.

\begin{intuitionbox}{Intuição Econômica}

\textbf{Em uma frase:} Uma externalidade existe quando o preço de um
produto não reflete todo o custo (ou benefício) que sua produção ou
consumo impõe à sociedade.

\textbf{Pense assim:} Quando uma fábrica em Cubatão polui o rio, o preço
do produto que ela vende não inclui o custo da água contaminada para os
pescadores e moradores rio abaixo. A fábrica ``exporta'' parte do seu
custo para a sociedade sem pagar por isso. Da mesma forma, quem se
vacina protege não só a si mesmo, mas também quem está ao redor --- um
benefício que não entra no cálculo individual.

\textbf{Por que isso importa:} Externalidades são a principal
justificativa econômica para intervenção do governo via impostos,
subsídios ou regulação --- do controle de emissões ao financiamento
público da vacinação.

\end{intuitionbox}

\begin{center}\rule{0.5\linewidth}{0.5pt}\end{center}

\section{20.2 O Preço Que Mente: Externalidades e Ineficiência
Alocativa}\label{202}

Definido o conceito, vem a pergunta que realmente importa: por que
externalidades fazem o mercado errar? A resposta é quase ofensivamente
simples. O preço --- aquele mecanismo que Adam Smith celebrou como o
grande coordenador da economia --- mente. Ou melhor, conta apenas metade
da história: reflete os custos e benefícios privados da transação, mas
ignora solenemente os efeitos sobre terceiros. Como resultado, as
decisões tomadas individualmente por agentes racionais conduzem a um
nível de produção ou consumo que difere do socialmente ótimo --- o
mercado ``erra'' sistematicamente.

É instrutivo conectar essa falha ao Primeiro Teorema do Bem-Estar.
Recordemos que esse teorema garante que todo equilíbrio competitivo é
Pareto-eficiente --- mas sob hipóteses específicas, entre as quais a de
\textbf{mercados completos}: deve existir um mercado para cada bem que
afeta a utilidade ou a produção de qualquer agente. Quando uma
externalidade está presente, \emph{falta} exatamente um mercado: não
existe mercado para a poluição, para o silêncio, para o ar limpo ou para
a imunidade de rebanho. Essa ``lacuna'' no sistema de mercados viola a
hipótese central do Primeiro Teorema e é a razão profunda pela qual
externalidades geram ineficiência. Em termos técnicos, a externalidade
cria uma \textbf{cunha} (wedge) entre os custos (ou benefícios) privados
e os custos (ou benefícios) sociais. Essa cunha distorce os sinais de
preço: o produtor vê um custo menor do que o verdadeiro custo social, ou
o consumidor percebe um benefício menor do que o verdadeiro benefício
social. A magnitude da ineficiência é proporcional ao tamanho dessa
cunha.

A presença de externalidades implica que as decisões privadas dos
agentes divergem do ótimo social. A razão fundamental é que os agentes
não incorporam em seus cálculos os efeitos que suas ações impõem sobre
terceiros (Browning \& Zupan, 2014, Cap. 20, apresentam aplicações
detalhadas).

\subsection{Externalidade negativa de
produção}\label{externalidade-negativa-de-produuxe7uxe3o}

Considere uma firma que produz um bem \(q\) com custo privado \(C(q)\) e
que gera poluição como subproduto. A poluição impõe um \textbf{custo
externo} \(E(q)\) sobre a sociedade (danos à saúde, degradação
ambiental). O \textbf{custo social} é:

{[} CS(q) = C(q) + E(q) \label{eq:20.3} \tag{20.3} {]}

A firma, buscando maximizar lucro, iguala o custo marginal privado ao
preço:

{[} P = C'(q) \quad \text{(CMg privado)} {]}

Mas a condição de eficiência social requer:

{[} P = C'(q) + E'(q) \quad \text{(CMg social)} \label{eq:20.4}
\tag{20.4} {]}

Como \(E'(q) > 0\), a firma produz \textbf{mais do que o ótimo social}:
\(q^{priv} > q^{soc}\).

\subsection{Externalidade positiva de
consumo}\label{externalidade-positiva-de-consumo}

No caso de uma externalidade positiva --- por exemplo, vacinação ---, o
\textbf{benefício social} excede o benefício privado:

{[} BS(q) = B\^{}\{priv\}(q) + B\^{}\{ext\}(q) {]}

O indivíduo consome até onde seu benefício marginal privado iguala o
preço:

{[} B'\^{}\{priv\}(q) = P \label{eq:20.5} \tag{20.5} {]}

Mas a eficiência requer:

{[} B'\^{}\{priv\}(q) + B'\^{}\{ext\}(q) = P \label{eq:20.6} \tag{20.6}
{]}

Como \(B'^{ext}(q) > 0\), o consumo privado fica \textbf{aquém do ótimo
social}: \(q^{priv} < q^{soc}\).

\begin{tipbox}{Regra Geral}

\begin{itemize}
\tightlist
\item
  \textbf{Externalidade negativa} → produção/consumo privado
  \textbf{excessivo} em relação ao ótimo social.
\item
  \textbf{Externalidade positiva} → produção/consumo privado
  \textbf{insuficiente} em relação ao ótimo social.
\item
  Em ambos os casos, o mercado \textbf{falha} em atingir eficiência de
  Pareto.
\end{itemize}

\end{tipbox}

\begin{boxmundobox}{Box Mundo 20.1 — Pedágios urbanos: precificando o congestionamento em Estocolmo, Londres e Singapura}

\textbf{Contexto:} O congestionamento viário é uma externalidade
negativa clássica: cada motorista que entra em uma via já congestionada
impõe custos a todos os demais usuários (tempo de viagem, consumo extra
de combustível, poluição) sem pagar por esse efeito. A solução
pigouviana natural é cobrar um pedágio que reflita o custo marginal
externo do congestionamento --- é exatamente o que três cidades
pioneiras fizeram.

\textbf{Dados:} Singapura introduziu o primeiro sistema de pedágio
urbano do mundo em 1975 (Electronic Road Pricing desde 1998), cobrando
tarifas que variam de SGD 0,50 a SGD 6,00 conforme o horário e o nível
de congestionamento. Londres implementou a \emph{Congestion Charge} em
2003, cobrando inicialmente £5 (atualizada para £15 em 2023) para
veículos que entram na zona central. Estocolmo adotou um sistema em 2006
após um período experimental com referendo: as tarifas variam entre SEK
15 e SEK 45 conforme o horário do dia.

\textbf{Análise:} Os resultados são consistentes nas três cidades. Em
Estocolmo, o tráfego na zona central caiu 20-25\% e as emissões de CO₂
na área taxada reduziram-se em 10-14\% (Eliasson et al., 2009). Em
Londres, o congestionamento diminuiu 30\% no primeiro ano, embora parte
do efeito tenha se dissipado com o tempo. O caso de Estocolmo é
especialmente instrutivo: a população inicialmente opunha-se ao pedágio,
mas após o período experimental, ao constatar a melhoria na qualidade do
tráfego, votou pela manutenção permanente do sistema em referendo. Essa
sequência --- resistência inicial seguida de apoio após a experiência
--- ilustra como a percepção pública dos benefícios de uma taxa
pigouviana pode mudar quando os cidadãos vivenciam a redução da
externalidade.

\textbf{Fonte:} Eliasson, J.; Hultkrantz, L.; Nerhagen, L.; Rosqvist, L.
S. (2009). The Stockholm congestion-charging trial 2006: overview of
effects. \emph{Transportation Research Part A}, 43(3), 240-250.

\end{boxmundobox}

\chapter{Modelo de Equilíbrio Parcial e
Soluções}\label{modelo-de-equiluxedbrio-parcial-e-soluuxe7uxf5es}

\section{20.3 Medindo a Mentira: Modelo de Equilíbrio Parcial de
Externalidades}\label{203}

Sabemos que o preço mente --- mas quanto ele mente? A intuição da seção
anterior nos disse que externalidades negativas levam à sobreprodução e
positivas à subprodução. Agora, formalizamos essa análise em um modelo
de equilíbrio parcial com duas firmas, o que nos permitirá medir
exatamente o tamanho da mentira e, na seção seguinte, calibrar os
instrumentos de política necessários para corrigi-la.

Formalizemos o problema usando um modelo de equilíbrio parcial com duas
firmas.

\textbf{Firma 1} (poluidora): produz \(q_1\) com custo \(C_1(q_1)\) e
gera emissões \(e = e(q_1)\), com \(e' > 0\).

\textbf{Firma 2} (afetada): produz \(q_2\) com custo \(C_2(q_2, e)\),
onde \(\partial C_2 / \partial e > 0\) --- a poluição de 1 eleva os
custos de 2.

\textbf{Equilíbrio privado}: cada firma maximiza seu próprio lucro,
ignorando o efeito sobre a outra.

Firma 1:

{[} \max\_\{q\_1\} ; p\_1 q\_1 - C\_1(q\_1) \implies p\_1 = C\_1'(q\_1)
{]}

\textbf{Ótimo social}: o planejador maximiza o lucro conjunto:

{[} \max\_\{q\_1, q\_2\} ; p\_1 q\_1 - C\_1(q\_1) + p\_2 q\_2 -
C\_2(q\_2, e(q\_1)) {]}

A condição de primeira ordem para \(q_1\) é:

{[} p\_1 = C\_1'(q\_1) + \frac{\partial C_2}{\partial e} \cdot e'(q\_1)
\label{eq:20.7} \tag{20.7} {]}

O termo \(\frac{\partial C_2}{\partial e} \cdot e'(q_1)\) é o
\textbf{custo marginal externo} (CME) --- o dano marginal imposto pela
firma 1 sobre a firma 2. No equilíbrio privado, esse termo é ignorado,
levando a \(q_1^{priv} > q_1^{soc}\).

A \textbf{perda de peso morto} (deadweight loss) associada à
externalidade é a área entre as curvas de custo marginal social e custo
marginal privado, entre \(q_1^{soc}\) e \(q_1^{priv}\):

{[} DWL = \int\_\{q\_1\textsuperscript{\{soc\}\}}\{q\_1\^{}\{priv\}\}
\left[ C_1'(q) + E'(q) - P \right] dq {]}

Geometricamente, essa integral corresponde à área do ``triângulo'' entre
a curva de custo marginal social (CMg + CME) e a reta de preço,
delimitado pelas quantidades \(q^{soc}\) e \(q^{priv}\). A interpretação
é direta: para cada unidade produzida além do ótimo social, o custo
social marginal excede o benefício marginal (preço), e a área acumulada
dessas ``perdas marginais'' constitui o peso morto. Note que a magnitude
do peso morto depende crucialmente de dois fatores: (i) o tamanho da
cunha entre CMg privado e CMg social (ou seja, a magnitude do custo
marginal externo) e (ii) a distância entre \(q^{priv}\) e \(q^{soc}\),
que por sua vez depende das \textbf{elasticidades} das curvas de oferta
e demanda. Se a oferta privada for muito inelástica (CMg privado muito
inclinado), a sobreprodução será pequena mesmo com um custo externo
grande, pois a quantidade responde pouco à divergência. Inversamente, se
a oferta for elástica, um custo externo modesto pode gerar uma distorção
quantitativa --- e um peso morto --- substancial. Essa relação entre
elasticidades e magnitude da ineficiência tem implicações diretas para a
política pública: setores com oferta elástica e externalidades
significativas são os candidatos prioritários à intervenção regulatória.

\begin{center}\rule{0.5\linewidth}{0.5pt}\end{center}

\figurePlaceholder{/micro-book/graficos/cap20/externalidade-negativa.html}

\textbf{Figura 20.1 --- Externalidade negativa e imposto pigouviano.} A
curva vermelha (CMg social) inclui o custo externo. A área sombreada
representa a perda de peso morto da sobreprodução. O imposto pigouviano
ótimo iguala o custo marginal externo no ótimo social.

\section{20.4 A Receita para a Fumaça: Soluções para Externalidades
Negativas}\label{204}

\subsection{20.4.1 Imposto Pigouviano}\label{imposto-pigouviano}

\begin{definitionbox}{Imposto Pigouviano}

Imposto por unidade de produção (ou de emissão) igual ao \textbf{custo
marginal externo} avaliado no ótimo social. Nomeado em homenagem a
Arthur Cecil Pigou (1920), que primeiro propôs essa correção.

\end{definitionbox}

O imposto \(t^*\) é definido como:

{[} t\^{}* = E'(q\^{}\{soc\}) = \left. \frac{\partial C_2}{\partial e}
\cdot e'(q\_1) \right\textbar\_\{q\_1 = q\_1\^{}\{soc\}\}
\label{eq:20.8} \tag{20.8} {]}

Com o imposto, a firma poluidora resolve:

{[} \max\_\{q\_1\} ; p\_1 q\_1 - C\_1(q\_1) - t\^{}* q\_1 {]}

Condição de primeira ordem:

{[} p\_1 = C\_1'(q\_1) + t\^{}* \label{eq:20.9} \tag{20.9} {]}

Como \(t^* = E'(q^{soc})\), a firma internaliza o custo externo e produz
\(q_1^{soc}\).

\begin{proofbox}{Demonstração: O Imposto Pigouviano Ótimo Restaura a Eficiência}

\textbf{Objetivo}: Mostrar que um imposto \(t^*\) igual ao custo
marginal externo no ótimo social induz a firma poluidora a escolher o
nível de produção socialmente eficiente.

\textbf{Configuração}: Considere uma economia com uma firma poluidora
que produz \(q\) unidades de um bem ao preço \(P\). O custo privado de
produção é \(C(q)\), com \(C'(q) > 0\) e \(C''(q) > 0\). A produção gera
um custo externo \(E(q)\), com \(E'(q) > 0\) e \(E''(q) \geq 0\). O
benefício social líquido é:

{[} W(q) = Pq - C(q) - E(q) {]}

\textbf{Passo 1 --- Ótimo social.} O planejador maximiza \(W(q)\):

{[} \frac{dW}{dq} = P - C'(q) - E'(q) = 0 {]}

Logo, o nível ótimo \(q^*\) satisfaz:

{[} P = C'(q\^{}\emph{) + E'(q\^{}}) \tag{i} {]}

\textbf{Passo 2 --- Equilíbrio privado sem imposto.} A firma maximiza
\(\pi(q) = Pq - C(q)\):

{[} P = C'(q\^{}\{priv\}) \tag{ii} {]}

Como \(E'(q) > 0\), comparando (i) e (ii), e dado que \(C''(q) > 0\),
temos \(q^{priv} > q^*\). A firma produz em excesso.

\textbf{Passo 3 --- Introdução do imposto pigouviano.} Defina o imposto
unitário:

{[} t\^{}* = E'(q\^{}*) {]}

Com o imposto, o lucro da firma é \(\pi^t(q) = Pq - C(q) - t^* q\). A
condição de maximização é:

{[} P = C'(q\^{}t) + t\^{}* = C'(q\^{}t) + E'(q\^{}*) \tag{iii} {]}

\textbf{Passo 4 --- Comparação.} Comparando (iii) com (i), ambas
requerem:

{[} P - E'(q\^{}*) = C'(q) {]}

Como \(C'(\cdot)\) é estritamente crescente (por \(C'' > 0\)), a solução
é única: \(q^t = q^*\).

\textbf{Conclusão.} O imposto \(t^* = E'(q^*)\) faz com que a firma
internalize o custo marginal externo e produza exatamente o nível
socialmente ótimo \(q^*\). A perda de peso morto é eliminada.
\(\blacksquare\)

\end{proofbox}

\begin{tipbox}{Subsídio Pigouviano}

De forma simétrica, para externalidades positivas, o instrumento
pigouviano é um \textbf{subsídio} por unidade igual ao benefício
marginal externo: \(s^* = B'^{ext}(q^{soc})\). Isso eleva a quantidade
produzida/consumida até o nível socialmente eficiente.

\end{tipbox}

\begin{exresolvidobox}{Exercício Resolvido 20.1 —Imposto pigouviano ótimo com custo quadrático}

\textbf{Enunciado.} Uma firma produz \(q\) unidades ao preço
\(P = 100\), com custo \(C(q) = 10q + q^2\). A produção gera custo
externo \(E(q) = 0{,}5q^2\). (a) Encontre \(q^{priv}\) e \(q^{soc}\).
(b) Calcule o imposto pigouviano ótimo. (c) Calcule o DWL eliminado.

\textbf{Solução.}

\textbf{(a)} Sem regulação:
\(P = C'(q) \implies 100 = 10 + 2q \implies q^{priv} = 45\).

Ótimo social:
\(P = C'(q) + E'(q) \implies 100 = 10 + 2q + q = 10 + 3q \implies q^{soc} = 30\).

\textbf{(b)} \(t^* = E'(q^{soc}) = q^{soc} = 30\).

Verificação: com imposto,
\(100 = 10 + 2q + 30 \implies q = 30 = q^{soc}\). \(\checkmark\)

\textbf{(c)}
\(DWL = \int_{30}^{45} (10 + 3q - 100)\,dq = \int_{30}^{45}(3q - 90)\,dq = \left[\frac{3q^2}{2} - 90q\right]_{30}^{45} = (3037{,}5 - 4050) - (1350 - 2700) = -1012{,}5 - (-1350) = 337{,}5\).

\end{exresolvidobox}

\begin{center}\rule{0.5\linewidth}{0.5pt}\end{center}

\figurePlaceholder{/micro-book/graficos/cap20/webr-pigou.html}

\textbf{WebR 20.1 --- Simulação de imposto pigouviano.} Altere os
parâmetros de custo e externalidade e observe como o imposto ótimo e o
DWL se ajustam.

\begin{center}\rule{0.5\linewidth}{0.5pt}\end{center}

\subsection{20.4.2 O Teorema de Coase}\label{o-teorema-de-coase}

\begin{theorembox}{Teorema de Coase (Coase, 1960)}

Se os \textbf{direitos de propriedade} estão claramente definidos e os
\textbf{custos de transação} são nulos, a barganha privada entre as
partes levará à alocação eficiente de recursos,
\textbf{independentemente} de qual parte detém o direito de propriedade.
A distribuição dos direitos afeta apenas a distribuição de riqueza, não
a eficiência.

\end{theorembox}

Formalmente, suponha que a firma poluidora e a firma afetada podem
negociar. Se a firma afetada tem o direito a um ambiente limpo, a
poluidora deve compensá-la pelo dano. Se a poluidora tem o direito de
poluir, a firma afetada paga para que ela reduza a produção. Em ambos os
casos, a produção converge para \(q^{soc}\) onde:

{[} C'(q) + E'(q) = P \label{eq:20.10} \tag{20.10} {]}

\textbf{Limites do Teorema de Coase:}

\begin{enumerate}
\def\labelenumi{\arabic{enumi}.}
\tightlist
\item
  \textbf{Custos de transação}: na prática, negociar é custoso,
  especialmente quando há muitas partes envolvidas.
\item
  \textbf{Externalidades difusas}: quando a poluição afeta milhões de
  pessoas, a barganha bilateral é inviável.
\item
  \textbf{Assimetria de informação}: as partes podem ter incentivos
  estratégicos para não revelar verdadeiramente seus custos e
  benefícios.
\item
  \textbf{Efeitos riqueza}: com utilidade marginal da renda decrescente,
  a atribuição inicial dos direitos pode afetar a alocação final.
\item
  \textbf{Problemas de hold-up}: investimentos específicos à relação
  podem gerar comportamento oportunista.
\end{enumerate}

\begin{infobox}{Nobel de Economia 1991 — Ronald Coase}

Ronald Coase recebeu o Prêmio Nobel ``por sua descoberta e elucidação do
significado dos custos de transação e dos direitos de propriedade para a
estrutura institucional e o funcionamento da economia.'' O artigo
seminal ``The Problem of Social Cost'' (1960) desafiou a prescrição
pigouviana dominante, mostrando que a atribuição de direitos de
propriedade e a barganha privada podem, sob certas condições, resolver o
problema das externalidades sem intervenção estatal. Coase também
introduziu o conceito de custos de transação como determinante da
fronteira entre firma e mercado em ``The Nature of the Firm'' (1937).

\end{infobox}

\begin{exresolvidobox}{Exercício Resolvido 20.2 —Negociação coaseana: fazendeiro vs. apicultor}

\textbf{Enunciado.} Um fazendeiro cria gado (\(n\) cabeças) com lucro
\(\pi_F(n) = 80n - 2n^2\). O gado danifica a lavoura vizinha em
\(D(n) = n^2\). (a) Encontre \(n^{priv}\) e \(n^{soc}\). (b) Mostre que
a negociação leva a \(n^{soc}\) independentemente de quem detém o
direito.

\textbf{Solução.}

\textbf{(a)} Sem negociação:
\(\pi_F'(n) = 80 - 4n = 0 \implies n^{priv} = 20\).

Ótimo social:
\(\max_n \pi_F(n) - D(n) = 80n - 2n^2 - n^2 = 80n - 3n^2\). CPO:
\(80 - 6n = 0 \implies n^{soc} = 40/3 \approx 13{,}33\).

\textbf{(b)} Se o fazendeiro tem o direito: parte de \(n = 20\). O
lavrador oferece compensação \(T\) para reduzir a \(n^{soc}\). Perda do
fazendeiro: \(\pi_F(20) - \pi_F(40/3) = 800 - 711{,}11 = 88{,}89\).
Ganho do lavrador: \(D(20) - D(40/3) = 400 - 177{,}78 = 222{,}22\). Como
\(222{,}22 > 88{,}89\), há espaço para acordo.

Se o lavrador tem o direito: parte de \(n = 0\). O fazendeiro oferece
\(T'\) para criar \(n^{soc}\) cabeças. Ganho do fazendeiro:
\(\pi_F(40/3) = 711{,}11\). Dano ao lavrador: \(D(40/3) = 177{,}78\).
Como \(711{,}11 > 177{,}78\), há espaço para acordo. Em ambos os casos,
\(n = n^{soc}\).

\end{exresolvidobox}

\begin{center}\rule{0.5\linewidth}{0.5pt}\end{center}

\figurePlaceholder{/micro-book/graficos/cap20/teorema-coase.html}

\textbf{Figura 20.2 --- Teorema de Coase.} Alterne entre atribuir o
direito de propriedade ao poluidor ou à vítima. Em ambos os casos, a
negociação leva ao mesmo nível eficiente de poluição onde BMg = DMg.
Apenas a direção da transferência muda.

\begin{center}\rule{0.5\linewidth}{0.5pt}\end{center}

\figurePlaceholder{/micro-book/graficos/cap20/webr-coase.html}

\textbf{WebR 20.2 --- Simulação do Teorema de Coase.} Varie os
parâmetros de lucro e dano e observe o intervalo de compensação viável.

\begin{center}\rule{0.5\linewidth}{0.5pt}\end{center}

\subsection{20.4.3 Regulação direta (command and
control)}\label{regulauxe7uxe3o-direta-command-and-control}

A regulação direta consiste em impor limites quantitativos de emissão ou
padrões tecnológicos às firmas poluidoras. Exemplos incluem:

\begin{itemize}
\tightlist
\item
  Limites máximos de emissão de poluentes (ex.: padrão CONAMA para
  emissões veiculares);
\item
  Obrigatoriedade de uso de tecnologia de controle (ex.: catalisadores);
\item
  Zoneamento industrial (restrições de localização).
\end{itemize}

A principal desvantagem da regulação direta é que ela tipicamente não é
\textbf{custo-efetiva}: ao impor o mesmo padrão a todas as firmas, não
explora as diferenças nos custos marginais de abatimento. Firmas com
custos de abatimento baixos deveriam reduzir mais emissões, e firmas com
custos altos, menos --- algo que um imposto ou mercado de permissões faz
automaticamente.

\subsection{20.4.4 Mercados de permissão de emissão (cap and
trade)}\label{mercados-de-permissuxe3o-de-emissuxe3o-cap-and-trade}

\begin{definitionbox}{Mercado de Permissões de Emissão (Cap and Trade)}

Sistema em que o regulador fixa um \textbf{teto} (cap) total de emissões
e distribui \textbf{permissões negociáveis} entre as firmas. As firmas
podem comprar e vender permissões no mercado. Em equilíbrio, o preço da
permissão iguala o custo marginal de abatimento entre todas as firmas,
garantindo \textbf{custo-efetividade}.

\end{definitionbox}

Se o teto é fixado no nível socialmente ótimo de emissões \(\bar{E}\), e
o mercado de permissões é competitivo, o preço de equilíbrio da
permissão é exatamente igual ao imposto pigouviano ótimo:

{[} p\_E = t\^{}* = E'(q\^{}\{soc\}) {]}

A equivalência entre imposto e cap-and-trade sob certeza (Weitzman,
1974) é um resultado fundamental. Sob \textbf{incerteza}, porém, os
instrumentos diferem: o imposto fixa o preço e deixa a quantidade
flutuar; o cap fixa a quantidade e deixa o preço flutuar. A escolha
ótima depende das elasticidades relativas dos custos de abatimento e dos
danos marginais.

\begin{exresolvidobox}{Exercício Resolvido 20.3 —Cap-and-trade com duas firmas}

\textbf{Enunciado.} Duas firmas emitem 100 ton cada. Custos de
abatimento: \(CA_1(a_1) = 3a_1^2\), \(CA_2(a_2) = a_2^2\). O regulador
impõe redução total de 80 ton. (a) Encontre a alocação custo-efetiva.
(b) Calcule o preço de equilíbrio da permissão. (c) Compare com
abatimento uniforme.

\textbf{Solução.}

\textbf{(a)} Equalização de custos marginais:
\(CMgA_1 = CMgA_2 \implies 6a_1 = 2a_2 \implies a_2 = 3a_1\). Restrição:
\(a_1 + 3a_1 = 80 \implies a_1 = 20\), \(a_2 = 60\).

\textbf{(b)}
\(p = CMgA_1(20) = 6 \times 20 = 120 = CMgA_2(60) = 2 \times 60 = 120\).
\(\checkmark\)

\textbf{(c)} Custo eficiente: \(3(20)^2 + (60)^2 = 1200 + 3600 = 4800\).
Uniforme (\(a_i = 40\)): \(3(40)^2 + (40)^2 = 4800 + 1600 = 6400\).
Economia: \(6400 - 4800 = 1600\) (25\%).

\end{exresolvidobox}

\begin{exresolvidobox}{Exercício Resolvido 20.4 —Imposto vs. cap-and-trade sob incerteza (Weitzman)}

\textbf{Enunciado.} O regulador estima o custo marginal de abatimento
como \(CMgA = 2a\), mas o verdadeiro é \(CMgA = 4a\). O dano marginal é
constante: \(D' = 60\). Compare os resultados sob imposto \(t = 60\) e
cap \(\bar{a} = 30\).

\textbf{Solução.}

Com \textbf{imposto} \(t = 60\): a firma abate até
\(4a = 60 \implies a = 15\) (abaixo da meta de 30). Custo de abatimento:
\(2(15)^2 = 450\). Dano residual extra: \(60 \times 15 = 900\). Mas o
custo marginal = dano marginal no nível escolhido.

Com \textbf{cap} \(\bar{a} = 30\): obriga abatimento de 30, com custo
marginal \(4 \times 30 = 120 \gg 60 = D'\). Custo de abatimento:
\(2(30)^2 = 1800\). A meta é atingida, mas a um custo marginal muito
acima do dano marginal.

Conclusão: com dano marginal plano e custos subestimados, o
\textbf{imposto domina} (resultado de Weitzman, 1974).

\end{exresolvidobox}

\begin{center}\rule{0.5\linewidth}{0.5pt}\end{center}

\figurePlaceholder{/micro-book/graficos/cap20/webr-cap-trade.html}

\textbf{WebR 20.3 --- Cap-and-trade e alocação custo-efetiva.} Ajuste os
custos de abatimento das firmas e a meta de redução para ver como o
mercado de permissões aloca o esforço.

\begin{center}\rule{0.5\linewidth}{0.5pt}\end{center}

\begin{boxmundobox}{Box Mundo 20.2 — O EU Emissions Trading System (EU ETS)}

\textbf{Contexto:} Lançado em 2005, o EU ETS é o maior mercado de
carbono do mundo, cobrindo cerca de 40\% das emissões de gases de efeito
estufa da União Europeia. O sistema opera por cap-and-trade: a UE fixa
um teto total de emissões que diminui anualmente, e as empresas dos
setores cobertos (energia, indústria pesada, aviação intra-EU) devem
possuir permissões correspondentes às suas emissões.

\textbf{Evolução:} O preço das permissões (EU Allowances, EUA) flutuou
enormemente: de menos de €5 na Fase I (2005-2007, com excesso de
permissões distribuídas gratuitamente) a mais de €100 em 2023 (Fase IV,
com teto mais apertado e reserva de estabilidade de mercado). O preço
reflete a escassez crescente à medida que o teto diminui.

\textbf{Lições:} O EU ETS demonstra que mercados de permissões funcionam
--- as emissões dos setores cobertos caíram \textasciitilde35\% entre
2005 e 2022 --- mas que o \emph{design} importa criticamente. A alocação
gratuita inicial gerou windfall profits; o excesso de permissões
deprimiu preços por uma década; e a introdução de mecanismos de
estabilidade (como o Market Stability Reserve em 2019) foi necessária
para corrigir o desequilíbrio.

\end{boxmundobox}

\begin{boxmundobox}{Box Mundo 20.3 — O programa americano de SO₂: cap-and-trade em ação}

\textbf{Contexto:} O Acid Rain Program, criado pelo Clean Air Act de
1990 nos EUA, foi o primeiro grande sistema de cap-and-trade do mundo,
voltado para a redução de emissões de dióxido de enxofre (SO₂) --- o
principal causador da chuva ácida. O programa impôs um teto nacional de
emissões e distribuiu permissões negociáveis entre usinas termelétricas.

\textbf{Resultados:} As emissões de SO₂ caíram de 15,7 milhões de
toneladas em 1990 para menos de 3 milhões em 2010 --- uma redução de
mais de 80\%. O custo total de abatimento ficou substancialmente abaixo
das estimativas ex ante, em grande parte porque o mercado de permissões
permitiu que as firmas com menores custos de abatimento reduzissem mais.

\textbf{Análise econômica:} Estudos estimam que os custos de compliance
com o mercado foram 15-50\% menores do que teriam sido sob regulação
uniforme (command-and-control). O programa é amplamente considerado um
dos casos de maior sucesso de política ambiental baseada em instrumentos
de mercado.

\end{boxmundobox}

\begin{boxbrasilbox}{Box Brasil 20.1 — REDD+ e o mercado voluntário de carbono florestal}

\textbf{Contexto:} O mecanismo REDD+ (Redução de Emissões por
Desmatamento e Degradação Florestal) busca atribuir valor econômico ao
carbono estocado nas florestas, criando incentivos para a conservação. O
Brasil é um dos principais participantes, com programas de REDD+ em
vários estados amazônicos. O princípio é coaseano: criar um ``direito de
propriedade'' sobre o carbono florestal e permitir transações de
mercado.

\textbf{Dados:} O Fundo Amazônia, criado em 2008, recebeu aportes
superiores a R\$ 3,4 bilhões da Noruega e Alemanha até 2023. O Brasil
aprovou em 2024 a regulamentação do Sistema Brasileiro de Comércio de
Emissões (SBCE), um sistema de cap-and-trade que abrange grandes
emissores.

\textbf{Lições:} O caso brasileiro ilustra a aplicação integrada dos
instrumentos discutidos: subsídio pigouviano (Fundo Amazônia), mercado
de permissões (SBCE, REDD+), regulação direta (Código Florestal) e os
limites do Teorema de Coase quando os custos de transação são elevados e
as externalidades são difusas e globais.

\end{boxbrasilbox}

\chapter{Bens Públicos: Atributos e
Alocação}\label{bens-puxfablicos-atributos-e-alocauxe7uxe3o}

\section{20.5 Quem Paga Pelo Poste de Luz?: Atributos dos Bens
Públicos}\label{205}

\begin{definitionbox}{Bem Público}

Um \textbf{bem público puro} é aquele que possui simultaneamente duas
propriedades:

\begin{enumerate}
\def\labelenumi{\arabic{enumi}.}
\item
  \textbf{Não-rivalidade}: o consumo do bem por um indivíduo não reduz a
  quantidade disponível para outros. Formalmente, se \(G\) é a
  quantidade provida, cada consumidor pode consumir \(G\) integralmente:
  \(g_i = G\) para todo \(i\).
\item
  \textbf{Não-exclusão}: não é possível (ou é excessivamente custoso)
  impedir indivíduos de consumir o bem, mesmo que não paguem por ele.
\end{enumerate}

\end{definitionbox}

\subsection{Classificação de bens}\label{classificauxe7uxe3o-de-bens}

A combinação dessas duas propriedades gera uma classificação de bens em
quatro categorias:

\textbf{Tabela 20.1 --- Classificação de bens por rivalidade e
excludência}

{\def\LTcaptype{none} % do not increment counter
\begin{longtable}[]{@{}
  >{\raggedright\arraybackslash}p{(\linewidth - 4\tabcolsep) * \real{0.3333}}
  >{\raggedright\arraybackslash}p{(\linewidth - 4\tabcolsep) * \real{0.3333}}
  >{\raggedright\arraybackslash}p{(\linewidth - 4\tabcolsep) * \real{0.3333}}@{}}
\toprule\noalign{}
\begin{minipage}[b]{\linewidth}\raggedright
\end{minipage} & \begin{minipage}[b]{\linewidth}\raggedright
\textbf{Excludente}
\end{minipage} & \begin{minipage}[b]{\linewidth}\raggedright
\textbf{Não-excludente}
\end{minipage} \\
\midrule\noalign{}
\endhead
\bottomrule\noalign{}
\endlastfoot
\textbf{Rival} & \textbf{Bem privado} (alimento, vestuário, combustível)
& \textbf{Recurso comum} (cardume no oceano, pasto comunitário, água de
aquífero) \\
\textbf{Não-rival} & \textbf{Bem de clube} (TV a cabo, pedágio, parque
com entrada paga) & \textbf{Bem público puro} (defesa nacional,
iluminação pública, ar limpo) \\
\end{longtable}
}

A tabela abaixo apresenta exemplos brasileiros para cada categoria:

\textbf{Tabela 20.2 --- Exemplos brasileiros de cada categoria de bem}

{\def\LTcaptype{none} % do not increment counter
\begin{longtable}[]{@{}
  >{\raggedright\arraybackslash}p{(\linewidth - 4\tabcolsep) * \real{0.3333}}
  >{\raggedright\arraybackslash}p{(\linewidth - 4\tabcolsep) * \real{0.3333}}
  >{\raggedright\arraybackslash}p{(\linewidth - 4\tabcolsep) * \real{0.3333}}@{}}
\toprule\noalign{}
\begin{minipage}[b]{\linewidth}\raggedright
Categoria
\end{minipage} & \begin{minipage}[b]{\linewidth}\raggedright
Propriedades
\end{minipage} & \begin{minipage}[b]{\linewidth}\raggedright
Exemplos Brasileiros
\end{minipage} \\
\midrule\noalign{}
\endhead
\bottomrule\noalign{}
\endlastfoot
\textbf{Bem privado} & Rival e excludente & Pão de queijo, gasolina,
corte de cabelo, ingresso de cinema \\
\textbf{Bem público puro} & Não-rival e não-excludente & Defesa nacional
(Forças Armadas), iluminação de via pública, sinal de rádio aberto,
conhecimento científico básico \\
\textbf{Bem de clube} & Não-rival (até congestionamento) e excludente &
Netflix, rodovia com pedágio (Via Dutra, Rodovia dos Bandeirantes),
clube recreativo, sistema de streaming de futebol \\
\textbf{Recurso comum} & Rival e não-excludente & Peixes no rio
Amazonas, pastagem no semiárido, água do Aquífero Guarani, vagas de
estacionamento público \\
\end{longtable}
}

\begin{boxbrasilbox}{Brasil na Prática — O SUS como bem público: universalidade, free-riding e congestionamento}

\textbf{Contexto.} O Sistema Único de Saúde (SUS) --- criado pela
Constituição de 1988 --- é o maior sistema público de saúde do mundo,
cobrindo \textbf{150 milhões de brasileiros} que dependem exclusivamente
dele (IBGE/PNS, 2019). O SUS opera sob o princípio da
\textbf{universalidade}: qualquer pessoa, independentemente de
contribuição, pode acessar seus serviços. Isso o torna, por desenho
institucional, \textbf{não-excludente}.

\textbf{É bem público?} Tecnicamente, o SUS é um caso limítrofe na
Tabela 20.1. Consultas e cirurgias são \textbf{rivais} --- se o médico
atende um paciente, não pode atender outro simultaneamente. Mas o
\emph{acesso ao sistema} tem características de bem público: o direito à
saúde é não-excludente e a cobertura universal gera
\textbf{externalidades positivas} (vacinação, vigilância epidemiológica,
controle de endemias). Na prática, o SUS funciona como um
\textbf{recurso comum congestionável} --- não-excludente, mas rival
quando a demanda excede a capacidade instalada.

\textbf{Dados.} O Brasil gasta \textasciitilde3,9\% do PIB em saúde
pública (2023) --- abaixo da média da OCDE (\textasciitilde6,5\%). O SUS
realizou \textbf{4,1 bilhões de procedimentos ambulatoriais} e
\textbf{11,3 milhões de internações} em 2022 (DATASUS). O tempo médio de
espera para cirurgia eletiva chega a \textbf{180 dias} em alguns estados
(TCU, 2023). Em 2019, 28,5\% da população possuía plano de saúde privado
--- um mecanismo de \textbf{auto-exclusão} que alivia a pressão sobre o
recurso comum.

\textbf{Conexão com a teoria.} O SUS ilustra o dilema central da Seção
20.6: a provisão pública de um bem (quase) público financiado por
impostos gerais enfrenta inevitavelmente o \textbf{problema do carona}
--- quem paga mais impostos financia desproporcionalmente o sistema, e
não há como condicionar acesso à contribuição. A subprovisão (filas,
falta de leitos, defasagem salarial de profissionais) é a manifestação
empírica da inequação \(G^{priv} < G^*\).

\textbf{Fonte:} IBGE, \emph{Pesquisa Nacional de Saúde}, 2019; DATASUS,
\emph{Informações de Saúde}, 2022; TCU, \emph{Relatório de Auditoria
Operacional --- Saúde}, 2023.

\end{boxbrasilbox}

\begin{notebox}{Bens Públicos Locais e Globais}

A provisão de bens públicos pode ter escala local (iluminação de uma
rua), regional (defesa costeira), nacional (sistema judiciário) ou
global (estabilidade climática, erradicação de doenças). A escala
determina qual nível de governo (ou cooperação internacional) é mais
adequado para a provisão.

\end{notebox}

\begin{exresolvidobox}{Exercício Resolvido 20.5 — Classificação de bens}

\textbf{Enunciado.} Classifique os seguintes bens como privado, público
puro, de clube ou recurso comum: (a) Wi-Fi aberto em praça pública; (b)
Vacina contra febre amarela; (c) Peixe no Rio São Francisco; (d) Pedágio
da Rodovia dos Bandeirantes; (e) Código-fonte do Linux.

\textbf{Solução.}

\textbf{(a)} Wi-Fi aberto: \textbf{bem público puro} (até
congestionamento) --- não-excludente (qualquer pessoa pode acessar) e
não-rival (até o ponto de saturação da rede). Com congestionamento,
torna-se recurso comum.

\textbf{(b)} Vacina: \textbf{bem privado} --- rival (cada dose é usada
por uma pessoa) e excludente (administrada individualmente). Porém, a
\emph{imunidade de rebanho} gerada é um \textbf{bem público}
(externalidade positiva), o que justifica o financiamento público.

\textbf{(c)} Peixe no Rio São Francisco: \textbf{recurso comum} ---
rival (peixe capturado por um pescador não está disponível para outro) e
não-excludente (difícil impedir acesso ao rio).

\textbf{(d)} Pedágio: \textbf{bem de clube} --- não-rival (até
congestionamento) e excludente (só entra quem paga a tarifa).

\textbf{(e)} Código-fonte do Linux: \textbf{bem público puro} ---
não-rival (copiar não reduz disponibilidade) e não-excludente (licença
open-source garante acesso livre).

\end{exresolvidobox}

\begin{center}\rule{0.5\linewidth}{0.5pt}\end{center}

\section{20.6 Some Verticalmente, Não Horizontalmente: Bens Públicos e
Alocação de Recursos}\label{206}

\subsection{Condição de eficiência (Samuelson,
1954)}\label{condiuxe7uxe3o-de-eficiuxeancia-samuelson-1954}

Para um bem privado, a eficiência requer que todos os consumidores se
deparem com o mesmo preço, que iguala o custo marginal. Para um bem
público, a condição de eficiência é fundamentalmente diferente.

\begin{theorembox}{Condição de Samuelson para Bens Públicos}

A provisão eficiente de um bem público requer que a \textbf{soma das
taxas marginais de substituição} de todos os indivíduos iguale a
\textbf{taxa marginal de transformação} (custo marginal de produção):

{[} \sum\_\{i=1\}\^{}\{N\} TMS\_i\^{}\{G,x\} = TMT\^{}\{G,x\}
\label{eq:20.11} \tag{20.11} {]}

onde
\(TMS_i^{G,x} = \frac{\partial U_i / \partial G}{\partial U_i / \partial x_i}\)
é a taxa marginal de substituição entre o bem público \(G\) e o bem
privado \(x\) para o indivíduo \(i\).

\end{theorembox}

\textbf{Intuição}: Como o bem público é não-rival, todos consomem a
mesma quantidade \(G\). O benefício social marginal de uma unidade
adicional é a soma dos benefícios marginais de todos os indivíduos.
Eficiência requer que esse benefício agregado iguale o custo marginal.

Diferentemente, para um bem privado, a eficiência requer \(TMS_i = TMT\)
para cada \(i\) individualmente (não a soma).

\begin{intuitionbox}{Intuição Econômica}

\textbf{Em uma frase:} Para bens privados, some as \emph{quantidades}
demandadas (soma horizontal); para bens públicos, some as
\emph{valorações marginais} (soma vertical).

\textbf{Pense assim:} Quando três amigos querem pizza (bem privado),
cada um come a sua --- a demanda de mercado é a soma horizontal das
demandas individuais. Quando três amigos querem iluminação na rua (bem
público), todos ``consomem'' a mesma luz --- o que importa é quanto cada
um valoriza marginalmente aquela luz. A demanda social é a soma vertical
das disposições a pagar.

\textbf{Por que isso importa:} Essa distinção é a razão pela qual
mercados funcionam bem para pizza mas mal para iluminação pública. O
mecanismo de preços revela a demanda individual por bens privados (cada
um compra sua quantidade ao preço de mercado), mas não revela a
valoração individual por bens públicos (todos consomem a mesma
quantidade independentemente do que pagam).

\end{intuitionbox}

\begin{center}\rule{0.5\linewidth}{0.5pt}\end{center}

\figurePlaceholder{/micro-book/graficos/cap20/bem-publico.html}

\textbf{Figura 20.3 --- Provisão ótima de bem público.} A soma vertical
dos benefícios marginais individuais determina o benefício marginal
social. O nível ótimo (Samuelson) ocorre onde a soma dos BMg iguala o
CMg. Os preços de Lindahl mostram a contribuição personalizada de cada
consumidor.

\subsection{Subprovisão pelo mercado}\label{subprovisuxe3o-pelo-mercado}

O mercado tende a \textbf{subprover} bens públicos porque cada
indivíduo, ao decidir sua contribuição voluntária, considera apenas
\textbf{seu} benefício marginal, ignorando o benefício que gera para os
demais.

Considere \(N\) indivíduos idênticos com utilidade
\(U_i = u(x_i) + v(G)\), onde \(G = \sum g_i\) é a contribuição total ao
bem público. Cada indivíduo escolhe \(g_i\) para maximizar:

{[} u(W\_i - g\_i) + v!\left(g\_i + \sum\_\{j \neq i\} g\_j\right) {]}

A condição de primeira ordem é:

{[} u'(W\_i - g\_i) = v'(G) \implies TMS\_i = 1 \label{eq:20.12}
\tag{20.12} {]}

Mas a condição de eficiência requer \(\sum TMS_i = 1\), ou seja,
\(TMS_i = 1/N\) para indivíduos idênticos. Como \(1 > 1/N\), cada
indivíduo demanda benefício marginal excessivamente alto, resultando em
\(G^{priv} < G^*\): o bem público é \textbf{subprovido}.

\begin{exresolvidobox}{Exercício Resolvido 20.6 — Condição de Samuelson com utilidade Cobb-Douglas}

\textbf{Enunciado.} Dois consumidores com utilidade
\(U_i = x_i^{1/2} \cdot G^{1/2}\) e renda \(W_i = 100\). Custo do bem
público: \(C(G) = G\). Derive a condição de Samuelson e encontre
\(G^*\).

\textbf{Solução.}

\(TMS_i^{G,x} = \frac{\partial U_i / \partial G}{\partial U_i / \partial x_i} = \frac{(1/2)x_i^{1/2} G^{-1/2}}{(1/2)x_i^{-1/2} G^{1/2}} = \frac{x_i}{G}\).

Condição de Samuelson:
\(\frac{x_1}{G} + \frac{x_2}{G} = 1 \implies x_1 + x_2 = G\).

Restrição de recursos: \(x_1 + x_2 + G = 200\). Substituindo:
\(G + G = 200 \implies G^* = 100\).

Cada consumidor consome \(x_i = 50\) do bem privado.

\end{exresolvidobox}

\begin{center}\rule{0.5\linewidth}{0.5pt}\end{center}

\figurePlaceholder{/micro-book/graficos/cap20/webr-samuelson.html}

\textbf{WebR 20.4 --- Condição de Samuelson e subprovisão.} Compare o
nível ótimo (Samuelson) com o nível de contribuição voluntária (Nash)
variando o número de consumidores.

\chapter{Preços de Lindahl e o Problema do
Carona}\label{preuxe7os-de-lindahl-e-o-problema-do-carona}

\section{20.7 Um Preço Para Cada Cidadão: Preços de Lindahl}\label{207}

\begin{definitionbox}{Equilíbrio de Lindahl}

Mecanismo hipotético de provisão de bens públicos em que cada indivíduo
paga um \textbf{preço personalizado} (\(\tau_i\)) pelo bem público,
igual à sua taxa marginal de substituição. A soma dos preços de Lindahl
iguala o custo marginal de produção:

{[} \sum\_\{i=1\}\^{}\{N\} \tau\_i = CMg(G) \label{eq:20.13} \tag{20.13}
{]}

Cada indivíduo, enfrentando seu preço personalizado, demanda a mesma
quantidade \(G^*\), e a condição de Samuelson é satisfeita.

\end{definitionbox}

O equilíbrio de Lindahl é eficiente por construção, mas enfrenta um
problema prático fundamental: para implementá-lo, o governo (ou o
mecanismo) precisa conhecer as preferências individuais de cada cidadão
--- informação que os cidadãos têm incentivo para \textbf{não revelar de
forma verdadeira}, conforme discutiremos na seção sobre o problema do
carona.

\begin{exresolvidobox}{Exercício Resolvido 20.7 — Preços de Lindahl com três consumidores}

\textbf{Enunciado.} Três consumidores com benefícios marginais
\(BMg_A = 50 - 2G\), \(BMg_B = 40 - G\), \(BMg_C = 30 - G\). O custo
marginal é \(CMg = 60\). Encontre o nível eficiente e os preços de
Lindahl.

\textbf{Solução.}

Condição de Samuelson: \(\sum BMg_i = CMg\).

\((50 - 2G) + (40 - G) + (30 - G) = 60 \implies 120 - 4G = 60 \implies G^* = 15\).

Preços de Lindahl: \(\tau_A = 50 - 2(15) = 20\),
\(\tau_B = 40 - 15 = 25\), \(\tau_C = 30 - 15 = 15\).

Verificação: \(20 + 25 + 15 = 60 = CMg\). \(\checkmark\)

O consumidor com maior valoração marginal (B, no nível ótimo) paga mais.
Cada consumidor, enfrentando seu preço personalizado, demanda exatamente
\(G^* = 15\).

\end{exresolvidobox}

\begin{center}\rule{0.5\linewidth}{0.5pt}\end{center}

\figurePlaceholder{/micro-book/graficos/cap20/tragedia-comuns.html}

\textbf{Figura 20.4 --- Tragédia dos comuns.} Cada usuário iguala o
produto médio ao custo (equilíbrio privado), enquanto o ótimo social
requer igualar o produto marginal ao custo. A área sombreada indica a
região de sobreuso do recurso comum.

\begin{center}\rule{0.5\linewidth}{0.5pt}\end{center}

\section{20.8 Eu Não Pago, Mas Uso: O Problema do Carona (Free
Rider)}\label{208}

\begin{definitionbox}{Problema do Carona}

Ocorre quando indivíduos racionais subinvestem na provisão de um bem
público (ou na revelação de suas preferências) porque esperam se
beneficiar das contribuições alheias sem pagar por elas. O bem público é
não-excludente: mesmo quem não paga pode consumir.

\end{definitionbox}

No contexto do equilíbrio de Lindahl, cada indivíduo tem incentivo para
subreportar sua valoração pelo bem público, pagando um preço de Lindahl
menor e deixando que outros financiem a provisão. Formalmente, se o
preço de Lindahl é determinado pela valoração reportada
\(\hat{\tau}_i\):

{[} \hat{\tau}\_i \textless{} \tau\_i\^{}\{verdadeiro\} = TMS\_i {]}

O comportamento de carona é mais severo quando:

\begin{itemize}
\tightlist
\item
  O grupo é \textbf{grande} (a contribuição individual tem efeito
  desprezível sobre \(G\));
\item
  O bem é \textbf{puramente não-excludente} (não há como punir quem não
  contribui);
\item
  Não há \textbf{interação repetida} ou mecanismos de reputação.
\end{itemize}

\begin{infobox}{Nobel de Economia 2009 — Elinor Ostrom}

Elinor Ostrom recebeu o Prêmio Nobel ``por sua análise da governança
econômica, especialmente dos bens comuns'' --- o primeiro Nobel para uma
mulher em Economia. O trabalho de Ostrom desafiou a visão convencional
de que recursos comuns (commons) inevitavelmente sofrem a ``tragédia dos
comuns'' de Hardin (1968). Através de extenso trabalho de campo e
análise comparativa, Ostrom demonstrou que comunidades frequentemente
desenvolvem \textbf{instituições locais} --- regras de uso,
monitoramento entre pares, sanções graduais --- que sustentam a
cooperação e previnem o esgotamento dos recursos, sem necessidade de
privatização ou regulação estatal centralizada. Seu \emph{Governing the
Commons} (1990) documentou exemplos de florestas comunitárias na Suíça,
sistemas de irrigação no Nepal e pesqueiros costeiros no Japão que foram
geridos com sucesso por séculos.

\end{infobox}

\begin{notebox}{Evidência Experimental}

Experimentos de bens públicos (jogos de contribuição voluntária) mostram
que as contribuições iniciais são tipicamente 40-60\% do ótimo social,
mas declinam ao longo do tempo para 10-20\%, convergindo para a previsão
teórica de subprovisão. Mecanismos de punição entre pares (peer
punishment) e comunicação face a face aumentam significativamente as
contribuições (Fehr \& Gächter, 2000; Ostrom et al., 1992).

\end{notebox}

\begin{boxbrasilbox}{Box Brasil 20.2 — Cobrança pelo uso da água no Brasil: o preço do recurso comum}

\textbf{Contexto:} A Política Nacional de Recursos Hídricos (Lei
9.433/1997) instituiu a cobrança pelo uso da água bruta no Brasil,
reconhecendo a água como \textbf{recurso comum} sujeito à tragédia dos
comuns. Sem precificação, o uso excessivo levaria ao esgotamento dos
mananciais.

\textbf{Mecanismo:} A cobrança é implementada pelos Comitês de Bacia
Hidrográfica --- instâncias de governança local que ecoam os princípios
de Ostrom. Usuários que captam água, consomem ou lançam efluentes pagam
pelo uso, com valores definidos coletivamente. A receita é reinvestida
na bacia.

\textbf{Dados:} Na bacia do rio Paraíba do Sul (RJ/SP/MG), a cobrança
arrecadou mais de R\$ 200 milhões entre 2003 e 2023. O valor cobrado por
metro cúbico varia entre R\$ 0,01 e R\$ 0,02 --- baixo em termos
absolutos, mas simbolicamente importante como sinal de escassez.

\textbf{Análise:} A cobrança combina dois instrumentos: um preço
pigouviano (internalizar o custo do uso sobre outros usuários) e
governança coaseana descentralizada (os próprios usuários definem as
regras via comitê). É um dos exemplos mais sofisticados de gestão de
recursos comuns no mundo em desenvolvimento.

\end{boxbrasilbox}

\begin{boxbrasilbox}{Brasil na Prática — Tragédia dos comuns na pesca artesanal do Nordeste}

\textbf{Contexto.} O litoral nordestino abriga mais de \textbf{500 mil
pescadores artesanais} --- a maior concentração do país (MPA/IBGE,
2023). A pesca artesanal opera em regime de \textbf{acesso aberto}:
qualquer pescador pode lançar redes no mesmo trecho de mar ou estuário,
sem direitos de propriedade definidos sobre o estoque pesqueiro. Essa é
a configuração clássica da \textbf{tragédia dos comuns} de Hardin
(1968).

\textbf{Dados.} A produção pesqueira extrativa marinha do Nordeste caiu
de \textasciitilde350 mil toneladas/ano (anos 2000) para
\textasciitilde230 mil toneladas/ano (2020) --- uma redução de 34\% em
duas décadas (IBGE, \emph{Pesquisa da Pecuária Municipal}). A lagosta
vermelha (\emph{Panulirus argus}), espécie de alto valor comercial, teve
sua captura reduzida em mais de 80\% desde os anos 1990. O ICMBio estima
que 15\% das espécies marinhas exploradas comercialmente estão
\textbf{sobreexplotadas}.

\textbf{O modelo em ação.} Cada pescador iguala seu benefício marginal
privado (receita da captura) ao custo marginal privado (combustível, mão
de obra), ignorando o \textbf{custo externo} que impõe sobre os demais:
a redução do estoque pesqueiro futuro. O equilíbrio privado ocorre onde
\(PMe = w/p\) (produto médio = custo por unidade de esforço), mas o
ótimo social requer \(PMg = w/p\) (produto marginal = custo) ---
exatamente como mostra a Figura 20.4. Como \(PMe > PMg\) para recursos
congestionáveis, o esforço de pesca no equilíbrio privado excede
sistematicamente o nível sustentável.

\textbf{Soluções tentadas.} O \emph{defeso} (período de proibição da
pesca, com seguro-defeso pago pelo governo) tenta limitar o esforço
total. Reservas extrativistas marinhas (RESEX) --- como a RESEX de
Canavieiras (BA) e a RESEX Acaú-Goiana (PB/PE) --- implementam
governança à la Ostrom, com regras comunitárias de uso. Evidências
sugerem que RESEX com gestão participativa efetiva conseguem manter
estoques pesqueiros 25-40\% acima de áreas não protegidas (Lopes et al.,
2013).

\textbf{Fonte:} IBGE, \emph{Pesquisa da Pecuária Municipal}, 2020;
ICMBio, \emph{Livro Vermelho da Fauna Brasileira}, 2018; Lopes, P.F.M.
et al.~(2013). Sugestões para melhorar a gestão de unidades de
conservação costeiras e marinhas no Brasil. \emph{Ambiente \&
Sociedade}, 16(4), 93-112.

\end{boxbrasilbox}

\chapter{Votação e Mecanismos de
Revelação}\label{votauxe7uxe3o-e-mecanismos-de-revelauxe7uxe3o}

\section{20.9 Enfiar 200 Milhões de Preferências Numa Urna: Votação e
Alocação de Recursos}\label{209}

Na ausência de mecanismos de mercado eficientes para bens públicos, as
sociedades recorrem a processos políticos --- especialmente a
\textbf{votação} --- para decidir o nível de provisão.

\subsection{Votação por maioria
simples}\label{votauxe7uxe3o-por-maioria-simples}

Considere uma comunidade de \(N\) cidadãos que deve decidir o nível de
gastos \(G\) com um bem público, financiado por imposto uniforme
\(T = CMg(G)/N\) por pessoa. Cada cidadão \(i\) tem nível preferido
\(G_i^*\) que maximiza \(U_i(G) - T\).

\begin{theorembox}{Teorema do Eleitor Mediano}

Se as preferências dos eleitores são \textbf{unimodais} (single-peaked)
e a escolha é unidimensional, a regra de maioria simples seleciona o
nível preferido pelo \textbf{eleitor mediano} --- aquele cujo nível
preferido \(G_m^*\) é tal que metade dos eleitores prefere mais e metade
prefere menos.

Formalmente, se \(G_1^* \leq G_2^* \leq \cdots \leq G_N^*\), o resultado
da votação por maioria é \(G_m^* = G_{(N+1)/2}^*\) (para \(N\) ímpar).

\end{theorembox}

\textbf{Relação com eficiência}: O nível escolhido pelo eleitor mediano
geralmente \textbf{não} coincide com o nível eficiente de Samuelson,
exceto por coincidência. A provisão pode ser excessiva ou insuficiente
dependendo da distribuição de preferências e renda na população.

\subsection{Paradoxo de Condorcet}\label{paradoxo-de-condorcet}

Quando as preferências não são unimodais ou a escolha é
multidimensional, a votação por maioria pode produzir \textbf{ciclos} (A
vence B, B vence C, C vence A), não existindo um vencedor de Condorcet.

\textbf{Exemplo:} Três eleitores (1, 2, 3) e três alternativas (A, B,
C).

{\def\LTcaptype{none} % do not increment counter
\begin{longtable}[]{@{}llll@{}}
\toprule\noalign{}
Eleitor & 1.ª preferência & 2.ª preferência & 3.ª preferência \\
\midrule\noalign{}
\endhead
\bottomrule\noalign{}
\endlastfoot
1 & A & B & C \\
2 & B & C & A \\
3 & C & A & B \\
\end{longtable}
}

\begin{itemize}
\tightlist
\item
  A vs.~B: A vence (eleitores 1 e 3).
\item
  B vs.~C: B vence (eleitores 1 e 2).
\item
  C vs.~A: C vence (eleitores 2 e 3).
\end{itemize}

Resultado: ciclo A ≻ B ≻ C ≻ A. Não existe vencedor de Condorcet.

\subsection{Teorema da Impossibilidade de
Arrow}\label{teorema-da-impossibilidade-de-arrow}

\begin{theorembox}{Teorema da Impossibilidade de Arrow (1951)}

Não existe regra de agregação de preferências (com três ou mais
alternativas e dois ou mais indivíduos) que satisfaça simultaneamente as
seguintes quatro condições:

\begin{enumerate}
\def\labelenumi{\arabic{enumi}.}
\tightlist
\item
  \textbf{Domínio irrestrito}: a regra funciona para qualquer perfil de
  preferências individuais.
\item
  \textbf{Princípio de Pareto (unanimidade)}: se todos preferem A a B, a
  escolha social deve preferir A a B.
\item
  \textbf{Independência de alternativas irrelevantes (IIA)}: a ordenação
  social entre A e B depende apenas das preferências individuais entre A
  e B, não das preferências sobre C.
\item
  \textbf{Não-ditadura}: não existe um indivíduo cuja preferência sempre
  determine a escolha social.
\end{enumerate}

\end{theorembox}

O teorema de Arrow demonstra que qualquer sistema de votação com três ou
mais alternativas necessariamente viola pelo menos uma dessas condições.
Em outras palavras: não existe sistema de votação perfeito.

\textbf{Regras de votação e suas violações:}

{\def\LTcaptype{none} % do not increment counter
\begin{longtable}[]{@{}ll@{}}
\toprule\noalign{}
Regra & Violação principal \\
\midrule\noalign{}
\endhead
\bottomrule\noalign{}
\endlastfoot
Maioria simples & Pode gerar ciclos (viola transitividade) \\
Ditadura & Viola não-ditadura \\
Unanimidade (veto) & Viola completude \\
Contagem de Borda & Viola IIA \\
Votação por aprovação & Viola IIA \\
\end{longtable}
}

\begin{infobox}{Nobel de Economia 1972 — Kenneth Arrow}

Kenneth Arrow recebeu o Prêmio Nobel (compartilhado com John Hicks)
``por suas contribuições pioneiras à teoria do equilíbrio geral e à
teoria do bem-estar.'' O Teorema da Impossibilidade, publicado em
\emph{Social Choice and Individual Values} (1951), é um dos resultados
mais profundos das ciências sociais: demonstra que a democracia, por
mais desejável que seja, não pode ser ``perfeita'' em sentido formal ---
todo sistema de votação envolve trade-offs fundamentais. O teorema
inspirou décadas de pesquisa em teoria da escolha social, design de
mecanismos e economia política.

\end{infobox}

\begin{center}\rule{0.5\linewidth}{0.5pt}\end{center}

\section{20.10 Fazendo Você Dizer a Verdade: Mecanismos de Revelação de
Preferências}\label{2010}

O desafio central na provisão de bens públicos é induzir os indivíduos a
\textbf{revelar honestamente} suas preferências. O mecanismo de
Vickrey-Clarke-Groves (VCG) fornece uma solução elegante.

\begin{definitionbox}{Mecanismo VCG (Vickrey-Clarke-Groves)}

Classe de mecanismos de revelação de preferências em que cada indivíduo
paga um imposto igual ao \textbf{custo externo} que sua participação
impõe sobre os demais membros do grupo. Sob esse mecanismo, é estratégia
dominante para cada indivíduo revelar sua verdadeira valoração.

\end{definitionbox}

\subsection{Funcionamento do mecanismo de Clarke (imposto
pivotal)}\label{funcionamento-do-mecanismo-de-clarke-imposto-pivotal}

\begin{enumerate}
\def\labelenumi{\arabic{enumi}.}
\tightlist
\item
  Cada indivíduo \(i\) reporta sua valoração \(\hat{v}_i\) pelo bem
  público.
\item
  O bem é provido se \(\sum_i \hat{v}_i \geq C\) (custo de provisão).
\item
  O indivíduo \(i\) é \textbf{pivotal} se sua presença altera a decisão
  (o bem é provido com ele mas não sem ele).
\item
  Se \(i\) é pivotal, paga um imposto igual a:
\end{enumerate}

{[} t\_i = C - \sum\_\{j \neq i\} \hat{v}\_j \label{eq:20.14}
\tag{20.14} {]}

Esse imposto corresponde ao ``custo'' que a presença de \(i\) impõe
sobre os demais (que precisam financiar a diferença).

\textbf{Por que funciona}: Se \(i\) reporta \(\hat{v}_i > v_i\), pode
tornar-se pivotal quando não deveria, e pagar um imposto superior ao seu
benefício real. Se reporta \(\hat{v}_i < v_i\), pode impedir a provisão
de um bem que lhe seria benéfico. Em ambos os casos, o desvio não é
lucrativo, e a verdade é estratégia dominante.

\begin{notebox}{Limitações do VCG}

O mecanismo VCG não é equilibrado orçamentariamente (os impostos
pivotais não cobrem o custo total), pode ser vulnerável a coalizões, e
requer que as valorações sejam \textbf{quase-lineares} (utilidade
transferível). Na prática, sua aplicação direta é rara, mas o princípio
inspira mecanismos em leilões (como o leilão de Vickrey) e em design de
mercados.

\end{notebox}

\begin{exresolvidobox}{Exercício Resolvido 20.8 — Mecanismo de Clarke para bem público discreto}

\textbf{Enunciado.} Três vizinhos decidem se constroem uma calçada
(custo = R\$900, dividido igualmente em R\$300 cada). Valorações:
\(v_A = 400\), \(v_B = 350\), \(v_C = 100\). (a) A calçada deve ser
construída? (b) Quem é pivotal? (c) Calcule os impostos pivotais.

\textbf{Solução.}

\textbf{(a)} \(\sum v_i = 400 + 350 + 100 = 850 < 900\). Pelo critério
de eficiência, a calçada \textbf{não} deve ser construída
(\(\sum v_i < C\)).

Mas calculemos os excedentes líquidos:
\(\sum (v_i - c_i) = (400 - 300) + (350 - 300) + (100 - 300) = 100 + 50 - 200 = -50 < 0\).
Confirmado: não construir.

\textbf{(b)} Sem A:
\(\sum_{j \neq A}(v_j - c_j) = 50 + (-200) = -150 < 0\). Decisão sem A:
não construir. Decisão com A: não construir. A \textbf{não é pivotal}.

Sem B: \(\sum_{j \neq B}(v_j - c_j) = 100 + (-200) = -100 < 0\). Decisão
sem B: não construir. Decisão com B: não construir. B \textbf{não é
pivotal}.

Sem C: \(\sum_{j \neq C}(v_j - c_j) = 100 + 50 = 150 > 0\). Decisão sem
C: \textbf{construir}. Decisão com C: não construir. C \textbf{é
pivotal} --- sua presença muda a decisão de ``construir'' para ``não
construir''.

\textbf{(c)} C é pivotal e impede um ganho líquido de 150 para os
demais. Contudo, como a decisão final é ``não construir'', não há
imposto a cobrar neste caso --- o imposto pivotal só se aplica quando o
agente pivotal altera a decisão \emph{para provisão}. Se reformularmos
com valorações mais altas (\(v_C = 200\)), a construção ocorreria e C
pagaria um imposto pelo dano imposto aos demais.

\end{exresolvidobox}

\begin{center}\rule{0.5\linewidth}{0.5pt}\end{center}

\figurePlaceholder{/micro-book/graficos/cap20/webr-vcg.html}

\textbf{WebR 20.5 --- Mecanismo VCG e imposto pivotal.} Insira as
valorações dos participantes e o custo do bem público para ver quem é
pivotal e os impostos correspondentes.

\begin{center}\rule{0.5\linewidth}{0.5pt}\end{center}

\begin{boxbrasilbox}{Box Brasil 20.3 — Desmatamento na Amazônia e precificação de carbono}

\textbf{Contexto:} O desmatamento na Amazônia brasileira é um caso
emblemático de \textbf{externalidade negativa de escala global}. A
remoção da floresta gera custos externos que incluem emissões de CO₂
(contribuindo para a mudança climática), perda de biodiversidade,
alteração do ciclo hidrológico (inclusive os ``rios voadores'' que
transportam umidade para o Centro-Sul do Brasil) e erosão do solo.

\textbf{Dimensão do problema:} Segundo dados do INPE/PRODES, o
desmatamento acumulado na Amazônia Legal superou 85 milhões de hectares
até 2023 (cerca de 17\% da floresta original). Após queda expressiva
entre 2004 e 2012 (de 27.772 km² para 4.571 km² anuais), as taxas
voltaram a subir, atingindo 13.235 km² em 2021, antes de recuar para
cerca de 9.001 km² em 2023 com o reforço de políticas de fiscalização.

\textbf{Análise econômica:} O desmatamento persiste porque os
\textbf{benefícios privados} (expansão da pecuária, agricultura,
extração madeireira) excedem os \textbf{custos privados} para os agentes
locais, embora os \textbf{custos sociais} --- incluindo danos climáticos
globais --- superem largamente os benefícios. Trata-se de uma
externalidade negativa clássica.

\textbf{Instrumentos de política:} O Fundo Amazônia (subsídio
pigouviano), REDD+ (mecanismo coaseano de direitos sobre carbono
florestal), o Sistema Brasileiro de Comércio de Emissões --- SBCE
(cap-and-trade, regulamentado em 2024) e o Código Florestal
(command-and-control). Este caso ilustra a aplicação integrada de todos
os instrumentos discutidos no capítulo.

\end{boxbrasilbox}

\begin{boxmundobox}{Box Mundo 20.4 — BBC e NHK: financiamento de radiodifusão pública}

\textbf{Contexto:} A radiodifusão é um bem público clássico (não-rival e
não-excludente via ondas de rádio/TV). Como financiar sua provisão sem
subprovisão?

\textbf{Modelos:} A BBC (Reino Unido) é financiada pela \emph{licence
fee} --- uma taxa anual obrigatória de £159 (2023) para qualquer
domicílio com aparelho de TV. A NHK (Japão) cobra uma taxa semelhante de
¥13.650/ano. Ambos os modelos evitam o problema do carona tornando o
pagamento compulsório.

\textbf{Alternativas:} Nos EUA, a PBS é financiada por uma combinação de
doações voluntárias, subsídio federal e patrocínio corporativo ---
resultando em orçamento muito menor per capita que BBC ou NHK. O modelo
americano ilustra a subprovisão prevista pela teoria quando o
financiamento é voluntário.

\textbf{Análise:} A taxa compulsória é economicamente equivalente a um
imposto de Lindahl uniforme --- não personalizado (todos pagam o mesmo),
mas que garante provisão. A discussão sobre o nível adequado da taxa e a
abrangência do serviço público é essencialmente um debate sobre a
condição de Samuelson na prática.

\end{boxmundobox}

\chapter{Exercícios e ANPEC --- Capítulo
20}\label{exercuxedcios-e-anpec-capuxedtulo-20}

\begin{tipbox}{Dados para exercicios empiricos}

\begin{itemize}
\tightlist
\item
  \textbf{SEEG (emissoes de GEE):}
  \href{https://seeg.eco.br/}{seeg.eco.br} --- Sistema de Estimativas de
  Emissoes de Gases de Efeito Estufa do Observatorio do Clima. Dados por
  setor e estado para estimar custos sociais de externalidades.
\item
  \textbf{INPE (desmatamento):}
  \href{http://terrabrasilis.dpi.inpe.br/}{terrabrasilis.dpi.inpe.br}
  --- Dados PRODES e DETER para analisar a tragedia dos comuns na
  Amazonia.
\item
  \textbf{ANA (recursos hidricos):}
  \href{https://www.snirh.gov.br/}{snirh.gov.br} --- Dados sobre
  outorgas e cobranca pelo uso da agua --- precificacao de externalidade
  em acao.
\item
  \textbf{Notebook:} N17 --- Modelo DICE simplificado (trajetorias
  otimas de emissao).
\end{itemize}

\end{tipbox}

\section{🧠 Revisão Rápida}\label{revisuxe3o-ruxe1pida-2}

\begin{infobox}{1. O que distingue uma externalidade tecnológica de uma externalidade pecuniária?}

Uma externalidade \textbf{tecnológica} (real) afeta diretamente a função
de produção ou utilidade de outro agente, sem passar pelo sistema de
preços (ex.: poluição). Uma externalidade \textbf{pecuniária} opera
\emph{via preços} (ex.: uma firma grande que reduz o preço de mercado e
prejudica concorrentes). Apenas externalidades tecnológicas geram
ineficiência alocativa e constituem falhas de mercado.

\end{infobox}

\begin{infobox}{2. Por que o imposto pigouviano restaura a eficiência?}

O imposto pigouviano \(t^* = E'(q^{soc})\) faz a firma pagar exatamente
o custo marginal externo que impõe à sociedade. Ao incluir esse custo no
cálculo de maximização de lucro, a firma ``internaliza'' a externalidade
e escolhe o nível de produção socialmente ótimo (onde CMg social =
preço).

\end{infobox}

\begin{infobox}{3. Qual é a diferença fundamental entre a condição de eficiência para bens privados e bens públicos?}

Para bens privados: \(TMS_i = TMT\) para cada indivíduo (igualdade
individual). Para bens públicos: \(\sum_i TMS_i = TMT\) (soma das TMS).
A diferença reflete a não-rivalidade: todos consomem a mesma quantidade
do bem público, então o benefício social é a \emph{soma} dos benefícios
individuais.

\end{infobox}

\begin{infobox}{4. Por que o Teorema de Coase tem aplicação prática limitada?}

Porque requer: (i) custos de transação nulos, (ii) direitos de
propriedade bem definidos, (iii) poucas partes envolvidas e (iv)
informação simétrica. Na prática, externalidades difusas (como poluição
do ar) envolvem milhões de partes, custos de negociação proibitivos e
assimetria de informação --- tornando a solução coaseana inviável.

\end{infobox}

\begin{infobox}{5. O que garante que o mecanismo VCG induz revelação verdadeira?}

No VCG, cada agente paga um imposto igual ao custo externo que sua
presença impõe sobre os demais. Se subreporta sua valoração, pode
impedir a provisão de um bem que lhe seria benéfico. Se sobrereporta,
pode tornar-se pivotal quando não deveria, pagando mais do que vale. Em
ambos os casos, desviar da verdade não é lucrativo --- a verdade é
estratégia dominante.

\end{infobox}

\begin{center}\rule{0.5\linewidth}{0.5pt}\end{center}

\section{📋 Resumo do Capítulo}\label{resumo-do-capuxedtulo-2}

\begin{enumerate}
\def\labelenumi{\arabic{enumi}.}
\tightlist
\item
  \textbf{Externalidades} surgem quando as ações de um agente afetam
  diretamente o bem-estar de outros fora do sistema de preços, gerando
  divergência entre custos/benefícios privados e sociais.
\item
  Externalidades negativas levam a \textbf{sobreprodução}; positivas, a
  \textbf{subprodução} em relação ao ótimo social.
\item
  As principais soluções são: \textbf{imposto/subsídio pigouviano}
  (internaliza o custo externo via preço), \textbf{Teorema de Coase}
  (negociação privada sob custos de transação nulos), \textbf{regulação
  direta} (limites quantitativos) e \textbf{cap-and-trade} (mercado de
  permissões).
\item
  \textbf{Bens públicos} (não-rivais e não-excludentes) são subprovidos
  pelo mercado devido ao \textbf{problema do carona}. A condição de
  Samuelson (\(\sum TMS = TMT\)) difere fundamentalmente da eficiência
  para bens privados.
\item
  \textbf{Preços de Lindahl} resolvem o problema em teoria, mas requerem
  informação que os agentes têm incentivo para não revelar. O
  \textbf{mecanismo VCG} torna a revelação verdadeira uma estratégia
  dominante.
\item
  O \textbf{Teorema da Impossibilidade de Arrow} mostra que nenhum
  sistema de votação satisfaz simultaneamente todos os critérios
  desejáveis de agregação de preferências.
\end{enumerate}

\begin{center}\rule{0.5\linewidth}{0.5pt}\end{center}

\section{🔑 Conceitos-Chave}\label{conceitos-chave-2}

\textbf{Tabela 20.3 --- Conceitos-chave do Capítulo 20}

{\def\LTcaptype{none} % do not increment counter
\begin{longtable}[]{@{}
  >{\raggedright\arraybackslash}p{(\linewidth - 2\tabcolsep) * \real{0.5000}}
  >{\raggedright\arraybackslash}p{(\linewidth - 2\tabcolsep) * \real{0.5000}}@{}}
\toprule\noalign{}
\begin{minipage}[b]{\linewidth}\raggedright
Conceito
\end{minipage} & \begin{minipage}[b]{\linewidth}\raggedright
Definição Resumida
\end{minipage} \\
\midrule\noalign{}
\endhead
\bottomrule\noalign{}
\endlastfoot
\textbf{Externalidade} & Efeito da ação de um agente sobre outro, fora
do sistema de preços \\
\textbf{Imposto pigouviano} & Imposto unitário = custo marginal externo
no ótimo social \\
\textbf{Teorema de Coase} & Barganha privada alcança eficiência se
custos de transação = 0 e direitos de propriedade definidos \\
\textbf{Cap-and-trade} & Teto de emissões + permissões negociáveis →
custo-efetividade \\
\textbf{Bem público} & Bem não-rival e não-excludente \\
\textbf{Condição de Samuelson} & \(\sum TMS_i = TMT\) para provisão
eficiente de bem público \\
\textbf{Preço de Lindahl} & Preço personalizado = TMS individual; soma =
CMg \\
\textbf{Problema do carona} & Incentivo a subinvestir na provisão de
bens públicos \\
\textbf{Teorema de Arrow} & Impossibilidade de regra de agregação
perfeita com ≥3 alternativas \\
\textbf{Mecanismo VCG} & Imposto pivotal torna revelação verdadeira
estratégia dominante \\
\end{longtable}
}

\begin{center}\rule{0.5\linewidth}{0.5pt}\end{center}

\section{✏️ Exercícios}\label{exercuxedcios-2}

\textbf{Exercício 20.1.} \{\#ex-20-1\} Uma fábrica de celulose produz
\(q\) toneladas com custo total \(C(q) = 10q + q^2\) e vende ao preço
\(P = 110\). A produção gera poluição com custo externo \(E(q) = 2q^2\).

\begin{enumerate}
\def\labelenumi{(\alph{enumi})}
\tightlist
\item
  Determine a quantidade produzida pela firma na ausência de regulação.
\item
  Determine a quantidade socialmente ótima.
\item
  Calcule o imposto pigouviano ótimo \(t^*\) e verifique que ele induz a
  firma a produzir a quantidade eficiente.
\item
  Calcule a perda de peso morto associada à ausência de regulação.
\end{enumerate}

🔑 Solução (Cap. 20)

\begin{center}\rule{0.5\linewidth}{0.5pt}\end{center}

\textbf{Exercício 20.2.} \{\#ex-20-2\} Considere três indivíduos (\(A\),
\(B\), \(C\)) com as seguintes valorações marginais por um bem público
\(G\) (medido em unidades):

{\def\LTcaptype{none} % do not increment counter
\begin{longtable}[]{@{}lllll@{}}
\toprule\noalign{}
\(G\) & \(BMg_A\) & \(BMg_B\) & \(BMg_C\) & \(CMg\) \\
\midrule\noalign{}
\endhead
\bottomrule\noalign{}
\endlastfoot
1 & 40 & 30 & 20 & 60 \\
2 & 35 & 25 & 15 & 60 \\
3 & 25 & 20 & 10 & 60 \\
4 & 15 & 10 & 5 & 60 \\
5 & 5 & 5 & 2 & 60 \\
\end{longtable}
}

\begin{enumerate}
\def\labelenumi{(\alph{enumi})}
\tightlist
\item
  Determine o nível eficiente de provisão do bem público.
\item
  Calcule os preços de Lindahl para cada indivíduo no nível eficiente.
\item
  Explique por que a provisão voluntária (contribuição individual)
  resultaria em \(G < G^*\).
\item
  Verifique se os preços de Lindahl são equilibrados orçamentariamente.
\end{enumerate}

🔑 Solução (Cap. 20)

\begin{center}\rule{0.5\linewidth}{0.5pt}\end{center}

\textbf{Exercício 20.3.} \{\#ex-20-3\} Duas firmas emitem poluentes. A
firma 1 tem custo de abatimento \(CA_1(a_1) = 2a_1^2\) e a firma 2,
\(CA_2(a_2) = a_2^2\), onde \(a_i\) é o abatimento em toneladas. Cada
firma emite 100 toneladas sem regulação. O regulador deseja reduzir as
emissões totais em 60 toneladas.

\begin{enumerate}
\def\labelenumi{(\alph{enumi})}
\tightlist
\item
  Qual é a alocação custo-efetiva de abatimento entre as duas firmas?
\item
  Qual seria o custo total se o regulador impusesse abatimento uniforme
  (30 toneladas cada)?
\item
  Compare com o custo da solução custo-efetiva e calcule a economia.
\item
  Se for utilizado um mercado de permissões (cap-and-trade), qual será o
  preço de equilíbrio da permissão?
\end{enumerate}

🔑 Solução (Cap. 20)

\begin{center}\rule{0.5\linewidth}{0.5pt}\end{center}

\textbf{Exercício 20.4.} \{\#ex-20-4\} Uma comunidade com 100 moradores
deve decidir se constrói um parque público (bem público discreto) ao
custo de R\$50.000. Cada morador \(i\) tem valoração \(v_i\) pela praça,
distribuída uniformemente entre R\$200 e R\$800.

\begin{enumerate}
\def\labelenumi{(\alph{enumi})}
\tightlist
\item
  O parque deve ser construído? (Compare a soma das valorações esperadas
  com o custo.)
\item
  Se a construção for financiada por contribuição voluntária uniforme de
  R\$500 por morador, quantos moradores estariam dispostos a contribuir?
\item
  Explique o problema do carona nesta situação.
\item
  Descreva como um mecanismo de imposto pivotal (Clarke) funcionaria
  neste contexto.
\end{enumerate}

🔑 Solução (Cap. 20)

\begin{center}\rule{0.5\linewidth}{0.5pt}\end{center}

\textbf{Exercício 20.5.} \{\#ex-20-5\} Considere uma economia com dois
bens --- um privado (\(x\)) e um público (\(G\)) --- e dois consumidores
com utilidades:

{[} U\_1 = \ln(x\_1) + 2\ln(G), \quad U\_2 = \ln(x\_2) + \ln(G) {]}

O preço do bem privado é 1 e o custo marginal do bem público é \(c\).
Cada consumidor tem renda \(W\).

\begin{enumerate}
\def\labelenumi{(\alph{enumi})}
\tightlist
\item
  Derive a condição de Samuelson para a provisão eficiente de \(G\).
\item
  Determine o nível eficiente \(G^*\) em função de \(W\) e \(c\).
\item
  Se cada consumidor contribui voluntariamente (equilíbrio de Nash em
  contribuições), determine o nível de \(G\) de equilíbrio.
\item
  Mostre que \(G^{Nash} < G^*\) e interprete economicamente.
\end{enumerate}

🔑 Solução (Cap. 20)

\begin{center}\rule{0.5\linewidth}{0.5pt}\end{center}

\textbf{Exercício 20.6.} \{\#ex-20-6\} Uma indústria siderúrgica produz
\(q\) toneladas de aço com custo total \(C(q) = 20q + 0{,}5q^2\), vende
ao preço \(P = 80\). Custo externo: \(E(q) = q^2\).

\begin{enumerate}
\def\labelenumi{(\alph{enumi})}
\tightlist
\item
  Encontre a quantidade sem regulação e a socialmente ótima.
\item
  Calcule o imposto pigouviano ótimo e verifique.
\item
  Calcule a receita fiscal e o peso morto eliminado pelo imposto.
\item
  Se o regulador fixar incorretamente \(t = 30\), qual será a quantidade
  e o peso morto residual?
\end{enumerate}

🔑 Solução (Cap. 20)

\begin{center}\rule{0.5\linewidth}{0.5pt}\end{center}

\textbf{Exercício 20.7.} \{\#ex-20-7\} Um fazendeiro cria \(n\) cabeças
de gado com lucro \(\pi_F(n) = 120n - 3n^2\). O gado causa dano à
lavoura vizinha de \(D(n) = 2n^2\).

\begin{enumerate}
\def\labelenumi{(\alph{enumi})}
\tightlist
\item
  Encontre o equilíbrio sem negociação.
\item
  Encontre o número socialmente ótimo de cabeças.
\item
  Se o fazendeiro tem o direito de criar gado livremente, mostre que a
  negociação coaseana leva ao ótimo social. Determine o intervalo de
  compensação.
\item
  Se o lavrador tem o direito a não sofrer danos, mostre que a
  negociação leva ao mesmo ótimo. Determine o intervalo de compensação.
\end{enumerate}

🔑 Solução (Cap. 20)

\begin{center}\rule{0.5\linewidth}{0.5pt}\end{center}

\textbf{Exercício 20.8.} \{\#ex-20-8\} \(N = 50\) consumidores
idênticos, renda \(W = 100\), utilidade \(U_i = x_i + 10\sqrt{G}\),
custo \(C(G) = G\).

\begin{enumerate}
\def\labelenumi{(\alph{enumi})}
\tightlist
\item
  Derive a condição de Samuelson e encontre \(G^*\).
\item
  Encontre o equilíbrio de Nash de contribuição voluntária \(G^{Nash}\).
\item
  Calcule a razão \(G^{Nash}/G^*\) e interprete.
\item
  O que acontece se \(N\) aumenta para 200? O problema do carona se
  agrava ou se atenua?
\end{enumerate}

🔑 Solução (Cap. 20)

\begin{center}\rule{0.5\linewidth}{0.5pt}\end{center}

\textbf{Exercício 20.9.} \{\#ex-20-9\} (\textbf{Tragédia dos comuns.})
Num lago de acesso aberto, a captura total é \(Q = E(100 - E)\), onde
\(E\) é o esforço total de pesca. Preço do peixe: \(p = 1\). Custo por
unidade de esforço: \(w = 10\). Há \(n = 10\) pescadores idênticos.

\begin{enumerate}
\def\labelenumi{(\alph{enumi})}
\tightlist
\item
  Encontre o esforço total de equilíbrio (acesso aberto).
\item
  Encontre o esforço socialmente eficiente.
\item
  Calcule o peso morto do acesso aberto.
\item
  Proponha duas políticas para alcançar o ótimo social.
\end{enumerate}

🔑 Solução (Cap. 20)

\begin{center}\rule{0.5\linewidth}{0.5pt}\end{center}

\textbf{Exercício 20.10.} \{\#ex-20-10\} (\textbf{Cap-and-trade
vs.~imposto sob incerteza.}) Três firmas com custos marginais de
abatimento: \(CMgA_1 = 2a_1\), \(CMgA_2 = 4a_2\), \(CMgA_3 = a_3\).
Meta: redução total de 100 toneladas.

\begin{enumerate}
\def\labelenumi{(\alph{enumi})}
\tightlist
\item
  Encontre a alocação custo-efetiva e o preço de equilíbrio das
  permissões. Mostre a equivalência com o imposto uniforme ótimo.
\item
  Compare o custo total da alocação custo-efetiva com abatimento
  uniforme.
\item
  Se os verdadeiros custos forem 50\% maiores que o estimado e o dano
  marginal for constante (\(D' = 40\)), compare imposto \(t = 40\) com
  cap = 100.
\item
  Relacione o resultado com Weitzman (1974).
\end{enumerate}

🔑 Solução (Cap. 20)

\begin{center}\rule{0.5\linewidth}{0.5pt}\end{center}

\section{🏆 Vem, ANPEC!}\label{vem-anpec-2}

\begin{infobox}{ANPEC 2017, Questão 10}

Sobre externalidades e bens públicos, classifique as afirmativas como
Verdadeiras (V) ou Falsas (F):

\begin{enumerate}
\def\labelenumi{(\arabic{enumi})}
\setcounter{enumi}{-1}
\item
  Na presença de externalidade negativa de produção, o equilíbrio
  competitivo gera quantidade superior à socialmente ótima.
\item
  O Teorema de Coase afirma que, havendo custos de transação nulos e
  direitos de propriedade bem definidos, a negociação privada sempre
  conduzirá à alocação eficiente, independentemente da atribuição
  inicial dos direitos.
\item
  A condição de Samuelson para provisão eficiente de bens públicos
  estabelece que cada consumidor deve igualar sua taxa marginal de
  substituição ao custo marginal de produção.
\item
  Em um sistema de cap-and-trade, o preço de equilíbrio da permissão de
  emissão iguala os custos marginais de abatimento entre as firmas
  participantes.
\item
  O problema do carona (free rider) tende a ser mais severo em grupos
  maiores.
\end{enumerate}

??? tip ``Gabarito'' (0) \textbf{V} --- Com externalidade negativa, CMg
social \textgreater{} CMg privado, levando a sobreprodução.

\begin{verbatim}
(1) **V** — Enunciado correto do Teorema de Coase.

(2) **F** — A condição de Samuelson exige que a **soma** das TMS iguale o CMg, não a TMS individual.

(3) **V** — No equilíbrio do mercado de permissões, os custos marginais de abatimento se equalizam.

(4) **V** — Em grupos maiores, a contribuição individual tem efeito menor sobre a provisão total, aumentando o incentivo ao free-riding (Olson, 1965).
\end{verbatim}

\end{infobox}

\begin{infobox}{ANPEC 2022, Questão 08}

Considere uma economia com dois consumidores (1 e 2), um bem privado
\(x\) e um bem público \(G\). As utilidades são
\(U_1 = x_1 + 8\sqrt{G}\) e \(U_2 = x_2 + 4\sqrt{G}\). O custo do bem
público é \(C(G) = G\). Cada consumidor tem renda \(W = 100\).

\begin{enumerate}
\def\labelenumi{(\arabic{enumi})}
\setcounter{enumi}{-1}
\item
  A condição de Samuelson para provisão eficiente é
  \(\frac{4}{\sqrt{G}} + \frac{2}{\sqrt{G}} = 1\).
\item
  O nível eficiente de bem público é \(G^* = 36\).
\item
  No equilíbrio de Nash de contribuição voluntária, \(G^{Nash} = 16\).
\item
  A provisão voluntária produz \(G^{Nash}/G^* = 4/9\) do nível
  eficiente.
\item
  A introdução de um mecanismo de Clarke resolveria o problema de
  subprovisão induzindo revelação verdadeira das preferências.
\end{enumerate}

??? tip ``Gabarito'' (0) \textbf{V} --- \(TMS_1 = 4/\sqrt{G}\),
\(TMS_2 = 2/\sqrt{G}\). Soma = \(6/\sqrt{G} = 1\)\ldots{} Espere:
\(TMS_1 = 8/(2\sqrt{G}) = 4/\sqrt{G}\),
\(TMS_2 = 4/(2\sqrt{G}) = 2/\sqrt{G}\). Soma:
\((4+2)/\sqrt{G} = 6/\sqrt{G}\). A questão diz
\(4/\sqrt{G} + 2/\sqrt{G}\), que é o mesmo que \(6/\sqrt{G} = 1\).
\textbf{V}.

\begin{verbatim}
(1) **V** — $6/\sqrt{G} = 1 \implies \sqrt{G} = 6 \implies G^* = 36$.

(2) **V** — No Nash, cada consumidor iguala sua TMS individual ao CMg. Consumidor 1: $4/\sqrt{G} = 1 \implies G = 16$. Se 1 contribui sozinho, $G = 16$. Consumidor 2: $2/\sqrt{G} = 1 \implies G = 4$. Mas 2 não contribui se 1 já provê $G = 16$ (pois $2/\sqrt{16} = 0{,}5 < 1$). Nash: $G = 16$. **V**.

(3) **V** — $16/36 = 4/9$. **V**.

(4) **V** — O mecanismo de Clarke torna a revelação verdadeira estratégia dominante, resolvendo o problema do carona. **V**.
\end{verbatim}

\end{infobox}

\begin{infobox}{ANPEC 2020, Questão 12}

Sobre o Teorema da Impossibilidade de Arrow:

\begin{enumerate}
\def\labelenumi{(\arabic{enumi})}
\setcounter{enumi}{-1}
\item
  O teorema se aplica quando há pelo menos três alternativas e dois
  indivíduos.
\item
  A condição de independência de alternativas irrelevantes (IIA) exige
  que a ordenação social entre A e B dependa apenas das preferências
  individuais entre A e B.
\item
  A contagem de Borda satisfaz todas as condições de Arrow.
\item
  O Teorema do Eleitor Mediano é válido quando as preferências são
  unimodais e a escolha é unidimensional.
\item
  O resultado do eleitor mediano sempre coincide com o nível eficiente
  de Samuelson para bens públicos.
\end{enumerate}

??? tip ``Gabarito'' (0) \textbf{V} --- Condição necessária para o
teorema.

\begin{verbatim}
(1) **V** — Definição correta de IIA.

(2) **F** — A contagem de Borda viola IIA (a ordenação entre A e B pode mudar com a introdução de C).

(3) **V** — Condições suficientes para o Teorema do Eleitor Mediano.

(4) **F** — O eleitor mediano reflete a preferência de um indivíduo específico, não a soma das TMS. Coincidência só ocorre em casos especiais.
\end{verbatim}

\end{infobox}

\chapter{Pesquisa em Ação --- Capítulo
20}\label{pesquisa-em-auxe7uxe3o-capuxedtulo-20}

\section{🔬 Pesquisa em Ação}\label{pesquisa-em-auxe7uxe3o-2}

\begin{pesquisabox}{[Eliasson, J.; Hultkrantz, L.; Nerhagen, L.; Rosqvist, L. S. (2009). The Stockholm Congestion-Charging Trial 2006: Overview of Effects. *Transportation Research Part A*, 43(3), 240–250.](https://doi.org/10.1016/j.tra.2008.09.007)}

\textbf{DOI:}
\href{https://doi.org/10.1016/j.tra.2008.09.007}{10.1016/j.tra.2008.09.007}

\textbf{Contexto.} O congestionamento viário é uma externalidade
negativa clássica. Em 2006, Estocolmo implementou um período
experimental de pedágio urbano (taxa pigouviana sobre congestionamento),
seguido de referendo. O paper avalia os efeitos do experimento.

\textbf{Método.} Comparação antes-depois das condições de tráfego na
zona tarifada, usando dados de fluxo veicular, tempos de viagem e
emissões. O período experimental de 7 meses permitiu avaliação
abrangente.

\textbf{Resultado.} O tráfego na zona central caiu 20-25\% durante o
período de cobrança. As emissões de CO₂ na área tarifada reduziram-se em
10-14\%. Os tempos de viagem diminuíram significativamente.
Crucialmente, a população, que inicialmente se opunha ao pedágio, votou
pela manutenção permanente no referendo --- um caso raro de preferência
revelada \emph{ex post} superando resistência \emph{ex ante}.

\textbf{Conexão com o capítulo.} O paper demonstra empiricamente a
eficácia do imposto pigouviano para externalidades de congestionamento
(Seção 20.4.1). A mudança de opinião pública ilustra como a experiência
concreta com a redução da externalidade pode alterar preferências
políticas.

\end{pesquisabox}

\begin{pesquisabox}{[Fehr, E.; Gächter, S. (2000). Cooperation and Punishment in Public Goods Experiments. *American Economic Review*, 90(4), 980–994.](https://doi.org/10.1257/aer.90.4.980)}

\textbf{DOI:}
\href{https://doi.org/10.1257/aer.90.4.980}{10.1257/aer.90.4.980}

\textbf{Contexto.} A teoria prevê que contribuições voluntárias a bens
públicos convergem para zero (free-riding). Mas isso ocorre na prática?
E mecanismos de punição podem sustentar a cooperação?

\textbf{Método.} Experimentos de contribuição voluntária com e sem
possibilidade de punição entre pares. Participantes decidem quanto
contribuir ao fundo público e, na condição com punição, podem gastar
para penalizar free-riders.

\textbf{Resultado.} Sem punição, contribuições declinam de
\textasciitilde50\% para \textasciitilde20\% do ótimo ao longo de 10
rodadas. Com punição, contribuições \emph{aumentam} para 60-90\% e se
sustentam. A punição é usada ativamente contra free-riders, mesmo sendo
custosa para quem pune (``punição altruísta''). O resultado sugere que
normas sociais de reciprocidade e justiça são forças poderosas na
provisão de bens públicos.

\textbf{Conexão com o capítulo.} O paper é a referência central do
experimento de sala de aula proposto na abertura do capítulo (Seções
20.5-20.8). Demonstra que o problema do carona é real, mas mitigável por
mecanismos sociais.

\end{pesquisabox}

\begin{pesquisabox}{[Schmalensee, R.; Stavins, R. N. (2013). The SO₂ Allowance Trading System: The Ironic History of a Grand Policy Experiment. *Journal of Economic Perspectives*, 27(1), 103–122.](https://doi.org/10.1257/jep.27.1.103)}

\textbf{DOI:}
\href{https://doi.org/10.1257/jep.27.1.103}{10.1257/jep.27.1.103}

\textbf{Contexto.} O programa americano de cap-and-trade para SO₂ (Acid
Rain Program, 1990) é frequentemente citado como o caso de maior sucesso
de política ambiental baseada em instrumentos de mercado. Mas qual é o
balanço completo?

\textbf{Método.} Revisão abrangente dos resultados econômicos e
ambientais do programa, incluindo custos de compliance, padrões de
comércio de permissões e reduções de emissão.

\textbf{Resultado.} Emissões caíram de 15,7 milhões de toneladas (1990)
para menos de 3 milhões (2010) --- redução de \textasciitilde80\%. Os
custos de compliance foram 15-50\% menores que estimativas ex ante. A
``ironia'' do título refere-se ao declínio do programa como modelo
político: apesar do sucesso, o instrumento de mercado perdeu apoio
político nos EUA, enquanto regulação direta voltou a ser preferida.

\textbf{Conexão com o capítulo.} O paper valida empiricamente a teoria
de cap-and-trade (Seção 20.4.4) e ilustra a equivalência imposto/cap sob
certeza. A ``ironia'' política conecta-se ao Teorema de Arrow (Seção
20.9): preferências sociais sobre instrumentos de política podem ser
intransitivas.

\end{pesquisabox}

\begin{pesquisabox}{[Ostrom, E. (1990). *Governing the Commons: The Evolution of Institutions for Collective Action*. Cambridge University Press.](https://doi.org/10.1017/CBO9780511807763)}

\textbf{DOI:}
\href{https://doi.org/10.1017/CBO9780511807763}{10.1017/CBO9780511807763}

\textbf{Contexto.} A ``tragédia dos comuns'' (Hardin, 1968) sugere que
recursos de acesso aberto serão inevitavelmente esgotados. A única
solução seria privatização ou regulação estatal. Ostrom desafiou essa
dicotomia.

\textbf{Método.} Análise comparativa de dezenas de casos de gestão
comunitária de recursos comuns em diferentes países e épocas: florestas
na Suíça e no Japão, sistemas de irrigação na Espanha e no Nepal,
pesqueiros costeiros na Turquia e no Canadá.

\textbf{Resultado.} Comunidades frequentemente desenvolvem instituições
locais eficazes --- regras de uso, monitoramento entre pares, sanções
graduais --- que sustentam a cooperação por séculos. Ostrom identificou
8 ``princípios de design'' para gestão bem-sucedida de commons,
incluindo fronteiras bem definidas, congruência entre regras e condições
locais, participação dos usuários na governança e mecanismos de
resolução de conflitos.

\textbf{Conexão com o capítulo.} O livro complementa as Seções
20.7-20.8, mostrando que entre a solução de mercado (Lindahl) e a
solução estatal (imposto pigouviano) existe um terceiro caminho: a
governança comunitária. O exemplo da cobrança pelo uso da água no Brasil
(Box Brasil 20.2) reflete os princípios de Ostrom.

\end{pesquisabox}

\begin{pesquisabox}{[Greenstone, M.; Jack, B. K. (2015). Envirodevonomics: A Research Agenda for an Emerging Field. *Journal of Economic Literature*, 53(1), 5–42.](https://doi.org/10.1257/jel.53.1.5)}

\textbf{DOI:}
\href{https://doi.org/10.1257/jel.53.1.5}{10.1257/jel.53.1.5}

\textbf{Contexto.} Países em desenvolvimento enfrentam simultaneamente
pobreza e degradação ambiental severa. Como a teoria das externalidades
se aplica quando instituições são fracas, direitos de propriedade são
mal definidos e capacidade regulatória é limitada?

\textbf{Método.} Survey da literatura emergente em ``envirodevonomics''
--- a interseção entre economia ambiental e economia do desenvolvimento.
Os autores revisam evidências sobre: valoração de serviços ambientais,
custo da poluição para saúde e produtividade, e eficácia de políticas
ambientais em países em desenvolvimento.

\textbf{Resultado.} A poluição do ar em Delhi equivale a fumar 3
cigarros por dia; a água contaminada em Bangladesh causa perdas
cognitivas mensuráveis em crianças. Porém, instrumentos clássicos
(imposto pigouviano, cap-and-trade) têm implementação limitada em países
com capacidade fiscal e regulatória fraca. Soluções adaptadas ---
pagamentos por serviços ambientais, auditoria ambiental com incentivos,
informação pública sobre qualidade do ar --- mostram resultados
promissores.

\textbf{Conexão com o capítulo.} O paper contextualiza os instrumentos
de política das Seções 20.4 no mundo real de países em desenvolvimento,
onde os pressupostos do modelo (custos de transação baixos, capacidade
regulatória, enforcement) frequentemente não se verificam.

\end{pesquisabox}

\begin{center}\rule{0.5\linewidth}{0.5pt}\end{center}

\begin{tipbox}{Exercício com IA — Externalidades e Bens Públicos}

\textbf{Prompt sugerido para ChatGPT/Claude:}

\emph{``Considere uma cidade com 500.000 habitantes que enfrenta
congestionamento viário severo. O custo externo do congestionamento é
estimado em R\$15 por viagem no horário de pico. A elasticidade-preço da
demanda por viagens de carro é -0,3. (a) Projete um sistema de pedágio
urbano pigouviano: qual deveria ser a tarifa no horário de pico? (b)
Estime a redução no número de viagens. (c) Que problemas distributivos
surgem? (d) Compare com alternativas: ampliação do transporte público,
rodízio de placas, home office obrigatório.''}

Use o modelo de IA como interlocutor para explorar as múltiplas
dimensões do problema --- mas verifique as contas e questione as
premissas!

\end{tipbox}

\begin{center}\rule{0.5\linewidth}{0.5pt}\end{center}

\section{📚 Referências do
Capítulo}\label{referuxeancias-do-capuxedtulo-2}

\begin{itemize}
\tightlist
\item
  Arrow, Kenneth J. 1951.
  \href{https://books.google.com/books/about/Social_Choice_and_Individual_Values.html?id=1W0PAQAAMAAJ}{\emph{Social
  Choice and Individual Values}}. New Haven: Yale University Press.
\item
  Coase, Ronald H. 1960. ``\href{https://doi.org/10.1086/466560}{The
  Problem of Social Cost}.'' \emph{Journal of Law and Economics} 3:
  1--44.
\item
  Eliasson, Jonas, Lars Hultkrantz, Lena Nerhagen, e Lena Smidfelt
  Rosqvist. 2009. ``\href{https://doi.org/10.1016/j.tra.2008.09.007}{The
  Stockholm Congestion-Charging Trial 2006: Overview of Effects}.''
  \emph{Transportation Research Part A} 43 (3): 240--250.
\item
  Fehr, Ernst, e Simon Gächter. 2000.
  ``\href{https://doi.org/10.1257/aer.90.4.980}{Cooperation and
  Punishment in Public Goods Experiments}.'' \emph{American Economic
  Review} 90 (4): 980--994.
\item
  Greenstone, Michael, e B. Kelsey Jack. 2015.
  ``\href{https://doi.org/10.1257/jel.53.1.5}{Envirodevonomics: A
  Research Agenda for an Emerging Field}.'' \emph{Journal of Economic
  Literature} 53 (1): 5--42.
\item
  Ostrom, Elinor. 1990.
  \href{https://doi.org/10.1017/CBO9780511807763}{\emph{Governing the
  Commons: The Evolution of Institutions for Collective Action}}.
  Cambridge: Cambridge University Press.
\item
  Pigou, Arthur Cecil. 1920.
  \href{https://www.econlib.org/library/NPDBooks/Pigou/pgEW.html}{\emph{The
  Economics of Welfare}}. London: Macmillan.
\item
  Samuelson, Paul A. 1954. ``\href{https://doi.org/10.2307/1925895}{The
  Pure Theory of Public Expenditure}.'' \emph{Review of Economics and
  Statistics} 36 (4): 387--389.
\item
  Schmalensee, Richard, e Robert N. Stavins. 2013.
  ``\href{https://doi.org/10.1257/jep.27.1.103}{The SO₂ Allowance
  Trading System: The Ironic History of a Grand Policy Experiment}.''
  \emph{Journal of Economic Perspectives} 27 (1): 103--122.
\item
  Weitzman, Martin L. 1974.
  ``\href{https://doi.org/10.2307/2296698}{Prices vs.~Quantities}.''
  \emph{Review of Economic Studies} 41 (4): 477--491.
\end{itemize}

\newpage

\end{document}
